%%%%%%%%%%%%%%%%%%%%%%%%%%%%%%%%%%%%%%%%%%%%%%%%%%%%%%%%%%%%%%%%%
\chapter{Unit 3: Order of Operations}
\label{chap:unit3}
%%%%%%%%%%%%%%%%%%%%%%%%%%%%%%%%%%%%%%%%%%%%%%%%%%%%%%%%%%%%%%%%%

%%===============================================================
\section{Order of Operations}
\label{sec:order-of-operations}
%%===============================================================

The \textbf{\textit{Order of Operations}} is used when doing expressions with more than one operation (e.g., ×, +, -). These are rules so you only get one answer all the time.



Example: When faced with $4+2 \times 3$, how do you proceed?



There are two ways:



$$\begin{aligned}

4 + 2 \times 3 &= (4 + 2) \times 3 \\\\ 

4 + 2 \times 3 &= 6 \times 3 \\\\ 

4 + 2 \times 3 &= 18

\end{aligned}$$





\textbf{or}



$$\begin{aligned}

4 + 2 \times 3 &= 4 + (2 \times 3) \\\\ 

4 + 2 \times 3 &= 4 + 6 \\\\ 

4 + 2 \times 3 &= 10

\end{aligned}$$



This is confusing, so which is correct? (Parentheses, "(" and ")" are used to show what to do first)



In order to communicate using mathematical expressions we must agree on an \textit{order of operations} so that each expression has only one value.



For the above example all mathematicians agree the correct answer is 10.



You're probably wondering what this order is.



\paragraph{The Standard Order of Operations}



Evaluate expressions in this order.



\begin{itemize}
    \item Parentheses or Brackets (evaluate what's inside them)
    \item Exponents
    \item Multiplication and/or division from left to right
    \item Addition and/or subtraction from left to right
\end{itemize}

\paragraph{An Easy Way of Remembering}



Use this memory tool to help remember the order!

 $$\begin{aligned}

 \textbf{P} &\text{lease} \hspace{44em} \\\\ 

 \textbf{E} &\text{xcuse} \\\\ 

 \textbf{M} &\text{y} \\\\ 

 \textbf{D} &\text{ear} \\\\ 

 \textbf{A} &\text{nnoying} \\\\ 

 \textbf{S} &\text{ister}

 \end{aligned}$$

 It is also commonly called by its acronym, \textbf{PEMDAS}.



An alternative form of this is \textbf{BIDMAS}

 $$ \begin{aligned}

 \textbf{B} &\text{rackets} \hspace{44em} \\\\ 

 \textbf{I} &\text{ndices} \\\\ 

 \textbf{D} &\text{ivision or} \\\\ 

 \textbf{M} &\text{ultiplication} \\\\ 

 \textbf{A} &\text{ddition or} \\\\ 

 \textbf{S} &\text{ubtraction}

 \end{aligned}$$



Yet another way of remembering this is \textbf{B}rackets \textbf{O}rders \textbf{D}ivision \textbf{M}ultiplication \textbf{A}ddition \textbf{S}ubtraction (BODMAS)or \textbf{B}ring \textbf{O}ur \textbf{D}ear \textbf{M}other \textbf{A}long \textbf{S}aturday



\paragraph{Examples}



$$\begin{array}{c|c|c}

\textbf{Expression} & \textbf{Evaluation} & \textbf{Operation} \\\\  \hline

 4 \times 2 +1 &

  \begin{aligned} & = \textbf{4 $\times$ 2} +1 \\\\ & = \textbf{8 + 1} \\\\ & = \textbf{9} \end{aligned} & \begin{aligned} & \text{Multiplication} \\\\  & \text{Addition} \\\\ & \end{aligned}  \\\\  \hline

 12 - 9 \div 3 &

  \begin{aligned} & = 12 - \bf{9 \div 3} \\\\ & = \mathbf{12-3} \\\\ & =\mathbf{9} \end{aligned} & \begin{aligned} & \text{Division} \\\\ & \text{Subtraction} \\\\ & \end{aligned}  \\\\  \hline

 2 \times 9 \div 3 &

  \begin{aligned} & = \textbf{2 $\times$ 9} \div 3 \\\\ & = \textbf{18 $\div$ 3} \\\\ & = \textbf{6} \end{aligned} & \begin{aligned} & \text{Left to Right} \\\\ & \text{Division} \\\\ & \end{aligned}  \\\\  \hline

 9 \div 3 \times 3 &

  \begin{aligned} & = \textbf{9 $\div$ 3} \times 3 \\\\ & = \textbf{3 $\times$ 3} \\\\ & = \textbf{9} \end{aligned} & \begin{aligned} & \text{Left to Right} \\\\ & \text{Multiplication} \\\\ & \end{aligned} \\\\  \hline

 3 + 12 \div (5-2) &

  \begin{aligned} & = 3+12 \div \textbf{(5-2)} \\\\ & = 3+ \textbf{12 $\div$ 3} \\\\ & = \textbf{3+4} \\\\ & = \textbf{7} \end{aligned} & \begin{aligned} & \text{Parentheses} \\\\ & \text{Division} \\\\\ & \text{Addition} \\\\ & \end{aligned}  \\\\  \hline

 7 \times 10 - (2 \times 4)^2 &

  \begin{aligned} & = 7 \times 10 - \textbf{(2 $\times$ 4)}^2 \\\\ & = 7 \times 10 - \mathbf{8^2} \\\\ & = \mathbf{7 \times 10} - 64 \\\\ & = \textbf{70 - 64} \\\\ & = \textbf{6} \end{aligned} & \begin{aligned} & \text{Parentheses} \\\\ & \text{Exponents} \\\\ & \text{Multiplication} \\\\ & \text{Subtraction} \\\\ & \end{aligned}  \\\\  \hline \end{array}$$





\paragraph{Practice Problems}



\textbf{Note:} the expressions in the following quiz, use an asterisk (\*) to indicate multiplication ($\times$) between adjacent factors. This use of the asterisk is nearly ubiquitous with the various computer languages, as the \textit{times} symbol is not an historically available keyboard character. Hand written expressions commonly use a small vertically centered dot (·) to indicate multiplication. Where unambiguous, multiplication is \textit{implied} between factors and a symbol is extraneous. 



\textbf{Quiz}



1. Evaluate the numerical expression: $2+4*3$  

2. Evaluate the numerical expression: $2*4+3$  

3. Evaluate the numerical expression: $(2+4)*3$  

4. Evaluate the numerical expression: $9^2 + 1 - 7*(8+4)/2$



\textbf{Answers}



1. $2+4*3 = 14$  

2. $2*4+3 = 11$  

3. $(2+4)*3 = 18$  

4. $9^2 + 1 - 7*(8+4)/2 = 40$



\paragraph{Playing with Mathematics}



To get yourself thinking about this, try this simple mathematical game:



Take the numbers 1 through 10 on the left side of an equation, and pick a number for the right side.



Example: $1\hspace{1em}2\hspace{1em}3\hspace{1em}4\hspace{1em}5\hspace{1em}6\hspace{1em}7\hspace{1em}8\hspace{1em}9\hspace{1em}10 = 1 $



Now put operators between those numbers. Only use parentheses when necessary.



Example: $1 + 2 - 3 + 4 - 5 + 6 + 7 + 8 - 9 - 10 = 1$



Change the number on the right-hand side. Can you generate an expression for this number? If not, can you prove why not?



Does this change if you change the order of the numbers?



\paragraph{Licensing}



Content obtained from \href{https://en.wikibooks.org/w/index.php?title=Algebra/Order_of_Operations\&oldid=3662722}{Order of Operations, Algebra, at Wikibooks} under a CC BY-SA license

\subsection{Learning Objectives}

\begin{enumerate}
    \item Here are three Student Learning Objectives (SLOs) related to the \textit{Order of Operations} concept:
    \item \textbf{Apply the Order of Operations:}
    \item \textbf{Outcome:} Students will correctly apply the order of operations (PEMDAS/BIDMAS) to simplify mathematical expressions.
    \item \textbf{Details:} This includes evaluating expressions by following the correct sequence: Parentheses/Brackets, Exponents/Indices, Multiplication and Division (left to right), and Addition and Subtraction (left to right). Students will demonstrate their understanding through the accurate simplification of complex expressions, ensuring that each expression yields only one correct answer.
    \item \textbf{Identify and Correct Mistakes in Order of Operations:}
    \item \textbf{Outcome:} Students will identify and correct errors in the simplification of mathematical expressions involving multiple operations.
    \item \textbf{Details:} Students will be presented with incorrectly simplified expressions and asked to diagnose the errors made. They will then re-evaluate the expressions using the correct order of operations to arrive at the correct solution. This objective emphasizes the importance of accuracy and the consequences of neglecting proper procedural order in mathematics.
    \item \textbf{Compare Different Methods of Simplification:}
    \item \textbf{Outcome:} Students will compare and contrast different methods of simplifying expressions, emphasizing the need for the standard order of operations.
    \item \textbf{Details:} Students will explore various methods of simplifying expressions and discuss why the standard order (PEMDAS/BIDMAS) ensures consistency and correctness. Through guided examples, students will understand that different approaches can lead to different results unless the standard order is followed, thus reinforcing the necessity of adhering to agreed-upon mathematical rules.
\end{enumerate}

\subsection{Assessment Instrument}

\begin{enumerate}
    \item \textbf{Assessment Instrument: Order of Operations Quiz} \textbf{Format:}
    \item \textbf{Type:} Mixed-format quiz (multiple-choice, short answer, and problem-solving questions)
    \item \textbf{Duration:} 45 minutes
    \item \textbf{Total Marks:} 30 points \textbf{Sections:}
    \item \textbf{Order of Operations (10 points):}
    \item \textbf{Question 1:} Simplify the following expressions using the correct order of operations:
    \item a. $7 + 4 \times 3^2 - 2$
    \item b. $(5 + 3) \times 2^2 + 7$
    \item c. $12 \div 4 + 6 \times (3 + 1)$
    \item d. $10 \times 2 + 6 \div 3 - 4$
    \item e. $(9 - 3) \times 2^3 - 8$
    \item \textbf{Problem Solving with Order of Operations (10 points):}
    \item \textbf{Question 2:} Solve the following word problems using the correct order of operations:
    \item a. Sarah buys 3 packs of pens and 2 notebooks. Each pack of pens costs `$`4 and each notebook costs `$`5. Calculate the total cost Sarah will pay. Use the expression: $3 \times 4 + 2 \times 5$
    \item b. A rectangle has a length of 7 units and a width of 3 units. The area of the rectangle is then multiplied by 2 and reduced by 5 units. Calculate the final result using the expression: $2 \times (7 \times 3) - 5$
    \item c. A recipe requires mixing 2 cups of flour with 3 times the number of cups of sugar. If you add another 4 cups of flour, what is the total number of cups used? Use the expression: $2 + 3 \times 1 + 4$
    \item \textbf{Application of Order of Operations in Complex Expressions (10 points):}
    \item \textbf{Question 3:} Evaluate the following expressions by applying the correct order of operations:
    \item a. $8 \times (5 + 2) - 9 \div 3$
    \item b. $15 - (4 \times 2) + 3^2$
    \item c. $(6 + 4) \div 2 \times (7 - 5) + 1$
    \item d. $5 + 3 \times (8 - 6) \div 2^2$
    \item e. $(4 \times 3^2) - (2 + 6) \times 3$ \textbf{Scoring Rubric:}
    \item \textbf{Order of Operations:} 2 points for each correctly simplified expression (total 10 points)
    \item \textbf{Problem Solving with Order of Operations:} 3.33 points for each correct solution (rounded down) (total 10 points)
    \item \textbf{Application of Order of Operations in Complex Expressions:} 2 points for each correctly simplified expression (total 10 points) This quiz is designed to assess students' understanding and application of the order of operations. It includes a mix of direct expression simplification, real-world problem-solving, and complex expression evaluation to ensure a comprehensive understanding of the concept.
\end{enumerate}

%%===============================================================
\section{Simplifying Exponents}
\label{sec:simplifying-exponents}
%%===============================================================

\paragraph{Properties of Exponents}



There are several properties of exponents which are frequently used to manipulate and simplify algebraic and arithmetic expressions.



$$\begin{align}

x^m x^n & = x^{m+n} \\\\ \dfrac{x^m}{x^n} & = x^{m-n} \\\\ (xy)^n & = x^n y^n \\\\ \left( \dfrac{x}{y} \right) ^n & = \dfrac{x^n}{y^n} \\\\ (x^m)^n & = x^{mn}

\end{align}$$





The first, third, and fifth properties may be extended to several factors, as follows:



$$\begin{aligned}

x^{m _1}x^{m _2}x^{m _3} \cdots x^{m _n} & = x^{m _1+m _2+m _3+ \, \cdots \, +m _n} \\\\  \\\\ (x _1x _2x _3 \cdots x _n)^m & = x _1^m x _2^m x _3^m \cdots x _n^m \\\\  \\\\ ( \cdots (((x^{m _1})^{m _2})^{m _3}) \cdots )^{m _n} & = x^{m _1m _2m _3 \cdots \, m _n}

\end{aligned}$$





\paragraph{Bases and exponents of one and zero}



Any number raised to an exponent of one equals itself. So, for example, $5^1 = 5$.



Any non-zero number raised to an exponent of zero equals one. So, for example, $5^0 = 1$.



Zero raised to any positive exponent is still zero. So, for example, $0^5 = 0$.



One raised to any exponent is still one. So, for example, $1^5 = 1$.



Zero raised to an exponent of zero is not defined.



\paragraph{Negative and noninteger powers}



A base may also be raised to a negative, fractional, or decimal power. These will be covered later in this lesson



\paragraph{Combining powers and roots}



The unit fraction notation used for roots previously may have given you the idea that roots are really the same as powers, only with a unit fraction (one over some number) instead of an integer as the exponent. Thus, the fractional notation is actually preferred in higher mathematics, although the root symbol is still used occasionally, especially for square roots.



\paragraph{Fractions as exponents}



Other (non-unit) fractions may also be used as exponents. In this case, the base number may be raised to the power of the numerator (top number in the fraction) then the denominator (bottom number) may be used to take the root. For example:



$8^{2/3} = \sqrt[3]{8^2} = \sqrt[3]{64} = 4$



Alternatively, you can take the root first and then apply the power:



$8^{2/3} = (\sqrt[3]{8})^2 = 2^2 = 4$



\paragraph{Decimal exponents}



Any fractional exponent can also be expressed as a decimal exponent. For example, a square root may also be written as:



$25^{1/2} = 25^{0.5} = 5 \;$



Also, decimals which can't be expressed as a fraction (irrational numbers) may be used as exponents:



$5^{\pi} = 5^{3.1415926...} = 156.99...\;$



Such problems aren't easy to solve by hand using basic math skills, but the answer can be estimated manually. In this case, since 3.1415926 is between 3 and 4 (and considerable closer to 3), we know that the answer will be between $5^3$ (or 125) and $5^4$ (or 625), and considerable closer to 125.



\paragraph{Negative exponents}



A negative exponent simply means you take the reciprocal (one over the number) of the base first, then apply the exponent:



$25^{-0.5} = 1/25^{0.5} = 1/5 \;$



Alternatively, you can first apply the exponent (ignoring the sign), then take the reciprocal:



$25^{-0.5} = 5^{-1} = 1/5 \;$



\paragraph{Fractions as bases}



When a fraction is raised to an exponent, both the numerator and denominator are raised to that exponent:



$(5/6)^2 = 5^2/6^2 = 25/36 \;$



Fractions may also be used for both the base and exponent:



$(25/36)^{1/2} = 25^{1/2}/36^{1/2} = 5/6 \;$



In addition, negative fractional exponents may be used, taking the reciprocal of the base, as always, to find the solution:



$(25/36)^{-1/2} = (36/25)^{1/2} = 36^{1/2}/25^{1/2} = 6/5 \;$



\paragraph{Negative bases}



Negative bases can be handled normally for integer powers:



$(-5)^2 = (-5) \times (-5) = 25 \;$



$(-5)^3 = (-5) \times (-5) \times (-5) = -125 \;$



$(-5)^4 = (-5) \times (-5) \times (-5) \times (-5) = 625 \;$



$(-5)^5 = (-5) \times (-5) \times (-5) \times (-5) \times (-5) = -3125 \;$



Note that negative bases raised to even powers produce positive results, while negative bases raised to odd powers produce negative results.



Be careful with negative signs. Since -5 = -1×5, there is a difference between $-5^2$ and $(-5)^2$. The former means the negative of 5 times 5, whereas the lattermeans -5 squared. In other words,



$-5^2 = -1\times5^2 = -1\times(5\times5) = -1\times25 = -25$



but



$(-5)^2=(-5)\times(-5)=25$



Roots and fractional/decimal powers are a bit trickier. Odd roots work out fine:



$(-125)^{1/3} = -5 \;$



Even roots, however, have no real solution:



$(-25)^{1/2} = ? \;$



$(-625)^{1/4} = ? \;$



Note that there is no real number, when multiplied by itself, which will produce -25, because 5×5 = 25 and -5×-5 = 25. There is actually a solution, called an \textbf{imaginary number}, but that won't be discussed until later lessons.



\paragraph{Principal root}



Note that, since both 5×5 = 25 and -5×-5 = 25, when we are asked to take the square root of 25 there are, in fact, two valid answers, 5 and -5. Actually, any even root of a positive number will have two solutions, with one being the negative of the other. This may seem unusual, but, in higher mathematics, problems often have multiple solutions.



However, for many problems, only the positive value seems to physically work. For example, if we are asked to figure the length of the sides of a square yard which has an area of 25 square units, only 5 units on a side works. If we said "each side can also have a length of -5 units", that doesn't make any sense. For this reason, the positive solution is called the \textbf{principal root}, and, depending on the question, may be the only desired answer. In cases where either answer is valid, it is sometimes written as ±5 (read as "plus or minus five"). However, the mathematical definition of the square root of x squared is the absolute value of x. Thus, square roots equations do not have two answers but two numbers can square to equal the same rational number.



\paragraph{Resources}



\begin{itemize}
    \item \href{https://www.youtube.com/watch?v=Zt2fdy3zrZU}{Simplifying Exponents}, The Organic Chemistry Tutor
\end{itemize}

\paragraph{Licensing}



Content obtained and/or adapted from:



\begin{itemize}
    \item \href{https://en.wikibooks.org/wiki/Primary_Mathematics/Powers,_roots,_and_exponents}{Powers, roots, and exponents; Wikibooks: Primary Mathematics} under a CC BY-SA license
\end{itemize}

\subsection{Learning Objectives}

\begin{enumerate}
    \item Here are three Student Learning Objectives (SLOs) for the topic "Properties of Exponents":
    \item \textbf{Apply the Product of Powers Property:}
    \item \textbf{Outcome:} Students will be able to apply the product of powers property to simplify expressions involving exponents.
    \item \textbf{Details:} This includes correctly using the rule $x^m \cdot x^n = x^{m+n}$ to combine exponential terms with the same base and simplifying algebraic expressions accordingly.
    \item \textbf{Utilize the Quotient of Powers Property:}
    \item \textbf{Outcome:} Students will be able to simplify expressions by applying the quotient of powers property.
    \item \textbf{Details:} This includes using the rule $$\frac{x^m}{x^n} = x^{m-n}$$ to divide exponential terms with the same base, and simplifying both numerical and algebraic expressions using this property.
    \item \textbf{Apply the Power of a Product Property:}
    \item \textbf{Outcome:} Students will be able to simplify expressions using the power of a product property.
    \item \textbf{Details:} This includes applying the rule $(xy)^n = x^n \cdot y^n$ to distribute exponents over products and simplify complex expressions involving multiple variables raised to a power.
\end{enumerate}

\subsection{Assessment Instrument}

\begin{enumerate}
    \item \textbf{Assessment Instrument: Simplifying Exponents Quiz} \textbf{Format:}
    \item \textbf{Type:} Mixed-format quiz (multiple-choice, short answer, and problem-solving questions)
    \item \textbf{Duration:} 45 minutes
    \item \textbf{Total Marks:} 30 points \textbf{Sections:}
    \item \textbf{Simplifying Exponents (10 points):}
    \item \textbf{Question 1 (Multiple-Choice):} Simplify the expression $x^3 \cdot x^5$ using the product of powers property.
    \item a. $x^{8}$
    \item b. $x^{15}$
    \item c. $x^{2}$
    \item d. $x^{5}$
    \item \textbf{Question 2 (Short Answer):} Simplify the expression $\frac{x^7}{x^4}$ using the quotient of powers property. Write your answer in terms of $x$.
    \item \textbf{Question 3 (Problem-Solving):} Simplify the expression $(2x^3y^2) \cdot (3x^2y^4)$ using the product of powers property for both bases $x$ and $y$. Show your work.
    \item \textbf{Combining Powers and Roots (10 points):}
    \item \textbf{Question 4 (Multiple-Choice):} Simplify the expression $(\sqrt[3]{8})^2$ using the properties of exponents.
    \item a. $4$
    \item b. $8$
    \item c. $64$
    \item d. $16$
    \item \textbf{Question 5 (Short Answer):} Express the value of $16^{3/4}$ as a simplified radical expression.
    \item \textbf{Question 6 (Problem-Solving):} Simplify the expression $\left(\frac{9}{16}\right)^{1/2}$ using the properties of exponents and fractions. Write your answer in simplest form.
    \item \textbf{Negative and Fractional Exponents (10 points):}
    \item \textbf{Question 7 (Multiple-Choice):} Simplify the expression $25^{-1/2}$.
    \item a. $5$
    \item b. $\frac{1}{25}$
    \item c. $\frac{1}{5}$
    \item d. $25$
    \item \textbf{Question 8 (Short Answer):} Simplify the expression $(-8)^{2/3}$ using the properties of exponents. Write your answer in simplest form.
    \item \textbf{Question 9 (Problem-Solving):} Simplify the expression $(\frac{27}{64})^{-2/3}$ using the properties of exponents and fractions. Show your work and provide the final answer in simplest form. \textbf{Scoring Rubric:}
    \item \textbf{Simplifying Exponents:} 3 points for each correct multiple-choice answer, 2.5 points for each correct short answer, 4 points for the problem-solving question (total 10 points)
    \item \textbf{Combining Powers and Roots:} 3 points for each correct multiple-choice answer, 2.5 points for each correct short answer, 4 points for the problem-solving question (total 10 points)
    \item \textbf{Negative and Fractional Exponents:} 3 points for each correct multiple-choice answer, 2.5 points for each correct short answer, 4 points for the problem-solving question (total 10 points) The quiz will help assess students' understanding and application of exponent properties, including simplification of expressions involving exponents, roots, and fractional powers. It ensures that they can accurately simplify expressions and perform the required operations with both positive and negative exponents.
\end{enumerate}

%%===============================================================
\section{Simplifying Radicals}
\label{sec:simplifying-radicals}
%%===============================================================

\paragraph{Radicals}



A \textbf{radical} is a special kind of number that is the root of a polynomial equation. First, let us look at one specific kind of radical, the "square root". It is a special kind of number related to squaring.



When we have a number, say 2, what positive number will give us, when we square it, the number 2?



\begin{itemize}
    \item 1 × 1 = 1, so the number must be higher than 1.
    \item 2 × 2 = 4, so the number must be lower than 2.
    \item 1.5 × 1.5 = 2.25, so the number must be lower than 1.5.
\end{itemize}

We could continue like this forever, and with each step get closer and closer to the answers (this kind of process for refining the answer is called iteration). The number we are looking for is approximately 1.41421\...



Obviously this is very difficult for us to work with, so we have a special notation. We write for a number \textit{a}, $\sqrt{a}$ to represent the number when squared will give us \textit{a} back.



Since this is the inverse operation from squaring, it can also be denoted as $a^{1/2}$ and $\newline$

$(a^{1/2})^2=a^1=a$.



We can extend this idea to other kinds of radicals. $\sqrt[3]{a}$ indicates the number $x$ such that $\newline$

$x^3=a$. For example, $\sqrt[3]{7}=1.91293...$



Note that it is not possible to find any real number whose square would be negative. Multiplying with a negative number changes the sign of the number being multiplied and two negative signs thus eliminate each other. For example, -7 × -7 = 49 and also 7 × 7 = 49 . Therefore the square root of a negative number is an undefined operation unless imaginary numbers are allowed as answers.



\paragraph{Simplifying Radicals}



To simplify a radical, you look for the largest perfect square factor of the number under the radical.



$${\sqrt{162}}$$ Once you find a perfect square factor, you can express the number under the radical as the product of two factors.



$${\sqrt{81*2}}$$ Then you can take the square root of the perfect square factor out of the radical and place it outside.



$$9{\sqrt{2}}$$ Once you have factored out all perfect square factors, you have simplified the radical.



\paragraph{Rationalizing the Denominator}



"Rationalizing the denominator" is simply taking the roots out of the denominator. This is needed because it is not proper to leave roots in the denominator. To rationalize the denominator you simply multiply the numerator and the denominator by the root that is in the denominator. For example:



${4\over\sqrt{5}}={4\sqrt{5}\over(\sqrt{5})^2}={4\sqrt{5}\over5}$



If you can simplify, do so.



$$\begin{aligned}

\frac{10}{\sqrt{5}} &= \frac{10\sqrt{5}}{(\sqrt{5})^2} &= \frac{10\sqrt{5}}{5} &= 2\sqrt{5}

\end{aligned}$$



$$\begin{aligned}

\frac{4}{\sqrt{8}} &= \frac{4}{\sqrt{4 \cdot 2}} &= \frac{4}{2\sqrt{2}} &= \frac{4\sqrt{2}}{2(\sqrt{2})^2} &= \frac{4\sqrt{2}}{4} &= \sqrt{2}

\end{aligned}$$





\paragraph{Applying Operations to Radicals}



\paragraph{Addition and Subtraction}



Adding and subtracting radicals can only be done when the numbers inside the radicals are the same. For instance, consider the following expression:



$${\sqrt{7}+\sqrt{28}}$$



You can not add those two terms as they are. However, the following equation simplifies the second term, and you get this:



$${\sqrt{7}+\sqrt{4*7}}={\sqrt{7}+2\sqrt{7}}$$



These terms can be added (or subtracted, for that matter). Also, coefficients of 1 are usually not written, so:



$${\sqrt{7}+2\sqrt{7}}={1\sqrt{7}+2\sqrt{7}}={3\sqrt{7}}$$



You can subtract radicals using the same procedure as adding, just instead of adding the coefficients, you subtract them.



The reason we can do this is the distributive property. Proving it is quite simple, and would look like this:



$${\sqrt{7}+2\sqrt{7}}$$



Remember that coefficients of 1 are not written, so this equation could also be written as:



$${1\sqrt{7}+2\sqrt{7}}$$



We then extract the term ${\sqrt{7}}$ from the expression:



$${\sqrt{7}*(1+2)}$$



We all know that ${1+2}={3}$, so:



$${\sqrt{7}*(3)}={3\sqrt{7}}$$



\paragraph{Multiplication and Division}



Multiplying and dividing radicals is quite simple. Multiplication requires very little work. When you are given two radicals to multiply, all you do is multiply the numbers inside the radicals and put the product under a radical. Consider the following equation:



$${\sqrt{11}*\sqrt{5}}={\sqrt{55}}$$



Division is a bit different, however. One way to divide radicals that uses a concept mentioned earlier on this page is to set up the division problem as a fraction. The following equation illustrates this concept using constants (regular numbers):



$${{1}\div{2}}={{1}\over{2}}={0.5}$$



The same concept applies to the following equation:



$${{\sqrt{10}}\div{\sqrt{3}}}={{\sqrt{10}}\over{\sqrt{3}}}$$



You can then rationalize the expression ${{\sqrt{10}}\over{\sqrt{3}}}$.



\paragraph{Resources}



\begin{itemize}
    \item \href{https://courses.lumenlearning.com/boundless-algebra/chapter/radicals/}{Introduction to Radicals}, Lumen Learning: Boundless Algebra
    \item \href{https://www.khanacademy.org/math/algebra/x2f8bb11595b61c86:rational-exponents-radicals/x2f8bb11595b61c86:simplifying-square-roots/v/simplifying-square-roots-1}{Simplifying Square Roots}, Khan Academy
    \item \href{https://www.mathwords.com/r/radical_rules.htm}{Useful Radical/Root Rules for Simplification}, Mathwords.com
    \item \href{https://math.libretexts.org/Bookshelves/Algebra/Book\%3A_Advanced_Algebra_(Redden}{Example Problems of Simplifying Radicals}/05\%3A_Radical_Functions_and_Equations/5.02\%3A_Simplifying_Radical_Expressions), LibreTexts
\end{itemize}

\paragraph{Licensing}



Content obtained from:



\begin{itemize}
    \item \href{https://en.wikibooks.org/wiki/Arithmetic/Radicals}{Radicals, Wikibooks: Arithmetic} under a CC BY-SA license
\end{itemize}

\subsection{Learning Objectives}

\begin{enumerate}
    \item \textbf{Simplifying Radicals:}
    \item \textbf{Outcome:} Students will be able to simplify radicals by factoring out the largest perfect square.
    \item \textbf{Details:} This includes identifying the perfect square factors of the radicand, expressing the radical as a product of square roots, and simplifying the expression by taking the square root of the perfect square factor.
    \item \textbf{Rationalizing the Denominator:}
    \item \textbf{Outcome:} Students will be able to rationalize the denominator of a fraction involving radicals.
    \item \textbf{Details:} This includes multiplying both the numerator and the denominator by the radical in the denominator to eliminate the radical, and simplifying the resulting expression.
    \item \textbf{Performing Operations with Radicals:}
    \item \textbf{Outcome:} Students will be able to perform addition, subtraction, multiplication, and division with radical expressions.
    \item \textbf{Details:} This includes recognizing when radicals can be added or subtracted by simplifying like terms, multiplying radicals by combining radicands, and dividing radicals by rationalizing the denominator if necessary.
\end{enumerate}

\subsection{Assessment Instrument}

\begin{enumerate}
    \item \textbf{Assessment Instrument: Simplifying Radicals Quiz} \textbf{Format:}
    \item \textbf{Type:} Mixed-format quiz (multiple-choice, short answer, and problem-solving questions)
    \item \textbf{Duration:} 45 minutes
    \item \textbf{Total Marks:} 30 points \textbf{Sections:}
    \item \textbf{Identifying Perfect Square Factors (10 points):}
    \item \textbf{Question 1:} Identify the largest perfect square factor of the following numbers and express each radical in its simplified form:
    \item a. $\sqrt{72}$
    \item b. $\sqrt{50}$
    \item c. $\sqrt{128}$
    \item d. $\sqrt{200}$
    \item e. $\sqrt{45}$
    \item \textbf{Scoring:} 2 points for each correct identification and simplification (total 10 points)
    \item \textbf{Simplifying Radicals (10 points):}
    \item \textbf{Question 2:} Simplify the following radicals by expressing the radicand as a product of perfect squares and taking the square root of the perfect square factors:
    \item a. $\sqrt{18}$
    \item b. $\sqrt{98}$
    \item c. $\sqrt{75}$
    \item d. $\sqrt{242}$
    \item e. $\sqrt{180}$
    \item \textbf{Scoring:} 2 points for each correctly simplified radical (total 10 points)
    \item \textbf{Problem-Solving with Radicals (10 points):}
    \item \textbf{Question 3:} Solve the following problems involving simplified radicals:
    \item a. Express $\sqrt{200}$ as a simplified radical and then add $\sqrt{200} + 3\sqrt{50}$.
    \item b. Simplify $\sqrt{128} \times \sqrt{2}$ and express your answer in simplified radical form.
    \item c. If $x = 2\sqrt{45}$, simplify the expression $3x - \sqrt{180}$.
    \item d. Simplify the expression $\sqrt{50} \times \sqrt{2}$ and then rationalize the denominator of the resulting expression $\frac{\sqrt{50} \times \sqrt{2}}{\sqrt{8}}$.
    \item \textbf{Scoring:} 2.5 points for each correct solution and simplification (total 10 points) \textbf{Scoring Rubric:}
    \item \textbf{Identifying Perfect Square Factors:} 2 points for correctly identifying the largest perfect square factor and expressing the radical in simplified form (total 10 points).
    \item \textbf{Simplifying Radicals:} 2 points for each correctly simplified radical (total 10 points).
    \item \textbf{Problem-Solving with Radicals:} 2.5 points for each correctly solved and simplified problem (total 10 points). This quiz assesses students' understanding of simplifying radicals by identifying perfect square factors, performing simplifications, and applying their knowledge in problem-solving scenarios.
\end{enumerate}

%%===============================================================
\section{Factoring Polynomials}
\label{sec:factoring-polynomials}
%%===============================================================

\paragraph{Useful Formulas}



Computing factors of polynomials requires knowledge of different formulas and some experience to find out which formula to be applied. Below, we give some important formulas:



${x^2-y^2=(x+y)(x-y)}$



${x^2+2xy+y^2=(x+y)^2}$



${x^2-2xy+y^2=(x-y)^2}$



${x^3-y^3=(x-y)(x^2+xy+y^2)}$



${x^3+y^3=(x+y)(x^2-xy+y^2)}$



${a^3+b^3+c^3-3abc=(a+b+c)(a^2+b^2+c^2-ab-ac-bc)}$



\paragraph{Methods}



Given the expression $x^2+3x+2$ , one may ask "what are the values of $x$ that make this expression 0?" If we factor we obtain



$$x^2+3x+2=(x+2)(x+1)$$



If $x=-1,-2$ , then one of the factors on the right becomes zero. Therefore, the whole must be zero. So, by factoring we have discovered the values of $x$ that render the expression zero. These values are termed "roots." In general, given a quadratic polynomial $px^2+qx+r$ that factors as



$$px^2+qx+r=(ax+c)(bx+d)$$



then we have that $x=-\frac{c}{a}$ and $x=-\frac{d}{b}$ are roots of the original polynomial.



A special case to be on the look out for is the difference of two squares, $a^2-b^2$ . In this case, we are always able to factor as



$$a^2-b^2=(a+b)(a-b)$$



For example, consider $4x^2-9$ . On initial inspection we would see that both $4x^2$ and $9$ are squares of $2x$ and $3$, respectively. Applying the previous rule we have



$$4x^2-9=(2x+3)(2x-3)$$



\paragraph{The AC method}



There is a way of simplifying the process of factoring using the AC method. Suppose that a quadratic polynomial has a formula of



$$ax^2+bx+c$$



If there are numbers $r$ and $q$ that satisfy both



$$r\cdot q=ac \text{ and } r+q=b$$



Then, we can rewrite the polynomial as



$$ax^2+bx+c=ax^2 + rx + qx + c$$



and factor out a common term from $(ax^2 + rx)$ and $(qx + c)$ to factor the polynomial.



For example, take the polynomial $2x^2 + 7x - 15$.



$$\begin{aligned}

ac &= 2(-15) = -30 \Leftrightarrow 10(-3) = -30 = ac \\\&\text{and } 10 + (-3) = 7 = b,

\end{aligned}$$

So, we can rewrite the polynomial as 

$$2x^2 + 10x - 3x - 15$$ 

From here, we can factor $2x$ from the first two terms, and $-3$ from the last two, which gives us 

$$2x(x + 5) + (-3)(x + 5) = (2x - 3)(x + 5)$$ 

Thus, the factored form of $2x^2 + 7x - 15$ is 

$$(2x - 3)(x + 5)$$ 

Note that if we reverse the order of $10x$ and $-3x$, we get 

$$\begin{aligned}

2x^2 - 3x + 10x - 15 

&= x(2x - 3) + 5(2x - 3) \\\&= (x + 5)(2x - 3) \ 

\end{aligned}$$

which is the same factored form that we got previously. Thus, the order of $rx$ and $qx$ should not matter when using the AC method to factor a quadratic polynomial.



\paragraph{The quadratic formula}



Given any quadratic equation $ax^2+bx+c=0\ ,\ a\ne0$, all solutions of the equation are given by the quadratic formula:



$$x=\frac{-b \pm \sqrt{b^2 - 4ac} }{2a}$$



Note that the value of $b^2-4ac$ will affect the number of \textit{real} solutions of the equation.



$$\begin{array}{|c|l|}

\hline

\textbf{If} & \textbf{Then} \\\\ 

\hline

b^2 - 4ac > 0 & \text{There are two real solutions for the equation} \\\\ 

\hline

b^2 - 4ac = 0 & \text{There is only one real solution for the equation} \\\\ 

\hline

b^2 - 4ac < 0 & \text{There are no real solutions for the equation} \\\\ 

\hline 

\end{array}$$



\paragraph{Remainder and Factor Theorem}



The polynomial division algorithm is as follows: suppose $d(x)$ and $p(x)$ are nonzero polynomials where the degree of $p(x)$ is greater than or equal to the degree of $d(x)$. Then there exist two unique polynomials, $q(x)$ and $r(x)$, such that $p(x) = d(x)q(x) + r(x)$, where either $r(x) = 0$ or the degree of $r(x)$ is strictly less than the degree of $d(x)$.



\paragraph{Remainder Theorem}



Suppose $p(x)$ is a polynomial of degree at least 1 and c is a real number. When $p(x)$ is divided by $(x - c)$ the remainder is $p(c)$.



> Proof: By the division algorithm, $p(x) = (x - c)q(x) + r$, where r must be a constant since $d(x) = x - c$ has a degree of 1. $p(x) = (x - c)q(x) + r$ must hold for all values of $x$, so we can set $x = c$ and get that $p(c) = (c - c)q(x) + r = r$. Thus the remainder $r = p(c)$.



\paragraph{Factor Theorem}



Suppose $p(x)$ is a nonzero polynomial. The real number $c$ is a zero of $p(x)$ if and only if $(x - c)$ is a factor of $p(x)$.



> By the division algorithm, $x - c$ is a factor of $p(x)$ if and only if $r = 0$. So, since $p(c) = r$ when $p(x)$ is divided by $x - c$, $x - c$ is a factor of $p(x)$ if and only if $p(c) = 0$; that is, if $c$ is a zero of $p(x)$.



\paragraph{Factor Theorem Example Problem}



> Determine if x + 2 is a factor of $2x^2 + 3x -2$.



Since c is positive instead of negative we need to use this basic identity:



$x + 2 = x - \left ( - 2 \right )$



Now we can use the factor theorem.



$2 \left (-2 \right )^2 + 3 \left (-2 \right ) -2 = 8 - 6 - 2 = 0$.



Since the resultant is 0, $(x+2)$ is a factor of $2x^2 + 3x -2$.



This means it is possible to re-state the polynomial in the form (x+2)( some linear expression of x).



So $2x^2 + 3x -2 = (x+2)(ax+b)$



Expanding the right hand side we get :



$2x^2 + 3x -2 = ax^2 + x( 2a+b) +2b$



Equating like terms we get :



$2= a$



$2a+b = 3$, and



$2b = -2$



Giving $a = 2$, $b= -1$ from the first and third equations and this works in the second, so



$2x^2 + 3x -2 = (x+2)(2x-1)$



\paragraph{Example Problems}



EXAMPLE 1: Find all the roots of $4x^2+7x-2$



> Finding the roots is equivalent to solving the equation $4x^2+7x-2=0$ . Applying the quadratic formula with $a=4\ ,\ b=7\ ,\ c=-2$, we have:



$$

\begin{aligned}x &= \frac{-7\pm\sqrt{7^2-4(4)(-2)}}{2(4)} \\\x &= \frac{-7\pm\sqrt{49 + 32}}{8} \\\x &= \frac{-7\pm\sqrt{81}}{8} \\\x &= \frac{-7\pm 9}{8} \\\x &= \frac{2}{8}\ ,\ x = \frac{-16}{8} \\\x &= \frac{1}{4}\ ,\ x = -2 \ 

\end{aligned}$$





> The quadratic formula can also help with factoring, as the next example demonstrates.



EXAMPLE 2: Factor the polynomial $4x^2+7x-2$



> We already know from the previous example that the polynomial has roots $x=\frac{1}{4}$ and $x=-2$ . Our factorization will take the form    $C(x+2)\left(x-\tfrac{1}{4}\right)$    .

> All we have to do is set this expression equal to our polynomial and solve for the unknown constant C:



$$\begin{aligned}

C(x+2)\left(x-\tfrac{1}{4}\right) &= 4x^2 + 7x - 2 \\\C\left(x^2 + \left(-\tfrac{1}{4} + 2\right)x - \tfrac{2}{4}\right) &= 4x^2 + 7x - 2 \\\C\left(x^2 + \tfrac{7}{4}x - \tfrac{1}{2}\right) &= 4x^2 + 7x - 2 \ 

\end{aligned}$$





> You can see that $C=4$ solves the equation. So the factorization is



$$\begin{aligned}

4x^2 + 7x - 2 &= 4(x+2)\left(x-\tfrac{1}{4}\right) \\\&= (x+2)(4x-1) \ 

\end{aligned}$$





EXAMPLE 3: Factor the polynomial $6x^2 + 7x + 2$.



> For this problem, let\'s use the AC method. $a = 6$, $b = 7$, and $c = 2$. So, $ac = 6(2) = 12$, and we need to find two numbers whose product equals 12, and whose sum equals 7. Both ac and b are positive, so we are only concerned with the positive factors of 12, which are 

$$1, 2, 3, 4, 6, \text{ and } 12$$

$$3(4) = 12 = ac \quad \text{ and } \quad 3 + 4 = 7 = b,$$ so we can rewrite the polynomial as 

$$\begin{aligned}

6x^2 + 3x + 4x + 2 &= 3x(2x + 1) + 2(2x + 1) \\\&= (3x + 2)(2x + 1) \ 

\end{aligned}$$



\paragraph{Resources}



\begin{itemize}
    \item \href{https://tutorial.math.lamar.edu/classes/alg/factoring.aspx}{Factoring Polynomials with Example Problems}, Paul\'s Online Notes
    \item \href{https://courses.lumenlearning.com/suny-beginalgebra/chapter/5-2-1-factor-trinomials/}{Factoring Polynomials}, Lumen Learning
    \item \href{https://www.youtube.com/watch?v=shHS-ERFVrc}{Difference of Squares}. Produced by TA Catherine Sporer, UTSA
    \item \href{http://www.algebralab.org/lessons/lesson.aspx?file=algebra_factoring.xml}{6 Methods of Factoring with Practive Problems}, AlgebraLAB
    \item \href{https://txwes.edu/media/twu/content-assets/images/academics/academic-success-center/A-C-Method.pdf}{AC Factoring Method}, Texas Wesleyan University
\end{itemize}

\paragraph{Licensing}



Content obtained and/or adapted from:



\begin{itemize}
    \item \href{https://en.wikibooks.org/wiki/Calculus/Algebra}{Algebra, Wikibooks: Calculus} under a CC BY-SA license
    \item \href{https://en.wikibooks.org/wiki/Algebra/Factoring_Polynomials}{Factoring Polynomials, Wikibooks: Algebra} under a CC BY-SA license
\end{itemize}

\subsection{Learning Objectives}

\begin{enumerate}
    \item Here are three Student Learning Objectives (SLOs) based on the material provided:
    \item \textbf{Apply Formulas for Polynomial Factoring:}
    \item \textbf{Outcome:} Students will be able to accurately apply key polynomial factoring formulas, such as the difference of squares, perfect square trinomials, and sum or difference of cubes.
    \item \textbf{Details:} This includes using formulas like $x^2-y^2=(x+y)(x-y)$ and $x^3-y^3=(x-y)(x^2+xy+y^2)$ to factor various polynomial expressions and demonstrate understanding through problem-solving.
    \item \textbf{Utilize the AC Method for Factoring Quadratics:}
    \item \textbf{Outcome:} Students will be able to factor quadratic polynomials using the AC method.
    \item \textbf{Details:} This involves identifying appropriate values for $r$ and $q$ that satisfy $r\cdot q=ac$ and $r+q=b$, then using these to rewrite and factor the quadratic expression.
    \item \textbf{Solve Quadratic Equations Using the Quadratic Formula:}
    \item \textbf{Outcome:} Students will be able to solve quadratic equations using the quadratic formula and interpret the results.
    \item \textbf{Details:} This includes calculating the discriminant ($b^2-4ac$) to determine the number of real solutions, then applying the quadratic formula to find the exact roots of the equation.
\end{enumerate}

\subsection{Assessment Instrument}

\begin{enumerate}
    \item \textbf{Assessment Instrument: Factoring Polynomials Quiz} \textbf{Format:}
    \item \textbf{Type:} Mixed-format quiz (multiple-choice, short answer, and problem-solving questions)
    \item \textbf{Duration:} 45 minutes
    \item \textbf{Total Marks:} 30 points \textbf{Sections:}
    \item \textbf{Recognizing and Applying Factoring Formulas (10 points):}
    \item \textbf{Question 1:} Identify the correct factorization of the following polynomials:
    \item a. $x^2 - 9$ (2 points)
    \item i. $(x - 3)(x + 3)$
    \item ii. $(x - 9)(x + 1)$
    \item iii. $(x + 9)(x - 1)$
    \item b. $x^3 - 8$ (2 points)
    \item i. $(x - 2)(x^2 + 2x + 4)$
    \item ii. $(x - 2)(x^2 - 2x + 4)$
    \item iii. $(x + 2)(x^2 + 4)$
    \item \textbf{Factoring Quadratics Using the AC Method (10 points):}
    \item \textbf{Question 2:} Factor the following quadratic polynomials using the AC method:
    \item a. $2x^2 + 7x + 3$ (5 points)
    \item b. $3x^2 + 14x + 8$ (5 points)
    \item \textbf{Solving Quadratic Equations Using the Quadratic Formula (10 points):}
    \item \textbf{Question 3:} Solve the following quadratic equations using the quadratic formula and provide the exact roots:
    \item a. $x^2 + 6x + 9 = 0$ (5 points)
    \item b. $2x^2 - 3x - 2 = 0$ (5 points) \textbf{Scoring Rubric:}
    \item \textbf{Recognizing and Applying Factoring Formulas:} 2 points for each correctly identified factorization (total 10 points).
    \item \textbf{Factoring Quadratics Using the AC Method:} 5 points for each correctly factored polynomial (total 10 points).
    \item \textbf{Solving Quadratic Equations Using the Quadratic Formula:} 5 points for each correctly solved equation with correct roots (total 10 points). This quiz assesses students' understanding of key factoring techniques and their ability to apply these methods to factor polynomials and solve quadratic equations accurately.
\end{enumerate}

%%===============================================================
\section{Linear Equations}
\label{sec:linear-equations}
%%===============================================================

\paragraph{What are Linear Equations?}



In the functions section we talk about how a function is like a box that takes an independent input value and uses a rule defined mathematically to create a unique output value. The value for the output is dependent on the value that is put in the box. We call the values that are going into the box the independent values or the domain. We call the values coming out of the box the dependent values or the range.



Unless we specify differently, on Cartesian graphs the domain is the real numbers. In the Cartesian Plane section we saw how running different values through a function to identify the points on the Cartesian plane by picking the first (x) value of the point the domain and the second (y) value from the range. To restate this: by convention the two variables for a function on the Cartesian Plane are x for the domain, the independent variable, and y for the range, the dependent variable. The variable y is the same as writing $f(x)$. Mathematicians recognize this equivalence but generally prefer to write y because its shorter. Because a function definition has an input and an output it must also contain an equal sign. The section in this book solving equations showed the various operations we can perform on both sides of an equal sign and still maintain the notion of equivalence. In this section we plug different values into the independent variable and solve to find the associated dependent variable. For instance if we start with the equation:



$$\begin{aligned}

\text{Given: } & y = x \\\\ 

\text{Subtract } x \text{ from both sides: } & y - x = x - x \\\\ 

\text{Simplify: } & y - x = 0 \\\\ 

\text{Alternatively, subtract } y \text{ from both sides: } & y - y = x - y \\\\ 

\text{Simplify: } & 0 = x - y \\\\ 

\text{Since } & y - x = 0 \equiv 0 = x - y, \text{ then: } \\\\ 

& y - x = 0 = x - y \\\\ 

\text{Using the transitive property: } & y - x = x - y \\\\ 

\text{Add } x + y \text{ to both sides: } & (y - x) + (y + x) = (x - y) + (y + x) \\\\ 

\text{Using the associative property: } & (y - y) + (x + x) = (x - x) + (y + y) \\\\ 

\text{Simplify: } & 2x = 2y \\\\ 

\text{Divide both sides by 2: } & \frac{2x}{2} = \frac{2y}{2} \\\\ 

\text{Simplify to show: } & x = y

\end{aligned}$$



We have not really proved anything mathematically above, but these operations allow us to manipulate equations to get the dependent variable by itself on one side of the equals sign. Then we can plug numbers into the independent variable to discover the function values for those numbers. Then we can draw these values as points on the Cartesian plane and get a feel for what the function would look like if we could see all the points defined by the function at once.



\paragraph{Horizontal Linear Equations}



Equations of the form $y = C _2$ are linear functions of the general form $y = m x + b$ where slope $m = 0$ and the constant $C _2$ \textbf{is} the y-intercept b (in the general form). The graph of this zero-slope function is a straight horizontal line, intercepting the y-axis at $C _2$, including zero and extends infinitely in the positive and negative directions for all $\mathbb{R}$ values of x (see the following diagram).



% [Image: y = c, where c = 3]



Figure: $y = c$, where $c = 3$



The domain for such functions is $\mathbb{R}$ covering all real numbers (unless otherwise specified), but the range is just the set $\\{ c \\}$. The equation $y = 0$ is the x-axis.



\paragraph{Vertical Linear Equations}



Equation $x = C _1$, x is one single value $C _1$ and y, being unrestricted, is every $\mathbb{R}$ number. The graph of $x = C _1$ is a straight vertical line where $x = C _1$, covering \textbf{all} positive, negative and zero values of y (see the following diagram).



% [Image: x = c, where c = 3]



Figure: $x = c$, where $c = 3$



Its domain is set $\\{ C _1 \\}$ and range is set $\mathbb{R}$ (unless otherwise specified). $x = C _1$ is technically not a function (there is more than one value of y for each value of x), but it's a relation. Vertical lines have \textbf{no} slope (m = divide by zero, undefined, plus and minus infinity). These are the only types of linear equations of the general form shown previously which are not linear functions. The equation $x = 0$ is \textbf{exactly} the y-axis. Lines with steepness approaching vertical have very large-magnitude slopes but are still functions.



\paragraph{General Linear Equations}



We are going to start by looking at simple functions called \textbf{linear equations}. When none of the instances of x and y in the algebraic expression defining the function rule have exponents then all the instances of x and y can be combined into just two occurrences. the graph of the expression can be represented as a straight line. The equation that expresses the function is considered a \textbf{linear equation} with two variables. The following equation is a simple example of such a linear equation:



$$y - x = 2 \,$$



Since y is the dependent variable it is standing in for the function. We can re-write the expression as $f(x) - x = 2$. If we add an x to both sides the equality property holds and we get the expression $f(x) - x + x = x + 2$. Simplifying we get $f(x) = x + 2$. In the following table we'll pick 3 values for x, and then calculate the dependent (y) values from $f(x)$.



$$\begin{array}{|c|c|c|}

\hline

\textbf{x value} & \textbf{y value (abscissa)} & \textbf{Coordinates (x,y)} \\\\ 

\hline

-1 & 1 & (-1,1) \\\\ 

\hline

0 & 2 & (0,2) \\\\ 

\hline

1 & 3 & (1,3) \\\\ 

\hline 

\end{array}$$



where x and y are variables to be plotted in a two-dimensional Cartesian coordinate graph as shown here:



% [Image: ]



% [Image: ]



This function is equivalent to the previous example of a linear equation, y - x = 2. The arrows at each end of the line indicate that the line extends infinitely in both directions. All linear functions of a single input variable have or can be algebraically arranged to have the general form:



$$y = f(x) = m x + b \,$$



where x and y are variables, $f(x)$ is the function of x, m is a constant called the \textbf{slope} of the line, and b is a constant which is the ordinate of the \textbf{y-intercept} (i. e. the value of y where the function line crosses the y-axis). The slope indicates the steepness of the line. In the previous example where y = x + 2, the slope m = 1 and the y-intercept ordinate $b = 2$. The $y = mx+b$ form of a linear function is called the \textbf{slope-intercept form}.



Unless a domain for x is otherwise stated, the domain for linear functions will be assumed to be all real numbers and so the lines in graphs of all linear functions extend infinitely in both directions. Also in linear functions with all real number domains, the range of a linear function will cover the entire set of real numbers for y, unless the slope m = 0 and the function equals a constant. In such cases, the range is simply the constant.



Conversely, if a function has the general form $y = mx + b$ or if it can be arranged to have that form, the function is linear. A linear equation with two variables has or can be algebraically rearranged to have the general form:



$$A x + B y = C \,$$



where x and y represent the linear variables, and the letters A, B, and C can represent any real constants, either positive or negative. Conversely, an equation with two variables x and y having that general form, or being able to be arranged in that form, would be linear as long as A and B are not both equal to 0. In the preceding equation, capital letters are to avoid confusion with other constants in this chapter and for consistency with Reference 1.



If one divides the preceding equation by B (when B is not 0) and solves for y, the following form can be obtained:



$$y = (-A/B) x + (C/B) \,$$



If one equates $-A/B$ to the slope m and $C/B$ to the y-intercept ordinate b, it can be seen that the general form for a linear equation and the slope-intercept form for a linear function are practically interconvertible except for the fact that, in a linear function, the B constant in the linear equation form cannot equal 0.



\paragraph{X and Y axis intercepts}



An axis intercept point is a point where the graph of a function, relation, or equation intersects the X or Y axes. This section is about finding out how a particular set of functions: linear functions cross the axes.



We know that the domains of most lines are infinite because they are defined at every value of X. The exception is lines that are defined as $X=c$ where c is a number we choose when we write the function. By definition this line is only defined for one value of X. Since the domain maps onto more than one value for the range this is actually a relationship and not a function. We've seen the graph of this relationships is a vertical line that passes through the point $(c,Y)$.



Lines with the equation $X=C$ intersect the X axis once and the Y axis 0 or infinite times.



We can also restrict the range of a function by simply writing $Y=c$ where c is again any number we choose. The graph of this line is a horizontal line that passes through the point $(X,c)$. When $Y=0$ this line is the same as the X axis. when $c \ne 0$, then the line can never intercept the X axis.



Lines with the equation $Y=C$ intersect the Y axis once and the X axis 0 or infinite times.



When looking at Cartesian graphs and linear equations we run into a mathematical axiom: "Two points determine a line.". We will see how this axiom affects the slope-intercept definition of a line $y= f(x) = m x + b$ in the next section. When two lines intersect they intersect at a point. If a line is not horizontal or perpendicular it will have to intersect the X and Y axes once, but only once.



In this page we are going to accept the statement"At most one line can be drawn through any point not on a given line parallel to the given line in a plane."



We've seen that for the equation $Y=mX + b$ the Y intercept will always be at b because that is where X=0.



Using Algebra we can subtract b from both sides: $Y - b = mX$



and multiply by $\frac {1}{m}$



$$\frac {Y}{m} - \frac {b}{m} = X$$



we can see that the X intercept is going to be $-\frac{b}{m}$.



An axis intercept may simply refer to the number value on the axis where the intersection occurs. For brevity we may say the line has an X intercept of 1 and a Y intercept of 2. After graphing just a few lines you will be able to tell this line points down and runs through quadrants II, I, and IV. With a little more practice you will be able to know that the equation for the line is Y=-2x + 2. We will see that by specifying the two points we are actually implying the slope of the line. There is an exception to this rule. If we say a line crosses the axes at 0 we know that the line will pass through 2 quadrants instead of 3, but we won't know which quadrants or how steep the line is. When we look at slope in the next section we will see why the equations above specify a point and a slope.



When you are trying to graph a linear equation finding the axes intercepts is often the easiest way to go about doing it. To find the x-intercept, set y = 0 and solve for x. To find the y-intercept, set $x = 0$ and solve for y. For most examples the intercepts are different points, and a line can be drawn through the two intercepts. If both intercepts are $(0,0)$, then another point must be determined to graph the line. If the equations is in the form $x = c$ or $y = c$, the horizontal or vertical lines are very simple to plot.



\paragraph{Slope}



\textbf{Algebra/Slope}



Slope is the change in the vertical distance of a line on a coordinate plane over the change in horizontal difference. In other words, it is the "rise" over the "run" or the steepness of a line. Slope is usually represented by the symbol $m$ like in the equation $y = mx + b$, m the coefficient of x represents the slope of the line.



Slope is computed by measuring the change in vertical distance divided by the change in horizontal difference, i.e.:



$$\ m = \frac{\Delta y}{\Delta x} = \frac{\text{rise}}{\text{run}}$$



The Greek uppercase letter $\Delta$ represents change, in this case change in the y-coordinates divided by the change in x-coordinates.



\textbf{Positive Slope/ Negative Slope}



If a line goes up from left to right, then the slope has to be positive. For example, a slope of `¾` would have a "rise" of 3, or go up 3; and a "run" of 4, or go right 4. Both numbers in the slope are either negative or positive in order to have a positive slope.



If a line goes down from left to right, then the slope has to be negative. For example, a slope of -3/4 would have a "rise" of -3, or go down 3; and a "run" of 4, or go right 4. Only one number in the slope can be negative for a line to have a negative slope.



\textbf{Other Types of Slope}



There are two special circumstances, no slope and slope of zero. A horizontal line has a slope of 0 and a vertical line has an undefined slope.



Horizontal lines have the form: $\ y = a$ ; where a is a constant, i.e. $a \in \mathbb{R}$Vertical lines have the form: $\ x = a$ ; where a is a constant, i.e. $a \in \mathbb{R}$



\paragraph{Determining Slope}



To determine the slope you need some information. This can include two (or more) coordinates, a parallel slope and a coordinate, a perpendicular slope and a coordinate, or the y-intercept and slope.



For completely horizontal lines, the difference in y coordinates between any two points is 0, so the slope $m = 0$, indicating no steepness in the line at all. If the line extends between right-upper $(+,+)$ and left-lower $( -, -)$ directions, then the slope is positive. As the slope increases, the line becomes steeper until the line is almost vertical when the slope is very large. When the slope $m = 1$, the line is diagonal with an angle halfway between the x and y axes. If the line extends between left-upper $(-,+)$ and right-lower $(+, -)$ directions, then the slope is negative. As the slope changes from 0 to very negative numbers, the steepness in the opposite direction increases. Compare the slope ( m ) values in the following graph of functions

$$y = 1 ( \text{where } m = 0), \quad y = \frac{1}{2} x + 1, \quad y = x + 1, \quad y = 2 x, \quad y = -\frac{1}{2} x + 1, \quad y = -x + 1, \quad \text{and } y = -2 x + 1$$

For all two-variable linear equations that can be converted to linear functions, the same calculation applies to slopes for those lines.



% [Image: ]



For the most part finding slope when given information is a simple matter. Simply take the slope equation y=mx+b and replace the variable with whatever information you know, and solve.



\paragraph{Two Coordinates}



To find the slope with two coordinates, you must first find the slope. Use the standard equation $\frac{y _2-y _1}{x _2-x _1}$. Put that into the equation as m, and replace x and y with x and y from one of the coordinates. Solve for b. Put that into the equation and you're done.



Example: (1,4) (4,8)



$m = \frac{y _2-y _1}{x _2-x _1}$



$m = \frac{4-1}{4-8}$



$m = \frac{3}{-4}$



Plug that right in.



$y = \frac{3}{-4}x+b$



$4 = \frac{3}{-4}(1)+b$



$4 = \frac{3}{-4}+b$



$4-\frac{3}{-4} = b$



$\frac{-19}{4} = b$



Put that in the equation and you're done.



$y = \frac{3}{-4}x + \frac{-19}{4}$



\paragraph{Parallel Lines}



If you have to find the slope of a line(Let's say AB) which is parallel to line (Let's say XY) then using the coordinates of line XY you can find the coordinates of slope of line AB As, Slope of line of a line parallel to another line is equal, i.e. Slope of AB = Slope of XY



\textbf{Example:}



AB and XY are PARALLEL Lines. Find the slope of AB.



Let line XY have the coordinates:-



$X(2,4)=(x1,y1)$ and $Y(3,6)=(x2,y2)$



Slope of XY



$m = (y2 - y1)/(x2-x1) = (6-4)/(3-2) = 2/1 \Rightarrow m = 2$



Slope of AB = Slope of XY, so Slope of AB = 2



\paragraph{Resources}



\begin{itemize}
    \item \href{https://courses.lumenlearning.com/beginalgebra/chapter/introduction-to-solving-linear-equations/}{Solving Linear Equations}, Lumen Learning
    \item \href{https://openstax.org/books/intermediate-algebra-2e/pages/2-1-use-a-general-strategy-to-solve-linear-equations}{General Strategy to Solving Linear Equations}, OpenStax
    \item \href{https://openstax.org/books/intermediate-algebra-2e/pages/2-3-solve-a-formula-for-a-specific-variable}{Solving Linear Equations for a Specific Variable}, OpenStax
    \item \href{https://www.youtube.com/watch?v=bsGyJ0GeRy8}{Solving Linear Equations: Example 1}. Produced by TA Catherine Sporer, UTSA
    \item \href{https://www.youtube.com/watch?v=nGZcfHQEdqo}{Solving Linear Equations: Example 2}. Produced by TA Catherine Sporer, UTSA
\end{itemize}

\paragraph{Licensing}



Content obtained and/or adapted from:



\begin{itemize}
    \item \href{https://en.wikibooks.org/wiki/Algebra/Linear_Equations_and_Functions}{Linear Equations and Functions, Wikibooks: Algebra} under a CC BY-SA license
    \item \href{https://en.wikibooks.org/wiki/Algebra/Intercepts}{Intercepts, Wikibooks: Algebra} under a CC BY-SA license
    \item \href{https://en.wikibooks.org/wiki/Algebra/Slope}{Slope, Wikibooks: Algebra} under a CC BY-SA license
\end{itemize}

\subsection{Learning Objectives}

\begin{enumerate}
    \item \textbf{Understanding the Concept of Linear Equations:}
    \item \textbf{Outcome:} Students will be able to define a linear equation and identify its basic components, including variables, constants, and the equal sign.
    \item \textbf{Details:} This includes understanding the general form of linear equations ($y = mx + b$), recognizing the role of the slope (m) and the y-intercept (b), and differentiating between linear equations and other types of equations.
    \item \textbf{Graphing Linear Equations:}
    \item \textbf{Outcome:} Students will be able to graph a linear equation on a Cartesian plane by identifying the slope and y-intercept and plotting the corresponding line.
    \item \textbf{Details:} This includes the ability to use the slope-intercept form ($y = mx + b$) to determine the steepness of the line and the point where the line crosses the y-axis, as well as drawing the line through calculated points.
    \item \textbf{Determining the Slope from Two Points:}
    \item \textbf{Outcome:} Students will be able to calculate the slope of a line given two points on the line.
    \item \textbf{Details:} This includes using the slope formula ($m = \frac{y _2 - y _1}{x _2 - x _1}$) to compute the slope and understanding the significance of the slope in determining the direction and steepness of the line.
\end{enumerate}

\subsection{Assessment Instrument}

\begin{enumerate}
    \item \textbf{Assessment Instrument: Linear Equations Quiz} \textbf{Format:}
    \item \textbf{Type:} Mixed-format quiz (multiple-choice, short answer, and problem-solving questions)
    \item \textbf{Duration:} 45 minutes
    \item \textbf{Total Marks:} 30 points \textbf{Sections:}
    \item \textbf{Understanding Linear Equations (10 points):}
    \item \textbf{Question 1:} Define a linear equation and identify its components. Provide an example.
    \item \textbf{Type:} Short answer (5 points)
    \item \textbf{Question 2:} Which of the following equations represents a linear equation? (Choose all that apply)
    \item a. $y = 2x + 3$
    \item b. $y = x^2 + 1$
    \item c. $y = \frac{1}{x}$
    \item d. $y = -4x + 7$
    \item \textbf{Type:} Multiple-choice (5 points)
    \item \textbf{Graphing Linear Equations (10 points):}
    \item \textbf{Question 3:} Given the equation $y = \frac{1}{2}x + 1$, plot the graph on a Cartesian plane and identify the slope and y-intercept.
    \item \textbf{Type:} Problem-solving (5 points)
    \item \textbf{Question 4:} If the equation $y = 3x - 4$ represents a line, what is the slope and what is the y-intercept? Explain how you would graph this line.
    \item \textbf{Type:} Short answer (5 points)
    \item \textbf{Determining Slope from Two Points (10 points):}
    \item \textbf{Question 5:} Calculate the slope of the line passing through the points (2, 5) and (4, 9). Show your work.
    \item \textbf{Type:} Problem-solving (5 points)
    \item \textbf{Question 6:} Two lines are given: Line A passes through the points (1, 2) and (3, 6), and Line B passes through the points (2, 3) and (4, 7). Are these lines parallel? Justify your answer.
    \item \textbf{Type:} Problem-solving (5 points) \textbf{Scoring Rubric:}
    \item \textbf{Understanding Linear Equations:} 5 points for the correct definition and identification of components, and 5 points for correct identification of linear equations (total 10 points).
    \item \textbf{Graphing Linear Equations:} 5 points for correctly plotting the graph and identifying slope/y-intercept, and 5 points for explaining the graphing process (total 10 points).
    \item \textbf{Determining Slope from Two Points:} 5 points for correct slope calculation with work shown, and 5 points for correct analysis of parallel lines with justification (total 10 points). This quiz will assess students' comprehension of linear equations, their ability to graph linear functions, and their skills in calculating and interpreting the slope of a line.
\end{enumerate}

%%===============================================================
\section{Solving Equations and Inequalities}
\label{sec:solving-equations-and-inequalities}
%%===============================================================

\paragraph{Introduction to Equations}



We've seen that an equation is like a balance. When two quantities are on either side of an equal sign in a mathematical statement, we are saying that the statement is only true under conditions that make the quantities the same. For instance, when we see the statement $x+2=3$ we now know that we are asking what number can we substitute for $x$ in the equation to make this statement true? One way you could work this out is by trying out different values of $x$ until you get one that works. This is called guess-and-check. Alternatively you might know the answer intuitively (by thinking \textit{What do I need to add to 2 to get 3?}).



However, if you have a more complicated problem such as $\frac{7x}{2} + 100 = 170$ you are likely to have trouble solving this problem intuitively or by guess-and-check. Because of this, mathematicians worked out a technique to solve this type of problem easily. This technique is the fundamental basis of algebra.



Since the equal sign means that both sides of the equation are the same (\textit{Same value, different appearance}), and that if you manipulate (using addition, multiplication, etc.) the values on both sides of the equal sign in the same way, then they will still be equal. This is also how a mechanical balance works. As long as you do the same thing to whatever is in both pans they stay level. If we have a value of 3 = in one pan and two things with the value of 2 and 1 in the other pan the balance is level. If we multiply both sides by 5 we get,



$$\begin{aligned}

3 \times 5 & = (2 + 1) \times 5 \\\15 & = (2 \times 5) + (1 \times 5) \\\15 & = 10 + 5 \\\15 & = 15

\end{aligned}$$





Notice how the equality still holds.



\paragraph{Neolithic Algebra}



\textbf{Neolithic Functions} Imagine you are the neolithic shepherd. What changes can you make to your bag of rocks?



\textbf{Possible Answers.}



\begin{itemize}
    \item Addition - When new sheep are added to your flock you would use addition. For instance in the spring you would have a "lamb born" function in which you add a rock to your bag. You might also have a "multiple birth" function in which you invoke the lamb born function once for each lamb born.
\end{itemize}

\begin{itemize}
    \item Subtraction - When you decide that a sheep is no longer part of your flock you would use a "remove sheep function" to keep your bag of rocks in balance with your herd. You might do different things with the rock for a missing sheep, a sheep you traded for something you needed, or because it became dinner one night. You might anticipate using these functions to manage your herd. If you know that you need to remove five rocks from your bag in the fall to buy firewood you might examine your bag carefully to know how many times you can invoke the "sheep dinner" function over the summer.
\end{itemize}

\begin{itemize}
    \item Multiplication - Imagine you are a shepherd just starting out. If your neighbors had a big enough flock they might promise you a number of sheep per month as payment for watching their flock. For instance if they promised you two sheep a month for tending their flock then you would know that after 1 month you would have two sheep, after two months you would have four sheep, and after an entire year you might have enough sheep to start your own operation.
\end{itemize}

\begin{itemize}
    \item Division - It is said that two things are inevitable: death and taxes. Division is the operation we use to make the inevitable fair. When a shepherd died his heirs could use division to divide up his flock fairly. They might use a function that allocated one sheep per heir until the flock was distributed. In what ways could the heirs distribute extra sheep if the flock was not evenly divisible? Similarly, when collecting taxes it is more fair to take a portion of the flock rather than a single amount. The clan leader would be better served taking 1/12th of a flock from each shepherd every year. If the clan leader asked for a specific number of sheep every year they might find that their clan became smaller every year just before tax time as new shepherds took their sheep to a new clan.
\end{itemize}

\paragraph{Using A Variable}



With algebra we use variables to represent things we haven't been able to count. Since we can't count the number the variable represents the number is called an "unknown". The most common variable is denoted $x$. To use a variable we write an expression which we know is true and has $x$ on its own on one side of an equals sign and a number on the other side of the equals sign.



Sometimes a variable can represent a set of numbers. In this case the number can be represented with set notation. We generally use a letter that reminds us what the variable represents.



Here are some examples of manipulating equations to get the $x$ on its own,



\textit{Example 1}: How many dollars do I need to see a movie and buy popcorn? If we let $x$ equal the dollars that I need to get into the theater, then we could use the variable set $A=\\{3,5,7,10\\}$ to represent a discount admission, a matinee admission, a regular admission, or an Imax admission. If we assume that popcorn costs 3 dollars at all the theaters we will go to then we can write the equation $x - A = 3$ to represent the money we need. We could add an A to both sides to get the equation $x = 3 + A$. Plugging the values for A into this equation we find that $x= \\{6, 8, 10, 13\\}$.



\textit{Example 2}: What is the value of $x$ in $2x = 4$?



\textit{Solution}: We can do the same in this case by dividing it by 2 (as $\frac{2x}{2}$ is $x$ (remember that $x$ is the same as 1$x$ because it has implied coefficient of 1)), but again to keep both sides of the equation equal we'll need to divide the other side by 2 as well to get,



$$\begin{aligned}

\frac{2x}{2} & = \frac{4}{2} \\\x & = \frac{4}{2} \\\x & = 2

\end{aligned}$$





\textit{Example 3}: What is the value of $x$ in $3x + 1 = 4$



\textit{Solution}: Here we first need to subtract 1 from both sides,



$$\begin{aligned}

3x + 1 - 1 & = 4 - 1 \\\3x & = 3

\end{aligned}$$





Then we divide both sides by 3 to get,



$$\begin{aligned}

\frac{3x}{3} & = \frac{3}{3} \\\x & = 1

\end{aligned}$$





Although in this case we chose to do the subtraction first and then the division, we could have done it the other way around, doing the division first followed by the subtraction, as follows,



$$\begin{aligned}

\frac{3x}{3} + \frac{1}{3} & = \frac{4}{3} \\\x + \frac{1}{3} & = \frac{4}{3} \\\x & = \frac{4}{3} - \frac{1}{3} \\\x & = \frac{3}{3} \\\x & = 1

\end{aligned}$$





When working problems it often helps to do additions/subtractions first and then multiplication/divisions, as this lets us avoid having to add or subtract fractions. However, both ways are equally valid.



\paragraph{Simplifying Equations}



Sometimes you'll come across equations which have a variable on both sides, for instance $5x-1=2x+2$, where x can be found on both sides of the equation.



We solve this type of equation in much the same way as we've solved the previous problems, but only this time you have to first make sure all of the variables are on the same side. The easiest way to see how to do this is by example:



Example 1: \textit{How do you find the value of $x$ in the equation, $5x-1=2x+2$?}



First of all you need to choose which side you want the variable to be on, the left or the right of the equals sign, in this case we'll choose to have the $x$ on the left hand side.



To do this we first have to look at where $x$ occurs on the right hand side; in this case it only appears in the term $2x$. As we don't want $x$ on the right hand side we need to get rid of it, and we can do this by subtracting $2x$ from the right side. Remember that for the equality to still be accurate we need to do the same on the left side as well.



$$\begin{aligned}

5x - 1 - 2x & = 2x + 2 - 2x \quad (\textbf{subtracting } 2x \textbf{ from both sides}) \\\\ 3x - 1 & = 2 \quad\quad\quad\quad\quad (\textbf{simplifying})

\end{aligned}$$





Now the equation is in a form which you are familiar with from the last chapter so hopefully you should now be able to solve this problem and get the answer $x=1$.



\paragraph{Manipulation of Radicals}



Let's say that we have a number $x$. The square root of $x$ is the number that, if multiplied by itself, equals $x$. Since there are two numbers which satisfy that condition, we usually specify the positive value. For example, the square root of 4 could be 2 (because $2 \times 2 = 4$) or it could be -2 (because $-2 \times -2 = 4$). We use the symbol $\sqrt{x}$ to indicate the positive square root of $x$.



The cube root of $x$ is the number that, if multiplied by itself three times, equals $x$. We use the symbol $\sqrt[3]{x}$ to indicate the cube root of $x$.



We use the symbol $\sqrt[n]{x}$ to indicate the number, which when multiplied $n$ times is equal to $x$. Or in symbols: if $y = \sqrt[n]{x}$ then $y^n = x$.



\paragraph{Rules}



$$\begin{aligned}

1.\ & \sqrt{x} \cdot \sqrt{x} = (\sqrt{x})^2 = x \hspace{100em}\\\\2.\ & \sqrt{\frac{x}{y}} = \frac{\sqrt{x}}{\sqrt{y}} \\\\3.\ & x^{\frac{m}{n}} = (\sqrt[n]{x})^m \\\\4.\ & \sqrt{x} \cdot \sqrt{y} = \sqrt{xy} \\\\5.\ & \sqrt[m]{\sqrt[n]{x}} = \sqrt[m \cdot n]{x} = x^{\frac{1}{m \cdot n}}

\end{aligned}$$





\paragraph{Solving for (Variable)}



When solving an equation, you usually solve for a specific variable. To do so, you have to get all instances of that variable on one side of the equals sign, and everything else on the other.



\paragraph{Properties of Equality}



The equal sign that depicts the fact that both sides of it are equal is a very strange symbol with many properties. It tells you various traits of each side, and it allows you to manipulate each side in specific ways. Here are the different properties of that sign:



$$\begin{array}{|c|c|c|} \hline \textbf{Property Name} & \textbf{Definition} & \textbf{Example} \\\\ \hline \text{Reflexive} & a = a & 8 = 8 \\\\ \hline \text{Symmetric} &  \begin{aligned} & \text{If } a = b, \\\\ & \text{then } b = a  \end{aligned} &  \begin{aligned} & \text{If } (3)(2) = 6, \\\\ & \text{then } 6 = (3)(2) \end{aligned} \\\\ \hline \text{Transitive} &  \begin{aligned} & \text{If } a = b \\\\  & \text{and } b = c, \\\\  & \text{then } a = c  \end{aligned} &  \begin{aligned} & \text{If } 8 = (4)(2) \\\\ & \text{and } (4)(2) = (2)(4), \\\\  & \text{then } 8 = (2)(4) \end{aligned} \\\\ \hline \text{Substitution} &  \begin{aligned} & \text{If } a = b, \\\\ & \text{then one can replace } \\\\ & a \text{ with } b \\\\  & \text{or vice versa}  \end{aligned} &  \begin{aligned} & \text{If } a = b \\\\  & \text{and } 1 + a = 3, \\\\  & \text{then } 1 + b = 3  \end{aligned} \\\\ \hline \text{Addition} &  \begin{aligned} & \text{You can add one number } \\\\ & \text{to both sides of the equation.}  \end{aligned} &  \begin{aligned}  & x - 6 = 14 \\\\  & x - 6 + 6 = 14 + 6 \\\\  & x = 20  \end{aligned} \\\\ \hline \text{Subtraction} &  \begin{aligned} & \text{You can subtract one number } \\\\ & \text{from both sides of the equation.} \end{aligned} &  \begin{aligned}  & x + 6 = 14 \\\\  & x + 6 - 6 = 14 - 6 \\\\  & x = 8  \end{aligned} \\\\ \hline \text{Multiplication} &  \begin{aligned} & \text{You can multiply both sides } \\\\ & \text{of the equation by a number.}  \end{aligned} &  \begin{aligned}  & \frac{x}{6} = 18 \\\\  & 6\left(\frac{x}{6}\right) = (18)(6) \\\\  & x = 108  \end{aligned} \\\\ \hline \text{Division} &  \begin{aligned} & \text{You can divide both sides } \\\\ & \text{of the equation by a number.}  \end{aligned} &  \begin{aligned}  & 6x = 18 \\\\  & 6x\left(\frac{1}{6}\right) = \frac{18}{6} \\\\  & x = 3  \end{aligned} \\\\ \hline \end{array}$$







\paragraph{Practice Problems}



\paragraph{Decide whether the following problems are expressions or equations:}



1. $2x + 6$:

\begin{itemize}
    \item [ ] Expression
    \item [ ] Equation
\end{itemize}

2. $3(x - 14)^2 = 12$:

\begin{itemize}
    \item [ ] Expression
    \item [ ] Equation
\end{itemize}

3. $4 + 6(2x + 16)$:

\begin{itemize}
    \item [ ] Expression
    \item [ ] Equation
\end{itemize}

\paragraph{Identify which properties are being used in the following problems:}



1. Given $a = b$ and $3 - a = 4$, so $3 - b = 4$:

\begin{itemize}
    \item [ ] Substitution Property of Equality
    \item [ ] Addition Property of Equality
    \item [ ] Multiplication Property of Equality
\end{itemize}

2. If $x + 9 = 12$, then $x = 3$:

\begin{itemize}
    \item [ ] Subtraction Property of Equality
    \item [ ] Division Property of Equality
    \item [ ] Addition Property of Equality
\end{itemize}

3. If $x - 9y = 4y$, then $x = 13y$:

\begin{itemize}
    \item [ ] Addition Property of Equality
    \item [ ] Substitution Property of Equality
    \item [ ] Multiplication Property of Equality
\end{itemize}

\paragraph{Basic Laws In Algebra}



In algebra we are working with the set of real numbers. We talked about the properties of real numbers with respect to mathematical operations in the section \href{Algebra/Real_Numbers "wikilink"}{Real Numbers}. If you don't remember the Commutative Property, the Associative Property, the Distributive Property, and the Identity Property go back to the real numbers section to review them.



There are several basic laws in algebra. Understanding these will help you to manipulate and solve equations, and to understand algebraic relationships.



\paragraph{Proportions or Ratios}



Ratios or proportions can be expressed as an equation of fractions



(for example $\frac {Q}{R} = \frac {S}{T}$ ),



or they can be expressed as a relationship Q \textbf{:} R = S \textbf{:} T , (expressed in words " 'Q' is to 'R' as 'S' is to 'T' ").



Using the words "is to" help us understand the physical relationship between the values, while the fractional representation can help us deal with the mathematical relationship.



For instance in America we know that 3 feet is to 1 yard as 6 feet is to 2 yards, but expressing this as a math equation: $\frac {3}{1} = \frac {6}{2} = \frac {2\*3}{2\*1}$ helps us see this is true because all we have done is double the original ratio.



Consider the relation 3\*4 = 2\*6 = 1\*12. Does this mean something more than the obvious fact that 12=12=12? How many ways could you package 12 items?



In general, these relationships can be described by a general equation like \textbf{$(Q)(T) = (S)(R)$}



Dividing each side of the equation by \textbf{$(T)(R)$} gives



$\frac{(Q)(T)}{(T)(R)} = \frac{(S)(R)}{(T)(R)}$ , which simplifies to ''$\frac{(Q)}{(R)} = \frac{(S)}{(T)}$



If we divide by \textbf{$(Q)(S)$} instead the results simplify to



\textbf{$\frac{(T)}{(S)} = \frac{(R)}{(Q)}$}



One needs to be careful since each term also has a proportional relation with the two other variables.



` In our examples all of the following are also valid`



` Q `\textbf{`:`}` S = R `\textbf{`:`}` T ,“ ‘Q’ is to ‘S’ as ‘R’ is to ‘T’ ”.`



` R `\textbf{`:`}` Q = T `\textbf{`:`}` S ,“ ‘R’ is to ‘Q’ as ‘T’ is to ‘S’ ”..  `



` S `\textbf{`:`}` Q = T `\textbf{`:`}` R ,“ ‘S’ is to ‘Q’ as ‘T’ is to ‘R’ ”.`



` Since 2\textit{6=3}4`



` 2:3 = 4:6, 2 is to 3 as 4 is to 6, and 2/3 = 4/6.`



` 4:2 = 6:3, 4 is to 2 as 6 is to 3, and 4/2 = 6/3.`



` 3:2 = 6:4, 3 is to 2 as 6 is to 4, and 3/2 = 6/4.`



` 2:4 = 3:6, 2 is to 4 as 3 is to 6, and 2/4 = 3/6.`



\paragraph{Solving Equations}



Although we have already solved a few equations, we will now discuss the formal idea of solving equations. To solve an equation, you are finding the value of any variables within the equation. To find the value of a variable, you have to manipulate the equation to state $*insert variable here* = *some number*$. Then you know the value of the variable! You will use the Properties of Equality to manipulate the equation into the desired form.



\paragraph{Practice Problems}



\paragraph{Solve for \( x \) in the following equations:}



1. $x + 7 = 12$



2. $-\frac{2}{3}x + 9 = 15$



3. $3x - 15 = 8x$



4. $c(x - 4m) = d(n - 9)$



5. $20x + 12 = 4(5x + 3)$



\paragraph{Introduction to Inequalities}



As opposed to an equation, an \textbf{inequality} is an expression that states that two quantities are unequal or not equivalent to one another. In most cases we use inequalities in real life more than equations (i.e. this shirt costs \$2 \textit{more than} that one).



Given that a and b are real numbers, there are four basic inequalities:



$\mathbf{a < b}$

\begin{itemize}
    \item a is "less than" b:
    \item Example: 2 < 4 ; -3 < 0; etc.
\end{itemize}

$\mathbf{a > b}$

\begin{itemize}
    \item $a$ is "greater than" $b$:
    \item Example: $-2 > -4$; $3 > 0$; etc.
\end{itemize}

$\mathbf{a \leq b}$

\begin{itemize}
    \item $a$ is "less than \textbf{or equal to}" $b$:
    \item Example: If we know that $x \leq 7$, then we can conclude that $x$ is equal to any value less than 7, including 7 itself.
\end{itemize}

$\mathbf{a \geq b}$

\begin{itemize}
    \item $a$ is "greater than \textbf{or equal to}" $b$:
    \item Example: Conversely, if $x \geq 7$, then $x$ is equal to any value greater than 7, including 7 itself.
\end{itemize}



\paragraph{Properties of Inequalities}



Just as there are four properties of equality, there are also four properties of inequality:



\textbf{Addition Property of Inequality}



If a, b, and c are real numbers such that $a > b$, then $a + c > b + c$. Conversely, if $a < b$, then $a + c < b + c$.



\textbf{Subtraction Property of Inequality}



If a, b, and c are real numbers such that $a > b$, then a - c > b - c. Conversely, if a < b, then a - c < b - c.



\textbf{Multiplication Property of Inequality}



If a, b, and c are real numbers such that a > b and \textbf{c > 0}, then ac (or a \* c) > bc (or b \* c). Conversely, if a < b and \textbf{c > 0}, then ac < bc. (Note that if c = 0, then both sides of the inequality are in fact equal.) We will also review cases where c is less than zero later in the lesson.



\textbf{Division Property of Inequality}



If a, b, and c are real numbers such that a > b, and \textbf{c > 0}, then $\frac {a}{c} > \frac {b}{c}$. Conversely under the same conditions, if a < b, then $\frac {a} {c} < \frac {b}{c}$. As with the Multiplication Property, there are special cases that will be discussed later when c < 0.



Note that all four properties also work with inequalities where $a \le b$ or $a \ge b$.



\paragraph{Trichotomy Property}



The statements above form the basis of Trichotomy Property:



\textit{Given any two real numbers a and b, then only `\textbf{`one`}` of the following statements must hold true:}



1.  $a < b$

2.  $a = b$

3.  $a > b$



So, if we are given any two unknown real-number values, then any one of the three statements will hold true.



\paragraph{Solving Inequalities}



Solving algebraic inequalities is more or less identical to solving algebraic equations. Consider the following example:



$2x + 4 \le 3x - 7$



Although it may be an inequality, we can use the Properties of Inequality stated above to solve. Start by subtracting 2x from both sides:



$2x - 2x + 4 \le 3x - 2x - 7$



$4 \le x - 7$



Finish by adding 7 to both sides.



$4 + 7 \le x - 7 + 7$



$11 \le x$



This can be rewritten as $x \ge 11$. To check, substitute any value greater than or equal to 11. However, in order to satisfy the Trichotomy Property, we'll substitute three different values: 10, 11, and 12.



$$\begin{aligned} & \begin{aligned} 2x + 4 &\leq 3x - 7 \\\\ 2(10) + 4 &\leq 3(10) - 7 \\\\ 24 &\leq 23 \end{aligned} \quad & \begin{aligned} 2x + 4 &\leq 3x - 7 \\\\ 2(11) + 4 &\leq 3(11) - 7 \\\\ 26 &\leq 26 \end{aligned} \quad\quad & \begin{aligned} 2x + 4 &\leq 3x - 7 \\\\ 2(12) + 4 &\leq 3(12) - 7 \\\\ 28 &\leq 29 \end{aligned} \end{aligned}$$



Ten is incorrect, whereas eleven and twelve satisfy the solution. Therefore, the \textbf{solution set} - the set of all answers which satisfy the original inequality - is $x \ge 11$. Written in set notation, the answer is {$x | x \ge 11$}. This is read as "the set of all x such that x is greater than or equal to 11".



\paragraph{Special Cases - A variable in the denominator}



For example, consider the inequality



$$\frac{2}{x-1}. The method for solving this kind of inequality involves four steps:



1.  Find out when the denominator is equal to 0. In this case it's when $x=1$.

2.  Pretend the inequality sign is an = sign and solve it as such: $\frac{2}{x-1}=2\,$, so $x=2$.

3.  Plot the points $x=1$ and $x=2$ on a number line with an unfilled circle because the original equation included < (it would have been a filled circle if the original equation included <= or >=). You now have three regions: $x2$.

4.  Test each region independently. in this case test if the inequality is true for 1x>2 (e.g. x=3). In this case the original inequation holds, and so the solution for the original inequation is 1>x>2.



Practice Problems



Determine if the number in parentheses satisfies the given inequality:



1. $x > 4 \quad (5)$

 - [ ] Yes

 - [ ] No



2. $2x \leq 9 \quad (6.5)$

 - [ ] Yes

 - [ ] No



3. $4x - 2 > 8 \quad (2.5)$

 - [ ] Yes

 - [ ] No



Special Cases



Let us suppose that we were given this inequality to solve:



$$-3x + 8 \geq 26$$



Using the same steps as above, start by subtracting 8 from both sides.



$-3x + 8 - 8 \ge 26 - 8$



$-3x \ge 18$



Now divide both sides by -3.



$-3x/-3 \ge 18/-3$



$x \ge -6$



Check this solution by substituting three numbers. We'll use -7, -6, and -5.



$$\begin{aligned} \& \begin{aligned} -3x + 8 \& \geq 26 \\\\ -3(-7) + 8 \& \geq 26 \\\\ 29 \& \geq 26 \end{aligned} \quad \& \begin{aligned} -3x + 8 \& \geq 26 \\\\ -3(-6) + 8 \& \geq 26 \\\\\ 26 \& \geq 26 \end{aligned} \quad \& \begin{aligned} -3x + 8 \& \geq 26 \\\\ -3(-5) + 8 \& \geq 26 \\\\ 23 \& \geq 26 \end{aligned} \end{aligned}$$



Wait, what happened? -7 and -6 satisfy the inequality, yet -7 is non-inclusive in the solution set! And -5, which is greater than -6, does not satisfy the inequality!



This is a special case involved in solving inequalities. Because the coefficient of the x-term was negative, the constant on the other side (26, which became 18 and then -6) switched signs. \textbf{In order to attain a valid solution if a negative number is divided, we need to \textit{switch the sign} in order to make sure that the solution set is correct.} Thus, $x \ge -6$ becomes $x \le -6$, or more specifically, $\\{x | x \le -6\\}$.





\paragraph{Practice Problems}



\paragraph{Solve the following inequalities and check each solution set}



1. $-3x - 5 > 22$



2. $4x - 2 \leq -6x - 7$



3. $\frac{-2x}{5} \geq 12 - 4x$



4. $11 - 7x < 39$



\paragraph{Graphing Solutions}



Because inequalities have multiple solutions, we need to be able to represent them graphically. In order to do so, we use the number line, depicted below:



                        <--|--|--|--|--|--|--|--|--|--|--|-->

                          -5 -4 -3 -2 -1  0  1  2  3  4  5





There are two ways of graphing solutions; however, each one is unique depending on the nature of the inequality. If, for example, the inequality contains the < or > sign, we use an open circle ("O") and place it on (or above) the corresponding position on the number line, then draw another line either left or right of the solution (based on the sign) to indicate the infinite number of solutions in the set. For example, if we were to graph x < 4:



                        <-----------------------------O

                        <--|--|--|--|--|--|--|--|--|--|--|-->

                          -5 -4 -3 -2 -1  0  1  2  3  4  5



The open circle indicates that 4 is \textbf{not included in the solution set}; however, all values less than 4 satisfy the solution.



If an inequality contains a $\le$ or $\ge$, then a closed circle (it will be depicted here with a \*) is placed on or above the corresponding position on the number line. Hence, the solution $x \ge -2$ is graphed as follows.



                                    *----------------------->

                        <--|--|--|--|--|--|--|--|--|--|--|-->

                          -5 -4 -3 -2 -1  0  1  2  3  4  5



The asterisk (\*) indicates that -2 is \textbf{inclusive in the solution set}, as are all values greater than -2.



\textbf{Practice Problems}



Graph the following inequalities on a number line:



1\. $x > 4$



2\. $x \le 1$



3\. $-1 > x$



\textbf{Solutions}



1\.



                                                      0----->

                        <--|--|--|--|--|--|--|--|--|--|--|-->

                          -5 -4 -3 -2 -1  0  1  2  3  4  5



2\.



                        <--------------------*

                        <--|--|--|--|--|--|--|--|--|--|--|-->

                          -5 -4 -3 -2 -1  0  1  2  3  4  5



(Note that your actual graph should have a closed circle in place of the dot. If in doubt, simply draw a circle and color it black.)



3\. Be careful here; you need to rearrange the inequality first before graphing. When rewritten, the inequality becomes x < -1:



                        <--------------O

                        <--|--|--|--|--|--|--|--|--|--|--|-->

                          -5 -4 -3 -2 -1  0  1  2  3  4  5



\paragraph{Lesson Review}



An inequality is a statement justifying that two quantities are not equal to each other. There are four cases of inequalities, two of which allow for equality ($a \le b$ and $a \ge b$). The four properties of inequality, which are more or less parallel to the properties of equality, can be used to solve simple inequalities. The only exceptions are in multiplication and division, where the signs must be reversed if both sides are multiplied or divided by a negative number. Finally, the solution set of an inequality can be graphed on the number line with either an open circle ("O") or a closed circle ("\*"), depending on the original sign used.



\paragraph{Resources and Examples}



See also: \href{Linear_Equations "wikilink"}{Linear Equations}



\paragraph{General Resources}



\begin{itemize}
    \item \href{https://tutorial.math.lamar.edu/classes/alg/Solving.aspx}{Solving Equations and Inequalities}, Paul's Online Notes
    \item \href{https://www.khanacademy.org/math/algebra/x2f8bb11595b61c86:solve-equations-inequalities}{Solving Equations and Inequalities}, Khan Academy
    \item \href{https://www.mathsisfun.com/algebra/inequality-solving.html}{Rules for Solving Inequalities}, Math Is Fun
\end{itemize}

\paragraph{Solving Different Types of Equations}



\begin{itemize}
    \item \href{https://courses.lumenlearning.com/ivytech-collegealgebra/}{College Algebra} from Lumen Learning:
    \item \href{https://courses.lumenlearning.com/ivytech-collegealgebra/chapter/solving-linear-equations-in-one-variable/}{Solving Linear Equations with One Variable}
    \item \href{https://courses.lumenlearning.com/ivytech-collegealgebra/chapter/solving-a-rational-equation/}{Solving Rational Equations}
    \item \href{https://courses.lumenlearning.com/ivytech-collegealgebra/chapter/solving-quadratic-equations-by-factoring/}{Solving Quadratic Equations by Factoring}
    \item \href{https://courses.lumenlearning.com/ivytech-collegealgebra/chapter/completing-the-square/}{Completing the Square to Solve Quadratic Equations}
    \item \href{https://courses.lumenlearning.com/ivytech-collegealgebra/chapter/using-the-quadratic-formula/}{Using the Quadratic Formula}
    \item \href{https://courses.lumenlearning.com/ivytech-collegealgebra/chapter/introduction-other-types-of-equations/}{Solving Other Types of Equations} (go through this section by clicking "next")
    \item \href{https://www.cliffsnotes.com/study-guides/algebra/algebra-i/inequalities-graphing-and-absolute-value/solving-equations-containing-absolute-value}{Solving Absolute Value Equations}, Cliff's Notes
\end{itemize}

\paragraph{Solving Different Types of Inequalities}



\begin{itemize}
    \item \href{https://www.youtube.com/watch?v=eNxruTEpY4o}{Solving Linear Inequalities}, patrickJMT on YouTube
    \item \href{https://www.youtube.com/watch?v=t54ccHYVhoo}{Solving Quadratic Inequalities}, patrickJMT on YouTube
    \item \href{https://www.youtube.com/watch?v=WZqZPLaK1cg}{Solving Quadratic Inequalities: More Examples}, patrickJMT on YouTube
    \item \href{https://www.youtube.com/watch?v=OzmqAZZd88g}{Solving Absolute Value Inequalities}, patrickJMT on YouTube
    \item \href{https://prep.math.lsa.umich.edu/cgi-bin/pmc/crtopic?sxn=8\&top=3\&crssxn=prep}{Solving Inequalities with Radicals}, University of Michigan Math Prep
\end{itemize}

NOTE: The same methods for solving equations can be used for solving inequalities. However, multiplying or dividing by a negative number on both sides of an inequality reverses the inequality sign. Examples of this:



\begin{itemize}
    \item $-x < 7 \implies (-1)(-x) > (-1)(7) \implies x > -7$
    \item $6 - y \ge x \implies y - 6 \le -x \implies y \le  6 - x$
    \item $-2x > 24 \implies (-2x)/(-2) < 24/(-2) \implies x < -12$
\end{itemize}

\paragraph{Licensing}



Content obtained and/or adapted from:



\begin{itemize}
    \item \href{https://en.wikibooks.org/wiki/Algebra/Solving_Equations}{Solving Equations, Wikibooks: Algebra} under a CC BY-SA license
    \item \href{https://en.wikibooks.org/wiki/Intermediate_Algebra/Solving_Inequalities}{Solving Inequalities, Wikibooks: Intermediate Algebra} under a CC BY-SA license
\end{itemize}

\subsection{Learning Objectives}

\begin{enumerate}
    \item Here are three Student Learning Objectives (SLOs) for the "Introduction to Equations" section:
    \item \textbf{Solve Basic Linear Equations:}
    \item \textbf{Outcome:} Students will be able to solve basic linear equations by applying addition, subtraction, multiplication, and division to both sides of the equation.
    \item \textbf{Details:} This includes understanding the concept of equality and correctly manipulating both sides of the equation to isolate the variable and find its value.
    \item \textbf{Simplify and Solve Equations with Variables on Both Sides:}
    \item \textbf{Outcome:} Students will be able to simplify and solve equations where the variable appears on both sides of the equal sign.
    \item \textbf{Details:} This includes the ability to identify and move terms involving the variable to one side of the equation and constant terms to the other side before solving.
    \item \textbf{Apply Properties of Equality:}
    \item \textbf{Outcome:} Students will be able to apply the properties of equality, including reflexive, symmetric, transitive, and substitution properties, to justify each step in solving an equation.
    \item \textbf{Details:} This involves using these properties to understand and explain the logical steps taken to maintain the balance of the equation during the solving process.
\end{enumerate}

\subsection{Assessment Instrument}

\begin{enumerate}
    \item \textbf{Assessment Instrument: Solving Equations and Inequalities Quiz} \textbf{Format:}
    \item \textbf{Type:} Mixed-format quiz (multiple-choice, short answer, and problem-solving questions)
    \item \textbf{Duration:} 45 minutes
    \item \textbf{Total Marks:} 30 points \textbf{Sections:}
    \item \textbf{Solving Linear Equations (10 points):}
    \item \textbf{Question 1:} Solve the following equations for $x$:
    \item a. $2x + 5 = 13$
    \item b. $\frac{7x}{2} - 4 = 10$
    \item c. $3x - 15 = 8x$
    \item d. $5x + 2 = 3x + 10$
    \item e. $\frac{2x - 4}{3} = 2$
    \item \textbf{Simplifying and Solving Equations with Variables on Both Sides (10 points):}
    \item \textbf{Question 2:} Simplify and solve the following equations:
    \item a. $5x - 1 = 2x + 2$
    \item b. $4(x - 3) = 2(2x + 5)$
    \item c. $7x + 2 = 3x - 6 + 4x$
    \item d. $6x + 4 - 3x = 2x + 10$
    \item \textbf{Solving Inequalities and Graphing Solutions (10 points):}
    \item \textbf{Question 3:} Solve the following inequalities and represent the solution on a number line:
    \item a. $2x + 3 > 7$
    \item b. $5 - 3x \leq 11$
    \item c. $\frac{4x - 2}{3} \geq 6$
    \item d. $-2x + 8 < 4$ \textbf{Scoring Rubric:}
    \item \textbf{Solving Linear Equations:} 2 points for each correctly solved equation (total 10 points)
    \item \textbf{Simplifying and Solving Equations with Variables on Both Sides:} 2.5 points for each correctly simplified and solved equation (total 10 points)
    \item \textbf{Solving Inequalities and Graphing Solutions:} 2.5 points for each correctly solved inequality and accurately graphed solution (total 10 points) This quiz assesses the students' ability to solve linear equations, simplify and solve equations with variables on both sides, and solve inequalities. It also evaluates their understanding of representing solutions graphically on a number line.
\end{enumerate}

%%===============================================================
\section{Graphs}
\label{sec:graphs}
%%===============================================================

% [Image: Graph of y=2x]



Figure: Graph of $y=2x$



It is sometimes difficult to understand the behavior of a function given only its definition; a visual representation or graph can be very helpful. A \textbf{graph} is a set of points in the Cartesian plane, where each point $(x,y)$ indicates that $f(x)=y$ . In other words, a graph uses the position of a point in one direction (the \textit{vertical-axis} or $y$\textit{-axis}) to indicate the value of $f$ for a position of the point in the other direction (the \textit{horizontal-axis} or $x$\textit{-axis}).



Functions may be graphed by finding the value of $f$ for various $x$ and plotting the points $(x,f(x))$ in a Cartesian plane. For the functions that you will deal with, the parts of the function between the points can generally be approximated by drawing a line or curve between the points. Extending the function beyond the set of points is also possible, but becomes increasingly inaccurate.



\paragraph{Linear functions}



Graphing linear functions are easy to understand and do. Because we know that two points can form a line, only two points are needed for us to graph a linear function if those two points are on the function. Oppositely, we can write down the equation of a linear function if we only know two points that are on the function.



The following section mainly talks about different forms of linear function notations so that you can easily identify or graph the function.



\paragraph{Introduction}



Plotting points like this is laborious. Fortunately, many functions' graphs fall into general patterns. For a simple case, consider functions of the form



$$f(x)=3x+2$$



The graph of $f$ is a single line, passing through the point $(0,2)$ with slope 3. Thus, after plotting the point, a straightedge may be used to draw the graph. This type of function is called \textit{linear} and there are a few different ways to present a function of this type.



\paragraph{Slope}



The slope is the backbone of linear functions because it shows how much the output of a function changes when the input changes. For example, if the slope of a function is 2, then it means when the input of a function increases by 1 unit, the output of the function increases by 2 units. Now, let's look at a more mathematical example.



> Consider this function: $f(x)=-5x+8$. What does the number $-5$ mean?

>

> It means that when $x$ increases by 1, $f(x)$ decreases by 5.

>

> Using mathematical terms:

>

> $f(x+1)-f(x)=(-5(x+1)+8)-(-5x+8)=-5$



It is easy to calculate the slope because the slope is like the speed of a vehicle. If we divide the change in distance and the corresponding change in time, we get the speed. Similarly, if we divide the change in $f(x)$ over the corresponding change in $x$, we get the slope.



If given two points, $(x _1,y _1)$ and $(x _2,y _2)$ , we may then compute the slope of the line that passes through these two points. Remember, the slope is determined as "rise over run." That is, the slope is the change in $y$-values divided by the change in $x$-values. In symbols:



$$\text{slope}~=\frac{\text{change in }y}{\text{change in }x}=\frac{\Delta y}{\Delta x}=\frac{y _2-y -1}{x _2-x _1}$$



Slope in a linear function: If two points $(x _1,y _1)$ and $(x _2,y _2)$ on a linear function, then the slope of the linear function is



$$m=\frac{y _2-y _1}{x _2-x _1}$$



Interestingly, there is a subtle relationship between the slope and the angle between the graph of the function and the positive $x$-axis, $\theta$. The relationship is:



> $m=\tan\theta$



It is an obvious relationship, but it can be ignored relatively easily.



\paragraph{Slope-intercept form}



% [Image: This is a linear function notated in the slope-intercept form. The slope here is a instead of m.]



Figure: This is a linear function notated in the slope-intercept form. The slope here is $a$ instead of $m$.



When we see a function presented as



> $y=f(x)=mx+b$



we call this presentation the \textbf{slope-intercept form}. This is because, not surprisingly, this way of writing a linear function involves the slope, $m$ , and the $y$-intercept, $b$ .



\textbf{Example 1:} Graph the function $f(x)=3x+2$.



The slope of the function is 3, and it intercepts the $y$-axis at point $(0,2)$. In order to graph the function, we need another point. Since the slope of the function is 3, then



> $f(x+1)=f(x)+3$



> $f(1)=f(0+1)=f(0)+3=5$



Knowing that the function goes through points $(0,2)\text{ and }(1,5)$, the function can be easily graphed.



\textbf{Example 2:} Now, consider another unknown linear function that goes through points $(-3,-4)\text{ and }(0,5)$. What is the equation for this function?



The slope can be calculate with the formula mentioned above.



$$m=\frac{y _2-y _1}{x _2-x _1}=\frac{5-(-4)}{0-(-3)}=3$$



And since the $y$-axis interception is $(0,5)$, we can know that



> $b=5$



Thus, the equation of this linear function should be



> $g(x)=3x+5$



\paragraph{Point-slope form}



If someone walks up to you and gives you one point and a slope, you can draw one line and only one line that goes through that point and has that slope. Said differently, a point and a slope uniquely determine a line. So, if given a point $(x _0,y _0)$ and a slope $m$ , we present the graph as



> $y-y _0=m(x-x _0)$



We call this presentation the \textbf{point-slope form}. The point-slope and slope-intercept form are essentially the same. In the point-slope form we can use any point the graph passes through. Where as, in the slope-intercept form, we use the $y$-intercept, that is the point $(0,b)$. The point-slope form is very important. Although it is not used as frequently as its counterpart the slope-intercept form, the concept of knowing a point and drawing the line in the direction of the slope will be encountered when we go into vector equations for lines and planes in future chapters.



\textbf{Example 1:} If a linear function goes through points $(3,-4)\text{ and }(1,5)$, what is the equation for this function?



The slope is:



$$m=\frac{y _2-y _1}{x _2-x _1}=\frac{5-(-4)}{1-3}=-\frac{9}{2}$$



Since we know two points, the following answers are all correct



$$y+4=-\frac{9}{2}(x-3)\text{ or }y-5=-\frac{9}{2}(x-1)$$



The two-point form is another form to write the equation for a linear function. It is similar to the point-slope form. Given points $(x _1,y _1)$ and $(x _2,y _2)$, we have the equation



$$y-y _1=\frac{y _2-y _1}{x _2-x _1}(x-x _1)$$



This presentation is in the \textbf{two-point form}. It is essentially the same as the point-slope form except we substitute the expression $\frac{y _2-y _1}{x _2-x _1}$ for $m$. However, this expression is not widely used in mathematics because in most situations, $(x _1,y _1)$ and $(x _2,y _2)$ are known coordinates. It would be redundant to write down a bulky $\frac{y _2-y _1}{x _2-x _1}$ instead of a simple expression of the slope.



\paragraph{Intercept form}



The intercept form looks like this:



$$\frac{x}{a}+\frac{y}{b}=1$$



By writing the function in the intercept form, we can quickly determine the $x,y$-axis intercepts.



> $x$-axis intercept: $(a,0)$

>

> $y$-axis intercept: $(0,b)$



When we discuss planes in 3D space, this form will be quite useful to determine the $x,y,z$-axis intercepts.



\paragraph{Quadratic functions}



To graph a quadratic function, there is the simple but work-heavy way, and there is the complicated but clever way. The simple way is to substitute the independent variable $x$ with various numbers and calculate the output $f(x)$. After some substitutions, plot those $(x,f(x))$ and connect those points with a curve. The complicated way is to find special points, such as intercepts and the vertex, and plot it out. The following section is a guide to find those special points, which will be useful in later chapters.



\paragraph{Standard form}



Quadratic functions are functions that look like this



> $y=f(x)=ax^2+bx+c$, where $a,b,c$ are constants



The constant $a$ determines the concavity of the function: if $a>0$, $f(x)$ concaves up; if $a0$, the vertex is the lowest point on the graph; if $a $y=f(x)=a(x-h)^2+k$, where $a,h,k$ are constants



The vertex of this function is $(h,k)$ because when $x=h$, $f(h)=k$. If $a>0$, $f(h)=k$ is the absolute minimum value that the function can achieve. If $a $y=ax^2+bx+c$ is the equation for the quadratic function. In this case, $a>0$, $c0$.



> Points $(\frac{-b-\sqrt{b^2-4ac}}{2a},0)\text{ and }(\frac{-b+\sqrt{b^2-4ac}}{2a},0)$



These are the coordinates for the two $x$-axis intercepts.



Knowing the coordinates, the function can be written in its factored form: $y=a(x-\frac{-b-\sqrt{b^2-4ac}}{2a})(x-\frac{-b+\sqrt{b^2-4ac}}{2a})$



> Point $(-\frac{b}{2a},-\frac{b^2-4ac}{4a})$ is the vertex for the quadratic function. Because $a>0$, the vertex is the lowest point on the graph.



Since the vertex is known, we can write the function in the vertex form: $y=a(x+\frac{b}{2a})^2-\frac{b^2-4ac}{4a}$



Although this does not look like the equation we've just discussed earlier, note that $-\frac{b^2-4ac}{4a}=+c-\frac{b^2}{4a}$.



Line $x=-\frac{b}{2a}$. The graph of the function is symmetric about this line. In other words,



$$f(-\frac{b}{2a}+x)=f(-\frac{b}{2a}-x)$$



If you can skillfully and quickly determine those special points, graphing quadratic functions will be less torturing.



\paragraph{Exponential and Logarithmic functions}



Exponential and logarithmic functions are inverse functions with each other. Take the exponential function $f(x)=a^x$ for example. The inverse function of $f(x)$, $f^{-1}(x)$, is



$$x=a^{f^{-1}(x)}$$



$$\Leftrightarrow f^{-1}(x)=\log _ax$$



which is a logarithmic function.



Since geometrically, the graph of the inverse function is flipping the graph of the original function over line $y=x$, we only need to know how to graph one of those functions.



\paragraph{Resources}



\paragraph{Basic Information}



\begin{itemize}
    \item \href{https://www.mathsisfun.com/data/cartesian-coordinates.html}{Introduction to Cartesian Coordinates and Graphs}, mathisfun.com
    \item \href{https://www.khanacademy.org/math/algebra/x2f8bb11595b61c86:foundation-algebra/x2f8bb11595b61c86:algebra-overview-history/v/descartes-and-cartesian-coordinates}{Introduction to the Coordinate Plane}, Khan Academy
    \item \href{https://courses.lumenlearning.com/ivytech-collegealgebra/chapter/plotting-ordered-pairs-in-the-cartesian-coordinate-system/}{Plotting Points on a Graph}, Lumen Learning
\end{itemize}

\paragraph{Graphing Linear Functions}



\begin{itemize}
    \item \href{https://courses.lumenlearning.com/waymakercollegealgebra/chapter/graph-linear-functions/}{Graphing Linear Functions}, Lumen Learning
    \item \href{https://www.khanacademy.org/math/algebra/x2f8bb11595b61c86:linear-equations-graphs}{Linear Equations and Graphs}, Khan Academy
    \item \href{https://www.palmbeachstate.edu/prepmathlw/Documents/graphing\%20linear\%20equations.pdf}{Graphing Linear Equations}, Palm Beach State College
    \item \href{https://www.youtube.com/watch?v=lkRyxnqOb8M}{Graphing a Line Using a Point and Slope}, patrickJMT
    \item \href{https://www.youtube.com/watch?v=mxBoni8N70Y}{Graphing a Line with X and Y Intercepts}, patrickJMT
\end{itemize}

\paragraph{Licensing}



Content obtained and/or adapted from:



\begin{itemize}
    \item \href{https://en.wikibooks.org/wiki/Calculus/Graphing_functions}{Graphing functions, Wikibooks: Calculus} under a CC BY-SA license
\end{itemize}

\subsection{Learning Objectives}

\begin{enumerate}
    \item Here are three Student Learning Objectives (SLOs) based on the graph of $y = 2x$:
    \item \textbf{Interpret the Slope of a Linear Function:}
    \item \textbf{Outcome:} Students will be able to determine and interpret the slope of a linear function from its graph.
    \item \textbf{Details:} This includes understanding that the slope represents the rate of change of the function and being able to calculate the slope from any two points on the line, such as recognizing that for $y = 2x$, the slope $m = 2$ indicates that for every unit increase in $x$, $y$ increases by 2 units.
    \item \textbf{Graph Linear Equations Using Points:}
    \item \textbf{Outcome:} Students will be able to graph a linear equation by identifying and plotting points on the Cartesian plane.
    \item \textbf{Details:} This includes selecting points on the graph of $y = 2x$ by substituting values for $x$ and finding the corresponding $y$ values, then plotting these points to form a straight line that represents the equation.
    \item \textbf{Analyze the Relationship Between a Function and its Graph:}
    \item \textbf{Outcome:} Students will be able to analyze and describe how the graphical representation of a linear function corresponds to its algebraic equation.
    \item \textbf{Details:} This involves understanding that the graph of $y = 2x$ is a straight line passing through the origin with a slope of 2, and explaining how changes in the slope or intercept affect the position and orientation of the line on the Cartesian plane.
\end{enumerate}

\subsection{Assessment Instrument}

\begin{enumerate}
    \item \textbf{Assessment Instrument: Graphs of Linear Functions Quiz} \textbf{Format:}
    \item \textbf{Type:} Mixed-format quiz (multiple-choice, short answer, and problem-solving questions)
    \item \textbf{Duration:} 45 minutes
    \item \textbf{Total Marks:} 30 points \textbf{Sections:}
    \item \textbf{Interpreting the Slope of a Linear Function (10 points):}
    \item \textbf{Question 1:} The graph of the function $ y = 2x $ is given. Answer the following:
    \item a. What is the slope of the function? (2 points)
    \item b. Explain in a sentence what the slope of this function means in terms of changes in $ x $ and $ y $. (3 points)
    \item c. If the slope were changed to 3, how would the graph of the function change? (5 points)
    \item \textbf{Graphing Linear Equations Using Points (10 points):}
    \item \textbf{Question 2:} Given the linear equation $ y = 2x - 3 $, do the following:
    \item a. Calculate the $ y $-value when $ x = 0 $ and when $ x = 2 $. (3 points)
    \item b. Plot the points $ (0, y) $ and $ (2, y) $ on the Cartesian plane and draw the line. (4 points)
    \item c. Describe the process you used to graph the equation. (3 points)
    \item \textbf{Analyzing the Relationship Between a Function and its Graph (10 points):}
    \item \textbf{Question 3:} Consider the function $ y = -\frac{1}{2}x + 4 $. Answer the following:
    \item a. Identify the slope and $ y $-intercept of the function. (4 points)
    \item b. Explain how the negative slope affects the direction of the line on the graph. (3 points)
    \item c. If the $ y $-intercept were 2 instead of 4, how would the graph change? (3 points) \textbf{Scoring Rubric:}
    \item \textbf{Interpreting the Slope of a Linear Function:} Points are awarded based on the correct identification of the slope, clear explanation of its meaning, and accurate description of changes in the graph (total 10 points).
    \item \textbf{Graphing Linear Equations Using Points:} Points are awarded for correctly calculating the $ y $-values, accurately plotting the points, drawing the line, and clearly describing the graphing process (total 10 points).
    \item \textbf{Analyzing the Relationship Between a Function and its Graph:} Points are awarded based on correctly identifying the slope and $ y $-intercept, explaining the impact of a negative slope, and accurately describing the effect of changes in the $ y $-intercept (total 10 points). This quiz assesses students' understanding of graphing linear functions, interpreting slopes, and analyzing the relationship between a function and its graph. It ensures that they can apply these concepts accurately and explain the underlying principles.
\end{enumerate}

%%===============================================================
\section{Functions}
\label{sec:functions}
%%===============================================================

\paragraph{Introduction}



Formally, a function f from a set X to a set Y is defined by a set G of ordered pairs (x, y) such that $x \in X$, $y \in Y$, and every element of X is the first component of exactly one ordered pair in G. In other words, for every x in X, there is exactly one element y such that the ordered pair (x, y) belongs to the set of pairs defining the function f. The set G is called the graph of the function. Formally speaking, it may be identified with the function, but this hides the usual interpretation of a function as a process. Therefore, in common usage, the function is generally distinguished from its graph.}} Whenever one quantity uniquely determines the value of another quantity, we have a function. That is, the set $X$ uniquely determines the set $Y$. You can think of a \textit{function} as a kind of machine. You feed the machine raw materials, and the machine changes the raw materials into a finished product.



\paragraph{Functions in everyday life}



Think about dropping a ball from a bridge. At each moment in time, the ball is a height above the ground. The height of the ball is a function of time. It was the job of physicists to come up with a formula for this function. This type of function is called \textit{real-valued} since the "finished product" is a number (or, more specifically, a real number).



Think about a wind storm. At different places, the wind can be blowing in different directions with different intensities. The direction and intensity of the wind can be thought of as a function of position. This is a function of two real variables (a location is described by two values - an $x$ and a $y$) which results in a vector (which is something that can be used to hold a direction and an intensity). These functions are studied in multivariable calculus (which is usually studied after a one year college level calculus course). This a vector-valued function of two real variables.



We will be looking at real-valued functions until studying multivariable calculus. Think of a real-valued function as an \textit{input-output machine}; you give the function an input, and it gives you an output which is a number (more specifically, a real number). For example, the squaring function takes the input 4 and gives the output value 16. The same squaring function takes the input -1 and gives the output value 1.



% [Image: This is an intuitive way to understand functions: a machine that makes the input x go through a transformation f into the output f(x)]



Figure: This is an intuitive way to understand functions: a machine that makes the input $x$ go through a transformation $f$ into the output $f(x)$



\paragraph{Notation}



Functions are used so much that there is a special notation for them. The notation is somewhat ambiguous, so familiarity with it is important in order to understand the intention of an equation or formula.



Though there are no strict rules for naming a function, it is standard practice to use the letters $f$ , $g$ , and $h$ to denote functions, and the variable $x$ to denote an independent variable. $y$ is used for both dependent and independent variables.



When discussing or working with a function $f$ , it's important to know not only the function, but also its independent variable $x$ . Thus, when referring to a function $f$, you usually do not write $f$, but instead $f(x)$ . The function is now referred to as "$f$ of $x$". The name of the function is adjacent to the independent variable (in parentheses). This is useful for indicating the value of the function at a particular value of the independent variable. For instance, if



$$f(x)=7x+1$$ , and if we want to use the value of $f$ for $x$ equal to $2$ , then we would substitute 2 for $x$ on both sides of the definition above and write



$$f(2)=7(2)+1=14+1=15$$



This notation is more informative than leaving off the independent variable and writing simply ''$f$' , but can be ambiguous since the parentheses next to $f$ can be misinterpreted as multiplication, $2f$. To make sure nobody is too confused, follow this procedure:



1.  Define the function $f$ by equating it to some expression.

2.  Give a sentence like the following: "At $x=c$, the function $f$ is\..."

3.  Evaluate the function.



\paragraph{Description}



There are many ways which people describe functions. In the examples above, a verbal descriptions is given (the height of the ball above the earth as a function of time). Here is a list of ways to describe functions. The top three listed approaches to describing functions are the most popular.



1.  A function is given a name (such as $f$) and a formula for the function is also given. For example, $f(x)=3x+2$ describes a function. We refer to the input as the \textbf{argument} of the function (or the \textbf{independent variable}), and to the output as the \textbf{value} of the function at the given argument.

2.  A function is described using an equation and two variables. One variable is for the input of the function and one is for the output of the function. The variable for the input is called the \textbf{independent variable}. The variable for the output is called the \textbf{dependent variable}. For example, $y=3x+2$ describes a function. The dependent variable appears by itself on the left hand side of equal sign.

3.  A verbal description of the function.



When a function is given a name (like in number 1 above), the name of the function is usually a single letter of the alphabet (such as $f$ or $g$). Some functions whose names are multiple letters (like the sine function $y=\sin(x)$



\paragraph{Plugging a value into a function}



If we write $f(x)=3x+2$ , then we know that



\begin{itemize}
    \item The function $f$ is a function of $x$ .
    \item To evaluate the function at a certain number, replace the $x$ with that number.
    \item Replacing $x$ with that number in the right side of the function will produce the function's output for that certain input.
    \item In English, the definition of $f$ is interpreted, "Given a number, $f$ will return \textit{two more than the triple of that number}."
\end{itemize}

How would we know the value of the function $f$ at 3? We would have the following three thoughts:



1.  $f(3)=3(3)+2$

2.  $3(3)+2=9+2$

3.  $9+2=11$



and we would write



$f(3)=3(3)+2=9+2=11$.



The value of $f$ at 3 is 11.



Note that $f(3)$ means the value of the dependent variable when $x$ takes on the value of 3. So we see that the number \textit{11} is the output of the function when we give the number \textit{3} as the input. People often summarize the work above by writing "the value of $f$ at three is eleven", or simply "$f$ of three equals eleven".



\paragraph{Basic concepts of functions}



The formal definition of a function states that a function is actually a \textit{mapping} that associates the elements of one set called the \textbf{domain} of the function, $A$, with the elements of another set called the \textit{range} of the function, $B$. For each value we select from the domain of the function, there exists \textit{exactly one} corresponding element in the range of the function. The definition of the function tells us which element in the range corresponds to the element we picked from the domain. An example of how is given below.



\paragraph{Example}



Let function $f(x)=5x+\frac{5x}{2}$ for all $x$. For what value of $x$ gives $f(x)=225\frac{1}{2}$ In mathematics, it is important to notice any repetition. If something seems to repeat, keep a note of that in the back of your mind somewhere.



Here, notice that $f(x)=5x+\frac{5x}{2}$ and $f(x)=225\frac{1}{2}$. Because $f(x)$ is equal to two different things, it must be the case that the other side of the equal sign to $f(x)$ is equal to the other. This property is known as the transitive property and could thus make the following equation below true:$$225\frac{1}{2}=5x+\frac{5x}{2}$$



Next, simplify --- make your life easier rather than harder. In this instance, since $5x+\frac{5x}{2}$ has $x$ as a like-term, and the two terms are fractions added to the other, put it over a common denominator and simplify. Similar, since $225\frac{1}{2}$ is a mixed fraction, $225\frac{1}{2}=225+\frac{1}{2}$.



$$\begin{aligned} 225\frac{1}{2} &= 5x + \frac{5x}{2} \\\\ \\\\ \Rightarrow \frac{225}{1} + \frac{1}{2} &= \frac{5x}{1} + \frac{5x}{2} \\\\ \\\\ \Rightarrow \frac{450}{2} + \frac{1}{2} &= \frac{10x}{2} + \frac{5x}{2} \\\\ \\\\ \Rightarrow \frac{451}{2} &= \frac{15x}{2} \end{aligned}$$





Multiply both sides by the reciprocal of the coefficient of $x$, $\frac{2}{15}$ from both sides by multiplying by it.

$$\begin{aligned} \frac{2}{15} \cdot \frac{451}{2} = x \\\\ \\\\ \Rightarrow \frac{451}{15} = x \end{aligned}$$





The value of $x$ that makes $f(x)=225\frac{1}{2}$ is $x=\frac{451}{15}$.



Classically, the element picked from the domain is pictured as something that is fed into the function and the corresponding element in the range is pictured as the output. Since we "pick" the element in the domain whose corresponding element in the range we want to find, we have control over what element we pick and hence this element is also known as the "independent variable". The element mapped in the range is beyond our control and is "mapped to" by the function. This element is hence also known as the "dependent variable", for it depends on which independent variable we pick. Since the elementary idea of functions is better understood from the classical viewpoint, we shall use it hereafter. However, it is still important to remember the correct definition of functions at all times.



\paragraph{Basic types of transformation}



To make it simple, for the function $f(x)$, all of the possible $x$ values constitute the domain, and all of the values $f(x)$ ($y$ on the x-y plane) constitute the range. To put it in more formal terms, a function $f$ is a mapping of some element $a\in A$, called the domain, to exactly one element $b\in B$, called the range, such that $f:A\to B$. The image below should help explain the modern definition of a function:



% [Image: A is the domain of the function while B is the range. This transformation from set A to B is an example of one-to-one function.]



Figure: $A$ is the domain of the function while $B$ is the range. This transformation from set $A$ to $B$ is an example of one-to-one function.



1.  A function is considered \textbf{one-to-one} if an element $a\in A$ from domain $A$ of function $f$, leads to exactly one element $b\in B$ from range $B$ of the function. By definition, since only one element $b$ is mapped by function $f$ from some element $a$, $f:A\to B$ implies that there exists only one element $b$ from the mapping. Therefore, there exists a one-to-one function because it complies with the definition of a function. This definition is similar to \textbf{\textit{Figure 1}}.

2.  A function is considered \textbf{many-to-one} if some elements $a _{1},a _{2},\ldots a _{n}\in A$ from domain $A$ of function $f$ maps to exactly one element $b\in B$ from range $B$ of the function. Since some elements $\left( a _{1},a _{2},\ldots a _{n} \right)\in A$ must map onto exactly one element $b\in B$, $f:A\to B$ must be compliant with the definition of a function. Therefore, there exists a many-to-one function.

3.  A function is considered \textbf{one-to-many} if exactly one element $a\in A$ from domain $A$ of function $f$ maps to some elements $\left( b _{1},b _{2},\ldots b _{n} \right)\in B$ from range $B$ of the function. If some element $a$ maps onto many distinct elements $b$, then $f:A\to B$ is non-functional since there exists many distinct elements $b$. Given many-to-one is non-compliant to the definition of a function, there exists no function that is one-to-many.



The modern definition describes a function sufficiently such that it helps us determine whether some new type of "function" is indeed one. Now that each case is defined above, you can now prove whether functions are one of these function cases. Here is an example problem:



\paragraph{One-to-One Example}



\textbf{Given}: $f(x)=ax+b$, where $a$ and $b$ are constant and $a\ne0$.



\textbf{Prove}: function $f$ is one-to-one.



Notice that the only changing element in the function $f$ is $x$. To prove a function is one-to-one by applying the definition may be impossible because although two random elements of domain set $A$ may not be many-to-one, there may be some elements $\left( a _{1},a _{2},\ldots a _{n} \right)\in A$ that may make the function many-to-one. To account for this possibility, we must prove that it is impossible to have some result like that.



Assume there exists two distinct elements $x _{1}\ne x _2$ that will result in $f\left( x _{1} \right) =f\left( x _{2} \right)$. This would make the function many-to-one. In consequence,



$$\begin{aligned} ax _{1} + b & = ax _{2} + b \quad & \text{Subtract } b \text{ from both sides:} \\\\ \Rightarrow ax _{1} & = ax _{2} \quad & \text{Subtract } ax _{2} \text{ from both sides:} \\\\ \Rightarrow ax _{1} - ax _{2} & = 0 \quad & \text{Factor } a \text{ from the left side:} \\\\ \Rightarrow a\left( x _{1} - x _{2} \right) & = 0 \quad & \text{Multiply } \frac{1}{a} \text{ to both sides:} \\\\ \Rightarrow \left( x _{1} - x _{2} \right) & = 0 \quad & \text{Add } x _{2} \text{ to both sides:} \\\\ \Rightarrow x _{1} & = x _{2} \end{aligned}$$



Notice that $x _{1}=x _{2}$. However, this is impossible because $x _{1}$ and $x _{2}$ are distinct. Ergo, $f\left( x _{1} \right)\ne f\left( x _{2} \right)$. No two distinct inputs can give the same output. Therefore, the function must be one-to-one.



\paragraph{Basic concepts}



\[\[, Range, Codomain.svg\|thumb\|230x230px\|The domain is all the elements in set $A$ that can be mapped to the elements in set $B$.



The range is those elements in set $B$ which are mapped with the domain.



The codomain is all the elements in set $B$.\]\]There are a few more important ideas to learn from the modern definition of the function, and it comes from knowing the difference between domain, range, and codomain. We already discussed some of these, yet knowing about sets adds a new definition for each of the following ideas. Therefore, let us discuss these based on these new ideas. Let $A$ and $B$ be a set. If we were to combine these two sets, it would be defined as the \textit{cartesian cross product} $A\times B$. The subset of this product is the function. The below definitions are a little confusing. The best way to interpret these is to draw an image. To the right of these definitions is the image that corresponds to it.



\paragraph{Definition of domain, range, and codomain of a function}



The \textbf{domain} is defined to be a set $A$ with all elements $a\in A$ mapping to at least one unique $b\in B$ .



The set of elements in $B$ is the \textbf{range} of the function mapping $f:A\to B$ in the cartesian cross product, whereby the set of all elements $a\in A$ maps to some element $b\in B$ .



The \textbf{codomain} is the set $B$, where it is not necessarily the case that all elements $b\in B$ was mapped by some $a\in A$.



Note that the codomain is not as important as the other two concepts.



Take $f(x)=\sqrt{1-x^2}$ for example:



% [Image: The domain of the function is the interval from -1 to 1]



Figure: The domain of the function is the interval from -1 to 1



Because of the square root, the content in the square root should be strictly not smaller than 0.



> $1-x^2\geq0$



> $\Leftrightarrow -1\leq x\leq 1$



Thus the domain is



> $x\in[-1,1]$



% [Image: The range of the function is the interval from 0 to 1]



Figure: The range of the function is the interval from 0 to 1



Correspondingly, the range will be



> $f(x)\in[0,1]$



\paragraph{Other types of transformation}



There is one more final aspect that needs to be defined. We already have a good idea of what makes a mapping a function (e.g. it cannot be one-to-many). However, three more definitions that you will often hear will be necessary to talk about: \textit{injective}, \textit{surjective}, \textit{bijective}.



% [Image: The function mapping on the left is an example of an injective function because it is one-to-one. However, it is not surjective because the range and the codomain are not the same.]



Figure: The function mapping on the left is an example of an injective function because it is one-to-one. However, it is not surjective because the range and the codomain are not the same.



\begin{itemize}
    \item A function is \textbf{injective} if it is one-to-one.
    \item A function is \textbf{surjective} if it is "onto." That is, the mapping $f:A\to B$ has $B$ as the \textit{range} of the function, where the codomain and the range of the function are the same.
    \item A function is \textbf{bijective} if it is both surjective and injective.
\end{itemize}

Again, the above definitions are often very confusing. Again, another image is provided to the right of the bullet points. Along with that, another example is also provided. Let us analyze the function $g(x)=ax^{2}$.



\textbf{Given}: $g(x)=ax^{2}$, where $a$ is constant and $a\ne0$.



\textbf{Prove}: function $g$ is non-surjective and non-injective.



Notice that the only changing element in the function $g$ is $x$. Let us check to see the conditions of this function.



\textbf{Is it injective}? Let $g(x)=y$, and solve for $x$. First, divide by $a$.



$$y=ax^{2}$$



$$\Rightarrow \frac{y}{a}=x^{2}$$



Then take the square root of $x^2$. By definition, $\sqrt{x^{2}}=\pm\sqrt{x^{2}}=\pm x$, so



$$\frac{y}{a}=x^{2}$$



$$\Rightarrow\pm\sqrt{\frac{y}{a}}=\pm x$$



Notice, however, what we learned from the above manipulation. As a result of solving for $x$, we found that there are two solutions for $x$. However, this resulted in two different values from $\frac{y}{a}$. This means that for some individual $x$ that gives $y$, there are two different inputs that result in the same value. Because $f(x _{1})=f(x _{2})$ when $x _{1}\ne x _{2}$, this is by definition non-injective.



\textbf{Is it surjective}? As a natural consequence of what we learned about inputs, let us determine the range of the function. After all, the only way to determine if something is surjective is to see if the range applies to all real numbers. A good way to determine this is by finding a pattern involving our domains. Let us say we input a negative number for the domain: $f(-x)=a(-x)^{2}=a\cdot(-x)\cdot(-x)=ax^2$. This seemingly trivial exercise tells us that negative numbers give us non-negative numbers for our range. This is huge information! For $x\in\mathbb{R}$, the function always results $y>0$ for our range. For the set $B$, we have elements in that set that have no mappings from the set $A$. As such, $y\le0$ is the codomain of set $B$. Therefore, this function is non-surjective!



% [Image: This is an example of an expression which fails the vertical line test.]



Figure: This is an example of an expression which fails the vertical line test.



\paragraph{Tests for equations}



\paragraph{The vertical line test}



The vertical line test is a systematic test to find out if an equation involving $x$ and $y$ can serve as a function (with $x$ the independent variable and $y$ the dependent variable). Simply graph the equation and draw a vertical line through each point of the $x$-axis. If any vertical line ever touches the graph at more than one point, then the equation is not a function; if the line always touches at most one point of the graph, then the equation is a function.



The circle, on the right, is not a function because the vertical line intercepts two points on the graph when $x=1$.



\paragraph{The horizontal line and the algebraic 1-1 test}



Similarly, the horizontal line test, though does not test if an equation is a function, tests if a function is injective (one-to-one). If any horizontal line ever touches the graph at more than one point, then the function is not one-to-one; if the line always touches at most one point on the graph, then the function is one-to-one.



The algebraic 1-1 test is the non-geometric way to see if a function is one-to-one. The basic concept is that:



Assume there is a function $f$. If:



> $f(a)=f(b)$, and $a=b$, then



function $f$ is one-to-one.



Here is an example: prove that $f(x)=\frac{1-2x}{1+x}$ is injective.



Since the notation is the notation for a function, the equation is a function. So we only need to prove that it is injective. Let $a$ and $b$ be the inputs of the function and that $f(a)=f(b)$. Thus,



$$ \begin{aligned} \frac{1-2a}{1+a} &= \frac{1-2b}{1+b} \\\\  \Leftrightarrow (1+b)(1-2a) &= (1+a)(1-2b) \\\\  \Leftrightarrow 1 - 2a + b - 2ab &= 1 - 2b + a - 2ab \\\\  \Leftrightarrow 1 - 2a + b &= 1 - 2b + a \\\\  \Leftrightarrow 1 - 2a + 3b &= 1 + a \\\\  \Leftrightarrow 1 + 3b &= 1 + 3a \\\\  \Leftrightarrow a &= b \end{aligned} $$



So, the result is $a=b$, proving that the function is injective.



Another example is proving that $g(x)=x^2$ is not injective.



Using the same method, one can find that $a=\pm b$, which is not $a=b$. So, the function is not injective.



\paragraph{Remarks}



The following arise as a direct consequence of the definition of functions:



1.  By definition, for each "input" a function returns only one "output", corresponding to that input. While the same output may correspond to more than one input, one input cannot correspond to more than one output. This is expressed graphically as the \textit{vertical line test}: a line drawn parallel to the axis of the dependent variable (normally vertical) will intersect the graph of a function only once. However, a line drawn parallel to the axis of the independent variable (normally horizontal) may intersect the graph of a function as many times as it likes. Equivalently, this has an algebraic (or formula-based) interpretation. We can always say if $a=b$ , then $f(a)=f(b)$ , but if we only know that $f(a)=f(b)$ then we can't be sure that $a=b$ .

2.  Each function has a set of values, the function's \textit{domain}, which it can accept as input. Perhaps this set is all positive real numbers; perhaps it is the set {pork, mutton, beef}. This set must be implicitly/explicitly defined in the definition of the function. You cannot feed the function an element that isn't in the domain, as the function is not defined for that input element.

3.  Each function has a set of values, the function's \textit{range}, which it can output. This may be the set of real numbers. It may be the set of positive integers or even the set $\\{0,1\\}$. This set, too, must be implicitly/explicitly defined in the definition of the function.



Functions are an important foundation of mathematics. This modern interpretation is a result of the hard work of the mathematicians of history. It was not until recently that the definition of the relation was introduced by René Descartes in \textit{Geometry} (1637). Nearly a century later, the standard notation ($f(x)=y$) was first introduced by Leonhard Euler in \textit{Introductio in Analysin Infinitorum} and \textit{Institutiones Calculi Differentialis}. The term function was also a new innovation during Euler's time as well, being introduced Gottfried Wilhelm Leibniz in a 1673 letter "to describe a quantity related to points of a curve, such as a coordinate or curve's slope." Finally, the modern definition of the function being the relation among sets was first introduced in 1908 by Godfrey Harold Hardy where there is a relation between two variables $x$ and $y$ such that "to some values of $x$ at any rate correspond values of $y$."



\paragraph{Important functions}



\paragraph{Polynomials}



Polynomial functions are the most common and most convenient functions in the world of calculus. Their frequent appearances and their applications on computer calculations have made them important. A polynomial in a single indeterminate x can always be written (or rewritten) in the form:



$$f(x)=a _n x^n+a _{n-1}x^{n-1}+a _{n-2}x^{n-2}+\cdot\cdot\cdot+a _2x^2+a _1x+a _0$$



To be more concise, it can also be written in the summation form:



$$\sum _{k=0}^n a _k x^k$$



\paragraph{Constant}



% [Image: Two linear functions are shown on the image. One is f(x)=x/2+2 and the other is g(x)=-2x+5]



Figure: Two linear functions are shown on the image. One is $f(x)=\frac{1}{2}x+2$ and the other is $g(x)=-2x+5$



When $n=0$, the polynomial can be rewritten into the following function:



> $f(x)=C$, where $C$ is a constant.



The graph of this function is a horizontal line passing the point $(0,C)$.



\paragraph{Linear}



When $n=1$, the polynomial can be rewritten into



> $f(x)=mx+b$, where $m\text{ and }b$ are constants.



The graph of this function is a straight line passing the point $(0,b)$ and $(-\frac{b}{m},0)$, and the slope of this function is $m$.



% [Image: This is the graph of a quadratic function.]



Figure: This is the graph of a quadratic function.



\paragraph{Quadratic}



When $n=2$, the polynomial can be rewritten into



> $f(x)=ax^2+bx+c$, where $a,b,c$ are constants.



The graph of this function is a parabola, like the trajectory of a basketball thrown into the hoop.



If there are questions about the quadratic formula and other basic polynomial concepts, please review the respective chapters in Algebra.



\paragraph{Trigonometric}



Trigonometric functions are also very important because it can connect algebra and geometry. Trigonometric functions are explained in detail here due to its importance and difficulty.



% [Image: The curve on the left is an exponential function while the curve on the right is a logarithmic one]



Figure: The curve on the left is an exponential function while the curve on the right is a logarithmic one.



\paragraph{Exponential and Logarithmic}



Exponential and logarithmic functions have a unique identity when calculating the derivatives, so this is a great time to review those functions. The exponential function is defined as:



$$f(x)=a^x$$, where $a$ is a constant. while the logarithmic function is defined as:



$$f(x)=\log _b x$$, where $b$ is a constant. A special number will be frequently seen in those functions: the Euler's constant, also known as the base of the natural logarithm. Notated as $e$, it is defined as

$$e=\sum _{n=0}^\infty \frac{1}{n!}$$



\paragraph{Signum}



The Signum (sign) function is simply defined as



$$\text{sgn}(x)= \begin{cases}-1 & \text{if} & x0 \end{cases}$$



\paragraph{Properties of functions}



Sometimes, a lot of function manipulations cannot be achieved without some help from basic properties of functions.



\paragraph{Domain and range}



This concept is discussed above.



\paragraph{Growth}



Although it seems obvious to spot a function increasing or decreasing, without the help of graphing software, we need a mathematical way to spot whether the function is increasing or decreasing, or else our human minds cannot immediately comprehend the huge amount of information.



Assume a function $f(x)$ with inputs $x _1,x _2$, and that $x _1\in[a,b]$, $x _2\in[a,b]$, and $x _2>x _1$ at all times.



> If for all $x _1$ and $x _2$, $f(x _2)-f(x _1)>0$, then

>

> $f(x)$ is increasing in $[a,b]$

>

> If for all $x _1$ and $x _2$, $f(x _2)-f(x _1)

> $f(x)$ is decreasing in $[a,b]$



\textbf{Example:} In which intervals is $f(x)=\frac{1}{1-x^2}$ increasing?



Firstly, the domain is important. Because the denominator cannot be 0, the domain for this function is



$x\in(-\infty,-1)\cup(-1,1)\cup(1,\infty)$



In $(-\infty,-1)$, the growth of the function is:



> Let $x _1,x _2\in[-\infty,-1]$ and $x _2>x _1$

>

> Thus,

>

> $$f(x _2)-f(x _1)=\frac{1}{1-x _2^2}-\frac{1}{1-x _1^2}=\frac{x _2^2-x _1^2}{(1-x _2^2)(1-x _1^2)}$$

>

> Because both $x _1,x _2\in[-\infty,-1]$

>

> We have $(1-x _2^2)(1-x _1^2)>0$

>

> And because $x _2 > x _1$ and $x _1 < x _2 < 0$

> $\newline$ we also have $x _2^2-x _1^2

> So, $$f(x _2)-f(x _1)=\frac{x _2^2-x _1^2}{(1-x _2^2)(1-x _1^2)}

> $f(x)$ is decreasing in $(-\infty,-1)$



In $(-1,1)$



> Let $x _1,x _2\in[-1,1]$ and $x _2>x _1$

>

> Thus,

>

> $$f(x _2)-f(x _1)=\frac{1}{1-x _2^2}-\frac{1}{1-x _1^2}=\frac{x _2^2-x _1^2}{(1-x _2^2)(1-x _1^2)}$$

>

> because both $x _1,x _2\in[-\infty,-1]$

> $\newline$ we have $(1-x _2^2)(1-x _1^2)>0$

>

> However, the sign of $x _2^2-x _1^2$ in $(-1,1)$ cannot be determined. It can only be determined in $(-1,0)\text{ and }(0,1)$.

>

> > In $(-1,0)$

> >

> > because $x _2 > x _1$ and $x _1 < x _2 < 0$

> > $\newline$ we have $x _2^2-x _1^2 < 0$

> >

> > In $(0,1)$

> >

> > And because $ x _2 > x _1 \text{ and }0 < x _1 < x _2$

> > $\newline$ we have $ x _2^2-x _1^2>0$

>

> As a result,

>

> $f(x)$ is decreasing in $(-1,0)$ and increasing in $(0,1)$.



In $(1,\infty)$



> Let $x _1,x _2\in[1,\infty]$ and $x _2>x _1$

>

> Thus,

>

> $$f(x _2)-f(x _1)=\frac{1}{1-x _2^2}-\frac{1}{1-x _1^2}=\frac{x _2^2-x _1^2}{(1-x _2^2)(1-x _1^2)}$$

>

> because both $x _1,x _2\in[-\infty,-1]$

> $\newline$ we have $(1-x _2^2)(1-x _1^2)>0$

>

> And because $x _2 > x _1 \text{ and } 0 < x _1 < x _2$

> $\newline$ we have $ x _2^2-x _1^2>0$

>

> > So, $$f(x _2)-f(x _1)=\frac{x _2^2-x _1^2}{(1-x _2^2)(1-x _1^2)} > 0$$

>

> $f(x)$ is increasing in $(1,\infty)$.



Therefore, the intervals in which the function is increasing are $(0,1)\And(1,\infty)$.



$\blacksquare$



After learning derivatives, there will be more ways to determine the growth of a function.



Parity



The properties odd and even are associated with symmetry. While even functions have a symmetry about the $y$-axis, odd functions are symmetric about the origin. In mathematical terms:



> A function is even when $f(-x)=f(x)$

>

> A function is odd when $f(-x)=-f(x)$



**Example:** Prove that $f(x)=\frac{1}{1-x^2}$ is an even function.



because $$f(-x)=\frac{1}{1-(-x)^2}=\frac{1}{1-x^2}=f(x)$$



we have that $ f(x)$ is an even function



Manipulating functions



Addition, Subtraction, Multiplication and Division of functions



For two real-valued functions, we can add the functions, multiply the functions, raised to a power, etc. For example:



-   If we add the functions $y=3x+2$ and $y=x^2$ , we obtain $y=x^2+3x+2$ .



-   If we subtract $y=3x+2$ from $y=x^2$ , we obtain $y=x^2-(3x+2)$ . We can also write this as $y=x^2-3x-2$ .



-   If we multiply the function $y=3x+2$ and the function $y=x^2$ , we obtain $y=(3x+2)x^2$ . We can also write this as $y=3x^3+2x^2$ .



-   If we divide the function $y=3x+2$ by the function $y=x^2$ , we obtain $y=\frac{3x+2}{x^2}$ .



If a math problem wants you to add two functions $f$ and $g$ , there are two ways that the problem will likely be worded:



1.  If you are told that $f(x)=3x+2$ , that $g(x)=x^2$ , that $h(x)=f(x)+g(x)$ and asked about $h$ , then you are being asked to add two functions. Your answer would be $h(x)=x^2+3x+2$ .

2.  If you are told that $f(x)=3x+2$ , that $g(x)=x^2$ and you are asked about $f+g$ , then you are being asked to add two functions. The addition of $f$ and $g$ is called $f+g$ . Your answer would be $(f+g)(x)=x^2+3x+2$ .



Similar statements can be made for subtraction, multiplication and division.



Let $f(x)=3x+2$ and: $g(x)=x^2$ . Let's add, subtract, multiply and divide.



$$\begin{aligned}

(f+g)(x)

    \& = f(x)+g(x) \\\    \& = (3x+2)+(x^2) \\\    \& = x^2+3x+2 \,

\end{aligned}$$



$$\begin{aligned}

(f-g)(x)

    \& = f(x)-g(x) \\\    \& = (3x+2)-(x^2) \\\    \& = -x^2+3x+2 \,

\end{aligned}$$



$$\begin{aligned}

 (f\times g)(x)

          \& = f(x)\times g(x) \\\          \& = (3x+2)\times(x^2)\\\          \& = 3x^3+2x^2 \,

\end{aligned}$$



$$\begin{aligned}

\left(\frac{f}{g}\right)(x)

            \& = \frac{f(x)}{g(x)} \\\            \& = \frac{3x+2}{x^2} \\\            \& = \frac{3}{x}+\frac{2}{x^2}

\end{aligned}$$



Composition of functions



We begin with a fun (and not too complicated) application of composition of functions before we talk about what composition of functions is.



Example: Dropping a ball If we drop a ball from a bridge which is 20 meters above the ground, then the height of our ball above the earth is a function of time. The physicists tell us that if we measure time in seconds and distance in meters, then the formula for height in terms of time is $h=-4.9t^2+20$ . Suppose we are tracking the ball with a camera and always want the ball to be in the center of our picture. Suppose we have $\theta=f(h)$ The angle will depend upon the height of the ball above the ground and the height above the ground depends upon time. So the angle will depend upon time. This can be written as $\theta=f(-4.9t^2+20)$ . We replace $h$ with what it is equal to. This is the essence of composition.



Composition of functions is another way to combine functions which is different from addition, subtraction, multiplication or division.



The value of a function $f$ depends upon the value of another variable $x$ ; however, that variable could be equal to another function $g$ , so its value depends on the value of a third variable. If this is the case, then the first variable is a function $h$ of the third variable; this function ($h$) is called the **composition** of the other two functions ($f$ and $g$).



Example: Composing two functions Let $f(x)=3x+2$ and: $g(x)=x^2$ . The composition of $f$ with $g$ is read as either "f composed with g" or "f of g of x."



Let



$$h(x)=f(g(x))$$ Then



$$\begin{aligned}

h(x) \& = f(g(x)) \\\     \& = f(x^2) \\\     \& = 3(x^2)+2 \\\     \& = 3x^2+2 \,

\end{aligned}$$



Sometimes a math problem asks you compute $(f\circ g)(x)$ when they want you to compute $f(g(x))$ ,



Here, $h$ is the composition of $f$ and $g$ and we write $h=f\circ g$ . Note that composition is not commutative:



$$f(g(x))=3x^2+2$$ , and



$$\begin{aligned}

g(f(x)) \& = g(3x+2) \\\        \& = (3x+2)^2 \\\        \& = 9x^2+12x+4 \,

\end{aligned}$$



> So, $f(g(x))\ne g(f(x))$.



Composition of functions is very common, mainly because functions themselves are common. For instance, squaring and sine are both functions:



$$\text{square}(x)=x^2$$



$$\text{sine}(x)=\sin(x)$$



Thus, the expression $\sin^2(x)$ is a composition of functions:



$$\begin{array}{|c|} \hline \sin^2(x) = \text{square}(\sin(x)) = \text{square}(\text{sine}(x)) \\\\ \hline \end{array}$$



(Note that this is *not* the same as $\text{sine}(\text{square}(x)) = \sin(x^2)$.) Since the function sine equals $\tfrac{1}{2}$ if $x = \tfrac{\pi}{6}$,



$$\text{square}\left(\sin\left(\tfrac{\pi}{6}\right)\right) = \text{square}\left(\tfrac{1}{2}\right)$$.



Since the function square equals $\tfrac{1}{4}$ if $x = \tfrac{\pi}{6}$,



$$\sin^2\left(\tfrac{\pi}{6}\right) = \text{square}\left(\sin\left(\tfrac{\pi}{6}\right)\right) = \text{square}\left(\tfrac{1}{2}\right) = \tfrac{1}{4}$$.



Transformations



Transformations are a type of function manipulation that are very common. They consist of multiplying, dividing, adding or subtracting constants to either the input or the output. Multiplying by a constant is called **dilation** and adding a constant is called **translation**. Here are a few examples:



$$\begin{aligned} f(2 \times x) \& \quad \text{Dilation} \\\\ f(x + 2) \& \quad \text{Translation} \\\\ 2 \times f(x) \& \quad \text{Dilation} \\\\ 2 + f(x) \& \quad \text{Translation} \end{aligned}$$



![Examples of horizontal and vertical translations](https://upload.wikimedia.org/wikipedia/commons/8/80/4_function_translations.jpg)



Figure: Examples of horizontal and vertical translations



![Examples of horizontal and vertical dilations](https://upload.wikimedia.org/wikipedia/commons/6/61/4_function_dilations.jpg)



Figure: Examples of horizontal and vertical dilations



Translations and dilations can be either horizontal or vertical. Examples of both vertical and horizontal translations can be seen at right. The red graphs represent functions in their 'original' state, the solid blue graphs have been translated (shifted) horizontally, and the dashed graphs have been translated vertically.



Dilations are demonstrated in a similar fashion. The function



$$f(2\times x)$$ has had its input doubled. One way to think about this is that now any change in the input will be doubled. If I add one to $x$, I add two to the input of $f$, so it will now change twice as quickly. Thus, this is a horizontal dilation by $\frac{1}{2}$ because the distance to the $y$-axis has been **halved**. A vertical dilation, such as



$$2\times f(x)$$ is slightly more straightforward. In this case, you double the output of the function. The output represents the distance from the $x$-axis, so in effect, you have made the graph of the function 'taller'. Here are a few basic examples where $a$ is any positive constant:



  ---------------------------------------------- ------------------------------------- ----------------------------------------------- ----------------

  Original graph                                 $f(x)$                                Rotation about origin                           $-f(-x)$

  Horizontal translation by $a$ units **left**   $f(x+a)$                              Horizontal translation by $a$ units **right**   $f(x-a)$

  Horizontal dilation by a factor of $a$         $f\left(x\times\tfrac{1}{a}\right)$   Vertical dilation by a factor of $a$            $a\times f(x)$

  Vertical translation by $a$ units **down**     $f(x)-a$                              Vertical translation by $a$ units **up**        $f(x)+a$

  Reflection about $x$-axis                      $-f(x)$                               Reflection about $y$-axis                       $f(-x)$

  ---------------------------------------------- ------------------------------------- ----------------------------------------------- ----------------



Inverse functions



We call $g(x)$ the inverse function of $f(x)$ if, for all $x$ :



$$g(f(x))=f(g(x))=x$$ A function $f(x)$ has an inverse function if and only if $f(x)$ is one-to-one. For example, the inverse of $f(x)=x+2$ is $g(x)=x-2$ . The function $f(x)=\sqrt{1-x^2}$ has no inverse because it is not injective. Similarly, the inverse functions of trigonometric functions have to undergo transformations and limitations to be considered as valid functions.



Notation



The inverse function of $f$ is denoted as $f^{-1}(x)$ . Thus, $f^{-1}(x)$ is defined as the function that follows this rule



$$f(f^{-1}(x))=f^{-1}(f(x))=x$$ To determine $f^{-1}(x)$ when given a function $f$ , substitute $f^{-1}(x)$ for $x$ and substitute $x$ for $f(x)$ . Then solve for $f^{-1}(x)$ , provided that it is also a function.



**Example:** Given $f(x)=2x-7$ , find $f^{-1}(x)$ .



Substitute $f^{-1}(x)$ for $x$ and substitute $x$ for $f(x)$ . Then solve for $f^{-1}(x)$ :



$$\begin{aligned} f(x) \&= 2x - 7 \\\\ \&\Leftrightarrow x = 2[f^{-1}(x)] - 7 \\\\ \&\Leftrightarrow x + 7 = 2[f^{-1}(x)] \\\\ \&\Leftrightarrow \frac{x + 7}{2} = f^{-1}(x) \end{aligned}$$



To check your work, confirm that $f^{-1}(f(x))=x$ :



$$\begin{aligned} f^{-1}(f(x)) \& = f^{-1}(2x - 7) \\\\ \& = \frac{(2x - 7) + 7}{2} \\\\ \& = \frac{2x}{2} \\\\ \& = x \end{aligned}$$



If $f$ isn't one-to-one, then, as we said before, it doesn't have an inverse. Then this method will fail.



**Example:** Given $f(x)=x^2$ , find $f^{-1}(x)$.



Substitute $f^{-1}(x)$ for $x$ and substitute $x$ for $f(x)$ . Then solve for $f^{-1}(x)$ :



$$\begin{aligned} f(x) \&= x^2 \\\\ x \&= \left(f^{-1}(x)\right)^2 \\\\ f^{-1}(x) \&= \pm\sqrt{x} \end{aligned}$$



Since there are two possibilities for $f^{-1}(x)$ , it's not a function. Thus $f(x)=x^2$ doesn't have an inverse. Of course, we could also have found this out from the graph by applying the Horizontal Line Test. It's useful, though, to have lots of ways to solve a problem, since in a specific case some of them might be very difficult while others might be easy. For example, we might only know an algebraic expression for $f(x)$ but not a graph.



\paragraph{Resources}



\begin{itemize}
    \item \href{https://www.khanacademy.org/math/algebra/x2f8bb11595b61c86:functions/x2f8bb11595b61c86:evaluating-functions/v/what-is-a-function}{What is a Function?}, Khan Academy
    \item \href{https://www.khanacademy.org/math/algebra2/functions_and_graphs}{Function and their graphs}, Khan Academy
    \item \href{https://www.cliffsnotes.com/study-guides/calculus/precalculus/functions/relations-vs-functions}{Relations vs. Functions}, Cliff's Notes
    \item \href{https://courses.lumenlearning.com/ivytech-collegealgebra/chapter/determine-whether-a-relation-represents-a-function/}{Determining if a Relation is a Function}, Lumen Learning
    \item \href{https://courses.lumenlearning.com/waymakercollegealgebra/chapter/identify-functions-using-graphs/}{Identifying Functions Using Graphs}, Lumen Learning
\end{itemize}

\paragraph{Licensing}



Content obtained and/or adapted from:



\begin{itemize}
    \item \href{https://en.wikibooks.org/wiki/Calculus/Functions}{Functions, Wikibooks: Calculus} under a CC BY-SA license
\end{itemize}

\subsection{Learning Objectives}

\begin{enumerate}
    \item \textbf{Understand the Definition of a Function:}
    \item \textbf{Outcome:} Students will be able to clearly define a function as a mapping from a set $ X $ to a set $ Y $, ensuring that each element in $ X $ corresponds to exactly one element in $ Y $.
    \item \textbf{Details:} This includes being able to explain the concept of a function both formally and intuitively, using examples like real-valued functions and vector-valued functions.
    \item \textbf{Evaluate Functions Using Function Notation:}
    \item \textbf{Outcome:} Students will be able to use function notation to evaluate the output of a function for a given input.
    \item \textbf{Details:} This includes substituting specific values into a function's formula and correctly interpreting the notation $ f(x) $ as the value of the function at $ x $.
    \item \textbf{Identify and Classify Functions Based on Injectivity and Surjectivity:}
    \item \textbf{Outcome:} Students will be able to determine whether a function is injective (one-to-one), surjective (onto), or bijective (both injective and surjective).
    \item \textbf{Details:} This includes using definitions and tests such as the vertical line test for functions, and horizontal line and algebraic tests for injectivity and surjectivity, with specific examples to solidify understanding.
\end{enumerate}

\subsection{Assessment Instrument}

\begin{enumerate}
    \item \textbf{Assessment Instrument: Properties of Functions Quiz} \textbf{Format:}
    \item \textbf{Type:} Mixed-format quiz (multiple-choice, short answer, and problem-solving questions)
    \item \textbf{Duration:} 45 minutes
    \item \textbf{Total Marks:} 30 points \textbf{Sections:}
    \item \textbf{Domain and Range (10 points):}
    \item \textbf{Question 1:} Identify the domain and range of the following functions:
    \item a. $ f(x) = \sqrt{x - 2} $
    \item b. $ g(x) = \frac{1}{x+3} $
    \item c. $ h(x) = x^2 - 4x + 3 $
    \item d. $ k(x) = \ln(x) $
    \item \textbf{Question 2:} Determine whether the following statements are true or false:
    \item a. The range of $ f(x) = \sqrt{x} $ is all non-negative real numbers.
    \item b. The domain of $ g(x) = \frac{1}{x-1} $ is all real numbers except $ x = 1 $.
    \item c. The domain of $ h(x) = \ln(x-1) $ is all real numbers.
    \item d. The range of $ k(x) = \frac{1}{x^2} $ is all positive real numbers.
    \item \textbf{Growth of Functions (10 points):}
    \item \textbf{Question 3:} For the function $ f(x) = \frac{1}{1 - x^2} $:
    \item a. Determine the intervals on which the function is increasing.
    \item b. Determine the intervals on which the function is decreasing.
    \item \textbf{Question 4:} Explain why the function $ f(x) = 2x + 3 $ is always increasing.
    \item \textbf{Parity of Functions (10 points):}
    \item \textbf{Question 5:} Determine whether the following functions are even, odd, or neither:
    \item a. $ f(x) = x^3 - x $
    \item b. $ g(x) = \cos(x) $
    \item c. $ h(x) = x^2 + 4x + 4 $
    \item d. $ k(x) = \sin(x) + x $
    \item \textbf{Question 6:} Prove that $ f(x) = \frac{1}{1 + x^2} $ is an even function. \textbf{Scoring Rubric:}
    \item \textbf{Domain and Range:} 2 points for each correct domain and range identification (total 8 points), and 0.5 points for each correct true/false response (total 2 points).
    \item \textbf{Growth of Functions:} 5 points for correctly identifying the intervals of increasing and decreasing (total 5 points), and 5 points for a clear explanation of the function's growth (total 5 points).
    \item \textbf{Parity of Functions:} 2 points for correctly determining the parity of each function (total 8 points), and 2 points for a correct proof of evenness (total 2 points). This quiz will help assess the students' understanding and application of the properties of functions, including domain and range, growth patterns, and the parity of functions. It ensures that they can correctly analyze and categorize functions based on these properties.
\end{enumerate}

%%===============================================================
\section{Domain}
\label{sec:domain}
%%===============================================================

\paragraph{Definition}



In mathematics, the domain or set of departure of a function is the set into which all of the input of the function is constrained to fall. It is the set X in the notation $f: X \rightarrow Y$, and is alternatively denoted as dom(f). Since a (total) function is defined on its entire domain, its domain coincides with its domain of definition. However, this coincidence is no longer true for a partial function, since the domain of definition of a partial function can be a proper subset of the domain.



A domain is part of a function f if f is defined as a triple (X, Y, G), where X is called the domain of f, Y its codomain, and G its graph.



A domain is not part of a function f if f is defined as just a graph. For example, it is sometimes convenient in set theory to permit the domain of a function to be a proper class X, in which case there is formally no such thing as a triple (X, Y, G). With such a definition, functions do not have a domain, although some authors still use it informally after introducing a function in the form $f: X \rightarrow Y$.



For instance, the domain of cosine is the set of all real numbers, while the domain of the square root consists only of numbers greater than or equal to 0 (ignoring complex numbers in both cases).



If the domain of a function is a subset of the real numbers and the function is represented in a Cartesian coordinate system, then the domain is represented on the x-axis. The domain of a function f can be thought of as the set of all x values that can be plugged into f(x) that return a valid output. For example, if we have a function g(x) in the Cartesian plane, the domain is all of the x values such that g(x) is a real number.



Examples:



\begin{itemize}
    \item Let $S$ be a set of ordered pairs such that $S = \\{(1,2), (2,3), (4, 7), (13, 9), (-20, 0)\\}$. The domain is the set of all x values of $S$, so the domain is $\\{-20, 1, 2, 4, 13\\}$.
    \item The domain of $g(x) = 1/x$ is all real numbers EXCEPT 0, since 1/0 is not defined.
    \item The domain of $h(x) = \sqrt{x}$ is $[0, \infty)$, since $\sqrt{x}$ is only defined when $x$ is nonnegative (that is, when $x$ is greater than or equal to 0).
\end{itemize}

\paragraph{Resources}



\begin{itemize}
    \item \href{https://www.youtube.com/watch?v=Q3NWljhiSJg}{Domain and Range: Basic Idea}, patrickJMT
    \item \href{https://courses.lumenlearning.com/ivytech-collegealgebra/chapter/find-the-domain-of-a-function-defined-by-an-equation/}{Finding Domain with an Equation}, Lumen Learning
    \item \href{https://courses.lumenlearning.com/ivytech-collegealgebra/chapter/find-domain-and-range-from-graphs/}{Finding Domain and Range with Graphs}, Lumen Learning
    \item \href{https://www.youtube.com/watch?v=w81y25anEOM}{Finding the Domain Algebraically}, patrickJMT
    \item \href{https://www.youtube.com/watch?v=BxaYyS6lsQ4}{Finding Domain and Range of a Piecewise Function}, patrickJMT
    \item \href{https://mathculus.com/how-to-find-the-domain-of-a-function-algebraically/}{How to Find Domain + Example Problems}, Math Culus
\end{itemize}

\paragraph{Licensing}



Content obtained and/or adapted from:



\begin{itemize}
    \item \href{https://en.wikipedia.org/wiki/Domain_of_a_function}{Domain of a Function, Wikipedia} under a CC BY-SA license
\end{itemize}

\subsection{Learning Objectives}

\begin{enumerate}
    \item Here are three Student Learning Objectives (SLOs) based on the provided definition of the domain of a function:
    \item \textbf{Identify the Domain of a Function:}
    \item \textbf{Outcome:} Students will be able to identify the domain of a function given its algebraic expression or graph.
    \item \textbf{Details:} This includes determining the set of all possible input values (x-values) that result in valid outputs for functions such as polynomials, rational functions, and square root functions.
    \item \textbf{Analyze the Domain of Complex Functions:}
    \item \textbf{Outcome:} Students will be able to analyze and describe the domain of more complex functions, including piecewise functions and functions involving absolute values.
    \item \textbf{Details:} This involves understanding the constraints on input values imposed by different components of the function, such as denominators, square roots, and conditional expressions in piecewise functions.
    \item \textbf{Compare Domains Across Different Functions:}
    \item \textbf{Outcome:} Students will be able to compare and contrast the domains of different functions, explaining how changes in function definitions affect their domains.
    \item \textbf{Details:} This includes exploring how modifications to a function's formula, such as shifting or scaling, impact the set of valid input values and understanding how the domain of one function relates to that of another.
\end{enumerate}

\subsection{Assessment Instrument}

\begin{enumerate}
    \item \textbf{Assessment Instrument: Domain of a Function Quiz} \textbf{Format:}
    \item \textbf{Type:} Mixed-format quiz (multiple-choice, short answer, and problem-solving questions)
    \item \textbf{Duration:} 45 minutes
    \item \textbf{Total Marks:} 30 points \textbf{Sections:}
    \item \textbf{Identifying the Domain (10 points):}
    \item \textbf{Question 1:} Determine the domain of the following functions:
    \item a. $ f(x) = \frac{1}{x-3} $
    \item b. $ g(x) = \sqrt{x+2} $
    \item c. $ h(x) = \frac{2x}{x^2 - 4} $
    \item d. $ k(x) = \ln(x) $
    \item e. $ m(x) = \sqrt{x^2 - 9} $
    \item \textbf{Analyzing the Domain of Complex Functions (10 points):}
    \item \textbf{Question 2:} Analyze the following functions and describe their domains:
    \item a. $ f(x) = \frac{1}{\sqrt{x-1}} $
    \item b. $ g(x) = \sqrt{4 - x^2} $
    \item c. $ h(x) = \frac{2x + 5}{x^2 - 1} $
    \item d. $ k(x) = \frac{\sqrt{x + 3}}{x - 2} $
    \item e. $ m(x) = |x - 3| - 4 $
    \item \textbf{Comparing Domains Across Different Functions (10 points):}
    \item \textbf{Question 3:} Compare the domains of the following pairs of functions and explain how the changes in the function's definition affect their domains:
    \item a. $ f(x) = \frac{1}{x} $ and $ g(x) = \frac{1}{x - 2} $
    \item b. $ h(x) = \sqrt{x} $ and $ k(x) = \sqrt{x - 4} $
    \item c. $ m(x) = \frac{1}{x^2 - 9} $ and $ n(x) = \frac{1}{x^2 + 1} $
    \item d. $ p(x) = \frac{1}{x^2} $ and $ q(x) = \frac{1}{x^2 - 1} $
    \item e. $ r(x) = \sqrt{x + 5} $ and $ s(x) = \sqrt{x - 5} $ \textbf{Scoring Rubric:}
    \item \textbf{Identifying the Domain:} 2 points for each correct domain (total 10 points)
    \item \textbf{Analyzing the Domain of Complex Functions:} 2 points for each correct analysis and domain description (total 10 points)
    \item \textbf{Comparing Domains Across Different Functions:} 2 points for each correct comparison and explanation (total 10 points) The quiz is designed to assess students' ability to identify, analyze, and compare the domains of various functions, ensuring a thorough understanding of how the domain affects the function and its valid input values.
\end{enumerate}

%%===============================================================
\section{Range}
\label{sec:range}
%%===============================================================

\paragraph{Definition}



In mathematics, the \textbf{range} of a function may refer to either of two closely related concepts:



\begin{itemize}
    \item The codomain of the function
    \item The image of the function
\end{itemize}

Given two sets X and Y, a binary relation f between X and Y is a (total) function (from X to Y) if for every x in X there is exactly one y in Y such that f relates x to y. The sets X and Y are called domain and codomain of f, respectively. The image of f is then the subset of Y consisting of only those elements y of Y such that there is at least one x in X with $f(x) = y$.



In algebra, the range (or codomain) of a function is all of the possible outputs of the function. That is, if x is any element of the domain of some function f, then $f(x)$ is in the range of the function f.



Examples:



\begin{itemize}
    \item Let $S$ be a set of ordered pairs such that $S = \\{(1,2), (2,3), (4, 7), (13, 9), (-20, 0)\\}$. The range is the set of all y values of $S$, so the range is $\\{0, 2, 3, 7, 9\\}$.
    \item The range of $g(x) = 1/x$ is all real numbers EXCEPT for 0. We know this because for all nonzero real numbers M, 1/M is a nonzero number and is in the domain of $g(x)$ (since the domain of this function is all nonzero numbers). So, we know that $1/(1/M) = M$ is in the range, where M is all nonzero numbers. There is no real number M such that $1/M = 0$ though, which is why 0 is not in the range of $g(x)$.
    \item The range of $h(x) = x^2 + 2$ is $[2,\infty)$. We can see this on the graph of $h(x)$ easily: the lowest point, or vertex, of the parabola is at $(0, 2)$, so 2 is in the range. The parabola extends up to infinity on either side of the vertex, so we know that the range must be all numbers from 2 to infinity.
\end{itemize}

\paragraph{Resources}



\begin{itemize}
    \item \href{https://www.intmath.com/functions-and-graphs/2a-domain-and-range.php}{Domain and Range}, Interactive Mathematics
    \item \href{https://www.youtube.com/watch?v=Q3NWljhiSJg}{Domain and Range: Basic Idea}, patrickJMT
    \item \href{https://courses.lumenlearning.com/ivytech-collegealgebra/chapter/find-domain-and-range-from-graphs/}{Finding Domain and Range with Graphs}, Lumen Learning
    \item \href{https://www.youtube.com/watch?v=BxaYyS6lsQ4}{Finding Domain and Range of a Piecewise Function}, patrickJMT
    \item \href{https://mathculus.com/how-to-find-the-range-of-a-function-algebraically/}{How to Find Range + Example Problems}, Math Culus
\end{itemize}

\paragraph{Licensing}



Content obtained and/or adapted from:



\begin{itemize}
    \item \href{https://en.wikipedia.org/wiki/Range_of_a_function}{Range of a function, Wikipedia} under a CC BY-SA license
\end{itemize}

\subsection{Learning Objectives}

\begin{enumerate}
    \item Here are three Student Learning Objectives (SLOs) based on the concept of the range of a function: ---
    \item \textbf{Identify the Range of a Function:}
    \item \textbf{Outcome:} Students will be able to determine the range of a given function by analyzing its algebraic expression.
    \item \textbf{Details:} This includes finding the set of all possible output values (range) for various types of functions, such as linear, quadratic, and rational functions, and justifying their answers through algebraic manipulation or by examining the function’s graph. ---
    \item \textbf{Distinguish Between Codomain and Image:}
    \item \textbf{Outcome:} Students will be able to differentiate between the codomain and image of a function and identify the image within a specific codomain.
    \item \textbf{Details:} This includes understanding that the codomain is the set of all potential outputs (as defined by the function’s definition), while the image is the actual set of outputs produced by applying the function to its domain. Students will practice with examples to clarify these distinctions. ---
    \item \textbf{Apply the Concept of Range to Real-World Problems:}
    \item \textbf{Outcome:} Students will be able to apply their understanding of the range of a function to solve real-world problems, such as determining the possible outputs in various scenarios (e.g., physics, economics).
    \item \textbf{Details:} This includes interpreting functions that model real-world phenomena, identifying constraints on outputs based on the situation, and justifying the range based on contextual factors such as physical limitations or initial conditions. --- These SLOs guide students in understanding and applying the concept of the range of a function, from basic identification to distinguishing related concepts and solving applied problems.
\end{enumerate}

\subsection{Assessment Instrument}

\begin{enumerate}
    \item \textbf{Assessment Instrument: Range of a Function Quiz} \textbf{Format:}
    \item \textbf{Type:} Mixed-format quiz (multiple-choice, short answer, and problem-solving questions)
    \item \textbf{Duration:} 45 minutes
    \item \textbf{Total Marks:} 30 points \textbf{Sections:}
    \item \textbf{Identifying the Range (10 points):}
    \item \textbf{Question 1:} Determine the range of the following functions:
    \item a. $ f(x) = x^2 + 3 $
    \item b. $ g(x) = \frac{1}{x} $
    \item c. $ h(x) = \sin(x) $
    \item d. $ k(x) = e^x $
    \item e. $ m(x) = \sqrt{x-2} $
    \item \textbf{Format:} Short answer (2 points each)
    \item \textbf{Distinguishing Codomain and Image (10 points):}
    \item \textbf{Question 2:} For each function below, identify the codomain and the image:
    \item a. $ f: \mathbb{R} \rightarrow \mathbb{R} $, $ f(x) = x^3 $
    \item b. $ g: \mathbb{R} \rightarrow [0, \infty) $, $ g(x) = x^2 $
    \item c. $ h: \mathbb{R} \rightarrow \mathbb{R} $, $ h(x) = \cos(x) $
    \item d. $ k: \mathbb{R} \rightarrow (0, \infty) $, $ k(x) = 2^x $
    \item \textbf{Format:} Short answer (2.5 points each)
    \item \textbf{Applying the Concept of Range (10 points):}
    \item \textbf{Question 3:} Solve the following problems by determining the range of the function in the context of the problem:
    \item a. A ball is thrown upwards with an initial velocity of 20 m/s. The height of the ball at time $ t $ is given by $ h(t) = -5t^2 + 20t + 1.5 $. Determine the range of the height $ h(t) $.
    \item b. A company models its revenue using the function $ R(x) = 5000(1 - e^{-0.1x}) $, where $ x $ represents the number of units sold. Determine the range of the revenue $ R(x) $.
    \item \textbf{Format:} Problem-solving (5 points each) \textbf{Scoring Rubric:}
    \item \textbf{Identifying the Range:} 2 points for each correctly identified range (total 10 points)
    \item \textbf{Distinguishing Codomain and Image:} 2.5 points for correctly identifying both the codomain and image of each function (total 10 points)
    \item \textbf{Applying the Concept of Range:} 5 points for correctly determining the range in each context and justifying the answer (total 10 points) This quiz assesses students' understanding of the range of a function, their ability to distinguish between related concepts, and their application of these concepts to real-world scenarios.
\end{enumerate}

%%===============================================================
\section{Toolkit Functions}
\label{sec:toolkit-functions}
%%===============================================================

\paragraph{Constant Function}



% [Image: Constant function]



For the constant function $f(x) = c$ the domain consists of all real numbers; there are no restrictions on the input. The only output value is the constant $c$, so the range is the set $\{c\}$ that contains this single element. In interval notation, this is written as $[c,c]$, the interval that both begins and ends with $c$.



Domain: $( -\infty ,\infty )$



Range: $[c, c]$



\paragraph{Identity Function}



% [Image: Identity function]



For the identity function $f(x) = x$ there is no restriction on $x$. Both the domain and range are the set of all real numbers.



Domain: $( -\infty ,\infty )$



Range: $( -\infty ,\infty )$



\paragraph{Absolute Value Function}



% [Image: Absolute value function]



For the absolute value function $f(x) = |x|$ there is no restriction on x. The outputs are $x$ for $x\ge 0$ and $-x$ for $x < 0$, so the range is all numbers greater than or equal to 0.



Domain: $( -\infty ,\infty )$



Range : $[0 ,\infty )$



\paragraph{Quadratic Function}



% [Image: Quadratic function]



For the quadratic function $f(x) = x^2$ the domain is all real numbers since the horizontal extent of the graph is the whole real number line. Because the graph does not include any negative values for the range, the range is only nonnegative real numbers.



Domain: $( -\infty ,\infty )$



Range : $[0 ,\infty )$



\paragraph{Cubic Function}



% [Image: Cubic function]



For the cubic function $f(x) = x^3$ the domain is all real numbers because the horizontal extent of the graph is the whole real number line. The same applies to the vertical extent of the graph, so the domain and range include all real numbers.



Domain: $( -\infty ,\infty )$



Range: $( -\infty ,\infty )$



\paragraph{Rational Function}



% [Image: Rational function]



For the rational function (also known as the reciprocal function) $f(x) = \frac{1}{x}$ we cannot divide by 0, so we must exclude 0 from the domain. Further, 1 divided by any value can never be 0, so the range also will not include 0.



Domain: $( -\infty , 0) \cup (0, \infty )$



Range: $( -\infty , 0) \cup (0, \infty )$



\paragraph{Squared Rational Function}



% [Image: Squared rational function]



For the squared rational function $f(x) = \frac{1}{x^2}$ we cannot divide by 0, so we must exclude 0 from the domain. Further, 1 divided by any value can never be 0, so the range also will not include 0. Also, since $x^2 > 0$ for all $x\neq 0$, the range will only consist of positive numbers.



Domain: $( -\infty , 0) \cup (0, \infty )$



Range: $(0, \infty )$



\paragraph{Square Root Function}



% [Image: Square root function]



For the square root function $f(x) = \sqrt{x}$ we cannot take the square root of a negative real number, so the domain must be 0 or greater. The range also excludes negative numbers because the square root of a positive number $x$ is defined to be positive, even though the square of the negative number $-\sqrt{x}$ also gives us $x$.



Domain: $[0, \infty )$



Range: $[0, \infty )$



\paragraph{Cube Root Function}



% [Image: Cube root function]



For the cube root function $f(x) = \sqrt[3]{x}$ the domain and range include all real numbers. Note that there is no problem taking a cube root, or any odd-integer root, of a negative number, and the resulting output is negative (it is an odd function).



Domain: $( -\infty ,\infty )$



Range: $( -\infty ,\infty )$



\paragraph{Resources}



\begin{itemize}
    \item \href{https://mathresearch.utsa.edu/wikiFiles/MAT1053/Domain_Range_and_Toolkit_Functions/MAT1053_M1.2Domain_Range_and_Toolkit_Functions.pdf}{Domain Range and Toolkit Functions}, Book Chapter
    \item \href{https://mathresearch.utsa.edu/wikiFiles/MAT1053/Domain_Range_and_Toolkit_Functions/MAT1053_M1.2Domain_Range_and_Toolkit_FunctionsGN.pdf}{Guided Notes}
    \item \href{https://courses.lumenlearning.com/cuny-hunter-collegealgebra/chapter/toolkit-functions/}{Toolkit Functions}, Lumen Learning
\end{itemize}

\paragraph{Licensing}



Content obtained and/or adapted from:



\begin{itemize}
    \item \href{https://courses.lumenlearning.com/cuny-hunter-collegealgebra/chapter/toolkit-functions/}{Toolkit Functions, Lumen Learning: Hunter College MATH101} under a CC0 (public domain) license
\end{itemize}

\subsection{Learning Objectives}

\begin{enumerate}
    \item Here are three Student Learning Objectives (SLOs) for the given functions:
    \item \textbf{Analyze the Domain and Range of Constant, Identity, and Quadratic Functions:}
    \item \textbf{Outcome:} Students will be able to determine and explain the domain and range for constant, identity, and quadratic functions.
    \item \textbf{Details:} This includes understanding how the domain and range are represented in interval notation and how they relate to the graph of the function.
    \item \textbf{Compare and Contrast the Behavior of Absolute Value and Square Root Functions:}
    \item \textbf{Outcome:} Students will be able to compare and contrast the domain, range, and graphical behavior of absolute value and square root functions.
    \item \textbf{Details:} This involves identifying how each function handles positive and negative inputs, as well as understanding the implications of their respective domains and ranges on the function's output.
    \item \textbf{Investigate the Properties of Rational and Squared Rational Functions:}
    \item \textbf{Outcome:} Students will be able to analyze and interpret the properties of rational and squared rational functions, specifically focusing on their domains and ranges.
    \item \textbf{Details:} This includes understanding the exclusions in the domain (such as where the function is undefined) and the implications for the range, as well as how these properties are reflected in the graphs of the functions.
\end{enumerate}

\subsection{Assessment Instrument}

\begin{enumerate}
    \item \textbf{Assessment Instrument: Toolkit Functions Quiz} \textbf{Format:}
    \item \textbf{Type:} Mixed-format quiz (multiple-choice, short answer, and problem-solving questions)
    \item \textbf{Duration:} 45 minutes
    \item \textbf{Total Marks:} 30 points \textbf{Sections:}
    \item \textbf{Analyzing Domain and Range (10 points):}
    \item \textbf{Question 1:} Determine the domain and range of the following functions:
    \item a. $f(x) = 5$ (Constant Function)
    \item b. $f(x) = |x|$ (Absolute Value Function)
    \item c. $f(x) = x^2$ (Quadratic Function)
    \item d. $f(x) = \frac{1}{x}$ (Rational Function)
    \item \textbf{Comparing Function Behaviors (10 points):}
    \item \textbf{Question 2:} Compare and contrast the behavior of the following functions:
    \item a. Explain the difference between the range of $f(x) = |x|$ (Absolute Value Function) and $f(x) = \sqrt{x}$ (Square Root Function).
    \item b. Describe how the domain of $f(x) = \frac{1}{x}$ (Rational Function) differs from the domain of $f(x) = x^3$ (Cubic Function).
    \item \textbf{Interpreting Graphs of Functions (10 points):}
    \item \textbf{Question 3:} Given the graphs of the following functions, identify the function type and describe the key properties such as domain, range, and any restrictions:
    \item a. A graph of $f(x) = x$ (Identity Function)
    \item b. A graph of $f(x) = \frac{1}{x^2}$ (Squared Rational Function)
    \item c. A graph of $f(x) = \sqrt[3]{x}$ (Cube Root Function) \textbf{Scoring Rubric:}
    \item \textbf{Analyzing Domain and Range:} 2.5 points for each correct identification and explanation of domain and range (total 10 points)
    \item \textbf{Comparing Function Behaviors:} 5 points for each correctly explained comparison (total 10 points)
    \item \textbf{Interpreting Graphs of Functions:} 3.33 points for each correct identification and explanation of the graph, domain, range, and restrictions (total 10 points) This quiz will assess students' understanding of domain and range analysis, comparison of different function behaviors, and interpretation of function graphs, focusing on the toolkit functions.
\end{enumerate}

%%===============================================================
\section{Intro to Power Functions}
\label{sec:intro-to-power-functions}
%%===============================================================

With practice the concept of slope for linear functions becomes intuitive. It makes sense that the line that fits the equation $y=2x$ has a steeper ascent then the line that fits the equation $y=\frac{1}{2}x$. You only have to move horizontally one unit to change your vertical direction two for the former when you graph $y=2x$. How many blocks do you need to move horizontally to change your vertical direction by one for the line $y=\frac{1}{2}x$?



> When we express concepts like $y=x^2$ the abstract behavior of what is being represented becomes a little harder to see.



A \textbf{monomial} of one variable, let's say x, is an algabraic expression of the form $c x^m $



where



\begin{itemize}
    \item $c$ is a constant, and
    \item $m$ is a non-negative integer (e.g., 0, 1, 2, 3, \...).
\end{itemize}

The integer $m$ is called the \textbf{degree} of the monomial.



The idea of a monomial of degree zero appears a bit mystical since it always represents one, except when the value of the variable is set equal to zero when the result is undefined. This idea allows us preserve the value of the constant in the monomial. We know that $cx^0$ is always equal to $c$ because even though we have 0 x's (somethings) we still have a c. When x = 0 things are difficult because the value we started with, 0, represents nothing.



For a monomial of power 1 we are multiplying C by one instance of our variable. When $x=0$ we get $c\*0=0$. When $x \ne 0$ we are multiplying c by 1 x. If x is less than 1 then c gets smaller, if x is more than 1 c gets bigger. When x is between 0 and -1 c gets smaller slower, when x is less than -1 c gets smaller faster.



A monomial with power two is one that "squares" the value of x. The reference to square is because using the multiplication operation once allows us to measure area. If you have something that is one unit on each side this is called a square unit. If you divide both sides of your square unit in half, you get 4 quarter units. We represent this with math by doing the multiplication $1/2 * 1/2 = 1/4$ Squaring something is a non-intuitive operation until you become comfortable with the graph of the function. We can see this with the story of the mathematician who was offered a reward by his king. The mathematician said he wanted a single grain of wheat, squared every day for 30 days. For the first seven days the king's servants delivered 1, 2, 4, 16, 256, 65, 536 grains of wheat to the mathematician. On the seventh day the value was 4,294,967,296 (4 gig in computer terms)\... Sometimes the story ends with the king re-negotiating, sometimes the story ends with the king executing the mathematician to preserve his kingdom, and sometimes the king is astute enough not to take the deal.



A monomial with power three is one that "cubes" the value of x. This is because we use the operation $x\*x\*x$ to measure the volume that a given area of x\*x takes up. If you have a cube that is 1 unit on each side and cut each side in half you will find that you have created 8 cubes. If the mathematician had asked to have the single grain of wheat cubed than the servants would have delivered $1$, $8$, $512$, $134 \times 10^6$, $242 \times 10^{22}$ grains of wheat and the kings deal would have needed to be re-negotiated two days earlier.



\paragraph{Power Function}



A power function is a function that can be represented in the form



$$f(x) = kx^p$$



where $k$ and $p$ are real numbers, and $k$ is known as the coefficient. We can also think of a power function as a monomial function; that is, a power function takes the form $y = f(x)$, where $f(x)$ is a single-variable monomial.



\paragraph{Resources}



\begin{itemize}
    \item \href{https://mathresearch.utsa.edu/wikiFiles/MAT1053/Intro_to_Polynomial_and_Power_Functions/MAT1053_M2.1Intro_to_Power_Functions_and_Polynomial_Functions.pdf}{Intro to Power Functions and Polynomial Functions}, Book Chapter
    \item \href{https://mathresearch.utsa.edu/wikiFiles/MAT1053/Intro_to_Polynomial_and_Power_Functions/MAT1053_M2.1Intro_to_Power_Functions_and_Polynomial_FunctionsGN.pdf}{Guided Notes}
\end{itemize}

\paragraph{Licensing}



Content obtained and/or adapted from:



\begin{itemize}
    \item \href{https://en.wikibooks.org/wiki/Algebra/Polynomials}{Polynomials, Wikibooks: Algebra} under a CC BY-SA license
\end{itemize}

\subsection{Learning Objectives}

\begin{enumerate}
    \item \textbf{Understanding Slope in Linear Functions:}
    \item \textbf{Outcome:} Students will be able to determine the slope of a linear function and understand the relationship between the slope and the steepness of the line.
    \item \textbf{Details:} This includes calculating the slope for given linear equations, interpreting the slope as a rate of change, and comparing the steepness of lines based on their slopes.
    \item \textbf{Identifying and Classifying Monomials:}
    \item \textbf{Outcome:} Students will be able to identify and classify monomials based on their degree and coefficient.
    \item \textbf{Details:} This involves recognizing monomials of varying degrees (e.g., degree 0, 1, 2, 3) and understanding the role of the constant and variable in the expression. Students will practice classifying examples and explaining how the degree affects the behavior of the function.
    \item \textbf{Exploring the Concept of Power Functions:}
    \item \textbf{Outcome:} Students will be able to describe and analyze power functions, particularly focusing on how changing the exponent and coefficient affects the graph and behavior of the function.
    \item \textbf{Details:} This includes understanding the definition of a power function, graphing various power functions (e.g., $f(x) = x^2$, $f(x) = x^3$), and analyzing the impact of different values of $k$ and $p$ on the shape and orientation of the graph.
\end{enumerate}

\subsection{Assessment Instrument}

\begin{enumerate}
    \item \textbf{Assessment Instrument: Intro to Power Functions Quiz} \textbf{Format:}
    \item \textbf{Type:} Mixed-format quiz (multiple-choice, short answer, and problem-solving questions)
    \item \textbf{Duration:} 45 minutes
    \item \textbf{Total Marks:} 30 points \textbf{Sections:}
    \item \textbf{Identifying Power Functions (10 points):}
    \item \textbf{Question 1:} Identify whether the following functions are power functions. If it is a power function, specify the values of $k$ and $p$:
    \item a. $f(x) = 3x^2 + 4$
    \item b. $g(x) = 7x^3$
    \item c. $h(x) = \dfrac{5}{x^2}$
    \item d. $j(x) = -4x^{1/2}$
    \item e. $k(x) = 2x + 1$
    \item \textbf{Graphing Power Functions (10 points):}
    \item \textbf{Question 2:} Sketch the graphs of the following power functions and describe their shape:
    \item a. $f(x) = x^2$
    \item b. $g(x) = -x^3$
    \item c. $h(x) = \dfrac{1}{x}$
    \item d. $j(x) = x^{1/3}$
    \item e. $k(x) = 2x^4$
    \item \textbf{Analyzing the Effect of Coefficients (10 points):}
    \item \textbf{Question 3:} For each of the following power functions, determine how changing the coefficient $k$ affects the graph. Provide a brief explanation:
    \item a. $f(x) = 3x^2$ vs. $g(x) = 6x^2$
    \item b. $h(x) = -2x^3$ vs. $j(x) = -5x^3$
    \item c. $k(x) = \dfrac{1}{2}x^4$ vs. $m(x) = 2x^4$
    \item d. $n(x) = x^{1/2}$ vs. $p(x) = 4x^{1/2}$ \textbf{Scoring Rubric:}
    \item \textbf{Identifying Power Functions:} 2 points for each correctly identified function and correct identification of $k$ and $p$ (total 10 points)
    \item \textbf{Graphing Power Functions:} 2 points for each correctly sketched graph and description of shape (total 10 points)
    \item \textbf{Analyzing the Effect of Coefficients:} 2.5 points for each correctly analyzed effect and explanation (total 10 points) This quiz will assess students' understanding of power functions, including identification, graphing, and analysis of how changes in coefficients impact the graph. It ensures that they can recognize and work with power functions effectively.
\end{enumerate}

%%===============================================================
\section{Intro to Polynomial Functions}
\label{sec:intro-to-polynomial-functions}
%%===============================================================

With practice the concept of slope for linear functions becomes intuitive. It makes sense that the line that fits the equation $y=2x$ has a steeper ascent then the line that fits the equation $y=\frac{1}{2}x$. You only have to move horizontally one unit to change your vertical direction two for the former when you graph $y=2x$. How many blocks do you need to move horizontally to change your vertical direction by one for the line $y=\frac{1}{2}x$?



When we express concepts like $y=x^2$ the abstract behavior of what is being represented becomes a little harder to see.



A \textbf{monomial} (or \textbf{power function}) of one variable, let's say x, is an algabraic expression of the form $c x^m $



where



\begin{itemize}
    \item $c$ is a constant, and
    \item $m$ is a non-negative integer (e.g., 0, 1, 2, 3, \...).
\end{itemize}

The integer $m$ is called the \textbf{degree} of the monomial.



The idea of a monomial of degree zero appears a bit mystical since it always represents one, except when the value of the variable is set equal to zero when the result is undefined. This idea allows us preserve the value of the constant in the monomial. We know that $cx^0$ is always equal to $c$ because even though we have 0 x's (somethings) we still have a c. When x = 0 things are difficult because the value we started with, 0, represents nothing.



For a monomial of power 1 we are multiplying C by one instance of our variable. When $x=0$ we get $c*0=0$. When $x \ne 0$ we are multiplying c by 1 x. If x is less than 1 then c gets smaller, if x is more than 1 c gets bigger. When x is between 0 and -1 c gets smaller slower, when x is less than -1 c gets smaller faster.



A monomial with power two is one that "squares" the value of x. The reference to square is because using the multiplication operation once allows us to measure area. If you have something that is one unit on each side this is called a square unit. If you divide both sides of your square unit in half, you get 4 quarter units. We represent this with math by doing the multiplication $1/2 * 1/2 = 1/4$ Squaring something is a non-intuitive operation until you become comfortable with the graph of the function. We can see this with the story of the mathematician who was offered a reward by his king. The mathematician said he wanted a single grain of wheat, squared every day for 30 days. For the first seven days the king's servants delivered 1, 2, 4, 16, 256, 65, 536 grains of wheat to the mathematician. On the seventh day the value was 4,294,967,296 (4 gig in computer terms)\... Sometimes the story ends with the king re-negotiating, sometimes the story ends with the king executing the mathematician to preserve his kingdom, and sometimes the king is astute enough not to take the deal.



A monomial with power three is one that "cubes" the value of x. This is because we use the operation x\*x\*x to measure the volume that a given area of x\*x takes up. If you have a cube that is 1 unit on each side and cut each side in half you will find that you have created 8 cubes. If the mathematician had asked to have the single grain of wheat cubed than the servants would have delivered $1$, $8$, $512$, $134 \times 10^6$, $242 \times 10^{22}$ grains of wheat and the kings deal would have needed to be re-negotiated two days earlier.



\paragraph{Polynomials}



A \textbf{polynomial} of one variable, x, is an algebraic expression that is a sum of one or more monomials. The \textbf{degree} of the \textbf{polynomial} is the highest degree of the monomials in the sum. An polynomial $P(x)$ can generically be expressed in the form



$$P(x) = a _n x^n + ... + a _i x^i + ... a _2 x^2 + a _1 x + a _0 $$



or



$$P(x) = \sum _{i=0}^{n}a _ix^i$$



The constants $a _i$ are called the \textbf{coefficients} of the polynomial.



Each of the individual monomials in the above sum, whose coefficient $a _i \neq 0$, is called a \textbf{term} of the polynomial. When i = 0, $x^i = 1$ and the corresponding term simply equals the constant $a _i$. Also when i = 1, the corresponding term equals $a _i x$.



A polynomial having two terms is called a \textbf{binomial}. A polynomial having three terms is called a \textbf{trinomial}.



\paragraph{Polynomial Equations}



We refer to all functions with one independent variable as $P(x)$. Each instance of $P(x)$ can be represented by an equation (either a monomial or a polynomial) which may have one or more places where the dependent variable is equal to zero. These places are called roots and they represent the number(s) whose value(s) for x make the function $P(x)=0$ true. These roots are called the \textbf{zeroes} of the polynomial (singular is \textbf{zero}). A polynomial of degree 1, will always look like a line when you graph it, and always has 1 real zero. A polynomial of degree 2, a quadratic function, can have 0, 1, or 2 real zeroes. A polynomial of degree 3 (a cubic function) can have 1 or 3 real zeroes. A polynomial of degree 4 can have 0, 2, or 4 real zeroes. Complex (unreal) zeroes, when present, always come in pairs. In general, a polynomial of degree n, where n is odd, can have from 1 to n real zeroes. A polynomial of degree n, where n is even, can have from 0 to n real zeroes.



When we graph polynomials each zero is a place where the polynomial crosses the x axis. A polynomial of degree one can be generically written as $P(x)=Mx + C$ where M and C can be any real number. We will see that quadratic functions are curves. The curve can bend before it ever touches the X axis in which case it has no zeroes, It can bend just as it touches the X axis, in which case it can have just one zero, or it can open up above or below the X axis in which case it will have two zeroes. If you think about this you will see that polynomials with an odd degree (1,3,5, \...) have to be positive and negative, so they have to cross the X axis at least once. Polynomials with an even degree (2,4,6,\....) might always be positive or negative and never have a zero.



Normally we represent a function in the form $P(x) = y$, but when we are looking for the roots of the function we want y to be equal to zero so we solve for the equation of $P(x)$ where $P(x) = 0$



$$\begin{array}{|c|c|c|c|} \hline \textbf{Order} & \textbf{Name} & \textbf{Number of bumps} & \textbf{Where found} \\\\ \hline 1 & \text{linear} & \text{no bumps - straight line} & \text{straight line equations} \\\\ \hline 2 & \text{quadratic} & \text{one bump} & \text{equations involving area and vibrations} \\\\ \hline 3 & \text{cubic} & \text{two bumps} & \text{equations involving volumes} \\\\ \hline 4 & \text{quartic} & \text{three bumps} & \text{some physics equations (melting ice)} \\\\ \hline n (5+) & & \text{n-1 bumps} & \text{very rare} \\\\ \hline \end{array}$$



\paragraph{Solving Polynomial Equations}



Some polynomial equations can be solved by factoring, and all equations of degrees 1-4 can be solved completely by formulae. Above degree 4, there are no formulae for solving completely, and you must rely on numerical analysis or factoring. This means that for polynomials of degree greater than 4 it is often impossible to find exact solutions.



\paragraph{Rational roots of polynomial equations}



Often we are interested in the rational roots of polynomials. A root is much like a factor of a number. For instance all even numbers have a factor of two. This means you can write the even numbers as two times another number. That is the numbers 2, 4, 6, 8 \... can be written as 2\*1, 2\*2, 2\*3, 2\*4 \... . This fact is helpful when you have a fraction of two even numbers. Given a fraction of two even numbers called N and M, N/M, you could reduce the fraction by re-writing it as $2\*n / 2\*m$. By keeping fractions in lowest terms it's easier to know when you can add or subtract them without looking for a common denominator.



\paragraph{An example of a use of a polynomial equation}



There is a story that in grade school the mathematician Gauss was asked to add the numbers 1 to 100 sequentially. He is said to have intuited the sum could be expressed with the formula n(n+1)/2 and quickly gave the answer 5050. The basis of this formula is that the numbers 1 through 49 added to the numbers 99 through 51 each yield 100. It is interesting to look at how this formula works for the values 9 and 10. For 10 we add the numbers 1+9, 2+ 8, 3+ 7, 4+ 6 to get 40 and we add the two remaining terms 5 and 10 to get 55. For 9 we add the terms 1 + 8, 2 + 7, 3 + 6, 4+ 5 to get 4\*9 = 36 + 9 = 45. In the first case the n + 1 is the odd number and represents adding the 10 and the middle number, the 5. In the second case the n is the odd number and the n+1 represents the sum for the preceding terms in the formula. You may or may not find stories like this intriguing based on how your personality reacts to what is known as the foundational crisis of mathematics. Learning mathematics is a lot like learning a foreign language. Some people seem more adept at learning languages than others, but with hard work learning a new language is something we can all do.



\paragraph{Multiplying polynomials together}



When we multiply polynomials together we rely heavily on the distributive property.



For instance when we multiply 67 by 5 we can divide the equation into (60 + 7)\*5 = (300 + 35) = 335. Additionally we can apply the commutative property to multiply multidigit numbers. 67\*25 = (60 + 7)(20 + 5) = ((60 + 7)\*20) + ((60 + 7) \*5) = (60\*20) + (7\*20) + (60\*5) + (7\*5) = 1200 + 140 + 300 + 35 = 1675. These properties are the foundation for the different forms of the mechanical calculating tool the \href{w:abacus "wikilink"}{abacus}.



When multiplying polynomials together we do similar operations. We use the commutative property to divide the multiplier into its component parts and multiply the multiplicant by each of these parts. For instance to multiply $x^2 + x$ by $x  + 1$ we first write the multiplicand and multiplier in terms of powers of x. This gives us $x^2 + x + 0x^0$ and $x + 1x^0$ The terms raised to the zero power represent constant integer terms in our equations. Next we apply the commutative property to rewrite the equations as $[(x^2 + x + 0x^0)\*1x^1] + [(x^2 + x + 0x^0)\*1x^0]$. We simplify these equations to be $[x^3 + x^2 + 0x^1] + [x^2 + x + 0x^0]$ (notice how our integer term drops out). Finally we combine like terms to get the answer $x^3 + 2x^2 + x +0x^0$. Let's repeat that in the more familiar columnar format of multiplication:



$$\begin{aligned} & &  & &1x^2&+&1x^1&+&0x^0 \\\\ &*&  & & & &1x^1&+&1x^0 \\\\ \hline & &  & &1x^2&+&1x^1&+&0x^0 \\\\ &+&1x^3 &+&1x^2&+&0x^1& & \\\\ \hline &=&1x^3 &+&2x^2&+&1x^1&+&0x^0 & = x^3 + 2x^2 + x \end{aligned}$$



By breaking a polynomial into its r



If we have a polynomial P(x)



$$P(x) = a _n x^n + ... + a _i x^i + ... a _2 x^2 + a _1 x + a _0 $$



The only possible rational roots (roots of the form p/q) are in the form



$$\frac{p}{q}=\frac{\text{factor}\ \text{of}\ a _0}{\text{factor}\ \text{ of }\ a _n}$$



\paragraph{Binomials}



A binomial is a sum or difference of two monomials. These can also be called polynomials, but to specify, these are binomials.



\paragraph{Examples}



$2x + 2$



$2y - 7$



\paragraph{How to factor}



To factor binomials, find the greatest common factor between the terms and factor.



\paragraph{Example}



$4x + 2$



The greatest common factor between these terms is 2 because both of the terms can be divided by it and the coefficient and constant is still an integer. The example factored would become:



$2(2x+1)$



\paragraph{Symmetrical Polynomials}



In mathematics a polynomial is considered to be symmetrical if you take the roots of the original polynomial and then interchange any root with another root, the polynomial will remain the same. For example the polynomial $8x^3 + 16x^2 - 22x - 30$ is symmetrical because its factorized form is $(2x - 3)(2x + 2)(2x + 5)$ and if you interchange the roots the resulting polynomial will be the same. However the polynomial $4x^3 + 12x^2 - 7x - 30$ is not symmetrical because it factorized form is $(2x - 3)(x + 2)(2x + 5)$ and if you interchange the roots the resulting polynomial will be $4x^3 + 18x^2 - 16x - 30$ if you switch the 2 and the 5 around.



\paragraph{Roots Of Quadratic Polynomial}



If we need to find the roots of a given quadratic function we have two formulae that can help us to find the roots of a quadratic equation.









Let $\alpha\,$ and $\beta\,$ be the roots of $ax^2+bx+c=0\,$. Then, $\alpha + \beta = - \frac{b}{a},\quad \alpha\beta = \frac{c}{a}$









\paragraph{Example}



Find the values of a and b of the equation $ax^2 + bx - 48\,$ if $\alpha + \beta = 6\,$ and $\alpha\beta = -16\,$.



1.  First we need to find the value of a and b, we use the relationships of the roots to find a and b.

    1.  $\alpha\beta = -16 = \frac{-48}{a}\,$ from this we can determine that a = 3

    2.  $\alpha + \beta = 6 = -\frac{b}{a}\,$

2.  Now that we have determined that a = 3 we can write the second relationship as:

    $$\alpha + \beta = 6 = -\frac{b}{3}\,$$ so we can determine that b = -18

3.  Now we can write the complete equation.

    $$3x^2 -18 x - 48\,$$



\paragraph{Roots Of Cubic Equations}



If we need to find the roots of a given cubic function we have three formulae that can help us to find the roots of a cubic equation.



Let $\alpha, \beta\,$ and $\gamma\,$ be the roots of $ax^3+bx^2+cx+d=0\,$. Then,

$$\sum\alpha = - \frac{b}{a},\quad \sum\alpha\beta = \frac{c}{a},\quad \alpha\beta\gamma = -\frac{d}{a}$$



Where: $\sum\alpha = \alpha + \beta + \gamma$



And: $\sum\alpha\beta = \alpha\beta + \alpha\gamma + \beta\gamma$



\paragraph{Example}



In this example we consider the special case of the cubic $x^3 + 21x^2 + cx + 280=0$, where c is to be determined and we are given the additional information that its 3 roots are in \href{A-level_Mathematics/C2/Sequences_and_Series#Arithmetic_Progression_.28AP.29 "wikilink"}{arithmetic progression}. Thus we can write the roots in the form p, p + q, p - q. Also factorize the equation.



1.  First to find p we use the $\sum\alpha$.

$$\begin{aligned} \sum\alpha & = p + (p + q) + (p - q) & = -\frac{21}{1} \\\\ \sum\alpha & = 3p & = -21 \\\\ p & = -7 \end{aligned}$$



2.  Then we need to find the value of q.

$$\begin{aligned} -7(-7+q)(-7-q) & = -\frac{280}{1} \\\\  7q^2 - 343 & = -280 \\\\  7q^2 & = 63 \\\\  q^2 & = 9 \\\\  q & = 3 \end{aligned}$$

3.  Now we can write out our roots.



> (-7 - 3),-7,(-7 + 3)

> -10,-7,-4



4.  We can now find c.

    $$\sum\alpha\beta = -7\times-4 + -7\times-10 + -10\times-4 = \frac{c}{1}$$



    $$138 = c\,$$

5.  The complete equation is $x^3 + 21x^2 + 138x + 280=0$

6.  Finally we write out the factorized equation.

7.  $(x + 7)(x+4)(x+10)=0\,$



\paragraph{Simple Substitution of Roots}



If you increase each root in a polynomial equation by the number n, you can calculate the resulting equation by replacing each x term in the original polynomial equation with (x - n). This leads to binomial expansion so make sure that you are well versed in it.



\paragraph{Example}



Suppose that the cubic equation $x^3+x^2-22x-40 = 0\,$ has roots $\alpha\,, \beta\,$ and $\gamma\,$. Find a cubic equation with the roots $\alpha - 2\,, \beta- 2$ and $\gamma -2\,$



1.  If $x=\alpha-2$ then $\alpha=x+2$. Since $\alpha$ is a root of the original equation you can replace each x term with x + 2:

    $$\left(x+2\right)^3+\left(x+2\right)^2-22\left(x+2\right)-40 = 0$$

2.  Using Binomial Expansion we can easily find the terms.

    $$\left(x^3+6x^2 + 12x + 8\right)+\left(x^2+4x+4\right)-22\left(x+2\right)-40 = 0$$

3.  Finally we combine all the terms and we have:



$$x^3 + 7x^2 -6x -72\,$$



\paragraph{Resources}



\begin{itemize}
    \item \href{https://en.wikibooks.org/wiki/Algebra/Polynomials}{Polynomials}, Wikibooks: Algebra
    \item \href{https://en.wikibooks.org/wiki/A-level_Mathematics/OCR/FP1/Roots_of_Polynomial_Equations}{Roots of Polynomial Equations}, Wikibooks: A-level Mathematics
    \item \href{https://mathresearch.utsa.edu/wikiFiles/MAT1053/Intro_to_Polynomial_and_Power_Functions/MAT1053_M2.1Intro_to_Power_Functions_and_Polynomial_Functions.pdf}{Intro to Power Functions and Polynomial Functions}, Book Chapter
    \item \href{https://mathresearch.utsa.edu/wikiFiles/MAT1053/Intro_to_Polynomial_and_Power_Functions/MAT1053_M2.1Intro_to_Power_Functions_and_Polynomial_FunctionsGN.pdf}{Guided Notes}
    \item \href{https://www.youtube.com/watch?v=nJLzFBsSX3o}{Multiplying Polynomials}. Produced by TA Catherine Sporer, UTSA
\end{itemize}

\paragraph{Licensing}



Content obtained and/or adapted from:



\begin{itemize}
    \item \href{https://en.wikibooks.org/wiki/Algebra/Polynomials}{Polynomials, Wikibooks: Algebra} under a CC BY-SA license
    \item \href{https://en.wikibooks.org/wiki/A-level_Mathematics/OCR/FP1/Roots_of_Polynomial_Equations}{Roots of Polynomial Equations, Wikibooks: A-level Mathematics} under a CC BY-SA license
\end{itemize}

\subsection{Learning Objectives}

\begin{enumerate}
    \item Here are three Student Learning Objectives (SLOs) based on the provided concepts:
    \item \textbf{Understand and Calculate Slope for Linear Functions:}
    \item \textbf{Outcome:} Students will be able to calculate and interpret the slope of linear functions, including identifying the steepness of a line given its equation.
    \item \textbf{Details:} This includes practice with equations of the form \(y = mx\), where \(m\) represents the slope. Students will be able to compare the steepness of lines with different slopes and understand how the slope affects the graph's direction.
    \item \textbf{Explore and Represent Monomials:}
    \item \textbf{Outcome:} Students will understand the concept of monomials, including their structure and behavior for different degrees.
    \item \textbf{Details:} This includes learning how to represent monomials of various degrees (e.g., \(cx^m\)) and exploring the special cases, such as monomials of degree zero, where \(m=0\) and \(cx^0 = c\), as well as understanding the implications when \(x=0\).
    \item \textbf{Graph and Analyze Polynomial Functions:}
    \item \textbf{Outcome:} Students will be able to graph polynomial functions and analyze the number and nature of their roots.
    \item \textbf{Details:} This includes understanding how the degree of a polynomial influences its graph, the possible number of real roots, and how these roots correspond to the points where the polynomial crosses the x-axis. Students will practice graphing polynomials and identifying roots, bumps, and other key features of the graph based on the polynomial's degree.
\end{enumerate}

\subsection{Assessment Instrument}

\begin{enumerate}
    \item \textbf{Assessment Instrument: Intro to Polynomial Functions Quiz} \textbf{Format:}
    \item \textbf{Type:} Mixed-format quiz (multiple-choice, short answer, and problem-solving questions)
    \item \textbf{Duration:} 45 minutes
    \item \textbf{Total Marks:} 30 points \textbf{Sections:}
    \item \textbf{Understanding Slope and Linear Functions (10 points):}
    \item \textbf{Question 1:} For the linear function $ y = \frac{1}{2}x $, how many units must you move horizontally to change the vertical direction by one unit?
    \item a. 0.5 units
    \item b. 1 unit
    \item c. 2 units
    \item d. 4 units
    \item \textbf{Question 2:} Given the function $ y = 3x $, what is the slope, and how does it compare to the slope of the function $ y = 1.5x $?
    \item \textbf{(3 points, short answer)}
    \item \textbf{Monomials and Their Behavior (10 points):}
    \item \textbf{Question 3:} Define a monomial and give an example of a monomial with degree zero. Explain why this monomial always equals the constant term.
    \item \textbf{(4 points, short answer)}
    \item \textbf{Question 4:} Evaluate the monomial $ 5x^2 $ when $ x = 3 $ and explain how the value changes when $ x = \frac{1}{2} $.
    \item \textbf{(6 points, problem-solving)}
    \item \textbf{Graphing and Analyzing Polynomial Functions (10 points):}
    \item \textbf{Question 5:} Consider the quadratic polynomial $ P(x) = x^2 - 4x + 4 $. How many real roots does this polynomial have, and what are they?
    \item a. No real roots
    \item b. One real root
    \item c. Two real roots
    \item d. Three real roots
    \item \textbf{Question 6:} Sketch the graph of the cubic polynomial $ P(x) = x^3 - 3x^2 + 2x $. Identify the points where the graph crosses the x-axis and determine the number of real roots.
    \item \textbf{(5 points, problem-solving)} \textbf{Scoring Rubric:}
    \item \textbf{Understanding Slope and Linear Functions:} 2 points for correct multiple-choice answer, 3 points for correct explanation of the slope comparison (total 5 points)
    \item \textbf{Monomials and Their Behavior:} 4 points for correct definition and example, 6 points for correct evaluation and explanation (total 10 points)
    \item \textbf{Graphing and Analyzing Polynomial Functions:} 2 points for correct identification of real roots, 3 points for correctly sketched graph and identification of roots (total 5 points) This quiz assesses the students' grasp of the concepts related to linear functions, monomials, and polynomial functions. It ensures they understand the slope, can work with monomials of various degrees, and can analyze and graph polynomial functions, focusing on the identification of roots and their implications.
\end{enumerate}

%%===============================================================
\section{Quadratic Functions}
\label{sec:quadratic-functions}
%%===============================================================

In algebra, a \textbf{quadratic function}, a \textbf{quadratic polynomial}, a \textbf{polynomial of degree 2}, or simply a \textbf{quadratic}, is a polynomial function with one or more variables in which the highest-degree term is of the second degree.



% [Image: A quadratic polynomial with two real roots (crossings of the x axis) and hence no complex roots. Some other quadratic polynomials have their minimum above the x axis, in which case there are no real roots and two complex roots.]



Figure: A quadratic polynomial with two real roots (crossings of the \textit{x} axis) and hence no complex roots. Some other quadratic polynomials have their minimum above the \textit{x} axis, in which case there are no real roots and two complex roots.



For example, a \textit{univariate} (single-variable) quadratic function has the form



$$f(x)=ax^2+bx+c,\quad a \ne 0$$

in the single variable \textit{x}. The graph of a univariate quadratic function is a parabola whose axis of symmetry is parallel to the \textit{y}-axis, as shown at right.



If the quadratic function is set equal to zero, then the result is a quadratic equation. The solutions to the univariate equation are called the roots of the univariate function.



The bivariate case in terms of variables \textit{x} and \textit{y} has the form



$$f(x,y) = a x^2 + by^2 + cx y+ d x+ ey + f \,$$



with at least one of \textit{a, b, c} not equal to zero, and an equation setting this function equal to zero gives rise to a conic section (a circle or other ellipse, a parabola, or a hyperbola).



A quadratic function in three variables \textit{x}, \textit{y,} and \textit{z} contains exclusively terms $x^2$, $y^2$, $z^2$, \textit{xy}, \textit{xz}, \textit{yz}, \textit{x}, \textit{y}, \textit{z}, and a constant:



$$f(x,y,z)=ax^2+by^2+cz^2+dxy+exz+fyz+gx+hy+iz +j,$$



with at least one of the coefficients \textit{a, b, c, d, e,} or \textit{f} of the second-degree terms being non-zero.



In general there can be an arbitrarily large number of variables, in which case the resulting surface of setting a quadratic function to zero is called a quadric, but the highest degree term must be of degree 2, such as $x^2$, \textit{xy}, \textit{yz}, etc.



\paragraph{Etymology}



The adjective \textit{quadratic} comes from the Latin word \textit{quadratum} ("square"). A term like $x^2$ is called a square in algebra because it is the area of a \textit{square} with side $x$.



\paragraph{Terminology}



\paragraph{Coefficients}



The coefficients of a polynomial are often taken to be real or complex numbers, but in fact, a polynomial may be defined over any ring.



\paragraph{Degree}



When using the term "quadratic polynomial", authors sometimes mean "having degree exactly 2", and sometimes "having degree at most 2". If the degree is less than 2, this may be called a "degenerate case". Usually the context will establish which of the two is meant.



Sometimes the word "order" is used with the meaning of "degree", e.g. a second-order polynomial.



\paragraph{Variables}



A quadratic polynomial may involve a single variable \textit{x} (the univariate case), or multiple variables such as \textit{x}, \textit{y}, and \textit{z} (the multivariate case).



\paragraph{The one-variable case}



Any single-variable quadratic polynomial may be written as



$$ax^2 + bx + c,\,$$

where \textit{x} is the variable, and \textit{a}, \textit{b}, and \textit{c} represent the coefficients. In elementary algebra, such polynomials often arise in the form of a quadratic equation $ax^2 + bx + c = 0$. The solutions to this equation are called the roots of the quadratic polynomial, and may be found through factorization, completing the square, graphing, Newton's method, or through the use of the quadratic formula. Each quadratic polynomial has an associated quadratic function, whose graph is a parabola.



\paragraph{Bivariate case}



Any quadratic polynomial with two variables may be written as



$$f(x,y) = a x^2 + b y^2 + cxy + dx+ e y + f, \,$$ where \textit{x} and \textit{y} are the variables and \textit{a}, \textit{b}, \textit{c}, \textit{d}, \textit{e}, and \textit{f} are the coefficients. Such polynomials are fundamental to the study of conic sections, which are characterized by equating the expression for \textit{f} (\textit{x}, \textit{y}) to zero. Similarly, quadratic polynomials with three or more variables correspond to quadric surfaces and hypersurfaces. In linear algebra, quadratic polynomials can be generalized to the notion of a quadratic form on a vector space.



\paragraph{Forms of a univariate quadratic function}



A univariate quadratic function can be expressed in three formats:



\begin{itemize}
    \item $f(x) = a x^2 + b x + c $ is called the \textbf{standard form},
    \item $f(x) = a(x - r _1)(x - r _2)$ is called the \textbf{factored form}, where $r _1$ and $r _2$ are the roots of the quadratic function and the solutions of the corresponding quadratic equation.
    \item $f(x) = a(x - h)^2 + k \,$ is called the \textbf{vertex form}, where $h$ and $k$ are the $x$ and $y$ coordinates of the vertex, respectively.
\end{itemize}

The coefficient \textit{a} is the same value in all three forms. To convert the \textbf{standard form} to \textbf{factored form}, one needs only the quadratic formula to determine the two roots $r _1$ and $r _2$. To convert the \textbf{standard form} to \textbf{vertex form}, one needs a process called completing the square. To convert the factored form (or vertex form) to standard form, one needs to multiply, expand and/or distribute the factors.



\paragraph{Graph of the univariate function}



% [Image: f(x) = ax^2]



Figure: $f(x) = ax^2 | _ {a=\\{0.1,0.3,1,3\\}} $



% [Image: f(x) = x^2 + bx]



Figure: $f(x) = x^2 + bx | _{b=\\{1,2,3,4\\}} $



% [Image: f(x) = x^2 + bx]



Figure: $f(x) = x^2 + bx | _{b=\\{-1,-2,-3,-4\\}} $



Regardless of the format, the graph of a univariate quadratic function $f(x) = ax^2 + bx + c$ is a parabola (as shown at the right). Equivalently, this is the graph of the bivariate quadratic equation $y = ax^2 + bx + c$.



\begin{itemize}
    \item If $a > 0$, the parabola opens upwards.
    \item If $a < 0$, the parabola opens downwards.
\end{itemize}

The coefficient $a$ controls the degree of curvature of the graph; a larger magnitude of $a$ gives the graph a more closed (sharply curved) appearance.



The coefficients $b$ and $a$ together control the location of the axis of symmetry of the parabola (also the `{{math|''x''}}`{=mediawiki}-coordinate of the vertex and the \textit{h} parameter in the vertex form) which is at



$$x = -\frac{b}{2a}.$$



The coefficient $c$ controls the height of the parabola; more specifically, it is the height of the parabola where it intercepts the \textit{y}-axis.



\paragraph{Vertex}



The \textbf{vertex} of a parabola is the place where it turns; hence, it is also called the \textbf{turning point}. If the quadratic function is in vertex form, the vertex is $(h, k)$. Using the method of completing the square, one can turn the standard form



$$f(x) = a x^2 + b x + c \,$$ into



$$\begin{aligned} f(x) &= ax^2 + bx + c \\\\ &= a(x - h)^2 + k \\\\ &= a\left(x - \frac{-b}{2a}\right)^2 + \left(c - \frac{b^2}{4a}\right) \end{aligned}$$





so the vertex, $(h, k)$, of the parabola in standard form is



$$\left(-\frac{b}{2a}, c - \frac{b^2}{4a}\right)$$



If the quadratic function is in factored form



$$f(x) = a(x - r _1)(x - r _2) \,$$ the average of the two roots, i.e.,



$$\frac{r _1 + r _2}{2} \,$$



is the \textit{x}-coordinate of the vertex, and hence the vertex $(h, k)$ is



$$\left(\frac{r _1 + r _2}{2}, f\left(\frac{r _1 + r _2}{2}\right)\right)$$



The vertex is also the maximum point if $a0$.



The vertical line



$$x=h=-\frac{b}{2a}$$



that passes through the vertex is also the \textbf{axis of symmetry} of the parabola.



\paragraph{Maximum and minimum points}



Using calculus, the vertex point, being a maximum or minimum of the function, can be obtained by finding the roots of the derivative:



$$f(x)=ax^2+bx+c \quad \Rightarrow \quad f'(x)=2ax+b \,$$



$x$ is a root of $f'(x)$ if $f'(x)=0$ resulting in



$$x=-\frac{b}{2a}$$ with the corresponding function value



$$f(x) = a \left (-\frac{b}{2a} \right)^2+b \left (-\frac{b}{2a} \right)+c = c-\frac{b^2}{4a} \,$$

So, again the vertex point coordinates, $(h, k)$, can be expressed as



$$\left (-\frac {b}{2a}, c-\frac {b^2}{4a} \right)$$



\paragraph{Roots of the univariate function}



% [Image: Graph of y = ax^2 + bx + c]



Figure: Graph of $y = ax^2 + bx + c$, where $a$ and the discriminant $b^2 - 4ac$ are positive, with

\begin{itemize}
    \item Roots and y-intercept in $\textcolor{red}{\text{red}}$
    \item Vertex and axis of symmetry in $\textcolor{blue}{\text{blue}}$
    \item Focus and directrix in $\textcolor{\text{pink}}$
\end{itemize}

% [Image: Quadratic function graph complex roots]



Figure: Visualisation of the complex roots of $y = ax^2 + bx + c$: the parabola is rotated $180^{\circ}$ about its vertex ($\textcolor{orange}{\text{orange}}$). Its x-intercepts are rotated $90^{\circ}$ around their mid-point, and the Cartesian plane is interpreted as the complex plane ($\textcolor{green}{\text{green}}$).



\paragraph{Exact roots}



The roots (or \textit{zeros}), $r _1$ and $r _2$, of the univariate quadratic function



$$ \begin{aligned} f(x) & = ax^2 + bx + c \\\\ & = a(x - r _1)(x - r _2), \end{aligned} $$



are the values of $x$ for which $f(x) = 0$.



When the coefficients $a$, $b$, and $c$, are real or complex, the roots are



$$r _1=\frac{-b - \sqrt{b^2-4ac}}{2a},$$



$$r _2=\frac{-b + \sqrt{b^2-4ac}}{2a}.$$



\paragraph{Upper bound on the magnitude of the roots}



The modulus of the roots of a quadratic $ax^2+bx+c\,$ can be no greater than $$\frac{\max(|a|, |b|, |c|)}{|a|}\times \phi \quad \text{where } \phi=\frac{1+\sqrt{5}}{2} = \text{ the golden ratio}$$



\paragraph{The square root of a univariate quadratic function}



The square root of a univariate quadratic function gives rise to one of the four conic sections, almost always either to an ellipse or to a hyperbola.



If $a>0\,$ then the equation $y = \pm \sqrt{a x^2 + b x + c}$ describes a hyperbola, as can be seen by squaring both sides. The directions of the axes of the hyperbola are determined by the ordinate of the minimum point of the corresponding parabola $y _p = a x^2 + b x + c \,$. If the ordinate is negative, then the hyperbola's major axis (through its vertices) is horizontal, while if the ordinate is positive then the hyperbola's major axis is vertical.



If $a0 \,$ the function has a minimum if $A>0$, and a maximum if $A0$ and a maximum if $A<0$; its graph forms a parabolic cylinder.



\paragraph{Resources}



\begin{itemize}
    \item \href{https://mathresearch.utsa.edu/wikiFiles/MAT1053/Quadratic_Functions/MAT1053_M2.2Quadratic_Functions.pdf}{Quadratic Functions}, Book Chapters
    \item \href{https://mathresearch.utsa.edu/wikiFiles/MAT1053/Quadratic_Functions/MAT1053_M2.2Quadratic_FunctionsGN.pdf}{Guided Notes}
\end{itemize}

\paragraph{Licensing}



Content obtained and/or adapted from:



\begin{itemize}
    \item \href{https://en.wikipedia.org/wiki/Quadratic_function}{Quadratic function, Wikipedia} under a CC BY-SA license
\end{itemize}

\subsection{Learning Objectives}

\begin{enumerate}
    \item \textbf{Identify the Standard Form of a Quadratic Function:}
    \item \textbf{Outcome:} Students will be able to recognize and write a quadratic function in its standard form, $ f(x) = ax^2 + bx + c $, where $ a \neq 0 $.
    \item \textbf{Details:} This includes understanding the significance of each coefficient in the equation and its effect on the graph of the function.
    \item \textbf{Solve Quadratic Equations Using the Quadratic Formula:}
    \item \textbf{Outcome:} Students will be able to solve quadratic equations using the quadratic formula $ x = \frac{-b \pm \sqrt{b^2 - 4ac}}{2a} $.
    \item \textbf{Details:} This includes identifying the discriminant, understanding its significance in determining the nature of the roots (real vs. complex), and applying the formula correctly to find the roots.
    \item \textbf{Graph Quadratic Functions and Identify Key Features:}
    \item \textbf{Outcome:} Students will be able to graph a quadratic function and identify its key features, including the vertex, axis of symmetry, and roots.
    \item \textbf{Details:} This includes determining whether the parabola opens upwards or downwards based on the sign of $ a $, finding the vertex using $ x = -\frac{b}{2a} $, and using this information to accurately sketch the graph.
\end{enumerate}

\subsection{Assessment Instrument}

\begin{enumerate}
    \item \textbf{Assessment Instrument: Quadratic Functions Quiz} \textbf{Format:}
    \item \textbf{Type:} Mixed-format quiz (multiple-choice, short answer, and problem-solving questions)
    \item \textbf{Duration:} 45 minutes
    \item \textbf{Total Marks:} 30 points \textbf{Sections:}
    \item \textbf{Identifying the Standard Form (10 points):}
    \item \textbf{Question 1:} Identify the standard form of the following quadratic functions and specify the values of $a$, $b$, and $c$:
    \item a. $ f(x) = 3x^2 - 5x + 2 $
    \item b. $ f(x) = -2x^2 + 4x - 1 $
    \item c. $ f(x) = x^2 + 7x $
    \item d. $ f(x) = 4x^2 - 8 $
    \item e. $ f(x) = -x^2 + 9 $
    \item \textbf{Solving Quadratic Equations (10 points):}
    \item \textbf{Question 2:} Solve the following quadratic equations using the quadratic formula:
    \item a. $ x^2 - 4x - 5 = 0 $
    \item b. $ 2x^2 + 3x - 2 = 0 $
    \item c. $ x^2 + 6x + 9 = 0 $
    \item d. $ 3x^2 - 7x + 2 = 0 $
    \item e. $ 4x^2 - 12x + 9 = 0 $
    \item \textbf{Graphing and Analyzing Quadratic Functions (10 points):}
    \item \textbf{Question 3:} For the following quadratic functions, graph the function, identify the vertex, axis of symmetry, and roots (if any):
    \item a. $ f(x) = x^2 - 4x + 3 $
    \item b. $ f(x) = -2x^2 + 4x - 1 $
    \item c. $ f(x) = 3x^2 + 6x + 2 $
    \item d. $ f(x) = x^2 + 2x + 1 $
    \item e. $ f(x) = -x^2 + 5x - 6 $ \textbf{Scoring Rubric:}
    \item \textbf{Identifying the Standard Form:} 2 points for each correct identification and specification of $a$, $b$, and $c$ (total 10 points)
    \item \textbf{Solving Quadratic Equations:} 2 points for each correctly solved equation (total 10 points)
    \item \textbf{Graphing and Analyzing Quadratic Functions:} 2 points for each correctly graphed function with accurate identification of vertex, axis of symmetry, and roots (total 10 points) This quiz will assess students' understanding of quadratic functions, including recognizing the standard form, solving quadratic equations, and analyzing and graphing quadratic functions. It ensures that students can apply the quadratic formula and interpret key features of the graph.
\end{enumerate}

%%===============================================================
\section{Rational Expressions}
\label{sec:rational-expressions}
%%===============================================================

Rational function is "any function which can be defined by a rational fraction, i.e. an algebraic fraction such that both the numerator and the denominator are polynomials".



\paragraph{Definition}



A function $f(x)$ is called a rational function if and only if it can be written in the form



$$f(x) = \frac{P(x)}{Q(x)}$$



where $P\,$ and $Q\,$ are polynomials in $x\,$ and $Q\,$ is not the zero polynomial. The domain of a function\|domain of $f\,$ is the set of all points $x\,$ for which the denominator $Q(x)\,$ is not zero.



However, if $\textstyle P$ and $\textstyle Q$ have a non constant polynomial greatest common divisor $\textstyle R$, then setting $\textstyle P=P _1R$ and $\textstyle Q=Q _1R$ produces a rational function



$$f _1(x) = \frac{P _1(x)}{Q _1(x)},$$ which may have a larger domain than $f(x)$, and is equal to $f(x)$ on the domain of $f(x).$ It is a common usage to identify $f(x)$ and $f _1(x)$, that is to extend "by continuity" the domain of $f(x)$ to that of $f _1(x).$ Indeed, one can define a rational fraction as an equivalence class of fractions of polynomials, where two fractions \textit{A}(\textit{x})/\textit{B}(\textit{x}) and \textit{C}(\textit{x})/\textit{D}(\textit{x}) are considered equivalent if \textit{A}(\textit{x})\textit{D}(\textit{x})=\textit{B}(\textit{x})\textit{C}(\textit{x}). In this case $\frac{P(x)}{Q(x)}$ is equivalent to $\frac{P _1(x)}{Q _1(x)}$.



\paragraph{Examples}



% [Image: y = (x^3-2x)/(2(x^2-5))]



Figure: $y = \frac{x^3-2x}{2(x^2-5)}$



% [Image: y = (x^2-3x-2)/(x^2-4)]



Figure: $y = \frac{x^2-3x-2}{x^2-4}$



The rational function $f(x) = \frac{x^3-2x}{2(x^2-5)}$ is not defined at $x^2=5 \Leftrightarrow x=\pm \sqrt{5}$. It is asymptotic to $\frac{x}{2}$ as \textit{x} approaches infinity.



The rational function $f(x) = \frac{x^2 + 2}{x^2 + 1}$ is defined for all real numbers, but not for all complex numbers, since if \textit{x} were a square root of $-1$ (i.e. the imaginary unit or its negative), then formal evaluation would lead to division by zero: $f(i) = \frac{i^2 + 2}{i^2 + 1} = \frac{-1 + 2}{-1 + 1} = \frac{1}{0}$, which is undefined.



A constant function such as $f(x) = \pi$ is a rational function since constants are polynomials. Note that the function itself is rational, even though the value (mathematics)\|value of $f(x)$ is irrational for all $x$.



Every polynomial function $f(x) = P(x)$ is a rational function with $Q(x) = 1$. A function that cannot be written in this form, such as $f(x) = \sin(x)$, is not a rational function. The adjective "irrational" is not generally used for functions.



The rational function $f(x) = \frac{x}{x}$ is equal to 1 for all \textit{x} except 0, where there is a removable singularity. The sum, product, or quotient (excepting division by the zero polynomial) of two rational functions is itself a rational function. However, the process of reduction to standard form may inadvertently result in the removal of such singularities unless care is taken. Using the definition of rational functions as equivalence classes gets around this, since \textit{x}/\textit{x} is equivalent to 1/1.



\paragraph{Sketch a graph of a rational function}



(1)Let's sketch the graph of $y=\frac{1}{x}$.

First, we must avoid $x=0$ because anything can not be divided by 0. Thus x should not be 0 in the equation. Now we just plug in some values of x. The result is as follows:

$$ \begin{array}{|c|c|} \hline \textbf{x} & \textbf{y} \\\\ \hline 1 & 1 \\\\ \hline 2 & 1/2 \\\\ \hline 3 & 1/3 \\\\ \hline -3 & -1/3 \\\\ \hline -2 & -1/2 \\\\ \hline -1 & -1 \\\\ \hline \end{array} $$ 



As x get large the function itself gets smaller and smaller. Here is the graph of $\frac{1}{x}$.



% [Image: ]



\paragraph{Resources}



\begin{itemize}
    \item \href{https://mathresearch.utsa.edu/wikiFiles/MAT1053/Rational_Functions/MAT1053_M4.2Rational_Functions.pdf}{Rational Functions}, Book Chapter
    \item \href{https://mathresearch.utsa.edu/wikiFiles/MAT1053/Rational_Functions/MAT1053_M4.2Rational_FunctionsGN.pdf}{Guided Notes}
\end{itemize}

\paragraph{Licensing}



Content obtained and/or adapted from:



\begin{itemize}
    \item \href{https://en.wikibooks.org/wiki/Intermediate_level_Mathematics/Rational_functions}{Rational Functions, Wikibooks: Intermediate Level Mathematics} under a CC BY-SA license
\end{itemize}

\subsection{Learning Objectives}

\begin{enumerate}
    \item Here are three Student Learning Objectives (SLOs) related to rational functions:
    \item \textbf{Identify and Classify Rational Functions:}
    \item \textbf{Outcome:} Students will be able to identify rational functions and classify them based on the degree of the numerator and denominator.
    \item \textbf{Details:} This includes recognizing the standard form of a rational function, understanding how the degrees of the numerator and denominator influence the function's characteristics, and distinguishing between proper and improper rational functions.
    \item \textbf{Analyze the Domain of Rational Functions:}
    \item \textbf{Outcome:} Students will be able to determine the domain of a given rational function by identifying the values of $x$ that cause the denominator to be zero.
    \item \textbf{Details:} This involves setting the denominator equal to zero, solving for $x$, and understanding how these values create restrictions on the function's domain. Additionally, students will explore how to extend the domain by simplifying the function where possible.
    \item \textbf{Graph Rational Functions:}
    \item \textbf{Outcome:} Students will be able to sketch the graph of a rational function, identifying key features such as asymptotes, intercepts, and discontinuities.
    \item \textbf{Details:} Students will practice plotting rational functions by first determining the vertical and horizontal asymptotes, then identifying any holes in the graph, and finally sketching the overall shape by evaluating the function at various points.
\end{enumerate}

\subsection{Assessment Instrument}

\begin{enumerate}
    \item \textbf{Assessment Instrument: Rational Expressions Quiz} \textbf{Format:}
    \item \textbf{Type:} Mixed-format quiz (multiple-choice, short answer, and problem-solving questions)
    \item \textbf{Duration:} 45 minutes
    \item \textbf{Total Marks:} 30 points \textbf{Sections:}
    \item \textbf{Identifying and Classifying Rational Functions (10 points):}
    \item \textbf{Question 1:} Determine if the following functions are rational. If they are, classify them as either proper or improper:
    \item a. $ f(x) = \frac{2x^2 + 3x + 1}{x^2 - 1} $
    \item b. $ g(x) = \frac{x^3 + 4x}{x - 2} $
    \item c. $ h(x) = \frac{3}{x^2 + 4x + 4} $
    \item d. $ j(x) = 5x + 3 $
    \item e. $ k(x) = \frac{x^2 + 1}{x^3 - x} $
    \item \textbf{Analyzing the Domain of Rational Functions (10 points):}
    \item \textbf{Question 2:} Find the domain of the following rational functions:
    \item a. $ f(x) = \frac{3x + 2}{x^2 - 9} $
    \item b. $ g(x) = \frac{x^2 - 4x}{x^2 + x - 6} $
    \item c. $ h(x) = \frac{2x^2 + x - 1}{x^2 + 3x + 2} $
    \item d. $ j(x) = \frac{5x + 7}{x^2 + x - 12} $
    \item e. $ k(x) = \frac{x^3 - x}{x^2 - 4} $
    \item \textbf{Graphing Rational Functions (10 points):}
    \item \textbf{Question 3:} Sketch the graph of the following rational functions, identifying key features such as asymptotes, intercepts, and any holes in the graph:
    \item a. $ f(x) = \frac{x^2 - 1}{x - 1} $
    \item b. $ g(x) = \frac{2x + 3}{x^2 - 4} $
    \item c. $ h(x) = \frac{1}{x^2} $
    \item d. $ j(x) = \frac{x^3 - 2x}{x^2 - 1} $
    \item e. $ k(x) = \frac{3x^2 + 2x - 1}{x - 2} $ \textbf{Scoring Rubric:}
    \item \textbf{Identifying and Classifying Rational Functions:} 2 points for each correct identification and classification (total 10 points).
    \item \textbf{Analyzing the Domain of Rational Functions:} 2 points for each correctly determined domain (total 10 points).
    \item \textbf{Graphing Rational Functions:} 2 points for each correctly sketched graph, including identification of key features (total 10 points). The quiz will assess students' understanding of rational functions, including their ability to identify, classify, analyze, and graph these functions. It ensures that students can apply their knowledge of rational expressions in various contexts and understand the implications of domain restrictions and graph features.
\end{enumerate}

%%===============================================================
\section{More on Polynomial Functions}
\label{sec:more-on-polynomial-functions}
%%===============================================================

\paragraph{Factoring}



\paragraph{Common Factor}



Factoring polynomials deals with picking out the common factor. Just like before when we factored a real number we apply the same idea to binomials, trinomials, and other polynomials. We assume that you don't know how to factor variables, so we will introduce the concept by analyzing the question into smaller and smaller pieces to \textit{further understanding} with the example problem below. For those who know how to do this, this may seem a little bit excessive, but it is important to fully understand what these type of problems means. Plus, in mathematics, when problems are very difficult, you break it down into easier pieces until you can finally tackle the big monster that is the problem.



\paragraph{Example Problem 1}



Factor the expression $x^8+x^5$. Notice that each term in the expression has an $x$ with no coefficients. If you have not memorized exponent rules, please do! The reason for that memorization is to simplify the process. When multiplying two terms such that $a^{p}a^{q}$, the result of the expression is $a^{p+q}$. Thus, you want to have some $x^{n}$ that when distributed onto every term in the parentheses, you get the same expression above. Let us set that up now:



$$x^{n}\left(x^{a}+x^{b}\right)$$



While this seems like a complicated way to look at it, a simple tasked is asked of you. In short, $x^{n}\left(x^{a}+x^{b}\right)=x^{n}x^{a}+x^{n}x^{b}$. Because $a^{p}a^{q}=a^{p+q}$, each term in $x^{n}x^{a}+x^{n}x^{b}$ has the same basic structure. Ergo,



$$x^{n+a}+x^{n+b}$$



Here we have what the question is actually asking us to figure out: $x^{n+a}+x^{n+b}=x^8+x^5$. Because $x$ is always equal to itself, yet the exponents are not, we may further rewrite the problem like so:



$$\begin{cases} n+a=8 \\\\ n+b=5 \end{cases}$$



Since $n$ is in the system of equations twice, it may be easier to eliminate $n$ from the system. In that case, multiply $-1$ to the equation (2). After, add equation (1) to equation (2):



> (3): $a-b=3$



You could keep going, but you are going to find out some hard truths, and you probably have no time to waste getting every thing done. If you want a proof that what you are doing is futile, read the note.^\[Note\ 1.\]^ What most problems mean by factoring an expression is to simply make it so that equation (3) is easy to solve. That is, let $b=0$. From there:



$$ \begin{aligned} a - 0 & = 3 \\\\ a & = 3 \end{aligned} $$





Because we allowed these two truths to exist, we can substitute these two values into the systems to find that $n=5$. This option was very arbitrary (as you would have learned if you read note 1). Nevertheless, this is what most problems want you to do when \textit{factoring an expression with no coefficients}. Allow $n$ to equal the smallest power value in the expression. Either way:



$$x^{5+3}+x^{5+0}=x^8+x^5$$



$$x^{5}x^{3}+x^{5}x^{0}=x^{5}\left(x^{3}+1\right)=x^8+x^5\blacksquare$$



The biggest truth you should have learned is when factoring an expression with no coefficients, you are really being asked to find the greatest common factor, and you do that by factoring the variable in the problem to the \textit{smallest power value} (highlighted in ${\color{Red}\text{red}}$). Here are a few examples:



\begin{itemize}
    \item $b^{18}+b^{3}-b^{{\color{Red}2}}=b^{2}\left(b^{16}+b-1\right)$
    \item $x^{8}+x^{5}+x^{{\color{Red}4}}=x^{4}\left(x^{4}+x+1\right)$
    \item $z^{6}+z^{{\color{Red}3}}=z^{3}\left(z^{3}+1\right)$
    \item $p^{\frac{1}{3}}+p^{{\color{Red}\frac{1}{4}}}=p^{\frac{1}{4}}\left(p^{\frac{4}{3}}+1\right)$
\end{itemize}

But what about when factoring with coefficients. This is a little bit more complicated, but it works the same either way.



\paragraph{Example Problem 2}



Factor the expression $42x^3+12x^2+24x$. Here, we should be focusing on the coefficients before the exponents. If we were to do that, we would analyze this small problem first:



$$\text{GCF}(42,12,24)$$



If you were to notice (or use a calculator), the largest value that is a factor of the numbers within is $6$. Let us go ahead and \textit{factor} that into our expression:



$$6\left(7x^3+2x^2+4x\right)$$



The coefficients should not matter if you factor because when multiplying the expression $\left(ca^{p}\right)a^{q}$, where $c$ is some constant, as per the associative property, $\left(ca^{p}\right)a^{q}=c\left(a^{p}a^{q}\right)$. From there, multiplying one term by a constant does not matter: $c\left(a^{p}a^{q}\right)=c\left(a^{p+q}\right)=ca^{p+q}$.^\*^



Therefore, we can apply the same general rule we discovered for \textbf{Example 1}: $42x^3+12x^2+24x=6x(7x^2+2x+4)\blacksquare$



\begin{itemize}
    \item Even if the other value is multiplied by another arbitrary constant $d$, so long as the commutative property is a thing, $\left(ca^{p}\right)\left(da^{q}\right)=ca^{p}da^{q}=cda^{p}a^{q}=cda^{p+q}$.
\end{itemize}

This procedure outlined above seems to work for any polynomial. Let's try it on a few:



\begin{itemize}
    \item $3b^{18}+12b^{3}-6b^{2}=3b^{2}\left(b^{16}+4b-2\right)$
    \item $12x^{8}+\frac{1}{16}x^{5}-\frac{1}{4}=4\left(3x^{8}+\frac{1}{64}x^{5}-\frac{1}{16}\right)$
    \item $6z^{6}4z^{3}+2=2\left(3z^{6}+2z^{3}+1\right)$
\end{itemize}

We may now synthesize the information we learned from these two problems to outline a general procedure:



Factoring any polynomial with a common term (procedure)



> Given a polynomial with coefficients and variables, the following procedure should be followed.

    1.  Find the greatest common factor for each of the coefficients.

    2.  If there is no constant (a value that has no variable attached to it), factor the variable to the lowest power value in the expression. If there is a constant, simply keep the greatest common factor outside the parentheses.



You are now ready to fully factor any expression with a common term.



\paragraph{Grouping}



When there is no common term that travels throughout the \textit{entire} polynomial you can factor by grouping. When factoring by grouping you need to split up the polynomial into two binomials.



\paragraph{Example Problem 3}



Factor the expression $x^3+x^2+2x+4$. Before we group, we need to do a quick check to see if we can factor the entire expression with just one term outside the parenthesis. If "4" was not a constant, we could factor the expression using the second step of the procedure. Another thing is that $\text{GCF}\left(1,1,2,4\right)=1$, which if we were to factor, that would mean that we are multiplying the same thing by 1, which is not very helpful.



What are we to do then? We could factor by putting parts of the expression that have something similar together. This is harder to explain than do, so follow along with us. First, we are going to rewrite the expression. However, instead of keeping it black, we are going to recolor the expression with two groups of colors, red and green:

$${\color{Red}x^3+x^2}+{\color{Green}2x+4}$$



First, we are going to factor the red bit, ${\color{Red}x^3+x^2}$. Notice that this part of the expression has a common factor $x^{2}$. Factor that bit out and leave it to the side:



$$x^3+x^2$$



$$\Rightarrow{\color{Red}x^{2}\left(x+1\right)}$$



From there we group the second binomial (in green) from the polynomial: ${\color{Green}2x+4}$. Notice the common factor is $2$, so factor that bit:



$$2x+4$$



$$\Rightarrow{\color{Green}2\left(x+4\right)}$$



With these two bits factored, put them together; after all, they were part of the same expression:



$${\color{Red}x^3+x^2}+{\color{Green}2x+4}$$



$$\Rightarrow{\color{Red}x^{2}\left(x+1\right)}+{\color{Green}2\left(x+4\right)}\blacksquare$$



You are now finished. The concept is simple, although some problems may be a little bit more difficult than others.



Notice that we factored two binomials. Would it work if we were to factor a trinomial and then a binomial? We will answer that in a little bit. For now, try the next problem below, where we group polynomials with different variables.



\paragraph{Example Problem 4}



Factor the expression $x^{2}y^{3}+2x^{2}y+4xy^{3}+8xy$. As before, we need to do a quick check to see if we can factor the entire expression with just one term outside the parenthesis. First, $\text{GCF}\left(1,2,4,8\right)=1$, which if we were to factor, that would mean that we are multiplying the same thing by 1, which is not very helpful. However, because there is an $xy$ inside the expression in some form, we could factor the entire expression like this:



$$xy\left(xy^{2}+2x+4y^{2}+8\right)$$



We are not done however. Notice that you could still group the polynomial inside the parentheses with common terms. The $xy$ may seem to complicate things; however, you can simply ignore it and say you are working on the inside of the parentheses. If you do get easily confused when too many things are on the paper or screen, it may be helpful to ignore them but always be sure to come back to them. Let us rewrite the expression inside of the parentheses in the same way we did in \textbf{\textit{Example 3}}:



$${\color{Red}xy^{2}+2x}+{\color{Green}4y^{2}+8}$$



Looking at the red term $\left({\color{Red}xy^{2}+2x}\right)$, there is a common factor of $x$ with both terms, and $\text{GCF}\left(1,2\right)=1$, so we will factor $1x=x$ outside the red term:



$$\left(xy^{2}+2x\right)$$



$${\color{Red}x\left(y^{2}+2\right)}$$



As before, do the green term ${\color{Green}4y^{2}+8}$. Look at the coefficients and constants first: $\text{GCF}\left(4,8\right)=4$. We found a factor! As per the definition of constant, there is no variable and thus nothing to factor. Therefore, we can now factor that bit:



$$4y^{2}+8$$



$${\color{Green}4\left(y^{2}+2\right)}$$



Finally, put everything together, including the $xy$ term we ignored for a large part of the problem. In the spirit of not confusing the reader, the red and green parts will be placed inside the brackets "\[\]":



$$xy\left[{\color{Red}x\left(y^{2}+2\right)}+{\color{Green}4\left(y^{2}+2\right)}\right]$$



For the more astute readers, you may have noticed that $y^{2}+2$ is multiplied to each term inside the square brackets. In case you do not see it, let $y^{2}+2=C$. If we were to replace the inside of the brackets like so (we are ignoring the colors since they are unnecessary for this part of the problem)$$xy\left(xC+4C\right)\text{,}$$we would notice that each term is multiplied by C, which is a factor of the expression and must be separated like so:



$$xy\left[C\left(x+4\right)\right]$$



Remembering that we let $y^{2}+2=C$, we substitute that back in and get



$$xy\left[\left(y^{2}+2\right)\left(x+4\right)\right]$$



Finally, because of the associative property, we may ignore the large square brackets (but keep the parentheses because anything inside the parentheses must be done first before multiplying them:



$$xy\left(y^{2}+2\right)\left(x+4\right)\blacksquare$$



This problem was a little bit more involved than expected. Is there an easier way to do the green and red bits? Yes, and it is a method called \textbf{F.O.I.L} (Firsts, Outside, Insides, Lasts). Here is what we mean by foil (see right):



% [Image: By Brandenads (talk) - Public Domain, owned and attributed towards Brandenade.]



Figure: By Brandenads (talk) - Public Domain, owned and attributed towards Brandenade.



By remembering the foil method we can \textit{move backwards} and solve this problem: $xy^{2}+2x+4y^{2}+8$.



We know that we multiply "firsts" with no coefficients. It is easier to :



$$\left(x+ \right)\left(y^{2}+\right)$$



From there we go "outsides". We see that the first term \textbf{$x$} is a multiple to the second term, $2x$, in the first binomial. So we get



$$\left(x+ \right)\left(y^{2}+2\right)$$



Next, "insides": we see \textbf{$y^{2}$} is a multiple of $4$. Therefore, we get



$(x+4)(y^{2}+2)$



In fact, we are finished. Notice that "lasts" is $4\cdot2$, which gives us $8$, the last term of the polynomial.



Keep in mind, this trick doesn't always work because the middle terms often end up having the same common variables and get simplified. It's helpful to start with the first terms and then go to the last terms. If you can get these two things you should be able to find the middle terms. Another thing to note is that we had to factor $xy$ to make this work. While it could be possible to do the original problem like this, it is much easier to use FOIL when the last term does not have any variables. Before we move on to the next section, we must answer one final pressing question that we asked at the beginning of this section.



\textbf{Factoring by grouping: a trinomial then a binomial?} This question is fundamental to the curious student. Let us try it out using \textbf{\textit{Example 3}}:



$$x^3+x^2+2x+4$$



Let us factor the trinomial: $x^3+x^2+2x=x\left(x^{2}+x+2\right)$. Next, factor the monomial: $4=2\left(2\right)$. Now put them together:



$$x\left(x^{2}+x+2\right)+2\left(2\right)$$



It seems to functionally work the same. In fact, it gives you the same answer. This should make sense because how you group anything does not change the answer (remember, associative property). However, it is standard mathematical syntax to factor a polynomial into two binomials. Because the CLEP exam is multiple choice, it should be easy to tell what the question is asking you to do. Either way, we have a procedure outlined for you again:



\textbf{Factoring any polynomial with $\underline{\text{no}}$ common term (procedure)}



>Given a polynomial with coefficients and variables, the following procedure should be followed.

    1.  Find the greatest common factor for each of the coefficients.

    2.  If there is no constant (a value that has no variable attached to it), factor the variable to the lowest power value in the expression. If there is a constant, simply keep the greatest common factor outside the parentheses.

    3.  Group the polynomial into two binomials, and apply steps 1 and 2 to the different groupings. Do this separately for the two groupings.

    4.  After completely factoring the two binomials. Put them together. If there is a common factor with the two combined binomials, factor it.



The procedure outlined above is best demonstrated in \textbf{Example 4}, and it does have a new: \textbf{factoring by grouping}. Simply look back at that example and apply each and every step onto that polynomial. Once you have done that, you now fully understand grouping polynomials.



\paragraph{Trinomials}



It may seem strange to start factoring larger expressions only to then move back to the trinomial. However, there is a reason why the book is arranged the way it is. Trinomial factoring may be the most difficult aspect of the journey to, ultimately, solve equations involving degrees higher than 1. Nevertheless, this difficult concept is necessary for the college algebra student, especially for the next two sections of the "factoring" section.



As before, we will introduce the concept by exploring the consequences and behaviors that a mathematician may have first discovered techniques to solve. First, let us apply our algorithm and see if it applies to problems like these.



\paragraph{Example 5}



Factor the expression $x^{2}+x-2$.



1.  $\text{GCF}(1,1,-2)=1$.

2.  We can only factor $x^{0}=1$ for the following expression.

3.  ???

4.  ???



Notice how our algorithm is confused at steps 3 and 4. Rightly, it should be because our algorithm assumes that we are factoring by grouping, which is not always applicable to all situations, such as this one. However, we can "force" the algorithm to fit into this situation. For example, we can try making the middle term split up into two terms. That is, we set

$$x=(f _{1}+f _{2})x=f _{1}x+f _{2}x \quad \text{ where } f _{1}+f _{2}=1$$



However, if we are to get anywhere with this, we want the middle term to have common factors. This is an interesting dilemma!



Our ultimate end goal is to use the algorithm we have used previously on the trinomial above to get a factored form such the following is true:



$$\begin{aligned} x^{2}+x-2&=x^{2}+f _{1}x+f _{2}x-2 \\\\ & =x(x+f _{1})+c(x+f _{1}) & \text{where }c\text{ is some constant multiple.} \\\\ & =(x+c)(x+f _{1}) \end{aligned}$$



However, notice how we got this result. If we want this to be factored, then $(x+f _{1})$ must be a factor of this trinomial. The question is, what is $c$? This is the ultimate question. To answer this, we need to analyze this application of the algorithm carefully. Since $f _{2}x-2$ is grouped together, we want to find a greatest common factor for $x$ and for $-2$. Ultimately, this means $f _{2}$ is the greatest factor of the grouped expression. Otherwise, it would be impossible to get $(x+f _{1})$ as a factor of the expression.



Therefore, $c=f _{2}$. This tells us a very important truth about factoring the above expression:

$f _{1}f _{2}=-2$



This is a result of applying our algorithm! Using what we know, let us rewrite the above problem in terms of what we know:

$$x^{2}+x-2=x^{2}+(f _{1}+f _{2})x-f _{1}f _{2} \text{, where} -2=f _{1}f _{2} \text{ and } 1=f _{1}+f _{2}$$



We can use a systems of equations!



$$\begin{cases} 1 & =f _{1}+f _{2} \\\\ -2 & =f _{1}f _{2} \end{cases}$$



For. the purposes of saving time, we will skip through the process of finding those factors. This can be trivially done with a "guess-and-check," but a technique for finding this factor will be explained later. For now, just thank us when we tell you the factors are $f _{1}=-1$ and $f _{2}=2$.



Using this result, we can now definitively use the same algorithm as before, but through some "fudging" or adjustments to the base expression we were solving for. Here is the process in action:



$\begin{aligned} x^{2}+x-2 & = x^{2}-x+2x-2 \\\\ & = x(x+1)-2(x+1) \\\\ & =(x-2)(x+1) \blacksquare \end{aligned}$



\textbf{Example 5} showed us a very important algorithm that is taught all around the world. It is called the \textbf{ac-method} or the \textbf{diamond method}, the steps for which we will outline below.



\textbf{AC-method of factoring trinomials}



> Given a trinomial of the form $ax^{2}+bx+c$, with $a\ne0$, the following procedure should be followed.

    1.  Find the greatest common factor for each of the coefficients.

    2.  If there is no constant (a value that has no variable attached to it), factor the variable to the lowest power value in the expression. If there is a constant, simply keep the greatest common factor outside the parentheses.

    3.  Find two values $f _{1}$ and $f _{2}$ whereby $ac=f _{1}f _{2}$ and $b=f _{1}+f _{2}$.

    4.  Rewrite the trinomial into the form $ax^{2}+f _{1}x+f _{2}x+c$; then, group the polynomial into two binomials, and apply steps 1 and 2 to the different groupings. Do this separately for the two groupings.

    5.  After completely factoring the two binomials. Put them together. If there is a common factor with the two combined binomials, factor it.



You should have realized by now that the ac-method is a modified version of the factoring by grouping method. This is why we decided to show factoring trinomials later, because the work involved is actually more than what you would expect from a carefully designed polynomial.



\paragraph{Difference of Two Squares}



So far in our factoring journey, we tried factoring general expressions that are not usually commonly seen. Despite that, there is a procedure that almost always works. Similarly, there are specialized cases in which a formula exists. It is the hope of this Wikibook that you understand why the formula is the way it is so that you may not simply blindly apply some seemingly unrelated relation.



\paragraph{Resources}



\begin{itemize}
    \item \href{https://mathresearch.utsa.edu/wikiFiles/MAT1053/More_on_Polynomial_and_Power_Functions/MAT1053_M4.1More_on_Power_Functions_and_Polynomial_Functions.pdf}{More on Power Functions and Polynomial Functions}, Book Chapter
    \item \href{https://courses.lumenlearning.com/ivytech-collegealgebra/chapter/use-arrow-notation/}{Arrow Notation for Polynomials}
    \item \href{https://mathresearch.utsa.edu/wikiFiles/MAT1053/More_on_Polynomial_and_Power_Functions/MAT1053_M4.1More_on_Power_Functions_and_Polynomial_FunctionsGN.pdf}{Guided Notes}
\end{itemize}

\paragraph{Licensing}



Content obtained and/or adapted from:



\begin{itemize}
    \item \href{https://en.wikibooks.org/wiki/CLEP_College_Algebra/Polynomials}{Polynomials, Wikibooks: CLEP College Algebra} under a CC BY-SA license
\end{itemize}

\subsection{Learning Objectives}

\begin{enumerate}
    \item Here are three Student Learning Objectives (SLOs) based on the topic of factoring, following the format provided:
    \item \textbf{Identify and Factor Common Factors:}
    \item \textbf{Outcome:} Students will be able to identify and factor out the greatest common factor (GCF) from polynomials, including those with coefficients and variables.
    \item \textbf{Details:} This includes understanding how to apply the GCF to simplify expressions by recognizing the smallest power value of a variable in an expression and factoring it out along with the GCF of the coefficients.
    \item \textbf{Factor by Grouping:}
    \item \textbf{Outcome:} Students will be able to factor polynomials by grouping terms when no common factor exists for the entire polynomial.
    \item \textbf{Details:} This involves dividing a polynomial into two binomials, factoring each group separately, and then combining the results. Emphasis is placed on recognizing when grouping is applicable and successfully applying the procedure to achieve the fully factored form.
    \item \textbf{Apply the AC-Method to Factor Trinomials:}
    \item \textbf{Outcome:} Students will be able to apply the AC-method (or diamond method) to factor trinomials, particularly when the trinomial cannot be factored by grouping directly.
    \item \textbf{Details:} This includes finding appropriate factors that satisfy the conditions $ac = f _1f _2$ and $b = f _1 + f _2$, rewriting the trinomial, and using the grouping method to achieve the factored form.
\end{enumerate}

\subsection{Assessment Instrument}

\begin{enumerate}
    \item \textbf{Assessment Instrument: Factoring Quiz} \textbf{Format:}
    \item \textbf{Type:} Mixed-format quiz (multiple-choice, short answer, and problem-solving questions)
    \item \textbf{Duration:} 45 minutes
    \item \textbf{Total Marks:} 30 points \textbf{Sections:}
    \item \textbf{Identifying and Factoring Common Factors (10 points):}
    \item \textbf{Question 1:} Identify the greatest common factor (GCF) and factor the following polynomials:
    \item a. $ 18x^5 + 12x^3 $
    \item b. $ 35y^4 - 21y^2 $
    \item c. $ 50a^6b^2 + 30a^4b^3 + 20a^2b^4 $
    \item \textbf{Scoring Rubric:} 3 points for correctly identifying the GCF and 2 points for correct factoring (total 10 points)
    \item \textbf{Factoring by Grouping (10 points):}
    \item \textbf{Question 2:} Factor the following polynomials by grouping:
    \item a. $ x^3 + 3x^2 + 2x + 6 $
    \item b. $ 4ab - 8a + 5b - 10 $
    \item \textbf{Scoring Rubric:} 5 points for correctly factoring each polynomial (total 10 points)
    \item \textbf{Applying the AC-Method to Factor Trinomials (10 points):}
    \item \textbf{Question 3:} Use the AC-method to factor the following trinomials:
    \item a. $ 2x^2 + 7x + 3 $
    \item b. $ 3x^2 - 5x - 2 $
    \item \textbf{Scoring Rubric:} 5 points for correctly applying the AC-method and achieving the factored form (total 10 points) \textbf{Scoring Rubric:}
    \item \textbf{Identifying and Factoring Common Factors:} 3 points for correctly identifying the GCF and 2 points for correct factoring (total 10 points)
    \item \textbf{Factoring by Grouping:} 5 points for correctly factoring each polynomial (total 10 points)
    \item \textbf{Applying the AC-Method to Factor Trinomials:} 5 points for correctly applying the AC-method and achieving the factored form (total 10 points) This quiz will assess the students' understanding and ability to factor polynomials by identifying common factors, grouping terms, and applying the AC-method to trinomials. The assessment is designed to ensure that students can apply these factoring techniques to a variety of polynomial expressions.
\end{enumerate}

%%===============================================================
\section{Exponential Functions}
\label{sec:exponential-functions}
%%===============================================================

\paragraph{Operations With Exponential Function}



An exponential function is a function where a constant base (b) is raised to a variable.



\paragraph{Multiplication}



Firstly, $b^x \times b^{2x}$ is $b^{\left(x + 2x\right)}\,$ which is $b^{3x}$. So when you multiply a base by the same base you add the variables. To clarify, here is an example with numbers:



$$\begin{array}{|c|c|c|c|c|} \hline x & 2^x & 2^{2x} & 2^x \times 2^{2x} & 2^{3x} \\\\ \hline 1 & 2 & 4 & 8 & 8 \\\\ \hline 2 & 4 & 16 & 64 & 64 \\\\ \hline \end{array}$$



\paragraph{Division}



Secondly $\frac{b^{2x}}{b^y}$ is $b^{\left( 2x - y \right)}\,$. So when a base is divided by the same base you subtract the variables.



Here is an example with numbers: $\frac{2^4}{2^2}=\frac{16}{4}=4=2^2$.



\paragraph{Base raised to two powers}



Thirdly $\left(b^{2x}\right)^{3x}$ is $b^{\left(2x\right) \times \left(3x\right)}$ which is $b^{6x^2}$. So when a base with a variable is raised to a variable you multiply the variables. Here is another example with numbers: (when x = 1) $\left(2^2\right)^3=4^3=64=2^6$.



\paragraph{Multiple bases}



Fourthly when $a^2 \times b^2 = ab \times ab$ it is the same as $\left(ab\right)^2\,$. Here is an example with numbers: $2^2 \times 3^2 = 36 = 6^2$. There is a similar situation with division: $\left(\frac{a}{b}\right)^2 = \frac{a}{b} \times \frac{a}{b} = \frac{a^2}{b^2}$. So when you multiply or divide two different bases raised to the same variable you can multiply or divide them first and then raise them to the variable.



\paragraph{Fractional exponents}



The last case is when $x$ is presented as a fraction, you can make a square root function, for example $b^\frac{1}{x}$ becomes $\pm \sqrt[x] b$. However it is customary to only use the positive root and so $b^\frac{1}{x}$ is defined as $\sqrt[x] b$. Another similar case is when the fraction has a constant (designated as $c$) other than 1 in the numerator , for example $b^ \frac {3}{x} = \left( \sqrt[x] b \right)^3$ so $b^ \frac {c}{x} = \left( \sqrt[x] b \right)^c$.



\paragraph{The Laws of Exponents}



The rules that have been suggested above are known as the laws of exponents and can be written as:



1.  $b^xb^y = b^{x+y}\,$

2.  $\frac{b^x}{b^y} = b^{x-y}$

3.  $\left(b^x\right)^y = b^{xy}$

4.  $a^n b^n = \left(ab\right)^n\,$

5.  $\left(\frac{a}{b}\right)^n = \frac{a^n}{b^n}$

6.  $b^{-n}=\frac{1}{b^n}$

7.  $b^ \frac {c}{x} = \left( \sqrt[x] b \right)^c$ where c is a constant

8.  $b^1=b\,$

9.  $b^0=1\,$



\paragraph{Graphing an Exponential Function}



When you graph an exponential function you use the same methods as with a regular function. There is a graph below that you can look at.



% [Image: ]



\paragraph{Solving Exponential Equations}



In order to solve an exponential equation you need to make sure that all the bases are the same. Then you can remove the base and solve for the variable. Here is an example:



Solve for x. $2^{\left(x-1\right)} = 16\,$



Now we convert 16 to a base 2 raised to a number.



$2^{\left(x-1\right)} = 2^4\,$



Now we can remove the base. So we have:



$x-1 = 4\,$



Finally solve for x.



$x = 5\,$



\paragraph{Resources}



\begin{itemize}
    \item \href{https://mathresearch.utsa.edu/wikiFiles/MAT1053/Exponential\%20Functions/MAT1053_M5.1Exponential_Functions.pdf}{Exponential Functions}, Book Chapter
    \item \href{https://mathresearch.utsa.edu/wikiFiles/MAT1053/Exponential\%20Functions/MAT1053_M5.1Exponential_FunctionsGN.pdf}{Guided Notes}
\end{itemize}

\paragraph{Licensing}



Content obtained and/or adapted from:



\begin{itemize}
    \item \href{https://en.wikibooks.org/wiki/A-level_Mathematics/OCR/C2/Logarithms_and_Exponentials}{Logarithms and Exponentials, Wikibooks: A-level Mathematics/OCR/C2} under a CC BY-SA license
\end{itemize}

\subsection{Learning Objectives}

\begin{enumerate}
    \item Here are three Student Learning Objectives (SLOs) for the topic "Operations With Exponential Functions":
    \item \textbf{Apply the Laws of Exponents:}
    \item \textbf{Outcome:} Students will be able to accurately apply the laws of exponents to simplify expressions involving multiplication, division, and powers of exponential functions.
    \item \textbf{Details:} This includes recognizing when to add, subtract, or multiply exponents, and applying these rules in both symbolic and numerical examples.
    \item \textbf{Solve Exponential Equations:}
    \item \textbf{Outcome:} Students will be able to solve exponential equations by expressing all terms with the same base and isolating the variable.
    \item \textbf{Details:} This involves converting exponential expressions to have a common base, removing the base, and solving the resulting linear or nonlinear equations.
    \item \textbf{Graph Exponential Functions:}
    \item \textbf{Outcome:} Students will be able to graph exponential functions and interpret their key features, such as intercepts, asymptotes, and growth or decay behavior.
    \item \textbf{Details:} This includes plotting exponential functions on a coordinate plane, identifying the base and exponent from the function’s equation, and describing the function’s behavior as the variable increases or decreases.
\end{enumerate}

\subsection{Assessment Instrument}

\begin{enumerate}
    \item \textbf{Assessment Instrument: Exponential Functions Quiz} \textbf{Format:}
    \item \textbf{Type:} Mixed-format quiz (multiple-choice, short answer, and problem-solving questions)
    \item \textbf{Duration:} 45 minutes
    \item \textbf{Total Marks:} 30 points \textbf{Sections:}
    \item \textbf{Laws of Exponents (10 points):}
    \item \textbf{Question 1:} Simplify the following expressions using the laws of exponents:
    \item a. $2^3 \times 2^4$
    \item b. $\frac{5^6}{5^2}$
    \item c. $(3^2)^4$
    \item d. $\frac{7^3}{7^{-2}}$
    \item \textbf{Question 2:} Which of the following is equivalent to $b^{\frac{3}{2}}$?
    \item a. $\sqrt{b^3}$
    \item b. $(\sqrt{b})^3$
    \item c. $\sqrt[3]{b^2}$
    \item d. $b^6$
    \item \textbf{Solving Exponential Equations (10 points):}
    \item \textbf{Question 3:} Solve the following exponential equations:
    \item a. $3^{2x} = 27$
    \item b. $4^{x-1} = 64$
    \item c. $2^{x+3} = 16$
    \item d. $5^{2x} = 125$
    \item \textbf{Question 4:} If $2^x = 32$, what is the value of $x$?
    \item \textbf{Graphing Exponential Functions (10 points):}
    \item \textbf{Question 5:} Given the function $f(x) = 2^x$, plot the graph of the function for $x = -2, -1, 0, 1, 2$. Label the key features such as the y-intercept and asymptote.
    \item \textbf{Question 6:} Describe how the graph of $f(x) = 3^x$ differs from the graph of $f(x) = 3^{-x}$. \textbf{Scoring Rubric:}
    \item \textbf{Laws of Exponents:} 2 points for each correctly simplified expression (total 8 points) and 2 points for the correct multiple-choice answer (total 10 points).
    \item \textbf{Solving Exponential Equations:} 2.5 points for each correctly solved equation (total 10 points).
    \item \textbf{Graphing Exponential Functions:} 5 points for a correct graph (total 5 points) and 5 points for a correct description (total 10 points). This quiz will assess the students' ability to apply the laws of exponents, solve exponential equations, and graph exponential functions, ensuring they have a comprehensive understanding of these concepts.
\end{enumerate}

%%===============================================================
\section{Logarithmic Functions}
\label{sec:logarithmic-functions}
%%===============================================================

\paragraph{Logarithmic Functions}



% [Image: ]



In mathematics you can find the inverse of an exponential function by switching x and y around: $y = b^x\,$ becomes $x = b^y\,$. The problem arises on how to find the value of y. The logarithmic function solved this problem. All conversions of logarithmic function into an exponential function follow the same pattern: $x = b^y\,$ becomes $y = \log _b x\,$. If a log is given without a written b then b=10. Also with logarithmic functions, $b > 0$ and $b \ne 1$. There are 2 cases where the log is equal to x: $\log _bb^X = X\,$ and $b^{\log _bX} = X\,$.



To recap, a \textbf{logarithm} is the inverse function of an exponent.



e.g. The inverse of the function $f(x) = 3^x$ is $f^{-1}(x) = \log _3 x$.



In general, $y = b^x \iff x = \log _b y$, given that $b > 0$.



\paragraph{Laws of Logarithmic Functions}



When X and Y are positive.



\begin{itemize}
    \item $\log _bXY = \log _bX + \log _bY\,$
    \item $\log _b \frac{X}{Y} = \log _bX - \log _bY\,$
    \item $\log _b X^k = k \log _bX\,$
\end{itemize}

\paragraph{Change of Base}



When x and b are positive real numbers and are not equal to 1. Then you can write $\log _a x\,$ as $\frac { \log _b x}{ \log _b a}$. This works for the natural log as well. here is an example:



\paragraph{Solving a Logarithmic Equation}



A \textbf{logarithmic equation} is an equation wherein one or more of the terms is a logarithm.



e.g. Solve $\log x + \log (x+2) = 2$ ($\log$ is another way of writing $\log _{10}$).



$$\begin{aligned} \log x + \log (x+2) & = 2 \\\\ \log (x(x+2)) & = 2 \\\\ x(x+2) & = 100 \\\\ x^2 + 2x & = 100 \\\\ (x + 1)^2 & = 101 \\\\ x+1 & = \sqrt{101} \\\\ x & = -1 \pm \sqrt{101} \end{aligned}$$



$\log _2 8 = \frac { \log 8}{ \log 2} = \frac {.9}{.3} = 3\,$ now check $2^3 = 8\,$



\paragraph{Resources}



\begin{itemize}
    \item \href{https://mathresearch.utsa.edu/wikiFiles/MAT1053/Logarithmic_Functions/MAT1053_M5.2Logarithmic_Functions.pdf}{Logarithmic Functions}, Book Chapter
    \item \href{https://mathresearch.utsa.edu/wikiFiles/MAT1053/Logarithmic_Functions/MAT1053_M5.2Logarithmic_FunctionsGN.pdf}{Guided Notes}
\end{itemize}

\paragraph{Licensing}



Content obtained and/or adapted from:



\begin{itemize}
    \item \href{https://en.wikibooks.org/wiki/A-level_Mathematics/OCR/C2/Logarithms_and_Exponentials}{Logarithms and Exponentials, Wikibooks: A-level Mathematics/OCR/C2} under a CC BY-SA license
\end{itemize}

\subsection{Learning Objectives}

\begin{enumerate}
    \item \textbf{Understand Logarithmic Functions:}
    \item \textbf{Outcome:} Students will be able to define logarithmic functions and explain their relationship to exponential functions.
    \item \textbf{Details:} This includes converting between exponential and logarithmic forms and identifying the components of each.
    \item \textbf{Apply Laws of Logarithmic Functions:}
    \item \textbf{Outcome:} Students will be able to utilize the laws of logarithms to simplify and manipulate logarithmic expressions.
    \item \textbf{Details:} This encompasses using properties such as the product, quotient, and power rules in various problem-solving contexts.
    \item \textbf{Solve Logarithmic Equations:}
    \item \textbf{Outcome:} Students will be able to solve equations that involve logarithmic expressions.
    \item \textbf{Details:} This involves applying techniques to isolate the variable, utilizing change of base formulas, and verifying solutions within the domain of the functions involved.
\end{enumerate}

\subsection{Assessment Instrument}

\begin{enumerate}
    \item \textbf{Assessment Instrument: Logarithmic Functions Quiz} \textbf{Format:}
    \item \textbf{Type:} Mixed-format quiz (multiple-choice, short answer, and problem-solving questions)
    \item \textbf{Duration:} 45 minutes
    \item \textbf{Total Marks:} 30 points \textbf{Sections:}
    \item \textbf{Understanding Logarithmic Functions (10 points):}
    \item \textbf{Question 1:} Convert the following exponential expressions to their equivalent logarithmic form:
    \item a. $2^5 = 32$
    \item b. $10^3 = 1000$
    \item c. $5^4 = 625$
    \item \textbf{Question 2:} Which of the following represents the inverse function of $f(x) = 4^x$?
    \item a. $f^{-1}(x) = \log _4 x$
    \item b. $f^{-1}(x) = 4x$
    \item c. $f^{-1}(x) = x^4$
    \item d. $f^{-1}(x) = \log _x 4$
    \item \textbf{Applying Laws of Logarithmic Functions (10 points):}
    \item \textbf{Question 3:} Simplify the following logarithmic expressions using the properties of logarithms:
    \item a. $\log _3 9 + \log _3 27$
    \item b. $\log _2 \frac{32}{8}$
    \item c. $3 \log _5 2$
    \item \textbf{Question 4:} Solve for $x$ in the equation $\log _4 (x^2) = 3$.
    \item \textbf{Solving Logarithmic Equations (10 points):}
    \item \textbf{Question 5:} Solve the following logarithmic equations:
    \item a. $\log _2 (x + 3) = 5$
    \item b. $\log _10 (3x - 7) = 2$
    \item \textbf{Question 6:} Use the change of base formula to evaluate $\log _3 81$. \textbf{Scoring Rubric:}
    \item \textbf{Understanding Logarithmic Functions:} 2.5 points for each correct answer (total 10 points)
    \item \textbf{Applying Laws of Logarithmic Functions:} 2.5 points for each correct simplification or solution (total 10 points)
    \item \textbf{Solving Logarithmic Equations:} 5 points for each correct solution (total 10 points) The quiz will assess students' understanding of logarithmic functions, their ability to apply the laws of logarithms, and their proficiency in solving logarithmic equations. This ensures that students can convert between exponential and logarithmic forms, simplify logarithmic expressions, and solve logarithmic equations accurately.
\end{enumerate}

%%===============================================================
\section{Logarithmic Properties}
\label{sec:logarithmic-properties}
%%===============================================================

Logarithmic Properties



$$ \begin{aligned} \text{Equality rule:} &\quad \text{If } \log _a(x) = \log _a(y), \text{ then } x = y \\\\ \text{Product rule:} &\quad \log _a(xy) = \log _a(x) + \log _a(y) \\\\ \text{Quotient rule:} &\quad \log _a\left(\frac{x}{y}\right) = \log _a(x) - \log _a(y) \\\\ \text{Power rule:} &\quad \log _a(x^n) = n\log _a(x) \\\\ \text{Change-of-base rule:} &\quad \log _a(x) = \frac{\log _b(x)}{\log _b(a)} \end{aligned} $$ 



\paragraph{Resources}



\begin{itemize}
    \item \href{https://mathresearch.utsa.edu/wikiFiles/MAT1053/Logarithmic\%20Properties/MAT1053_M6.1Logarithmic_Properties.pdf}{Logarithmic Properties}, Book Chapter
    \item \href{https://mathresearch.utsa.edu/wikiFiles/MAT1053/Logarithmic\%20Properties/MAT1053_M6.1Logarithmic_PropertiesGN.pdf}{Guided Notes}
\end{itemize}

\subsection{Learning Objectives}

\begin{enumerate}
    \item Here are three Student Learning Objectives (SLOs) related to logarithmic properties:
    \item \textbf{Apply the Equality Rule of Logarithms:}
    \item \textbf{Outcome:} Students will be able to solve equations involving logarithms by applying the equality rule, which states that if $ \log _a(x) = \log _a(y) $, then $ x = y $.
    \item \textbf{Details:} This includes identifying when logarithmic expressions can be equated and solving for the unknown variable by recognizing and applying the equality rule.
    \item \textbf{Use the Product, Quotient, and Power Rules of Logarithms:}
    \item \textbf{Outcome:} Students will be able to simplify logarithmic expressions using the product rule, quotient rule, and power rule of logarithms.
    \item \textbf{Details:} This involves understanding how to expand or condense logarithmic expressions. Students will demonstrate the ability to rewrite expressions like $ \log _a(xy) $, $ \log _a\left(\frac{x}{y}\right) $, and $ \log _a(x^n) $ using these rules and apply these rules in solving logarithmic equations.
    \item \textbf{Apply the Change-of-Base Rule:}
    \item \textbf{Outcome:} Students will be able to convert logarithms from one base to another using the change-of-base rule, $ \log _a(x) = \frac{\log _b(x)}{\log _b(a)} $.
    \item \textbf{Details:} This includes demonstrating the ability to use the change-of-base formula to calculate logarithms using different bases and solving related problems on calculators that only accept logarithms of base 10 or base $ e $.
\end{enumerate}

\subsection{Assessment Instrument}

\begin{enumerate}
    \item \textbf{Assessment Instrument: Logarithmic Properties Quiz} \textbf{Format:}
    \item \textbf{Type:} Mixed-format quiz (multiple-choice, short answer, and problem-solving questions)
    \item \textbf{Duration:} 45 minutes
    \item \textbf{Total Marks:} 30 points \textbf{Sections:}
    \item \textbf{Applying the Equality Rule (10 points):}
    \item \textbf{Question 1:} Solve the following logarithmic equations using the equality rule:
    \item a. If $ \log _3(x) = \log _3(9) $, find the value of $ x $.
    \item b. If $ \log _{10}(2x - 5) = \log _{10}(11) $, solve for $ x $.
    \item c. Given that $ \log _2(y + 4) = \log _2(12) $, determine the value of $ y $.
    \item \textbf{Simplifying Logarithmic Expressions (10 points):}
    \item \textbf{Question 2:} Simplify the following logarithmic expressions using the product, quotient, and power rules:
    \item a. $ \log _5(25) + \log _5(4) $
    \item b. $ \log _7\left(\dfrac{49}{7}\right) $
    \item c. $ 3\log _2(8) $
    \item d. $ \log _3(81) - \log _3(9) $
    \item \textbf{Using the Change-of-Base Rule (10 points):}
    \item \textbf{Question 3:} Convert the following logarithms to base 10 using the change-of-base rule:
    \item a. $ \log _2(32) $
    \item b. $ \log _5(125) $
    \item c. $ \log _4(16) $
    \item d. $ \log _3(27) $ \textbf{Scoring Rubric:}
    \item \textbf{Applying the Equality Rule:} 3.33 points for each correctly solved equation (total 10 points)
    \item \textbf{Simplifying Logarithmic Expressions:} 2.5 points for each correctly simplified expression (total 10 points)
    \item \textbf{Using the Change-of-Base Rule:} 2.5 points for each correctly converted logarithm (total 10 points) The quiz will assess the students' understanding and application of logarithmic properties, including solving logarithmic equations, simplifying expressions, and converting between logarithmic bases. This ensures that students can accurately apply the rules of logarithms in various contexts.
\end{enumerate}

%%===============================================================
\section{Exponential Properties}
\label{sec:exponential-properties}
%%===============================================================

\paragraph{Introduction}



Exponential properties can be used to manipulate equations involving exponential expressions and/or functions. Here are some important exponential properties:



$$ \begin{aligned} \text{Negative exponent property:} & \quad a^{-n} = \frac{1}{a^{n}} \quad \text{and} \quad \frac{1}{a^{-n}} = a^{n} \\\\ \text{Product of like bases:} & \quad a^m a^n = a^{m+n} \\\\ \text{Quotient of like bases:} & \quad \frac{a^m}{a^n} = a^m a^{-n} = a^{m-n} \\\\ \text{Multiple powers:} & \quad (a^m)^n = a^{mn} = (a^n)^m \\\\ \text{Product to a power:} & \quad (ab)^n = a^n b^n \\\\ \text{Quotient to a power:} & \quad \left(\frac{a}{b}\right)^n = \frac{a^n}{b^n} \end{aligned} $$ 



Special cases involving 0:



\begin{itemize}
    \item For any nonzero number $a$, $a^0 = 1$.
    \item For any positive number $m$, $0^m = 0$.
    \item $0^m$ does not exist if m is negative (since $0^{-n} = 1/0^n = 1/0$).
    \item $0^0$ is indeterminate or undefined depending on the context).
\end{itemize}

\paragraph{Resources}



\begin{itemize}
    \item \href{https://tutoring.asu.edu/sites/default/files/exponentialandlogrithmicproperties.pdf}{Exponential and Logarithmic Properties}, Arizona State University
    \item \href{https://www.khanacademy.org/math/cc-eighth-grade-math/cc-8th-numbers-operations/cc-8th-exponent-properties/a/exponent-properties-review}{Exponent Properties Review}, Khan Academy
\end{itemize}

\subsection{Learning Objectives}

\begin{enumerate}
    \item Here are three Student Learning Objectives (SLOs) based on the introduction to exponential properties:
    \item \textbf{Apply Negative Exponent Properties:}
    \item \textbf{Outcome:} Students will be able to apply the negative exponent property to simplify expressions and solve equations involving negative exponents.
    \item \textbf{Details:} This includes rewriting expressions with negative exponents as their reciprocal positive exponents and manipulating these expressions to solve equations.
    \item \textbf{Manipulate Exponential Expressions Using the Product and Quotient Properties:}
    \item \textbf{Outcome:} Students will be able to manipulate and simplify exponential expressions using the properties of the product and quotient of like bases.
    \item \textbf{Details:} This involves understanding and applying the rules $a^m \cdot a^n = a^{m+n}$ and $\frac{a^m}{a^n} = a^{m-n}$ in various algebraic contexts, including simplifying complex expressions.
    \item \textbf{Evaluate Expressions Using Multiple Powers and Power Properties:}
    \item \textbf{Outcome:} Students will be able to evaluate and simplify expressions involving powers raised to additional powers and products or quotients raised to a power.
    \item \textbf{Details:} This includes applying the rules $(a^m)^n = a^{mn}$, $(ab)^n = a^n b^n$, and $\left(\frac{a}{b}\right)^n = \frac{a^n}{b^n}$ to simplify and evaluate exponential expressions, as well as solve equations where these properties are applicable.
\end{enumerate}

\subsection{Assessment Instrument}

\begin{enumerate}
    \item \textbf{Assessment Instrument: Exponential Properties Quiz} \textbf{Format:}
    \item \textbf{Type:} Mixed-format quiz (multiple-choice, short answer, and problem-solving questions)
    \item \textbf{Duration:} 45 minutes
    \item \textbf{Total Marks:} 30 points \textbf{Sections:}
    \item \textbf{Negative Exponents (10 points):}
    \item \textbf{Question 1:} Simplify the following expressions using the negative exponent property:
    \item a. $3^{-2}$
    \item b. $\dfrac{1}{4^{-3}}$
    \item c. $x^{-5} \cdot x^2$
    \item d. $5x^{-2} \cdot 2x^3$
    \item e. $\left(\dfrac{2}{5}\right)^{-2}$
    \item \textbf{Product and Quotient of Like Bases (10 points):}
    \item \textbf{Question 2:} Simplify the following expressions using the product and quotient of like bases properties:
    \item a. $7^3 \cdot 7^4$
    \item b. $\dfrac{9^5}{9^2}$
    \item c. $x^4 \cdot x^{-2}$
    \item d. $\dfrac{5^8}{5^3}$
    \item e. $\dfrac{a^6 \cdot a^{-3}}{a^2}$
    \item \textbf{Multiple Powers and Power Properties (10 points):}
    \item \textbf{Question 3:} Simplify the following expressions using multiple powers and power properties:
    \item a. $(2^3)^2$
    \item b. $\left(\dfrac{3}{4}\right)^3$
    \item c. $(x^2y^3)^4$
    \item d. $\left(\dfrac{5x^2}{2}\right)^3$
    \item e. $\left(\dfrac{a^2b^{-3}}{c^4}\right)^2$ \textbf{Scoring Rubric:}
    \item \textbf{Negative Exponents:} 2 points for each correct simplification (total 10 points)
    \item \textbf{Product and Quotient of Like Bases:} 2 points for each correct simplification (total 10 points)
    \item \textbf{Multiple Powers and Power Properties:} 2 points for each correct simplification (total 10 points) The quiz will assess students' understanding and application of exponential properties, including negative exponents, product and quotient of like bases, and multiple powers. This ensures they can correctly simplify and manipulate expressions involving these properties.
\end{enumerate}

%%===============================================================
\section{Exponential and Logarithmic Equations}
\label{sec:exponential-and-logarithmic-equations}
%%===============================================================

\paragraph{Logarithms and Exponents}



A \textbf{logarithm} is the inverse function of an exponent.



e.g. The inverse of the function $f(x) = 3^x$ is $f^{-1}(x) = \log _3 x$.



In general, $y = b^x \iff x = \log _b y$, given that $b > 0$.



\paragraph{Laws of Logarithms}



The laws of logarithms can be derived from the laws of exponentiation:



$$\begin{aligned} x^{a+b} = x^a \times x^b & \iff \log a + \log b = \log ab \\\\ x^{a-b} = x^a \div x^b & \iff \log a - \log b = \log a/b \\\\ (x^a)^b = x^{ab} & \iff \log a^b = b \log a \end{aligned}$$



These laws apply to logarithms of any given base



\paragraph{Natural Logarithms}



The \textbf{natural logarithm} is a logarithm with base $e$, where $e$ is a constant such that the function $e^x$ is its own derivative.



The natural logarithm has a special symbol: $\ln x$



The graph $y = e^{kx}$ exhibits exponential growth when $k > 0$ and exponential decay when $k < 0$. The inverse graph is $y = \frac{1}{k} \ln x$. \href{https://www.desmos.com/calculator/azbv4kz38z}{Here} is an interactive graph which shows the two functions as inverses of one another.



\paragraph{Solving Logarithmic and Exponential Equations}



An \textbf{exponential equation} is an equation in which one or more of the terms is an exponential function. e.g. $5^x = 2^{x+2}$. Exponential equations can be solved with logarithms.



e.g. Solve $3^{x+1} = 4^{2x-1}$



$$\begin{aligned} 3^{x+1} & = 4^{2x-1} \\\\ (x+1)\ln 3 & = (2x-1)\ln 4 \\\\ x\ln 3 + \ln 3 & = 2x\ln 4 - \ln 4 \\\\ \ln 3 + \ln 4 & = x(2\ln 4 - \ln 3) \\\\ x & = \frac{\ln 3 + \ln 4}{2\ln 4 - \ln 3} \\\\ x & \approx 1.4844 \end{aligned}$$



A \textbf{logarithmic equation} is an equation wherein one or more of the terms is a logarithm.



e.g. Solve $\lg x + \lg (x+2) = 2$



$$\begin{aligned} \lg x + \lg (x+2) & = 2 \\\\ \lg (x(x+2)) & = 2 \\\\ x(x+2) & = 100 \\\\ x^2 + 2x & = 100 \\\\ (x + 1)^2 & = 101 \\\\ x+1 & = \sqrt{101} \\\\ x & = -1 \pm \sqrt{101} \end{aligned}$$



\paragraph{Converting Relationships to a Linear Form}



In maths and science, it is easier to deal with linear relationships than non-linear relationships. Logarithms can be used to convert some non-linear relationships into linear relationships.



\paragraph{Exponential Relationships}



An \textbf{exponential relationship} is of the form $y = ab^x$. If we take the natural logarithm of both sides, we get $\ln y = \ln a + x \ln b$. We now have a linear relationship between $\ln y$ and $x$.



e.g. The following data is related with an exponential relationship. Determine this exponential relationship, then convert it to linear form.



$$ \begin{array}{|c|c|} \hline x & y \\\\ \hline 0 & 5 \\\\ \hline 2 & 45 \\\\ \hline 4 & 405 \\\\ \hline \end{array} $$



$$\begin{aligned} \text{Exponential relationship } \implies y & = ab^x \\\\ 5 & = ab^0 = a(1) \\\\ a & = 5 \\\\ y & = 5b^x \\\\ 45 & = 5b^2 \\\\ 9 & = b^2 \\\\ b & = 3 \\\\ y & = 5(3^x) \end{aligned}$$



Now convert it to linear form by taking the natural logarithm of both sides:



$$\begin{aligned} y & = 5(3^x) \\\\ \ln y & = \ln 5 + x \ln 3 \end{aligned}$$



\paragraph{Power Relationships}



A \textbf{power relationship} is of the form $y = ax^b$. If we take the natural logarithm of both sides, we get $\ln y = \ln a + b \ln x$. This is a linear relationship between $\ln y$ and $\ln x$.



e.g. The amount of time that a planet takes to travel around the sun (its orbital period) and its distance from the sun are related by a power law. Use the following data to deduce this power law:



$$ \begin{array}{|c|c|c|} \hline \textbf{Planet} & \textbf{Distance from Sun} \,(\times 10^6 \, \text{km}) & \textbf{Orbital Period (days)} \\\\ \hline \text{Earth} & 149.6 & 365.2 \\\\ \hline \text{Mars} & 227.9 & 687.0 \\\\ \hline \text{Jupiter} & 778.6 & 4331 \\\\ \hline \end{array} $$



$$\begin{aligned} \text{Power law}\implies T & = aR^b \\\\ \text{Use Earth data}\implies 365.2 & = a(149.6^b) \\\\ \ln 365.2 & = \ln a + b \ln 149.6 \\\\ \text{Use Mars data}\implies 687.0 & = a(227.9^b) \\\\ \ln 687.0 & = \ln a + b \ln 227.9 \\\\ \ln 687.0 - \ln 365.2 & = \ln a - \ln a + b\ln 227.9 - b\ln 149.6 \\\\ \ln \frac{687.0}{365.2} & = 0 + b(\ln\frac{227.9}{149.6}) \\\\ b & = \frac{\ln \tfrac{687.0}{365.2}}{\ln\tfrac{227.9}{149.6}} \approx 1.5011 \\\\ \ln 365.2 & = \ln a + 1.5011 \ln 149.6 \\\\ \ln a & = \ln 365.2 - \ln 1839.9 \\\\ \ln a & = \ln 0.1985 \\\\ a & = 0.1985 \\\\ \therefore T & = 0.1985R^{1.5011} \end{aligned}$$



\paragraph{Change of Base Formula}



$\log _{y}x=\frac{\log _{a}x} {\log _{a}y}$ where \textit{a} is any positive number, distinct from 1. Generally, \textit{a} is either 10 (for common logs) or \textit{e} (for natural logs).



Proof:

$$ \begin{aligned} \log _{y}x & = b \\\\ y^b & = x \\\\ \log _{a}y^b & = \log _{a}x & \quad \text{Putting both sides to } \log _{a}\\\\ b\log _{a}y & = \log _{a}x \\\\ b & = \frac{\log _{a}x}{\log _{a}y} \\\\ \log _{y}x & = \frac{\log _{a}x}{\log _{a}y} & \quad \text{Replacing } b \text{ from first line} \end{aligned} $$



\paragraph{Swap of Base and Exponent Formula}



$a^{\log _{b}c}=c^{\log _{b}a}$ where a or c must not be equal to 1.



Proof:



$$ \begin{aligned} \log _{a}b &= \frac{1}{\log _{b}a} \quad \text{by the change of base formula above.} \\\\ \text{Note that } & a = c^{\log _{c}a}. \text{ Then } a^{\log _{b}c} \text{ can be rewritten as} \\\\ & \left(c^{\log _{c}a}\right)^{\log _{b}c} \quad \text{or by the exponential rule as} \\\\ & c^{\left(\log _{c}a \cdot \log _{b}c\right)} \\\\ & c^{\left(\log _{c}a \cdot \frac{1}{\log _{c}b}\right)} \quad \text{using the inverse rule noted above} \\\\ & c^{\log _{b}a} \quad \text{by the change of base formula.} \end{aligned} $$



\paragraph{Resources}



\begin{itemize}
    \item \href{https://mathresearch.utsa.edu/wikiFiles/MAT1093/Logarithmic\%20and\%20Exponential\%20Equations/Esparza\%201093\%20Notes\%207.4.pdf}{Logarithmic and Exponential Equations}. Written notes created by Professor Esparza, UTSA.
    \item \href{https://www.khanacademy.org/math/algebra2/x2ec2f6f830c9fb89:logs/x2ec2f6f830c9fb89:exp-eq-log/a/solving-exponential-equations-with-logarithms}{Solving Exponential Equations with Logarithms}, Khan Academy
    \item \href{https://www.mcckc.edu/tutoring/docs/br/math/expon_logar/Solving_Exponential_and_Logarithmic_Equations.pdf}{Solving Exponential and Logarithmic Equations}, Metropolitan Community College Kansas City
    \item \href{https://tutorial.math.lamar.edu/classes/calci/explogeqns.aspx}{Exponential and Logarithmic Equations}, Paul\'s Online Notes
\end{itemize}

\paragraph{Licensing}



Content obtained and/or adapted from:



\begin{itemize}
    \item \href{https://en.wikibooks.org/wiki/A-level_Mathematics/CIE/Pure_Mathematics_2/Logarithmic_and_Exponential_Functions}{Logarithmic and Exponential Functions, Wikibooks: A-level Mathematics} under a CC BY-SA license
    \item \href{https://en.wikibooks.org/wiki/Algebra/Logarithms}{Logarithms, Wikibooks: Algebra} under a CC BY-SA license
\end{itemize}

\subsection{Learning Objectives}

\begin{enumerate}
    \item Here are three Student Learning Objectives (SLOs) for the topic of Logarithms and Exponents, formatted similarly to your example:
    \item \textbf{Understand the Relationship between Logarithms and Exponents:}
    \item \textbf{Outcome:} Students will be able to explain the inverse relationship between logarithms and exponents, including converting between exponential and logarithmic forms of an equation.
    \item \textbf{Details:} This includes demonstrating that for any base $ b > 0 $, the equation $ y = b^x $ can be rewritten as $ x = \log _b y $ and vice versa, and applying this knowledge to solve related problems.
    \item \textbf{Apply the Laws of Logarithms:}
    \item \textbf{Outcome:} Students will be able to apply the laws of logarithms to simplify logarithmic expressions and solve logarithmic equations.
    \item \textbf{Details:} This includes using the product, quotient, and power rules of logarithms, such as $ \log(ab) = \log a + \log b $, $ \log\left(\frac{a}{b}\right) = \log a - \log b $, and $ \log(a^b) = b\log a $, and solving equations that require the application of these laws.
    \item \textbf{Solve Exponential and Logarithmic Equations:}
    \item \textbf{Outcome:} Students will be able to solve exponential and logarithmic equations by applying logarithmic transformations and the properties of exponents.
    \item \textbf{Details:} This includes solving equations of the form $ a^x = b^y $ by taking logarithms on both sides, as well as solving logarithmic equations by exponentiating both sides to eliminate the logarithm.
\end{enumerate}

\subsection{Assessment Instrument}

\begin{enumerate}
    \item \textbf{Assessment Instrument: Exponential and Logarithmic Equations Quiz} \textbf{Format:}
    \item \textbf{Type:} Mixed-format quiz (multiple-choice, short answer, and problem-solving questions)
    \item \textbf{Duration:} 45 minutes
    \item \textbf{Total Marks:} 30 points \textbf{Sections:}
    \item \textbf{Multiple Choice Questions (10 points):}
    \item \textbf{Question 1:} Which of the following is the correct solution for the equation $ 2^x = 16 $?
    \item a. $ x = 2 $
    \item b. $ x = 4 $
    \item c. $ x = 8 $
    \item d. $ x = 16 $
    \item \textbf{Question 2:} What is the logarithmic form of the equation $ 5^3 = 125 $?
    \item a. $ \log _5 3 = 125 $
    \item b. $ \log _3 5 = 125 $
    \item c. $ \log _5 125 = 3 $
    \item d. $ \log _125 5 = 3 $
    \item \textbf{Short Answer Questions (10 points):}
    \item \textbf{Question 3:} Solve the following exponential equation for $ x $: $ 3^{2x} = 81 $.
    \item \textbf{Question 4:} Express the equation $ \log _b x = y $ in its exponential form.
    \item \textbf{Question 5:} If $ \log _2 x = 5 $, what is the value of $ x $?
    \item \textbf{Problem-Solving Questions (10 points):}
    \item \textbf{Question 6:} Solve the logarithmic equation $ \log _3 (x - 2) + \log _3 4 = \log _3 12 $ and find the value of $ x $.
    \item \textbf{Question 7:} Solve the exponential equation $ 4^{x+1} = 2^{2x+3} $ for $ x $. \textbf{Scoring Rubric:}
    \item \textbf{Multiple Choice Questions:} 2.5 points for each correct answer (total 10 points).
    \item \textbf{Short Answer Questions:} 2 points for each correct answer (total 10 points).
    \item \textbf{Problem-Solving Questions:} 5 points for each correctly solved problem, with partial credit for correct steps even if the final answer is incorrect (total 10 points). This quiz will assess students' understanding of exponential and logarithmic equations, including their ability to convert between forms, apply logarithmic rules, and solve both types of equations accurately.
\end{enumerate}

%%===============================================================
\section{Simple Interest}
\label{sec:simple-interest}
%%===============================================================

\textbf{Simple interest} is calculated only on the principal amount, or on that portion of the principal amount that remains. It excludes the effect of compounding. Simple interest can be applied over a time period other than a year, for example, every month.



Simple interest is calculated according to the following formula:



$$\frac {r \cdot B \cdot m}{n}$$



where



\begin{itemize}
    \item \textit{r} is the simple annual interest rate
    \item \textit{B} is the initial balance
    \item \textit{m} is the number of time periods elapsed and
    \item \textit{n} is the frequency of applying interest.
\end{itemize}

For example, imagine that a credit card holder has an outstanding balance of $\text{\$}$2500 and that the simple annual interest rate is 12.99\% \textit{per annum}, applied monthly, so the frequency of applying interest is 12 per year. Over one month,



$$\frac{0.1299 \times \$2500}{12} = \$27.06$$



interest is due (rounded to the nearest cent).



Simple interest applied over 3 months would be



$$\frac{0.1299 \times \$2500 \times 3}{12} = \$81.19$$



If the card holder pays off only interest at the end of each of the 3 months, the total amount of interest paid would be



$$\frac{0.1299 \times \$2500}{12} \times 3 = \$27.06\text{ per month} \times 3\text{ months} =\$81.18$$



which is the simple interest applied over 3 months, as calculated above. (The one cent difference arises due to rounding to the nearest cent.)



\paragraph{Resources}



\begin{itemize}
    \item \href{https://mathresearch.utsa.edu/wikiFiles/MAT1053/Simple_and_Compound_Interest/MAT1053_M7.1Simple_and_Compound_Interest.pdf}{Simple and Compound Interest}, Book Chapter
    \item \href{https://mathresearch.utsa.edu/wikiFiles/MAT1053/Simple_and_Compound_Interest/MAT1053_M7.1Simple_and_Compound_InterestGN.pdf}{Guided Notes}
\end{itemize}

\paragraph{Licensing}



Content obtained and/or adapted from:



\begin{itemize}
    \item \href{https://en.wikipedia.org/wiki/Interest}{Interest, Wikipedia} under a CC BY-SA license
\end{itemize}

\subsection{Learning Objectives}

\begin{enumerate}
    \item Here are three Student Learning Objectives (SLOs) related to simple interest:
    \item \textbf{Calculate Simple Interest:}
    \item \textbf{Outcome:} Students will be able to accurately calculate simple interest using the formula $\frac{r \cdot B \cdot m}{n}$.
    \item \textbf{Details:} This includes understanding the variables in the formula (simple annual interest rate, initial balance, number of time periods, and frequency of interest application) and applying the formula to different financial scenarios, such as calculating interest for various periods and interest rates.
    \item \textbf{Analyze the Effects of Interest Frequency on Total Interest:}
    \item \textbf{Outcome:} Students will be able to analyze and compare the effects of changing the frequency of interest application (e.g., monthly, quarterly, annually) on the total simple interest accrued.
    \item \textbf{Details:} This involves calculating and interpreting the differences in simple interest amounts when the frequency of interest application varies, demonstrating how more frequent interest application impacts the total interest paid over a given period.
    \item \textbf{Interpret Simple Interest in Financial Contexts:}
    \item \textbf{Outcome:} Students will be able to interpret and explain the concept of simple interest in the context of financial decisions, such as credit card payments and loans.
    \item \textbf{Details:} This includes explaining how simple interest is calculated, its implications for short-term and long-term financial planning, and understanding how to manage debt or investments that accrue simple interest, using real-world examples to illustrate key concepts.
\end{enumerate}

\subsection{Assessment Instrument}

\begin{enumerate}
    \item \textbf{Assessment Instrument: Simple Interest Quiz} \textbf{Format:}
    \item \textbf{Type:} Mixed-format quiz (multiple-choice, short answer, and problem-solving questions)
    \item \textbf{Duration:} 45 minutes
    \item \textbf{Total Marks:} 30 points \textbf{Sections:}
    \item \textbf{Understanding Simple Interest (10 points):}
    \item \textbf{Question 1:} Define simple interest and describe how it differs from compound interest. (3 points)
    \item \textbf{Question 2:} Identify the components of the simple interest formula $\frac{r \cdot B \cdot m}{n}$. (3 points)
    \item \textbf{Question 3:} Explain what happens to the total simple interest if the interest is applied semi-annually instead of annually, assuming all other factors remain constant. (4 points)
    \item \textbf{Calculating Simple Interest (10 points):}
    \item \textbf{Question 4:} A loan of $\text{\$}$1500 is given at a simple annual interest rate of 8\%, applied quarterly. Calculate the interest accrued after 6 months. (5 points)
    \item \textbf{Question 5:} A savings account has a balance of $\text{\$}$2000, and the bank offers a simple interest rate of 5\% per annum, applied monthly. How much interest will be earned after 3 months? (5 points)
    \item \textbf{Analyzing Financial Scenarios (10 points):}
    \item \textbf{Question 6:} A credit card has an outstanding balance of $\text{\$}$5000 with a simple annual interest rate of 15\%, applied monthly. Calculate the interest due after one month if no payments are made. (4 points)
    \item \textbf{Question 7:} A borrower takes a loan of $\text{\$}$3000 at a simple interest rate of 10\% per annum. If the borrower pays off the interest each quarter, how much total interest will be paid after 9 months? (6 points) \textbf{Scoring Rubric:}
    \item \textbf{Understanding Simple Interest:} Points will be awarded based on the accuracy and completeness of definitions, identification of formula components, and explanations (total 10 points).
    \item \textbf{Calculating Simple Interest:} Full points for correct calculations and correct application of the simple interest formula (total 10 points).
    \item \textbf{Analyzing Financial Scenarios:} Points awarded for accurate problem-solving and correct interpretation of financial scenarios involving simple interest (total 10 points). This quiz is designed to assess students' understanding of simple interest, their ability to perform accurate calculations, and their capability to analyze and interpret financial scenarios involving simple interest.
\end{enumerate}

%%===============================================================
\section{Compound Interest}
\label{sec:compound-interest}
%%===============================================================

\textbf{Compound interest} includes interest earned on the interest that was previously accumulated.



Compare, for example, a bond paying 6 percent semiannually (that is, coupons of 3 percent twice a year) with a certificate of deposit (GIC) that pays 6 percent interest once a year. The total interest payment is $\text{\$}$6 per $\text{\$}$100 par value in both cases, but the holder of the semiannual bond receives half the $\text{\$}$6 per year after only 6 months (time preference), and so has the opportunity to reinvest the first $\text{\$}$3 coupon payment after the first 6 months, and earn additional interest.



For example, suppose an investor buys $\text{\$}$10,000 par value of a US dollar bond, which pays coupons twice a year, and that the bond's simple annual coupon rate is 6 percent per year. This means that every 6 months, the issuer pays the holder of the bond a coupon of 3 dollars per 100 dollars par value. At the end of 6 months, the issuer pays the holder:



$$\frac {r \cdot B \cdot m}{n} = \frac {6\% \times \$10\,000 \times 1}{2} = \$300$$



Assuming the market price of the bond is 100, so it is trading at par value, suppose further that the holder immediately reinvests the coupon by spending it on another $\text{\$}$300 par value of the bond. In total, the investor therefore now holds:



$$\$10\,000 + \$300 = \left(1 + \frac{r}{n}\right) \cdot B = \left(1 + \frac{6\%}{2}\right) \times \$10\,000$$



and so earns a coupon at the end of the next 6 months of:



$$\begin{aligned}\frac {r \cdot B \cdot m}{n} & = \frac {6\% \times \left(\$10\,000 + \$300\right)}{2} \\\\ & = \frac {6\% \times \left(1 + \frac{6\%}{2}\right) \times \$10\,000}{2} \\\\ & = \$309\end{aligned}$$



Assuming the bond remains priced at par, the investor accumulates at the end of a full 12 months a total value of:



$$\begin{aligned}\$10,000 + \$300 + \$309 & = \$10\,000 + \frac {6\% \times \$10,000}{2} + \frac {6\% \times \left( 1 + \frac {6\%}{2}\right) \times \$10\,000}{2} \\\\ & = \$10\,000 \times \left(1 + \frac{6\%}{2}\right)^2\end{aligned}$$



and the investor earned in total:



$$\begin{aligned}\$10\,000 \times \left(1 + \frac {6\%}{2}\right)^2 - \$10\,000 \\\\ = \$10\,000 \times \left( \left( 1 + \frac {6\%}{2}\right)^2 - 1\right)\end{aligned}$$



The formula for the \textbf{annual equivalent compound interest rate} is:



$$\left(1 + \frac{r}{n}\right)^n - 1$$



where



\begin{itemize}
    \item r is the simple annual rate of interest
    \item n is the frequency of applying interest
\end{itemize}

For example, in the case of a 6\% simple annual rate, the annual equivalent compound rate is:



$$\left(1 + \frac{6\%}{2}\right)^2 - 1 = 1.03^2 - 1 = 6.09\%$$



\paragraph{Resources}



\begin{itemize}
    \item \href{https://mathresearch.utsa.edu/wikiFiles/MAT1053/Simple_and_Compound_Interest/MAT1053_M7.1Simple_and_Compound_Interest.pdf}{Simple and Compound Interest}, Book Chapter
    \item \href{https://mathresearch.utsa.edu/wikiFiles/MAT1053/Simple_and_Compound_Interest/MAT1053_M7.1Simple_and_Compound_InterestGN.pdf}{Guided Notes}
\end{itemize}

\paragraph{Licensing}



Content obtained and/or adapted from:



\begin{itemize}
    \item \href{https://en.wikipedia.org/wiki/Interest}{Interest, Wikipedia} under a CC BY-SA license
\end{itemize}

\subsection{Learning Objectives}

\begin{enumerate}
    \item Here are three student learning objectives (SLOs) related to the concept of compound interest, using the example format you provided:
    \item \textbf{Calculate Compound Interest:}
    \item \textbf{Outcome:} Students will be able to calculate the future value of an investment using the compound interest formula, given the principal, interest rate, compounding frequency, and time period.
    \item \textbf{Details:} This includes understanding how to apply the formula $ A = P \left(1 + \frac{r}{n}\right)^{nt} $, where $ A $ is the amount of money accumulated after n years, including interest, $ P $ is the principal amount (the initial money), $ r $ is the annual interest rate (in decimal), $ n $ is the number of times that interest is compounded per year, and $ t $ is the time the money is invested for.
    \item \textbf{Compare Investment Strategies:}
    \item \textbf{Outcome:} Students will be able to compare different investment strategies by calculating and contrasting the compound interest accrued over time in scenarios with varying compounding frequencies (e.g., semiannual vs. annual).
    \item \textbf{Details:} This includes analyzing how more frequent compounding (e.g., semiannual vs. annual) results in a higher effective interest rate and a greater accumulation of interest over time, using examples such as bonds and certificates of deposit.
    \item \textbf{Derive the Effective Annual Rate (EAR):}
    \item \textbf{Outcome:} Students will be able to derive and interpret the Effective Annual Rate (EAR) from a given nominal interest rate and compounding frequency.
    \item \textbf{Details:} This includes using the formula $ \text{EAR} = \left(1 + \frac{r}{n}\right)^n - 1 $, where $ r $ is the nominal annual interest rate and $ n $ is the number of compounding periods per year. Students will understand how the EAR provides a more accurate measure of interest earned or paid annually when interest is compounded more frequently than once a year.
\end{enumerate}

\subsection{Assessment Instrument}

\begin{enumerate}
    \item \textbf{Assessment Instrument: Compound Interest Quiz} \textbf{Format:}
    \item \textbf{Type:} Mixed-format quiz (multiple-choice, short answer, and problem-solving questions)
    \item \textbf{Duration:} 45 minutes
    \item \textbf{Total Marks:} 30 points \textbf{Sections:}
    \item \textbf{Understanding Compound Interest (10 points):}
    \item \textbf{Question 1 (Multiple Choice, 2 points each):}
    \item a. Which of the following statements correctly describes compound interest?
    \item A) Interest earned on the original principal only
    \item B) Interest earned on both the original principal and previously accumulated interest
    \item C) Interest earned only after a specific period of time
    \item D) Interest earned at a fixed rate regardless of the compounding frequency
    \item b. If a bond pays an annual interest rate of 5\% compounded quarterly, which of the following is the correct expression for calculating the total amount accumulated after 1 year?
    \item A) $P(1 + 0.05)$
    \item B) $P(1 + \frac{0.05}{4})^4$
    \item C) $P(1 + \frac{0.05}{12})^{12}$
    \item D) $P(1 + 0.05)^4$
    \item c. Which factor does NOT directly affect the amount of compound interest earned?
    \item A) The principal amount
    \item B) The interest rate
    \item C) The frequency of compounding
    \item D) The type of investment
    \item d. What is the formula for the Effective Annual Rate (EAR) if the nominal interest rate is 8\% compounded monthly?
    \item A) $\left(1 + \frac{0.08}{12}\right)^{12} - 1$
    \item B) $\left(1 + 0.08\right)^{12} - 1$
    \item C) $\left(1 + \frac{0.08}{2}\right)^2 - 1$
    \item D) $\left(1 + 0.08\right) - 1$
    \item \textbf{Calculating Compound Interest (10 points):}
    \item \textbf{Question 2 (Problem-Solving, 5 points each):}
    \item a. Calculate the amount accumulated after 3 years if $1,000 is invested at an annual interest rate of 5\% compounded annually.
    \item b. An investor deposits $2,000 into a savings account that earns 4\% interest compounded quarterly. How much will the investment be worth after 2 years?
    \item \textbf{Comparing Investment Strategies (10 points):}
    \item \textbf{Question 3 (Short Answer, 5 points each):}
    \item a. Explain why an investment with a 6\% annual interest rate compounded semiannually might yield a higher return than the same investment with 6\% interest compounded annually.
    \item b. Given the nominal interest rate of 6\%, calculate and compare the Effective Annual Rate (EAR) for semiannual and quarterly compounding. Which compounding frequency offers a better return? \textbf{Scoring Rubric:}
    \item \textbf{Understanding Compound Interest:} 2 points for each correct answer (total 10 points)
    \item \textbf{Calculating Compound Interest:} Full marks for correctly applying the compound interest formula and arriving at the correct amount (5 points each, total 10 points)
    \item \textbf{Comparing Investment Strategies:} Full marks for correctly explaining the concept and accurately calculating the EAR (5 points each, total 10 points) The quiz will assess students' understanding of compound interest concepts, their ability to calculate compound interest, and their ability to compare different investment strategies based on compounding frequency. This ensures that students grasp both the theoretical and practical aspects of compound interest.
\end{enumerate}

%%===============================================================
\section{Annuities}
\label{sec:annuities}
%%===============================================================

% [Image: The present value of 1,000, 100 years into the future. Curves represent constant discount rates of 2\%, 3\%, 5\%, and 7\%.]



Figure: The present value of $\text{\$}$1,000, 100 years into the future. Curves represent constant discount rates of 2\%, 3\%, 5\%, and 7\%.



The \textbf{time value of money} is the widely accepted conjecture that there is greater benefit to receiving a sum of money now rather than an identical sum later. It may be seen as an implication of the later-developed concept of time preference.



The time value of money is among the factors considered when weighing the opportunity costs of spending rather than saving or investing money. As such, it is among the reasons why interest is paid or earned: interest, whether it is on a bank deposit or debt, compensates the depositor or lender for the loss of their use of their money. Investors are willing to forgo spending their money now only if they expect a favorable net return on their investment in the future, such that the increased value to be available later is sufficiently high to offset both the preference to spending money now and inflation (if present); see required rate of return.



\paragraph{History}



The Talmud (\textasciitilde{}500 CE) recognizes the time value of money. In Tractate Makkos page 3a the Talmud discusses a case where witnesses falsely claimed that the term of a loan was 30 days when it was actually 10 years. The false witnesses must pay the difference of the value of the loan "in a situation where he would be required to give the money back (within) thirty days\..., and that same sum in a situation where he would be required to give the money back (within) 10 years\...The difference is the sum that the testimony of the (false) witnesses sought to have the borrower lose; therefore, it is the sum that they must pay."



The notion was later described by Martín de Azpilcueta (1491–1586) of the School of Salamanca.



\paragraph{Calculations}



Time value of money problems involve the net value of cash flows at different points in time.



In a typical case, the variables might be: a balance (the real or nominal value of a debt or a financial asset in terms of monetary units), a periodic rate of interest, the number of periods, and a series of cash flows. (In the case of a debt, cash flows are payments against principal and interest; in the case of a financial asset, these are contributions to or withdrawals from the balance.) More generally, the cash flows may not be periodic but may be specified individually. Any of these variables may be the independent variable (the sought-for answer) in a given problem. For example, one may know that: the interest is 0.5\% per period (per month, say); the number of periods is 60 (months); the initial balance (of the debt, in this case) is 25,000 units; and the final balance is 0 units. The unknown variable may be the monthly payment that the borrower must pay.



For example, £100 invested for one year, earning 5\% interest, will be worth £105 after one year; therefore, £100 paid now \textit{and} £105 paid exactly one year later \textit{both} have the same value to a recipient who expects 5\% interest assuming that inflation would be zero percent. That is, £100 invested for one year at 5\% interest has a \textit{future value} of £105 under the assumption that inflation would be zero percent.



This principle allows for the valuation of a likely stream of income in the future, in such a way that annual incomes are discounted and then added together, thus providing a lump-sum "present value" of the entire income stream; all of the standard calculations for time value of money derive from the most basic algebraic expression for the present value of a future sum, "discounted" to the present by an amount equal to the time value of money. For example, the future value sum $FV$ to be received in one year is discounted at the rate of interest $r$ to give the present value sum $PV$:



$$PV = \frac{FV}{(1+r)}$$



Some standard calculations based on the time value of money are:



\begin{itemize}
    \item \textit{Present value}: The current worth of a future sum of money or stream of cash flows, given a specified rate of return. Future cash flows are "discounted" at the \textit{discount rate;} the higher the discount rate, the lower the present value of the future cash flows. Determining the appropriate discount rate is the key to valuing future cash flows properly, whether they be earnings or obligations.
    \item \textit{Present value of an annuity}: An annuity is a series of equal payments or receipts that occur at evenly spaced intervals. Leases and rental payments are examples. The payments or receipts occur at the end of each period for an ordinary annuity while they occur at the beginning of each period for an annuity due.
\end{itemize}

> \textit{Present value of a perpetuity} is an infinite and constant stream of identical cash flows.



\begin{itemize}
    \item \textit{Future value}: The value of an asset or cash at a specified date in the future, based on the value of that asset in the present.
    \item \textit{Future value of an annuity (FVA)}: The future value of a stream of payments (annuity), assuming the payments are invested at a given rate of interest.
\end{itemize}

There are several basic equations that represent the equalities listed above. The solutions may be found using (in most cases) the formulas, a financial calculator or a spreadsheet. The formulas are programmed into most financial calculators and several spreadsheet functions (such as PV, FV, RATE, NPER, and PMT).



For any of the equations below, the formula may also be rearranged to determine one of the other unknowns. In the case of the standard annuity formula, there is no closed-form algebraic solution for the interest rate (although financial calculators and spreadsheet programs can readily determine solutions through rapid trial and error algorithms).



These equations are frequently combined for particular uses. For example, bonds can be readily priced using these equations. A typical coupon bond is composed of two types of payments: a stream of coupon payments similar to an annuity, and a lump-sum return of capital at the end of the bond's maturity–-that is, a future payment. The two formulas can be combined to determine the present value of the bond.



An important note is that the interest rate \textit{i} is the interest rate for the relevant period. For an annuity that makes one payment per year, \textit{i} will be the annual interest rate. For an income or payment stream with a different payment schedule, the interest rate must be converted into the relevant periodic interest rate. For example, a monthly rate for a mortgage with monthly payments requires that the interest rate be divided by 12 (see the example below). See compound interest for details on converting between different periodic interest rates.



The rate of return in the calculations can be either the variable solved for, or a predefined variable that measures a discount rate, interest, inflation, rate of return, cost of equity, cost of debt or any number of other analogous concepts. The choice of the appropriate rate is critical to the exercise, and the use of an incorrect discount rate will make the results meaningless.



For calculations involving annuities, it must be decided whether the payments are made at the end of each period (known as an ordinary annuity), or at the beginning of each period (known as an annuity due). When using a financial calculator or a spreadsheet, it can usually be set for either calculation. The following formulas are for an ordinary annuity. For the answer for the present value of an annuity due, the PV of an ordinary annuity can be multiplied by (1 + \textit{i}).



\paragraph{Formula}



The following formula use these common variables:



\begin{itemize}
    \item \textit{PV} is the value at time zero (present value)
    \item \textit{FV} is the value at time \textit{n} (future value)
    \item \textit{A} is the value of the individual payments in each compounding period
    \item \textit{n} is the number of periods (not necessarily an integer)
    \item \textit{i} is the interest rate at which the amount compounds each period
    \item \textit{g} is the growing rate of payments over each time period
\end{itemize}

\paragraph{Future value of a present sum}



The future value (\textit{FV}) formula is similar and uses the same variables.



$$FV   \ = \  PV \cdot (1+i)^n$$



\paragraph{Present value of a future sum}



The present value formula is the core formula for the time value of money; each of the other formulae is derived from this formula. For example, the annuity formula is the sum of a series of present value calculations.



The present value (\textit{PV}) formula has four variables, each of which can be solved for by numerical methods:



$$PV \ = \ \frac{FV}{(1+i)^n}$$



The cumulative present value of future cash flows can be calculated by summing the contributions of $FV _t$, the value of cash flow at time \textit{t}:



$$PV \ = \ \sum _{t=1}^{n} \frac{FV _{t}}{(1+i)^t}$$



Note that this series can be summed for a given value of \textit{n}, or when \textit{n} is 8. This is a very general formula, which leads to several important special cases given below.



\paragraph{Present value of an annuity for n payment periods}



In this case the cash flow values remain the same throughout the \textit{n} periods. The present value of an annuity (PVA) formula has four variables, each of which can be solved for by numerical methods:



$$PV(A) \,=\,\frac{A}{i} \cdot \left[ {1-\frac{1}{\left(1+i\right)^n}} \right]$$



To get the PV of an annuity due, multiply the above equation by (1 + \textit{i}).



\paragraph{Present value of a growing annuity}



In this case each cash flow grows by a factor of (1+\textit{g}). Similar to the formula for an annuity, the present value of a growing annuity (PVGA) uses the same variables with the addition of \textit{g} as the rate of growth of the annuity (A is the annuity payment in the first period). This is a calculation that is rarely provided for on financial calculators.



\textbf{Where i ? g :}



$$PV(A)\,=\,{A \over (i-g)}\left[ 1- \left({1+g \over 1+i}\right)^n \right]$$ \textbf{Where i = g :}



$$PV(A)\,=\,{A \times n \over 1+i}$$



To get the PV of a growing annuity due, multiply the above equation by (1 + \textit{i}).



\paragraph{Present value of a perpetuity}



A perpetuity is payments of a set amount of money that occur on a routine basis and continue forever. When \textit{n} ? 8, the \textit{PV} of a perpetuity (a perpetual annuity) formula becomes a simple division.



$$PV(P) \ = \ { A \over i }$$



\paragraph{Present value of a growing perpetuity}



When the perpetual annuity payment grows at a fixed rate (\textit{g}, with \textit{g} \< \textit{i}) the value is determined according to the following formula, obtained by setting \textit{n} to infinity in the earlier formula for a growing perpetuity:



$$PV(A)\,=\,{A \over i-g}$$



In practice, there are few securities with precise characteristics, and the application of this valuation approach is subject to various qualifications and modifications. Most importantly, it is rare to find a growing perpetual annuity with fixed rates of growth and true perpetual cash flow generation. Despite these qualifications, the general approach may be used in valuations of real estate, equities, and other assets.



This is the well known Gordon growth model used for stock valuation.



\paragraph{Future value of an annuity}



The future value (after \textit{n} periods) of an annuity (FVA) formula has four variables, each of which can be solved for by numerical methods:



$$FV(A) \,=\,A\cdot\frac{\left(1+i\right)^n-1}{i}$$



To get the FV of an annuity due, multiply the above equation by (1 + i).



\paragraph{Future value of a growing annuity}



The future value (after \textit{n} periods) of a growing annuity (FVA) formula has five variables, each of which can be solved for by numerical methods:



\textbf{Where i ? g :}



$$FV(A) \,=\,A\cdot\frac{\left(1+i\right)^n-\left(1+g\right)^n}{i-g}$$ \textbf{Where i = g :}



$$FV(A) \,=\,A\cdot n(1+i)^{n-1}$$



\paragraph{Formula table}



The following table summarizes the different formulas commonly used in calculating the time value of money. These values are often displayed in tables where the interest rate and time are specified.



$$ \begin{array}{|c|c|c|} \hline \textbf{Find} & \textbf{Given} & \textbf{Formula} \\\\ \hline \text{Future value (F)} & \text{Present value (P)} & F = P \cdot (1+i)^n \\\\ \hline \text{Present value (P)} & \text{Future value (F)} & P = F \cdot (1+i)^{-n} \\\\ \hline \text{Repeating payment (A)} & \text{Future value (F)} & A = F \cdot \frac{i}{(1+i)^n-1} \\\\ \hline \text{Repeating payment (A)} & \text{Present value (P)} & A = P \cdot \frac{i(1+i)^n}{(1+i)^n-1} \\\\ \hline \text{Future value (F)} & \text{Repeating payment (A)} & F = A \cdot \frac{(1+i)^n-1}{i} \\\\ \hline \text{Present value (P)} & \text{Repeating payment (A)} & P = A \cdot \frac{(1+i)^n-1}{i(1+i)^n} \\\\ \hline \text{Future value (F)} & \text{Initial gradient payment (G)} & F = G \cdot \frac{(1+i)^n-in-1}{i^2} \\\\ \hline \text{Present value (P)} & \text{Initial gradient payment (G)} & P = G \cdot \frac{(1+i)^n-in-1}{i^2(1+i)^n} \\\\ \hline \text{Fixed payment (A)} & \text{Initial gradient payment (G)} & A = G \cdot \left[\frac{1}{i}-\frac{n}{(1+i)^n-1}\right] \\\\ \hline \text{Future value (F)} & \begin{aligned} & \text{Initial exponentially} \\\\ & \text{increasing payment (D)} \end{aligned} & F = D \cdot \frac{(1+g)^n-(1+i)^n}{g-i} \, \text{(for } i \neq g\text{)} \\\\ \text{} & \text{Increasing percentage (g)} & F = D \cdot \frac{n(1+i)^n}{1+g} \, \text{(for } i = g\text{)} \\\\ \hline \text{Present value (P)} & \begin{aligned} & \text{Initial exponentially} \\\\ & \text{increasing payment (D)} \end{aligned} & P = D \cdot \frac{\left(\frac{1+g}{1+i}\right)^n-1}{g-i} \, \text{(for } i \neq g\text{)} \\\\ \text{} & \text{Increasing percentage (g)} & P = D \cdot \frac{n}{1+g} \, \text{(for } i = g\text{)} \\\\ \hline \end{array} $$



Notes:



\begin{itemize}
    \item \textit{A} is a fixed payment amount, every period
    \item \textit{G} is the initial payment amount of an increasing payment amount, that starts at \textit{G} and increases by \textit{G} for each subsequent period.
    \item \textit{D} is the initial payment amount of an exponentially (geometrically) increasing payment amount, that starts at \textit{D} and increases by a factor of (1+\textit{g}) each subsequent period.
\end{itemize}

\paragraph{Derivations}



\paragraph{Annuity derivation}



The formula for the present value of a regular stream of future payments (an annuity) is derived from a sum of the formula for future value of a single future payment, as below, where \textit{C} is the payment amount and \textit{n} the period.



A single payment C at future time \textit{m} has the following future value at future time \textit{n}:



$$FV \ = C(1+i)^{n-m}$$



Summing over all payments from time 1 to time n, then reversing t



$$FVA \ = \sum _{m=1}^n C(1+i)^{n-m} \ = \sum _{k=0}^{n-1} C(1+i)^k$$



Note that this is a geometric series, with the initial value being \textit{a} = \textit{C}, the multiplicative factor being 1 + \textit{i}, with \textit{n} terms. Applying the formula for geometric series, we get



$$FVA \ = \frac{ C ( 1 - (1+i)^n )}{1 - (1+i)} \ = \frac{ C ( 1 - (1+i)^n )}{-i}$$



The present value of the annuity (PVA) is obtained by simply dividing by $(1+i)^n$:



$$PVA \ = \frac{FVA}{(1+i)^n} = \frac{C}{i} \left( 1 - \frac{1}{(1+i)^n} \right)$$



Another simple and intuitive way to derive the future value of an annuity is to consider an endowment, whose interest is paid as the annuity, and whose principal remains constant. The principal of this hypothetical endowment can be computed as that whose interest equals the annuity payment amount:



$$\text{Principal} \times i = C$$



$$\text{Principal} = \frac{C}{i} + \text{goal}$$



Note that no money enters or leaves the combined system of endowment principal + accumulated annuity payments, and thus the future value of this system can be computed simply via the future value formula:



$$FV = PV(1+i)^n$$



Initially, before any payments, the present value of the system is just the endowment principal, $PV = \frac{C}{i}$. At the end, the future value is the endowment principal (which is the same) plus the future value of the total annuity payments ($FV = \frac{C}{i} + FVA$). Plugging this back into the equation:



$$\frac{C}{i} + FVA = \frac{C}{i} (1+i)^n$$



$$FVA = \frac{C}{i} \left[ \left(1+i \right)^n - 1 \right]$$



\paragraph{Perpetuity derivation}



Without showing the formal derivation here, the perpetuity formula is derived from the annuity formula. Specifically, the term:



$$\left({1 - {1 \over { (1+i)^n } }}\right)$$ can be seen to approach the value of 1 as \textit{n} grows larger. At infinity, it is equal to 1, leaving ${C \over i}$ as the only term remaining.



\paragraph{Continuous compounding}



Rates are sometimes converted into the continuous compound interest rate equivalent because the continuous equivalent is more convenient (for example, more easily differentiated). Each of the formula above may be restated in their continuous equivalents. For example, the present value at time 0 of a future payment at time \textbf{t} can be restated in the following way, where \textbf{e} is the base of the natural logarithm and \textbf{r} is the continuously compounded rate:



$$\text{PV}  = \text{FV}\cdot e^{-rt}$$ This can be generalized to discount rates that vary over time: instead of a constant discount rate \textit{r,} one uses a function of time \textit{r}(\textit{t}). In that case the discount factor, and thus the present value, of a cash flow at time \textit{T} is given by the integral of the continuously compounded rate \textit{r}(\textit{t}):



$$\text{PV}  = \text{FV}\cdot \exp\left(-\int_0^T r(t)\,dt\right)$$ Indeed, a key reason for using continuous compounding is to simplify the analysis of varying discount rates and to allow one to use the tools of calculus. Further, for interest accrued and capitalized overnight (hence compounded daily), continuous compounding is a close approximation for the actual daily compounding. More sophisticated analysis includes the use of differential equations, as detailed below.



\paragraph{Examples}



Using continuous compounding yields the following formulas for various instruments:



Annuity:

$$\ PV \ = \ {A(1-e^{-rt}) \over e^r -1}$$



Perpetuity:

$$\ PV  \ = \  {A \over e^r - 1}$$



Growing annuity:

$$\ PV  \ = \  {Ae^{-g}(1-e^{-(r-g)t}) \over e^{(r-g)} - 1}$$



Growing perpetuity:

$$\ PV  \ = \  {Ae^{-g} \over e^{(r-g)} - 1}$$



Annuity with continuous payments:

$$\ PV  \ = \  { 1 - e^{(-rt)} \over r }$$



These formulas assume that payment A is made in the first payment period and annuity ends at time t.



\paragraph{Resources}



\begin{itemize}
    \item \href{https://mathresearch.utsa.edu/wikiFiles/MAT1053/Annuities/MAT1053_M7.2Annuities.pdf}{Annuities}, Book Chapter
    \item \href{https://mathresearch.utsa.edu/wikiFiles/MAT1053/Annuities/MAT1053_M7.2AnnuitiesGN.pdf}{Guided Notes}
\end{itemize}

\paragraph{Licensing}



Content obtained and/or adapted from:



\begin{itemize}
    \item \href{https://en.wikipedia.org/wiki/Time_value_of_money}{Time value of money, Wikipedia} under a CC BY-SA license
\end{itemize}

\subsection{Learning Objectives}

\begin{enumerate}
    \item Here are three Student Learning Objectives (SLOs) related to the time value of money concept, using the format provided:
    \item \textbf{Calculate Present Value:}
    \item \textbf{Outcome:} Students will be able to calculate the present value of future cash flows using various discount rates.
    \item \textbf{Details:} This includes applying the present value formula $$PV = \frac{FV}{(1 + r)^n}$$ to determine the current worth of a future sum under different interest rates, such as 2\%, 3\%, 5\%, and 7\%, and interpreting the impact of varying discount rates on the present value.
    \item \textbf{Analyze Time Value of Money:}
    \item \textbf{Outcome:} Students will be able to analyze and explain the concept of the time value of money and its implications for financial decision-making.
    \item \textbf{Details:} This involves understanding how interest rates, time periods, and future cash flows influence present value, and applying this knowledge to assess investment opportunities and evaluate financial decisions, such as whether to invest a sum of money now or later.
    \item \textbf{Compare Future and Present Values:}
    \item \textbf{Outcome:} Students will be able to compare and contrast future values and present values of annuities, lump-sum payments, and perpetuities.
    \item \textbf{Details:} This includes solving problems involving the future value formula $$FV = PV \times (1 + r)^n$$ and understanding how the choice of discount rate and time period affects the future and present values of financial assets, particularly in the context of long-term investments. These SLOs are designed to ensure students understand and can apply the fundamental concepts of the time value of money, critical for making informed financial decisions.
\end{enumerate}

\subsection{Assessment Instrument}

\begin{enumerate}
    \item \textbf{Assessment Instrument: Annuities Quiz} \textbf{Format:}
    \item \textbf{Type:} Mixed-format quiz (multiple-choice, short answer, and problem-solving questions)
    \item \textbf{Duration:} 45 minutes
    \item \textbf{Total Marks:} 30 points \textbf{Sections:}
    \item \textbf{Concept of Annuities (10 points):}
    \item \textbf{Question 1:} Explain what an annuity is and how it differs from a lump-sum payment.
    \item \textbf{(Short Answer, 5 points)}
    \item \textbf{Question 2:} Which of the following is an example of an annuity?
    \item a. A one-time payment of $10,000 received today.
    \item b. Monthly mortgage payments.
    \item c. A single $5,000 payment received 10 years from now.
    \item d. Quarterly dividends received from stock investments.
    \item \textbf{(Multiple Choice, 5 points)}
    \item \textbf{Calculating Present Value of an Annuity (10 points):}
    \item \textbf{Question 3:} Calculate the present value of an annuity that pays `$`1,000 per year for 5 years at an interest rate of 5\%. Use the formula: $$PV(A) = \frac{A}{i} \cdot \left[1 - \frac{1}{(1+i)^n}\right]$$
    \item \textbf{(Problem-solving, 10 points)}
    \item \textbf{Comparing Ordinary Annuities and Annuities Due (10 points):}
    \item \textbf{Question 4:} Explain the difference between an ordinary annuity and an annuity due. Provide an example of each.
    \item \textbf{(Short Answer, 5 points)}
    \item \textbf{Question 5:} If an ordinary annuity pays $500 at the end of each year for 4 years, what would be the present value if the interest rate is 3\%? How does this value compare to an annuity due with the same payment and interest rate?
    \item \textbf{(Problem-solving, 5 points)} \textbf{Scoring Rubric:}
    \item \textbf{Concept of Annuities:} 5 points for a clear and accurate explanation; 5 points for the correct choice in the multiple-choice question (total 10 points).
    \item \textbf{Calculating Present Value of an Annuity:} 10 points for correctly solving the present value of the annuity using the provided formula (total 10 points).
    \item \textbf{Comparing Ordinary Annuities and Annuities Due:} 5 points for correctly explaining the difference with examples; 5 points for solving the present value and comparing the two types of annuities (total 10 points). The quiz will assess students' understanding of the concept of annuities, their ability to calculate the present value, and their knowledge of the differences between ordinary annuities and annuities due. It ensures that students can apply the necessary formulas and clearly articulate the differences in these financial concepts.
\end{enumerate}

%%===============================================================
\section{Payout Annuities}
\label{sec:payout-annuities}
%%===============================================================

\textbf{Payout annuities} are typically used after retirement. With a payout annuity, you start with money in the account, and pull money out of the account on a regular basis. Any remaining money in the account earns interest. After a fixed amount of time, the account will end up empty. Perhaps you have saved $\text{\$}$500,000 for retirement, and want to take money out of the account each month to live on. You want the money to last you 20 years. This is a payout annuity. The formula is derived in a similar way as we did for savings annuities. The details are omitted here.



$$P _{0} = \frac{d \left( 1 - \left( 1 + \frac{r}{k} \right)^{-Nk} \right)}{\left( \frac{r}{k} \right)}$$



where:



\begin{itemize}
    \item $P _{0}$ is the balance in the account at the beginning (starting amount, or principal).
    \item $d$ is the regular withdrawal (the amount you take out each year, each month, etc.).
    \item $r$ is the annual interest rate (in decimal form. Example: 5\% = 0.05).
    \item $k$ is the number of compounding periods in one year.
    \item $N$ is the number of years we plan to take withdrawals.
\end{itemize}

Like with annuities, the compounding frequency is not always explicitly given, but is determined by how often you take the withdrawals.



\textbf{When do you use this}

Payout annuities assume that you take money from the account on a regular schedule (every month, year, quarter, etc.) and let the rest sit there earning interest.

\begin{itemize}
    \item Compound interest: One deposit
    \item Annuity: Many deposits.
    \item Payout Annuity: Many withdrawals
\end{itemize}

\paragraph{Example}



After retiring, you want to be able to take $\text{\$}$1000 every month for a total of 20 years from your retirement account. The account earns 6\% interest. How much will you need in your account when you retire?







In this example,

\begin{itemize}
    \item d = $\text{\$}$1000 (the monthly withdrawal)
\end{itemize}

\begin{itemize}
    \item r = 0.06 (6\% annual rate)
\end{itemize}

\begin{itemize}
    \item k = 12 (since we're doing monthly withdrawals, we'll compound monthly)
\end{itemize}

\begin{itemize}
    \item N = 20 (since were taking withdrawals for 20 years)
\end{itemize}

We're looking for $P _0$; how much money needs to be in the account at the beginning.



Putting this into the equation:

$$\begin{aligned} & {{P} _{0}}=\frac{1000\left(1-{{\left(1+\frac{0.06}{12}\right)}^{-20(12)}}\right)}{\left(\frac{0.06}{12}\right)} \\\\ & {{P} _{0}}=\frac{1000\times\left(1-{{\left(1.005\right)}^{-240}}\right)}{\left(0.005\right)} \\\\ & {{P} _{0}}=\frac{1000\times\left(1-0.302\right)}{\left(0.005\right)}=\text{\$}139,600 \\    \end{aligned}$$



You will need to have $\text{\$}$139,600 in your account when you retire.



Notice that you withdrew a total of $\text{\$}$240,000 ($\text{\$}$1000 a month for 240 months). The difference between what you pulled out and what you started with is the interest earned. In this case it is $\text{\$}$240,000 – $\text{\$}$139,600 = $\text{\$}$100,400 in interest.



You know you will have $\text{\$}$500,000 in your account when you retire. You want to be able to take monthly withdrawals from the account for a total of 30 years. Your retirement account earns 8\% interest. How much will you be able to withdraw each month?



\textbf{Example}



In this example, we're looking for d.

\begin{itemize}
    \item r = 0.08 (8\% annual rate)
    \item k = 12 (since we're withdrawing monthly)
    \item N = 30 (30 years)
    \item $P _0$ = $\text{\$}$500,000 (we are beginning with $\text{\$}$500,000)
\end{itemize}

In this case, we're going to have to set up the equation, and solve for d.

$$\begin{aligned} & 500,000=\frac{d\left(1-{{\left(1+\frac{0.08}{12}\right)}^{-30(12)}}\right)}{\left(\frac{0.08}{12}\right)} \\\\ & 500,000=\frac{d\left(1-{{\left(1.00667\right)}^{-360}}\right)}{\left(0.00667\right)} \\\\ & 500,000=d(136.232) \\\\  &d=\frac{500,000}{136.232}=\$3670.21  \\\\    \end{aligned}$$

You would be able to withdraw $\text{\$}$3,670.21 each month for 30 years.



\paragraph{Licensing}



Content obtained and/or adapted from:



\begin{itemize}
    \item \href{https://courses.lumenlearning.com/math4libarts/chapter/payout-annuities/}{Payout Annuities, Lumen Learning} under a CC BY-SA license
\end{itemize}

\subsection{Learning Objectives}

\begin{enumerate}
    \item Here are three Student Learning Objectives (SLOs) related to the concept of payout annuities:
    \item \textbf{Calculate the Initial Principal Required for Payout Annuities:}
    \item \textbf{Outcome:} Students will be able to determine the initial amount of money needed in a retirement account to ensure regular withdrawals over a specified period, considering the interest rate and compounding frequency.
    \item \textbf{Details:} This includes applying the formula for payout annuities to calculate the required starting principal ($P _0$) based on given variables such as the withdrawal amount, interest rate, compounding periods, and the duration of withdrawals.
    \item \textbf{Determine the Regular Withdrawal Amount from a Payout Annuity:}
    \item \textbf{Outcome:} Students will be able to calculate the regular withdrawal amount ($d$) that can be sustained from a retirement account over a specified period, given an initial principal, interest rate, and compounding frequency.
    \item \textbf{Details:} This involves manipulating the payout annuity formula to solve for the withdrawal amount ($d$), ensuring that the account balance reaches zero at the end of the withdrawal period.
    \item \textbf{Analyze the Effect of Interest Rates and Compounding Frequency on Payout Annuities:}
    \item \textbf{Outcome:} Students will be able to analyze how different interest rates and compounding frequencies impact the sustainability of regular withdrawals from a retirement account.
    \item \textbf{Details:} This includes comparing scenarios with varying interest rates and compounding frequencies to understand their effects on both the initial principal needed and the regular withdrawal amount, helping students make informed decisions about their retirement planning.
\end{enumerate}

\subsection{Assessment Instrument}

\begin{enumerate}
    \item \textbf{Assessment Instrument: Payout Annuities Quiz} \textbf{Format:}
    \item \textbf{Type:} Mixed-format quiz (multiple-choice, short answer, and problem-solving questions)
    \item \textbf{Duration:} 45 minutes
    \item \textbf{Total Marks:} 30 points \textbf{Sections:}
    \item \textbf{Understanding the Payout Annuity Formula (10 points):}
    \item \textbf{Question 1 (Multiple Choice):} Which of the following best describes the variable $P _0$ in the payout annuity formula?
    \item a. The annual interest rate
    \item b. The total amount withdrawn over the annuity period
    \item c. The initial principal or starting amount in the account
    \item d. The regular withdrawal amount
    \item \textbf{Question 2 (Short Answer):} Explain the significance of the compounding frequency ($k$) in the payout annuity formula.
    \item \textbf{Question 3 (Multiple Choice):} If the annual interest rate ($r$) increases, what happens to the initial principal $P _0$ required for a fixed withdrawal amount?
    \item a. Increases
    \item b. Decreases
    \item c. Stays the same
    \item d. Not enough information to determine
    \item \textbf{Calculating the Initial Principal (10 points):}
    \item \textbf{Question 4 (Problem-Solving):} You plan to withdraw $\text{\$}$1,500 per month for 25 years from your retirement account, which earns 5% interest compounded monthly. Calculate the initial amount ($P _0$) needed in the account at the beginning. (Show all steps of your calculation.)
    \item \textbf{Question 5 (Multiple Choice):} If the number of years for withdrawals ($N$) increases while all other factors remain the same, what is the effect on the initial principal $P _0$?
    \item a. Increases
    \item b. Decreases
    \item c. Stays the same
    \item d. Not enough information to determine
    \item \textbf{Determining Regular Withdrawal Amount (10 points):}
    \item \textbf{Question 6 (Problem-Solving):} You have saved $\text{\$}$200,000 for retirement and want to withdraw money monthly for 15 years. Your account earns 7% interest compounded monthly. Calculate the amount you can withdraw each month ($d$). (Show all steps of your calculation.)
    \item \textbf{Question 7 (Multiple Choice):} If the interest rate ($r$) decreases, what happens to the monthly withdrawal amount $d$ for a fixed principal $P _0$?
    \item a. Increases
    \item b. Decreases
    \item c. Stays the same
    \item d. Not enough information to determine \textbf{Scoring Rubric:}
    \item \textbf{Understanding the Payout Annuity Formula:} 3 points for each correct multiple-choice answer (total 6 points), 4 points for a correct short answer explanation (total 10 points)
    \item \textbf{Calculating the Initial Principal:} 6 points for a correct problem-solving answer (including all steps shown), 4 points for a correct multiple-choice answer (total 10 points)
    \item \textbf{Determining Regular Withdrawal Amount:} 6 points for a correct problem-solving answer (including all steps shown), 4 points for a correct multiple-choice answer (total 10 points) The quiz will assess the students' understanding of the payout annuity formula, their ability to calculate the necessary initial principal for a specified withdrawal plan, and their ability to determine the regular withdrawal amount based on given financial conditions.
\end{enumerate}

%%===============================================================
\section{Loans}
\label{sec:loans}
%%===============================================================

In finance, a \textbf{loan} is the lending of money by one or more individuals, organizations, or other entities to other individuals, organizations etc. The recipient (i.e., the borrower) incurs a debt and is usually liable to pay interest on that debt until it is repaid as well as to repay the principal amount borrowed.



The document evidencing the debt (e.g., a promissory note) will normally specify, among other things, the principal amount of money borrowed, the interest rate the lender is charging, and the date of repayment. A loan entails the reallocation of the subject asset(s) for a period of time, between the lender and the borrower.



The interest provides an incentive for the lender to engage in the loan. In a legal loan, each of these obligations and restrictions is enforced by contract, which can also place the borrower under additional restrictions known as loan covenants. Although this article focuses on monetary loans, in practice, any material object might be lent.



Acting as a provider of loans is one of the main activities of financial institutions such as banks and credit card companies. For other institutions, issuing of debt contracts such as bonds is a typical source of funding.



\paragraph{Personal loan}



\paragraph{Secured}



A secured loan is a form of debt in which the borrower pledges some asset (i.e., a car, a house) as collateral.



A mortgage loan is a very common type of loan, used by many individuals to purchase residential property. The lender, usually a financial institution, is given security – a lien on the title to the property – until the mortgage is paid off in full. In the case of home loans, if the borrower defaults on the loan, the bank would have the legal right to repossess the house and sell it, to recover sums owing to it.



Similarly, a loan taken out to buy a car may be secured by the car. The duration of the loan is much shorter – often corresponding to the useful life of the car. There are two types of auto loans, direct and indirect. In a direct auto loan, a bank lends the money directly to a consumer. In an indirect auto loan, a car dealership (or a connected company) acts as an intermediary between the bank or financial institution and the consumer.



Other forms of secured loans include loans against securities - such as shares, mutual funds, bonds, etc. This particular instrument issues customers a line of credit based on the quality of the securities pledged. Gold loans are issued to customers after evaluating the quantity and quality of gold in the items pledged. Corporate entities can also take out secured lending by pledging the company's assets, including the company itself. The interest rates for secured loans are usually lower than those of unsecured loans. Usually, the lending institution employs people (on a roll or on a contract basis) to evaluate the quality of pledged collateral before sanctioning the loan.



\paragraph{Unsecured}



Unsecured loans are monetary loans that are not secured against the borrower's assets. These may be available from financial institutions under many different guises or marketing packages:



\begin{itemize}
    \item Credit cards
    \item Personal loans
    \item Bank overdrafts
    \item Credit facilities or lines of credit
    \item Corporate bonds (may be secured or unsecured)
    \item Peer-to-peer lending
\end{itemize}

The interest rates applicable to these different forms may vary depending on the lender and the borrower. These may or may not be regulated by law. In the United Kingdom, when applied to individuals, these may come under the Consumer Credit Act 1974.



Interest rates on unsecured loans are nearly always higher than for secured loans because an unsecured lender's options for recourse against the borrower in the event of default are severely limited, subjecting the lender to higher risk compared to that encountered for a secured loan. An unsecured lender must sue the borrower, obtain a money judgment for breach of contract, and then pursue execution of the judgment against the borrower's unencumbered assets (that is, the ones not already pledged to secured lenders). In insolvency proceedings, secured lenders traditionally have priority over unsecured lenders when a court divides up the borrower's assets. Thus, a higher interest rate reflects the additional risk that in the event of insolvency, the debt may be uncollectible.



\paragraph{Demand}



Demand loans are short-term loans that typically do not have fixed dates for repayment. Instead, demand loans carry a floating interest rate, which varies according to the prime lending rate or other defined contract terms. Demand loans can be "called" for repayment by the lending institution at any time. Demand loans may be unsecured or secured.



\paragraph{Subsidized}



A subsidized loan is a loan on which the interest is reduced by an explicit or hidden subsidy. In the context of college loans in the United States, it refers to a loan on which no interest is accrued while a student remains enrolled in education.



\paragraph{Concessional}



A concessional loan, sometimes called a "soft loan", is granted on terms substantially more generous than market loans either through below-market interest rates, by grace periods, or a combination of both. Such loans may be made by foreign governments to developing countries or may be offered to employees of lending institutions as an employee benefit (sometimes called a \textit{perk}).



\paragraph{Target markets}



Loans can also be categorized according to whether the debtor is an individual person (consumer) or a business.



\paragraph{Personal}



Common personal loans include mortgage loans, car loans, home equity lines of credit, credit cards, installment loans, and payday loans. The credit score of the borrower is a major component in and underwriting and interest rates (APR) of these loans. The monthly payments of personal loans can be decreased by selecting longer payment terms, but overall interest paid increases as well. A personal loan can be obtained from banks, alternative (non-bank) lenders, online loan providers and private lenders.



\paragraph{Commercial}



Loans to businesses are similar to the above but also include commercial mortgages and corporate bonds and government guaranteed loans. Underwriting is not based upon credit score but rather credit rating.



\paragraph{Loan payment}



The most typical loan payment type is the fully amortizing payment in which each monthly rate has the same value over time.



The fixed monthly payment \textit{P} for a loan of \textit{L} for \textit{n} months and a monthly interest rate \textit{c} is:



$$P = L \cdot \frac{c\,(1 + c)^n}{(1 + c)^n - 1}$$



\paragraph{Abuses in lending}



Predatory lending is one form of abuse in the granting of loans. It usually involves granting a loan in order to put the borrower in a position that one can gain advantage over them; subprime mortgage-lending and payday-lending are two examples, where the moneylender is not authorized or regulated, the lender could be considered a loan shark.



Usury is a different form of abuse, where the lender charges excessive interest. In different time periods and cultures, the acceptable interest rate has varied, from no interest at all to unlimited interest rates. Credit card companies in some countries have been accused by consumer organizations of lending at usurious interest rates and making money out of frivolous "extra charges".



Abuses can also take place in the form of the customer defrauding the lender by borrowing without intending to repay the loan.



\paragraph{United States taxes}



Most of the basic rules governing how loans are handled for tax purposes in the United States are codified by both Congress (the Internal Revenue Code) and the Treasury Department (Treasury Regulations – another set of rules that interpret the Internal Revenue Code).



1\. A loan is not gross income to the borrower. Since the borrower has the obligation to repay the loan, the borrower has no accession to wealth.



2\. The lender may not deduct (from own gross income) the amount of the loan. The rationale here is that one asset (the cash) has been converted into a different asset (a promise of repayment). Deductions are not typically available when an outlay serves to create a new or different asset.



3\. The amount paid to satisfy the loan obligation is not deductible (from own gross income) by the borrower.



4\. Repayment of the loan is not gross income to the lender. In effect, the promise of repayment is converted back to cash, with no accession to wealth by the lender.



5\. Interest paid to the lender is included in the lender's gross income. Interest paid represents compensation for the use of the lender's money or property and thus represents profit or an accession to wealth to the lender. Interest income can be attributed to lenders even if the lender doesn't charge a minimum amount of interest.



6\. Interest paid to the lender may be deductible by the borrower. In general, interest paid in connection with the borrower's business activity is deductible, while interest paid on personal loans are not deductible. The major exception here is interest paid on a home mortgage.



\paragraph{Income from discharge of indebtedness}



Although a loan does not start out as income to the borrower, it becomes income to the borrower if the borrower is discharged of indebtedness. Thus, if a debt is discharged, then the borrower essentially has received income equal to the amount of the indebtedness. The Internal Revenue Code lists "Income from Discharge of Indebtedness" in Section 61(a)(12) as a source of gross income.



Example: X owes Y $\text{\$}$50,000. If Y discharges the indebtedness, then X no longer owes Y $\text{\$}$50,000. For purposes of calculating income, this is treated the same way as if Y gave X $\text{\$}$50,000.



For a more detailed description of the "discharge of indebtedness", look at Section 108 (Cancellation of Debt (COD) Income) of the Internal Revenue Code.



\paragraph{Resources}



\begin{itemize}
    \item \href{https://mathresearch.utsa.edu/wikiFiles/MAT1053/Loans/MAT1053_M8.1Loans.pdf}{Loans}, Book Chapter
    \item \href{https://mathresearch.utsa.edu/wikiFiles/MAT1053/Loans/MAT1053_M8.1LoansGN.pdf}{Guided Notes}
\end{itemize}

\paragraph{Licensing}



Content obtained and/or adapted from:



\begin{itemize}
    \item \href{https://en.wikipedia.org/wiki/Loan}{Loan, Wikipedia} under a CC BY-SA license
\end{itemize}

\subsection{Learning Objectives}

\begin{enumerate}
    \item Here are three Student Learning Objectives (SLOs) based on the provided information about loans:
    \item \textbf{Differentiate Between Secured and Unsecured Loans:}
    \item \textbf{Outcome:} Students will be able to distinguish between secured and unsecured loans by identifying the types of collateral involved and the implications of these distinctions for both the lender and the borrower.
    \item \textbf{Details:} This includes understanding the risks and benefits associated with each type of loan, the interest rates typically applied, and the legal consequences in the event of default.
    \item \textbf{Analyze the Structure of Loan Repayment:}
    \item \textbf{Outcome:} Students will be able to calculate loan repayment schedules using the amortization formula, recognizing the impact of different interest rates and loan terms on the overall cost of the loan.
    \item \textbf{Details:} This includes the ability to interpret the formula for fully amortizing payments, apply it to various loan scenarios, and explain the financial implications of different loan terms (e.g., short-term vs. long-term loans).
    \item \textbf{Identify and Evaluate Loan Abuses:}
    \item \textbf{Outcome:} Students will be able to identify common forms of abuse in lending practices, such as predatory lending and usury, and evaluate their legal and ethical implications.
    \item \textbf{Details:} This includes understanding the regulatory frameworks that govern lending practices, recognizing signs of unethical behavior by lenders, and analyzing case studies where loan abuse has led to legal action or reform.
\end{enumerate}

\subsection{Assessment Instrument}

\begin{enumerate}
    \item \textbf{Assessment Instrument: Loans Quiz} \textbf{Format:}
    \item \textbf{Type:} Mixed-format quiz (multiple-choice, short answer, and problem-solving questions)
    \item \textbf{Duration:} 45 minutes
    \item \textbf{Total Marks:} 30 points \textbf{Sections:}
    \item \textbf{Types of Loans (10 points):}
    \item \textbf{Question 1:} (Multiple Choice) Which of the following is an example of a secured loan?
    \item a. Credit card
    \item b. Personal loan without collateral
    \item c. Mortgage
    \item d. Bank overdraft
    \item \textbf{Question 2:} (Short Answer) Explain the difference between a secured loan and an unsecured loan. Provide one example of each.
    \item \textbf{Loan Repayment Calculation (10 points):}
    \item \textbf{Question 3:} (Problem-Solving) Calculate the monthly payment for a loan of `$`20,000 with an interest rate of 5% per annum, to be repaid over 10 years. Use the amortization formula: $P = L \cdot \frac{c\,(1 + c)^n}{(1 + c)^n - 1}$ where $P$ is the monthly payment, $L$ is the loan amount, $c$ is the monthly interest rate, and $n$ is the number of months.
    \item \textbf{Question 4:} (Short Answer) If a borrower takes out a loan of `$`10,000 with a fixed interest rate of 6\% per annum for 5 years, explain how the loan's total cost would change if the repayment period was extended to 10 years.
    \item \textbf{Ethical and Legal Considerations in Lending (10 points):}
    \item \textbf{Question 5:} (Multiple Choice) Which of the following is a characteristic of predatory lending?
    \item a. Offering loans at a below-market interest rate
    \item b. Charging excessive interest rates and fees
    \item c. Providing transparent loan terms and conditions
    \item d. Offering grace periods for loan repayment
    \item \textbf{Question 6:} (Short Answer) Define usury and explain why it is considered unethical. How do regulatory frameworks typically address this issue?
    \item \textbf{Question 7:} (Problem-Solving) A borrower takes out a payday loan of `$`500 with a 15\% interest rate, payable in two weeks. Calculate the total amount the borrower must repay at the end of the two weeks. Discuss why payday loans might be considered predatory. \textbf{Scoring Rubric:}
    \item \textbf{Types of Loans:} 5 points for each correct answer/explanation (total 10 points)
    \item \textbf{Loan Repayment Calculation:} 5 points for each correct calculation/explanation (total 10 points)
    \item \textbf{Ethical and Legal Considerations in Lending:} 2.5 points for each correct multiple-choice answer, 5 points for the short answer, 5 points for the problem-solving question (total 10 points) The quiz assesses students' understanding of different types of loans, their ability to calculate loan repayments using the amortization formula, and their comprehension of the ethical and legal considerations in lending. It ensures that students can apply their knowledge in practical scenarios related to finance.
\end{enumerate}

%%===============================================================
\section{Promissory Notes}
\label{sec:promissory-notes}
%%===============================================================

% [Image: A 1926 promissory note from the Imperial Bank of India, Rangoon, Burma for 20,000 rupees plus interest]



Figure: A 1926 promissory note from the Imperial Bank of India, Rangoon, Burma for 20,000 rupees plus interest. 



A \textbf{promissory note}, sometimes referred to as a \textbf{note payable}, is a legal instrument (more particularly, a financing instrument and a debt instrument), in which one party (the \textit{maker} or \textit{issuer}) promises in writing to pay a determinate sum of money to the other (the \textit{payee}), either at a fixed or determinable future time or on demand of the payee, under specific terms.



\paragraph{Overview}



The terms of a note usually include the principal amount, the interest rate if any, the parties, the date, the terms of repayment (which could include interest) and the maturity date. Sometimes, provisions are included concerning the payee's rights in the event of a default, which may include foreclosure of the maker's assets. In foreclosures and contract breaches, promissory notes under CPLR 5001 allow creditors to recover prejudgement interest from the date interest is due until liability is established. For loans between individuals, writing and signing a promissory note are often instrumental for tax and record keeping. A promissory note alone is typically unsecured.



\paragraph{Terminology}



The term note payable is commonly used in accounting (as distinguished from accounts payable) or commonly as just a "note", it is internationally defined by the \textit{Convention providing a uniform law for bills of exchange and promissory notes}, but regional variations exist. A banknote is frequently referred to as a promissory note, as it is made by a bank and payable to bearer on demand. Mortgage notes are another prominent example.



If the promissory note is unconditional and readily saleable, it is called a negotiable instrument.



\textbf{Demand promissory notes} are notes that do not carry a specific maturity date, but are due on demand of the lender. Usually the lender will only give the borrower a few days' notice before the payment is due.



Promissory notes may be used in combination with security agreements. For example, a promissory note may be used in combination with a mortgage, in which case it is called a mortgage note.



\paragraph{Loan contracts}



In common speech, other terms, such as "loan", "loan agreement", and "loan contract" may be used interchangeably with "promissory note". The term "loan contract" is often used to describe a contract that is lengthy and detailed.



A promissory note is very similar to a loan. Each is a legally binding contract to unconditionally repay a specified amount within a defined time frame. However, a promissory note is generally less detailed and less rigid than a loan contract. For one thing, loan agreements often require repayment in installments, while promissory notes typically do not. Furthermore, a loan agreement usually includes the terms for recourse in the case of default, such as establishing the right to foreclose, while a promissory note does not.



\paragraph{Difference from IOU}



Promissory notes differ from IOUs in that they contain a specific promise to pay along with the steps and timeline for repayment as well as consequences if repayment fails. IOUs only acknowledge that a debt exists.



\paragraph{Negotiability}



Negotiable instruments are unconditional and impose few to no duties on the issuer or payee other than payment. In the United States, whether a promissory note is a negotiable instrument can have significant legal impacts, as only negotiable instruments are subject to Article 3 of the Uniform Commercial Code and the application of the holder in due course rule. The negotiability of mortgage notes has been debated, particularly due to the obligations and "baggage" associated with mortgages; however, in mortgages notes are often determined to be negotiable instruments.



In the United States, the Non-Negotiable Long Form Promissory Note is not required.



\paragraph{Use as financial instruments}



Promissory notes are a common financial instrument in many jurisdictions, employed principally for short time financing of companies. Often, the seller or provider of a service is not paid upfront by the buyer (usually, another company), but within a period of time, the length of which has been agreed upon by both the seller and the buyer. The reasons for this may vary; historically, many companies used to balance their books and execute payments and debts at the end of each week or tax month; any product bought before that time would be paid only then. Depending on the jurisdiction, this deferred payment period can be regulated by law; in countries like France, Italy or Spain, it usually ranges between 30 and 90 days after the purchase.



When a company engages in many of such transactions, for instance by having provided services to many customers all of whom then deferred their payment, it is possible that the company may be owed enough money that its own liquidity position (i.e., the amount of cash it holds) is hampered, and finds itself unable to honour their own debts, despite the fact that by the books, the company remains solvent. In those cases, the company has the option of asking the bank for a short-term loan, or using any other such short-term financial arrangements to avoid insolvency. However, in jurisdictions where promissory notes are commonplace, the company (called the \textit{payee} or \textit{lender}) can ask one of its debtors (called the \textit{maker}, \textit{borrower} or \textit{payor}) to \textit{accept} a promissory note, whereby the maker signs a legally binding agreement to honour the amount established in the promissory note (usually, part or all its debt) within the agreed period of time. The lender can then take the promissory note to a financial institution (usually a bank, albeit this could also be a private person, or another company), that will exchange the promissory note for cash; usually, the promissory note is cashed in for the amount established in the promissory note, less a small discount.



Once the promissory note reaches its \textit{maturity date}, its current holder (the bank) can execute it over the emitter of the note (the debtor), who would have to pay the bank the amount \textit{promised} in the note. If the maker fails to pay, however, the bank retains the right to go to the company that cashed the promissory note in, and demand payment. In the case of unsecured promissory notes, the lender accepts the promissory note based solely on the maker's ability to repay; if the maker fails to pay, the lender must honour the debt to the bank. In the case of a secured promissory note, the lender accepts the promissory note based on the maker's ability to repay, but the note is secured by a thing of value; if the maker fails to pay and the bank reclaims payment, the lender has the right to execute the security.



\paragraph{Use as private money}



Thus, promissory notes can work as a form of private money. In the past, particularly during the 19th century, their widespread and unregulated use was a source of great risk for banks and private financiers, who would often face the insolvency of both debtors, or simply be scammed by both.



\paragraph{Resources}



\begin{itemize}
    \item \href{https://mathresearch.utsa.edu/wikiFiles/MAT1053/Promisory_Notes/MAT1053_M8.2Promisory_Notes.pdf}{Promissory Notes}, Book Chapter
    \item \href{https://mathresearch.utsa.edu/wikiFiles/MAT1053/Promisory_Notes/MAT1053_M8.2Promisory_NotesGN.pdf}{Guided Notes}
\end{itemize}

\paragraph{Licensing}



Content obtained and/or adapted from:



\begin{itemize}
    \item \href{https://en.wikipedia.org/wiki/Promissory_note}{Promissory note, Wikipedia} under a CC BY-SA license
\end{itemize}

\subsection{Learning Objectives}

\begin{enumerate}
    \item \textbf{Understand the Purpose of Promissory Notes:}
    \item \textbf{Outcome:} Students will be able to explain the primary purpose of a promissory note as a legal instrument used for financial transactions.
    \item \textbf{Details:} This includes identifying the key elements of a promissory note, such as the principal amount, interest rate, parties involved, repayment terms, and maturity date, and explaining how these elements function within the context of a financial agreement.
    \item \textbf{Differentiate Between Promissory Notes and Similar Instruments:}
    \item \textbf{Outcome:} Students will be able to differentiate a promissory note from similar financial instruments, such as loan contracts and IOUs.
    \item \textbf{Details:} This includes understanding the specific legal obligations and characteristics that distinguish promissory notes from other instruments, such as the presence of a specific promise to pay, the inclusion of repayment terms, and the potential for negotiability.
    \item \textbf{Analyze the Legal Implications of Negotiable and Non-Negotiable Promissory Notes:}
    \item \textbf{Outcome:} Students will be able to analyze the legal implications of whether a promissory note is negotiable or non-negotiable.
    \item \textbf{Details:} This includes understanding how negotiability affects the enforceability of the note under different jurisdictions, particularly under the Uniform Commercial Code in the United States, and explaining the potential risks and benefits for the parties involved in the transaction.
\end{enumerate}

\subsection{Assessment Instrument}

\begin{enumerate}
    \item \textbf{Assessment Instrument: Promissory Notes Quiz} \textbf{Format:}
    \item \textbf{Type:} Mixed-format quiz (multiple-choice, short answer, and problem-solving questions)
    \item \textbf{Duration:} 45 minutes
    \item \textbf{Total Marks:} 30 points \textbf{Sections:}
    \item \textbf{Understanding Promissory Notes (10 points):}
    \item \textbf{Question 1:} Multiple Choice (2 points each)
    \item a. Which of the following best defines a promissory note?
    \item i. A promise to provide a service in the future
    \item ii. A legal instrument in which one party promises to pay a determinate sum of money to another party
    \item iii. A contract for the sale of goods
    \item iv. A receipt of payment
    \item b. What is typically not included in a promissory note?
    \item i. Principal amount
    \item ii. Terms of repayment
    \item iii. Payee's rights in the event of default
    \item iv. Inventory of goods
    \item c. In the context of promissory notes, what does "maturity date" refer to?
    \item i. The date the note was issued
    \item ii. The date the repayment is due
    \item iii. The date the interest begins to accrue
    \item iv. The date the borrower signs the note
    \item d. A promissory note that can be transferred and sold to another party is known as:
    \item i. A negotiable instrument
    \item ii. A mortgage note
    \item iii. A demand note
    \item iv. A loan agreement
    \item e. Which of the following best describes the difference between a promissory note and an IOU?
    \item i. A promissory note is less formal than an IOU
    \item ii. An IOU does not include a specific promise to pay or a repayment schedule
    \item iii. An IOU is legally binding, but a promissory note is not
    \item iv. A promissory note does not involve financial transactions
    \item \textbf{Analyzing Legal Implications (10 points):}
    \item \textbf{Question 2:} Short Answer (5 points each)
    \item a. Explain the significance of a promissory note being a negotiable instrument under U.S. law. What are the legal implications for the issuer and the holder?
    \item b. Describe how the use of a secured promissory note can protect the lender in the event of the borrower’s default. Provide an example to illustrate your explanation.
    \item \textbf{Application of Concepts (10 points):}
    \item \textbf{Question 3:} Problem-Solving (5 points each)
    \item a. A company issues an unsecured promissory note to a lender for `$`10,000, due in 90 days. After 60 days, the lender needs immediate funds and decides to sell the note to a bank at a 5\% discount. How much will the lender receive from the bank, and what will the bank collect at maturity?
    \item b. A borrower defaults on a secured promissory note. The security was a piece of equipment valued at `$`15,000. The outstanding debt is `$`12,000. Explain the process that the lender can follow to recover the debt and any potential financial outcomes. \textbf{Scoring Rubric:}
    \item \textbf{Understanding Promissory Notes:} 2 points for each correct answer (total 10 points)
    \item \textbf{Analyzing Legal Implications:} 5 points for each well-explained answer, demonstrating understanding of the legal concepts (total 10 points)
    \item \textbf{Application of Concepts:} 5 points for each correct and clearly explained problem-solving process (total 10 points) The quiz will assess the students' understanding of promissory notes, their legal implications, and their ability to apply these concepts to real-world scenarios.
\end{enumerate}

%%===============================================================
\section{Systems of Equations in Two Variables}
\label{sec:systems-of-equations-in-two-variables}
%%===============================================================

\paragraph{Introduction}



A system of equations consists of two or more equations made up of two or more variables such that all equations in the system are considered simultaneously. To find the unique solution to a system of equations, we must find a numerical value for each variable in the system that will satisfy all equations in the system at the same time. Some systems may not have a solution (for example, two distinct lines that are parallel to one another) and others may have an infinite number of solutions (for example, $y = 0$ and $y = \sin(x)$, since $\sin(x) = 0$ for an infinite number of x-values). In order for a linear system to have a unique solution, there must be at least as many equations as there are variables. Even so, this does not guarantee a unique solution.



Systems of equations can be solved in a number of ways. We can use elimination and/or substitution to solve for points of intersection shared by all equations in the system. We can also sometimes graph the equations to find the points shared by all equations, given that our system is graphable.



Here are some examples of systems of equations with two variables:



\begin{itemize}
    \item One solution: $y = 2x$ and $y = x + 5$. By use of substitution, elimination, or graphing, we can find that the solution to this system is $(5, 10)$, as this is the only point shared by both equations in the system.
    \item Multiple solutions: $y = x^2$ and $y = 4$. Since $x^2 = 4$ when x = 2 or -2, we know there are two solutions to this system: (2, 4) and (-2, 4). This is also clear when we graph these two equations, as $y = 4$ intersects the parabola $y = x^2$ when x = -2, and when x = 2.
    \item No solutions: $y = x^2$ and $y = -3$. Since $x^2$ is positive for all values of x, it cannot equal -3. When graphed, we can see that these two equations never intersect. Thus there are no solutions to this system of equations.
    \item Infinite solutions: $y = \sin(x)$ and $y = 0$. $\sin(x) = 0$ for all $x = \pi k$, where k is some integer. So, this system has an infinite number of solutions of the form $(\pi k, 0)$; that is, the solution set for this system of equations is {\.....$(-2\pi, 0)$, $(-\pi, 0)$, $(0, 0)$, $(\pi, 0)$, $(2\pi, 0)$\.....}.
\end{itemize}

\paragraph{Linear Equations with Two Variables}



In the previous module, linear equations with two variables were discussed. A single linear equation having two unknown variables is practically insufficient to solve or even narrow down the solutions for the two variables, although it does establish a relationship between them. The relationship is shown graphically as a line. Another linear equation with the same two variables may be enough to narrow down the solution to the two equations to one value for the first variable and one value for the second variable, i. e. to solve the system of two simultaneous linear equations. Let's see how two linear equations with the same two unknowns might be related to each other. Since we said it was given that both equations were linear, the graphs of both equations would be lines in the same two-dimensional coordinate plane (for a system with two variables). The lines could be related to each other in the following three ways:



\textbf{1.} The graphs of both equations could \textbf{coincide} giving the \textbf{same line}. This means that the two equations are providing the same information about how the variables are related to each other. The two equations are basically the same, perhaps just different versions or forms of each other. Either one could be mathematically manipulated to produce the other one. Both lines would have the same slope and the same y-intercept. Such equations are considered \textbf{dependent} on each other. Since no new information is provided, the addition of the second equation does not solve the problem by narrowing the solution set down to one solution.



\textbf{Example}: Dependent linear equations



$$6x - 3y = 12 $$

$$y = 2x - 4 $$



The above two equations provide the same information and result is the same graph, i. e. lines which \textbf{coincide} as shown in the following image.



% [Image: ]



Let's see how these equations can be mathematically manipulated to show they are basically the same.



Divide both sides of the first equation $6x - 3y = 12 $ by 3 to give



$$2x -y = 4 $$



Now add y to both sides



$$2x = 4 + y $$

Now subtract 4 from both sides



$$y = 2x - 4 $$



This is the same as the second equation in the example. This is the slope-intercept form of the equation, from which a slope and a y-intercept unique to the line can be compared with any other equations in the slope-intercept form.



\textbf{2.} The graphs of two lines could be \textbf{parallel} although not the same. The two lines do not intersect each other at any point. This means there is no solution which satisfies both equations simultaneously, i. e. at the same time. The solution set for this system of simultaneous linear equations is the empty set. Such equations are considered \textbf{inconsistent} with each other and actually give contradictory information if it is claimed they are both true at the same time in the same problem. The parallel lines have equal slopes but different y-intercepts.



Sets of equations which have at least one common point which might provide a solution set are \textbf{consistent} with each other. For example, the dependent equations mentioned previously are consistent with each other.



\textbf{Example}: Inconsistent linear equations



$$3x - 2y = -2 $$



$$3x -2y = 2 $$



To compare slopes and y-intercepts for these two linear equations, we place them in the slope-intercept forms. Subtract 3x from both sides of both equations.



$3x - 2y = -2 \qquad  \qquad  3x -2y = 2$ 



$-2y = -3x - 2 \qquad \qquad -2y = -3x + 2$



Divide both sides of both equations by -2 and simplify to get slope-intercept forms for comparison.



$(-2y)/(-2) = (-3x - 2)/(-2) \qquad  (-2y)/(-2) = (-3x + 2)/(-2)$



> $y = \frac{3}{2} x + 1 \qquad  \qquad \qquad  \qquad \qquad y = \frac{3}{2} x - 1$



Now, both slopes are equal at 3/2, but the y-intercepts at 1 and -1 are different.The lines are parallel. The graphs are shown here:



% [Image: ]



\textbf{3.} If the two lines are not the same and are not parallel, then they would intersect at one point because they are graphed in the same two-dimensional coordinate plane. The one point of intersection is the ordered pair of numbers which is the solution to the system of two linear equations and two unknowns. The two equations provide enough information to solve the problem and further equations are not needed. Such equations intersecting at a point providing a solution to the problem are considered \textbf{independent} of each other. The lines have different slopes but may or may not have the same y-intercept. Because such equations provide at least one solution point, they are \textbf{consistent} with each other.



\textbf{Example}: Consistent independent linear equations



$$y = 3x - 5 $$



$$y = -x - 1 $$



Both of these equations are given in the slope-intercept, so it is easy to compare slopes and y-intercepts. For these two linear functions, both slopes are different and both y-intercepts are different. This means the lines are neither dependent nor inconsistent, so on a two-dimensional graph they must intersect at some point. In fact, the graph shows the lines intersecting at (1,-2), which is the ordered pair solution to this system of independent simultaneous equations. Visual inspection of a graph cannot be relied on to give perfectly accurate coordinates every time, so either the point is tested with both equations or one of the following two methods is used to determine accurate coordinates for the intersection point.



% [Image: ]



\paragraph{Resources}



\begin{itemize}
    \item \href{https://tutorial.math.lamar.edu/classes/alg/systemstwovrble.aspx}{Linear Systems with Two Variables}, Paul's Online Notes
    \item \href{https://openstax.org/books/college-algebra/pages/7-1-systems-of-linear-equations-two-variables}{Systems of Linear Equations with Two Variables}, OpenStax
    \item \href{https://openstax.org/books/college-algebra/pages/7-3-systems-of-nonlinear-equations-and-inequalities-two-variables}{Systems of Nonlinear Equations with Two Variables}, OpenStax
    \item \href{https://www.youtube.com/watch?v=ej25myhYcSg}{Solving Systems of Equations with Elimination}, patrickJMT
    \item \href{https://www.youtube.com/watch?v=cwHR_B9zK7k}{Solving Systems of Equations with Substitution}, patrickJMT
    \item \href{https://www.youtube.com/watch?v=WNkPKv0OTGI}{Solving Systems of Equations with Graphing}, patrickJMT
\end{itemize}

\paragraph{Licensing}



Content obtained and/or adapted from:



\begin{itemize}
    \item \href{https://openstax.org/books/college-algebra/pages/7-1-systems-of-linear-equations-two-variables}{Systems of Linear Equations: Two Variables, OpenStax: College Algebra} under a CC BY license
    \item \href{https://en.wikibooks.org/wiki/Algebra/Systems_of_Equations}{Systems of Equations, Wikibooks: Algebra} under a CC BY-SA license
\end{itemize}

\subsection{Learning Objectives}

\begin{enumerate}
    \item Here are three Student Learning Objectives (SLOs) based on the provided content:
    \item \textbf{Solve Systems of Linear Equations:}
    \item \textbf{Outcome:} Students will be able to solve systems of linear equations with two variables using substitution, elimination, and graphical methods.
    \item \textbf{Details:} This includes understanding how to apply different methods to find unique solutions, multiple solutions, or determine when no solution exists. Students will also interpret the graphical representation of linear systems.
    \item \textbf{Classify Systems of Linear Equations:}
    \item \textbf{Outcome:} Students will be able to classify systems of linear equations as consistent, inconsistent, dependent, or independent.
    \item \textbf{Details:} This involves analyzing the relationship between the lines represented by the equations in the system, including identifying whether they coincide, are parallel, or intersect at a single point.
    \item \textbf{Analyze Solution Sets of Linear Systems:}
    \item \textbf{Outcome:} Students will be able to analyze the solution set of a system of linear equations and determine whether it has a single solution, no solution, or infinitely many solutions.
    \item \textbf{Details:} This includes understanding the conditions under which each type of solution occurs and being able to justify the classification using algebraic methods and graphical representations.
\end{enumerate}

\subsection{Assessment Instrument}

\begin{enumerate}
    \item \textbf{Assessment Instrument: Systems of Equations in Two Variables Quiz} \textbf{Format:}
    \item \textbf{Type:} Mixed-format quiz (multiple-choice, short answer, and problem-solving questions)
    \item \textbf{Duration:} 45 minutes
    \item \textbf{Total Marks:} 30 points \textbf{Sections:}
    \item \textbf{Classifying Systems of Equations (10 points):}
    \item \textbf{Question 1:} For each of the following systems of equations, classify them as consistent and independent, consistent and dependent, or inconsistent. Justify your classification.
    \item a. $2x + 3y = 6$ and $4x + 6y = 12$
    \item b. $x - y = 1$ and $2x + y = 7$
    \item c. $3x + 2y = 5$ and $3x + 2y = -4$
    \item \textbf{Solving Systems of Equations (10 points):}
    \item \textbf{Question 2:} Solve the following systems of equations using the substitution or elimination method:
    \item a. $x + y = 7$ and $2x - y = 3$
    \item b. $3x - 2y = 4$ and $x + 3y = 7$
    \item \textbf{Question 3:} Solve the system graphically and identify the solution (if any):
    \item $y = 2x - 5$ and $y = -x + 2$
    \item \textbf{Analyzing Solution Sets (10 points):}
    \item \textbf{Question 4:} Determine the number of solutions for each of the following systems and explain your reasoning:
    \item a. $y = x^2 + 1$ and $y = 3x + 4$
    \item b. $y = 2x + 3$ and $y = 2x - 1$
    \item \textbf{Question 5:} Consider the system $y = \sin(x)$ and $y = 0$. How many solutions does this system have? Explain how you arrived at your answer. \textbf{Scoring Rubric:}
    \item \textbf{Classifying Systems of Equations:} 3.33 points for each correct classification and justification (total 10 points)
    \item \textbf{Solving Systems of Equations:} 3.33 points for each correctly solved system (total 10 points)
    \item \textbf{Analyzing Solution Sets:} 3.33 points for each correct analysis and explanation (total 10 points) The quiz is designed to evaluate students' understanding of systems of equations, including their ability to classify, solve, and analyze solution sets for different types of systems.
\end{enumerate}

%%===============================================================
\section{Systems of Inequalities in Two Variables}
\label{sec:systems-of-inequalities-in-two-variables}
%%===============================================================

In mathematics a \textbf{linear inequality} is an inequality which involves a linear function. A linear inequality contains one of the symbols of inequality. It shows the data which is not equal in graph form.



\begin{itemize}
    \item $<$ less than
    \item $>$ greater than
    \item $\leq$ less than or equal to
    \item $\geq$ greater than or equal to
    \item $\neq$ not equal to
    \item $=$ equal to
\end{itemize}

A linear inequality looks exactly like a linear equation, with the inequality sign replacing the equality sign.



\paragraph{Linear inequalities of real numbers}



\paragraph{Two-dimensional linear inequalities}



% [Image: Graph of linear inequality: x + 3y < 9]



Figure: Graph of linear inequality: $x + 3y < 9$



Two-dimensional linear inequalities are expressions in two variables of the form:



$$ax + by < c \text{ and } ax + by \geq c,$$ where the inequalities may either be strict or not. The solution set of such an inequality can be graphically represented by a half-plane (all the points on one "side" of a fixed line) in the Euclidean plane. The line that determines the half-planes (\textit{ax} + \textit{by} = \textit{c}) is not included in the solution set when the inequality is strict. A simple procedure to determine which half-plane is in the solution set is to calculate the value of \textit{ax} + \textit{by} at a point $(x _0, y _0)$ which is not on the line and observe whether or not the inequality is satisfied.



For example, to draw the solution set of $x + 3y < 9$, one first draws the line with equation $x + 3y = 9$ as a dotted line, to indicate that the line is not included in the solution set since the inequality is strict. Then, pick a convenient point not on the line, such as $(0,0)$. Since $0 + 3(0) = 0 < 9$, this point is in the solution set, so the half-plane containing this point (the half-plane "below" the line) is the solution set of this linear inequality.



\paragraph{Linear inequalities in general dimensions}



In $\mathbb{R} ^n$ linear inequalities are the expressions that may be written in the form



$$f(\bar{x}) < b$$ or $f(\bar{x}) \leq b,$ where \textit{f} is a linear form (also called a \textit{linear functional}), $\bar{x} = (x _1,x _2,\ldots,x _n)$ and \textit{b} a constant real number.



More concretely, this may be written out as



$$a _1 x _1 + a _2 x _2 + \cdots + a _n x _n < b$$ or



$$a _1 x _1 + a _2 x _2 + \cdots + a _n x _n \leq b.$$



Here $x _1, x _2,...,x _n$ are called the unknowns, and $a _{1}, a _{2},..., a _{n}$ are called the coefficients.



Alternatively, these may be written as



$$g(x) < 0 \,$$ or $g(x) \leq 0,$ where \textit{g} is an affine function.



That is



> $a _0 + a _1 x _1 + a _2 x _2 + \cdots + a _n x _n < 0$



or



> $a _0 + a _1 x _1 + a _2 x _2 + \cdots + a _n x _n \leq 0.$



Note that any inequality containing a "greater than" or a "greater than or equal" sign can be rewritten with a "less than" or "less than or equal" sign, so there is no need to define linear inequalities using those signs.



\paragraph{Systems of linear inequalities}



A system of linear inequalities is a set of linear inequalities in the same variables:



$$ \begin{aligned} a _{11} x _1 &\; + \;&& a _{12} x _2 &\; + \cdots + \;&& a _{1n} x _n &\; \leq \;&& b _1 \\\\ a _{21} x _1 &\; + \;&& a _{22} x _2 &\; + \cdots + \;&& a _{2n} x _n &\; \leq \;&& b _2 \\\\ \vdots & && \vdots & && \vdots & && \vdots \\\\ a _{m1} x _1 &\; + \;&& a _{m2} x _2 &\; + \cdots + \;&& a _{mn} x _n &\; \leq \;&& b _m \\\\ \end{aligned} $$

 Here $x _1,\ x _2,...,x _n$ are the unknowns, $a _{11},\ a _{12},...,\ a _{mn}$ are the coefficients of the system, and $b _1,\ b _2,...,b _m$ are the constant terms.



This can be concisely written as the matrix inequality



$$Ax \leq b,$$



where \textit{A} is an $m \times n$ matrix, \textit{x} is an $n \times 1$ column vector of variables, and \textit{b} is an $m \times 1$ column vector of constants.



In the above systems both strict and non-strict inequalities may be used.



\begin{itemize}
    \item Not all systems of linear inequalities have solutions.
\end{itemize}

Variables can be eliminated from systems of linear inequalities using Fourier–Motzkin elimination.



\paragraph{Applications}



\paragraph{Polyhedra}



The set of solutions of a real linear inequality constitutes a half-space of the 'n'-dimensional real space, one of the two defined by the corresponding linear equation.



The set of solutions of a system of linear inequalities corresponds to the intersection of the half-spaces defined by individual inequalities. It is a convex set, since the half-spaces are convex sets, and the intersection of a set of convex sets is also convex. In the non-degenerate cases this convex set is a convex polyhedron (possibly unbounded, e.g., a half-space, a slab between two parallel half-spaces or a polyhedral cone). It may also be empty or a convex polyhedron of lower dimension confined to an affine subspace of the \textit{n}-dimensional space $\mathbb{R} ^n$.



\paragraph{Linear programming}



A linear programming problem seeks to optimize (find a maximum or minimum value) a function (called the objective function) subject to a number of constraints on the variables which, in general, are linear inequalities. The list of constraints is a system of linear inequalities.



\paragraph{Resources}



\begin{itemize}
    \item \href{https://courses.lumenlearning.com/cuny-hunter-collegealgebra/chapter/outcome-graphing-linear-equations/}{Systems of Linear Inequalities in Two Variables}, Lumen Learning
    \item \href{https://www.youtube.com/watch?v=PAb1lCtpF8w}{Solving Systems of Inequalities}, Textbook Tactics on YouTube
    \item \href{https://www.mathplanet.com/education/algebra-1/systems-of-linear-equations-and-inequalities/systems-of-linear-inequalities}{Graphing Systems of Linear Inequalities with Two Variables}, Math Planet
\end{itemize}

\paragraph{Licensing}



Content obtained and/or adapted from:



\begin{itemize}
    \item \href{https://en.wikipedia.org/wiki/Linear_inequality}{Linear inequality, Wikipedia} under a CC BY-SA license
\end{itemize}

\subsection{Learning Objectives}

\begin{enumerate}
    \item Here are three Student Learning Objectives (SLOs) based on the topic of linear inequalities:
    \item \textbf{Understand and Graph Linear Inequalities:}
    \item \textbf{Outcome:} Students will be able to graph linear inequalities in two variables and identify the solution set as a half-plane on the coordinate plane.
    \item \textbf{Details:} This includes plotting the boundary line (solid for $\leq$ or $\geq$, dashed for $<$ or $>$), selecting a test point to determine the correct half-plane, and shading the appropriate region.
    \item \textbf{Solve Systems of Linear Inequalities:}
    \item \textbf{Outcome:} Students will be able to solve systems of linear inequalities and represent the solution set as the intersection of half-planes.
    \item \textbf{Details:} This involves solving multiple linear inequalities, graphing each inequality on the same coordinate plane, and identifying the region where all inequalities are satisfied simultaneously.
    \item \textbf{Apply Linear Inequalities to Real-World Problems:}
    \item \textbf{Outcome:} Students will be able to formulate and solve linear inequalities based on real-world scenarios, and interpret the results within the context of the problem.
    \item \textbf{Details:} This includes translating a given situation into a system of linear inequalities, solving the system graphically or algebraically, and explaining the meaning of the solution in terms of the original problem context.
\end{enumerate}

\subsection{Assessment Instrument}

\begin{enumerate}
    \item \textbf{Assessment Instrument: Systems of Inequalities in Two Variables Quiz} \textbf{Format:}
    \item \textbf{Type:} Mixed-format quiz (multiple-choice, short answer, and problem-solving questions)
    \item \textbf{Duration:} 45 minutes
    \item \textbf{Total Marks:} 30 points \textbf{Sections:}
    \item \textbf{Understanding Linear Inequalities (10 points):}
    \item \textbf{Question 1:} Identify the correct graph for the inequality $2x - y > 3$ from the options below:
    \item a. A graph where the line $2x - y = 3$ is solid, and the region above it is shaded.
    \item b. A graph where the line $2x - y = 3$ is dashed, and the region below it is shaded.
    \item c. A graph where the line $2x - y = 3$ is dashed, and the region above it is shaded.
    \item d. A graph where the line $2x - y = 3$ is solid, and the region below it is shaded.
    \item \textbf{Solving Systems of Linear Inequalities (10 points):}
    \item \textbf{Question 2:} Solve the following system of inequalities and graph the solution set on the coordinate plane: $$ \begin{aligned} x + 2y &\leq 4 \\\\ x - y &> -1 \end{aligned} $$
    \item a. Graph the boundary lines and indicate the solution region.
    \item b. Determine if the point (1, 1) is within the solution set. Justify your answer.
    \item \textbf{Application of Systems of Linear Inequalities (10 points):}
    \item \textbf{Question 3:} A company produces two products, A and B. The profit from each product is `$`10 for A and `$`15 for B. The company must meet the following constraints: $$ \begin{aligned} 2A + 3B &\leq 60 \quad (\text{resource constraint}) \\\\ A &\geq 10 \quad (\text{minimum production of A}) \end{aligned} $$
    \item a. Write the system of inequalities that represents this situation.
    \item b. Graph the system of inequalities.
    \item c. Identify the feasible region and determine the maximum possible profit based on the given constraints. \textbf{Scoring Rubric:}
    \item \textbf{Understanding Linear Inequalities:} 5 points for correct identification of the graph, 5 points for correctly explaining the choice (total 10 points).
    \item \textbf{Solving Systems of Linear Inequalities:} 5 points for accurate graphing of inequalities, 2.5 points for correct boundary lines and shading, 2.5 points for evaluating the point (total 10 points).
    \item \textbf{Application of Systems of Linear Inequalities:} 3 points for writing the correct system of inequalities, 3 points for accurate graphing, 4 points for identifying the feasible region and calculating the maximum profit (total 10 points). This quiz assesses students' ability to graph linear inequalities, solve systems of linear inequalities, and apply these concepts to real-world problems, ensuring that they understand and can effectively communicate their reasoning.
\end{enumerate}

%%===============================================================
\section{Linear Programming}
\label{sec:linear-programming}
%%===============================================================

Linear Programming (LP) is a mathematical modelling technique useful for allocation of limited resources such as material, machines etc to several competing activities such as projects, services etc. A typical linear programming problem consists of a linear objective function which is to be maximized or minimized subject to a finite number of linear constraints. (By a linear function we mean a function of the form $a _1x _1+a _2x _2+\cdots a _nx _n$ where $x _1,x _2\cdots x _n$ are all variables.)



The founders of LP are George B. Dantzig, who published the simplex method in 1947, John von Neumann, who developed the theory of the duality in the same year, and Leonid Kantorovich, a Russian mathematician who used similar techniques in economics before Dantzig and won the Nobel prize in 1975 in economics. The linear programming problem was first shown to be solvable in polynomial time by Leonid Khachiyan in 1979, but a larger major theoretical and practical breakthrough in the field came in 1984 when Narendra Karmarkar introduced a new interior point method for solving linear programming problems.



\paragraph{The Two Variable LP model}



% [Image: Example of graphical solution to a LP with two variables X_1 and X_2] 



Figure: Example of graphical solution to a LP with two variables $X _1$ and $X _2$



Two \textbf{variable} LP models don't usually exist in practice but their study can provide us with an example as to how a general model can be handled.



We will explain the model by an example.



Consider a chemical company which produces two salts: X and Y. Suppose an acid and a base are the only two chemicals required by the company to produce these two salts. Further suppose that by past experience the owner of the company has obtained the following data:



To make 1 unit of the salt X, 6 units of the acid and 1 unit of the base is required. To make 1 unit of the salt Y, 4 units of the acid and 2 units of the base are required. At most 24 units of the acid and 6 units of the base are available daily. 1 unit of the salt X gets him a profit of 5 (in whatever currency) and 1 unit of the salt Y gets him a profit of 4. In addition due to a market survey he knows that the daily demand for the salt Y cannot exceed that for X by more than 1 unit, and that the maximum daily demand for the salt Y is 2 units.



The company owner wants the best \textit{product mix}. That is to say, he wants to know the amount of X and Y he should make daily so that he gets the highest profit.



To formulate the LP model for this problem we first need to identify the \textit{decision variables}. These are the variables which will represent the entities about which we have to actually make a decision about. In our problem it is clear that the entities are the salts X and Y and so the variables should represent the amount they should be each made. So let,



$x _1$= Units produced daily of salt X



$x _2$= Units produced daily of salt Y



The next step is to decide what should be the objective function for our model. It is clear from the name itself that the objective function represents our objective i.e. maximizing our total profit. Since the total profit is usually the sum of the profits obtained from X and Y it should be equal to $5x _1+4x _2$. This is assuming that if 1 unit of X (or Y) gives a profit of 5 (or 4) then $x _1$ (or $x _2$) units would give a profit of $5x _1$ (or $4x _2$). Hence if z represents the total daily profit the objective is:



Maximize z = $5x _1+4x _2$.



Now we consider the constraints that restrict the usage of the acids and the salts. We know that the total amount of acid which we use in the day can't exceed 24. Since $6x _1$ and $4x _2$ are reasonably the amounts of the acid used in making the salt X and Y we may conclude that our constraint is:



$6x _1+4x _2\le 24$



Similarly, for the base:



$x _1+2x _2\le 6$



Now the market dictates that the excess of the daily production of Y over X which is $x _2-x _1$ can't exceed 1 which means that:



$x _2-x _1\le 1$



Also the maximum daily demand for the salt Y is 2 units so that:



$x _2\le 2$



An implicit assumption is that the variables $x _1$ and $x _2$ can't assume negative values (Why?). The \textit{nonnegativity restrictions}, $x _1,x _2\ge 0$ account for this requirement.



The complete model can be now stated as:



Maximize $z = 5x _1+4x _2$



subject to,



$6x _1+4x _2\le 24$,



$x _1+2x _2\le 6$,



$x _2-x _1\le 1$,



$x _2\le 2$,



$x _1,x _2\ge 0$.



Any values of $x _1$ and $x _2$ that satisfy all the five constraints constitute a \textit{feasible solution}. Otherwise the solution is infeasible. For example, the solution $x _1=3$ and $x _2=1$ is feasible because on substitution in the constraints none of the inequalities are violated. On the other hand the solution, the solution $x _1=4$ and $x _2=1$ is infeasible because it does not satisfy constraint (1), namely 6\*4+4\*1=28 which is larger then the right hand side (=24).



Our goal is to find the \textit{optimum solution}, i.e. one which maximizes the objective function besides being feasible. The procedure to do that will be discussed in the succeeding sections.



\paragraph{Properties of the LP model}



LP models have the following properties:



\begin{itemize}
    \item \textbf{Proportionality}: the contribution of each decision variable is directly proportional to its value in both the objective function and in the constraints. For example, the contribution of the first decision variable in the first constraint is $6x _1$. This is directly proportional to $x _1$ with 6 as the constant of proportionality. Roughly speaking, the rule of three is followed.
\end{itemize}

\begin{itemize}
    \item \textbf{Additivity}: the contribution of all the variables in the objective function and the constraints to be the sum of the individual contributions of each variable. For example, the total profit is $5x _1+4x _2$ which is the sum of the individual profits $5x _1$ and $4x _2$.
\end{itemize}

\begin{itemize}
    \item \textbf{Certainty}: all the objective and constraint coefficients are deterministic, that is all the data about profit, availability and requirements is constant.
\end{itemize}

\paragraph{Resources}



\begin{itemize}
    \item \href{https://mathresearch.utsa.edu/wikiFiles/MAT1053/Linear_Programming/MAT1053_M9.2Linear_Programming.pdf}{Linear Programming}, Book Chapter
    \item \href{https://mathresearch.utsa.edu/wikiFiles/MAT1053/Linear_Programming/MAT1053_M9.2Linear_ProgrammingGN.pdf}{Guided Notes}
\end{itemize}

\paragraph{Licensing}



Content obtained and/or adapted from:



\begin{itemize}
    \item \href{https://en.wikibooks.org/wiki/Operations_Research/Linear_Programming}{Linear programming, Wikibooks} under a CC BY-SA license
\end{itemize}

\subsection{Learning Objectives}

\begin{enumerate}
    \item Here are three Student Learning Objectives (SLOs) based on the topic of Linear Programming (LP):
    \item \textbf{Formulate Linear Programming Problems:}
    \item \textbf{Outcome:} Students will be able to formulate a linear programming problem by identifying decision variables, defining the objective function, and establishing relevant linear constraints.
    \item \textbf{Details:} This includes setting up a mathematical model for real-world scenarios, such as resource allocation or optimization problems, by converting practical constraints into linear inequalities.
    \item \textbf{Solve Linear Programming Problems Graphically:}
    \item \textbf{Outcome:} Students will be able to solve two-variable linear programming problems using graphical methods, including identifying feasible regions, evaluating corner points, and determining the optimal solution.
    \item \textbf{Details:} This involves plotting constraints on a graph, identifying the feasible region, and applying the objective function to find the optimal solution, while understanding the implications of different solutions on the objective.
    \item \textbf{Analyze the Properties of Linear Programming Models:}
    \item \textbf{Outcome:} Students will be able to analyze the key properties of linear programming models, including proportionality, additivity, and certainty, and explain how these properties impact the formulation and solution of LP problems.
    \item \textbf{Details:} This involves understanding the mathematical assumptions underlying LP models, evaluating the validity of these assumptions in different contexts, and recognizing the limitations that these properties impose on the model's application in real-world scenarios.
\end{enumerate}

\subsection{Assessment Instrument}

\begin{enumerate}
    \item \textbf{Assessment Instrument: Linear Programming Quiz} \textbf{Format:}
    \item \textbf{Type:} Mixed-format quiz (multiple-choice, short answer, and problem-solving questions)
    \item \textbf{Duration:} 45 minutes
    \item \textbf{Total Marks:} 30 points \textbf{Sections:}
    \item \textbf{Formulating Linear Programming Problems (10 points):}
    \item \textbf{Question 1:} (Multiple Choice) Which of the following correctly represents the objective function for a linear programming problem where the profit from product A is `$`3 per unit and from product B is `$`4 per unit, with $x _1$ and $x _2$ representing the number of units produced of products A and B, respectively?
    \item a. Maximize $z = 3x _1 + 4x _2$
    \item b. Maximize $z = 4x _1 + 3x _2$
    \item c. Maximize $z = 3x _1 - 4x _2$
    \item d. Maximize $z = 4x _1 - 3x _2$
    \item \textbf{Question 2:} (Short Answer) Identify the decision variables and formulate the objective function for a company that wants to minimize costs, where the cost per unit of product A is `$`10 and product B is `$`15, and $x _1$ and $x _2$ represent the quantities of products A and B, respectively.
    \item \textbf{Solving Linear Programming Problems Graphically (10 points):}
    \item \textbf{Question 3:} (Problem-Solving) A company produces two products, X and Y. The profit from product X is `$`5 per unit, and the profit from product Y is `$`4 per unit. The production process for these products requires two resources: Resource 1 and Resource 2. Each unit of product X requires 2 units of Resource 1 and 3 units of Resource 2, while each unit of product Y requires 4 units of Resource 1 and 1 unit of Resource 2. The company has 16 units of Resource 1 and 9 units of Resource 2 available. Formulate the linear programming model for this problem, and then graphically determine the feasible region.
    \item \textbf{Question 4:} (Multiple Choice) Which of the following points is a feasible solution for the linear programming model defined by the constraints $2x _1 + 3x _2 \le 12$ and $x _1 + 2x _2 \le 8$, with $x _1, x _2 \ge 0$?
    \item a. $(2, 3)$
    \item b. $(3, 2)$
    \item c. $(4, 1)$
    \item d. $(1, 4)$
    \item \textbf{Analyzing the Properties of Linear Programming Models (10 points):}
    \item \textbf{Question 5:} (Short Answer) Explain the concept of proportionality in the context of linear programming. Provide an example of how this property is reflected in the constraints of a linear programming problem.
    \item \textbf{Question 6:} (Problem-Solving) Consider the linear programming model with the objective function $\text{Maximize  }  z = 7x _1 + 5x _2$ subject to the constraints $x _1 + x _2 \le 10$, $2x _1 + 3x _2 \le 18$, and $x _1, x _2 \ge 0$. Analyze whether the solution $(x _1, x _2) = (4, 5)$ satisfies the property of certainty and whether it is feasible. \textbf{Scoring Rubric:}
    \item \textbf{Formulating Linear Programming Problems:} 5 points for each correct answer (total 10 points)
    \item \textbf{Solving Linear Programming Problems Graphically:} 5 points for each correctly solved problem (total 10 points)
    \item \textbf{Analyzing the Properties of Linear Programming Models:} 5 points for each correct and well-explained answer (total 10 points) The quiz will help assess the students' understanding of linear programming concepts, including formulating problems, solving them graphically, and analyzing the key properties of LP models. It ensures that they can accurately apply the theoretical knowledge to practical scenarios.
\end{enumerate}

%%===============================================================
\section{Matrices}
\label{sec:matrices}
%%===============================================================

\paragraph{Matrices}



A matrix is an array of numbers arranged into rows and columns. Some examples of matrices are,



$$A=

\begin{bmatrix}

2 & 4 & 6 \\\\ 0 & -5 & 7.105 \\\\ 1 & -3 & 2

\end{bmatrix}

,\quad

B=

\begin{bmatrix}

	-3 & -4 \\\\ \pi & \sqrt{2}

\end{bmatrix}

,\text{ and }\,\,\,

C=

\begin{bmatrix}

	-3 & -5 & 0 & 0 \\\\ -1 & 4.56 & 3.28 & 19

\end{bmatrix}$$



When describing matrices we indicate the number of rows first, then the number of columns. For example, the matrix $C$ with two rows and four columns is said to be a $2\times 4$ matrix.



It is standard notation to name matrices with capital letters and to use lower case letters with subscripts to identify particular entries in a matrix.



For example, to identify the entry in row 1 and column 3 of matrix $A$ we would write $a _{13}$. To indicate that this entry is a six we would write the equation $a _{13}=6$.



Two matrices are considered to be \textbf{equal} only if they are the same size and every pair of corresponding elements are equal.



A \textbf{column matrix} is a matrix with only one column. Similarly, a \textbf{row matrix} has only one row.



\paragraph{Vectors}



A \textbf{vector} is an object often defined by a long list of properties. However, for now we will avoid the more complicated definition, and just say that a vector is an ordered list of numbers. Later we will see that vectors can really be much more.



An ordered pair, $(x, y)$, that is used to identify a point in the plane can be considered to be a vector.



Similarly, an ordered triple, $(x, y, z)$ is a vector.



Obviously, row and column matrices can also be considered to be a vector.



It is common to name vectors using variables with arrows above.



For example, we might write $\vec{v}=(2, 3, 5, -4),\text{ or } \vec{w}=

\begin{bmatrix} 4 \\\\ 5 \\\\ 0 \end{bmatrix}$.



For the most part, it will convenient to think of vectors as column matrices.



\paragraph{Resources}



\begin{itemize}
    \item \href{https://www.unf.edu/\textasciitilde{}mzhan/chapter3.pdf}{Brief Introduction to Vectors and Matrices}, University of North Florida
    \item \href{https://www.usna.edu/Users/math/uhan/sm286a/lessons/02.pdf}{Introduction to Matrices and Vectors}, United States Naval Academy
\end{itemize}

\paragraph{Licensing}



Content obtained and/or adapted from:



\begin{itemize}
    \item \href{https://en.wikibooks.org/wiki/Linear_Algebra/Matrices_and_Vectors/}{Matrices and Vectors, Wikibooks: Linear Algebra} under a CC BY-SA license
\end{itemize}

\subsection{Learning Objectives}

\begin{enumerate}
    \item \textbf{Matrix Classification and Identification:}
    \item \textbf{Outcome:} Students will be able to classify matrices based on their dimensions as row matrices, column matrices, square matrices, or rectangular matrices.
    \item \textbf{Details:} This includes understanding the structure of a matrix, identifying the number of rows and columns, and distinguishing between different types of matrices using specific examples.
    \item \textbf{Matrix Equality:}
    \item \textbf{Outcome:} Students will be able to determine the equality of two matrices by comparing their dimensions and corresponding elements.
    \item \textbf{Details:} This includes understanding the conditions necessary for two matrices to be considered equal and applying this knowledge to specific examples where students compare matrices element-wise.
    \item \textbf{Vector Representation Using Matrices:}
    \item \textbf{Outcome:} Students will be able to represent vectors as row or column matrices and identify vectors in various forms, such as ordered pairs or triples.
    \item \textbf{Details:} This includes converting vectors from one form to another, recognizing the relationship between vectors and matrices, and using correct notation to represent vectors as matrices in different scenarios.
\end{enumerate}

\subsection{Assessment Instrument}

\begin{enumerate}
    \item \textbf{Assessment Instrument: Matrix Concepts Quiz} \textbf{Format:}
    \item \textbf{Type:} Mixed-format quiz (multiple-choice, short answer, and problem-solving questions)
    \item \textbf{Duration:} 45 minutes
    \item \textbf{Total Marks:} 30 points \textbf{Sections:}
    \item \textbf{Matrix Classification and Identification (10 points):}
    \item \textbf{Question 1:} Identify the type of the following matrices as either row matrix, column matrix, square matrix, or rectangular matrix:
    \item a. $A = \begin{bmatrix} 1 & 2 & 3 \end{bmatrix}$
    \item b. $B = \begin{bmatrix} 4 \\\\ 5 \end{bmatrix}$
    \item c. $C = \begin{bmatrix} 1 & 2 \\\\ 3 & 4 \end{bmatrix}$
    \item d. $D = \begin{bmatrix} -1 & 0 & 2 \\\\ 3 & 4 & 5 \end{bmatrix}$
    \item \textbf{Matrix Equality (10 points):}
    \item \textbf{Question 2:} Determine whether the following pairs of matrices are equal. Justify your answer by comparing their dimensions and corresponding elements:
    \item a. $A = \begin{bmatrix} 2 & 3 \\\\ 4 & 5 \end{bmatrix}, B = \begin{bmatrix} 2 & 3 \\\\ 4 & 5 \end{bmatrix}$
    \item b. $C = \begin{bmatrix} 1 & -1 \end{bmatrix}, D = \begin{bmatrix} 1 & -1 \end{bmatrix}$
    \item c. $E = \begin{bmatrix} 0 & 2 & 4 \end{bmatrix}, F = \begin{bmatrix} 0 \\\\ 2 \\\\ 4 \end{bmatrix}$
    \item \textbf{Vector Representation Using Matrices (10 points):}
    \item \textbf{Question 3:} Convert the following vectors into matrix form, either as row or column matrices:
    \item a. $\vec{v} = (3, -2, 7)$
    \item b. $\vec{w} = (-1, 0, 4, 8)$
    \item c. $\vec{u} = (5, -3)$ \textbf{Scoring Rubric:}
    \item \textbf{Matrix Classification and Identification:} 2.5 points for each correctly identified matrix type (total 10 points)
    \item \textbf{Matrix Equality:} 3.33 points for each correct determination and justification (total 10 points)
    \item \textbf{Vector Representation Using Matrices:} 3.33 points for each correctly converted vector (total 10 points) The quiz is designed to assess students' understanding of matrices, their ability to classify matrices, determine matrix equality, and convert vectors into matrix form.
\end{enumerate}

%%===============================================================
\section{Matrix Operations}
\label{sec:matrix-operations}
%%===============================================================

\paragraph{Matrix Operations}



\paragraph{Adding and subtracting matrices}



In order to add or subtract two matrices, they must be of the same dimension; that is, the two matrices must have the same number of rows and the same number of columns. To add two matrices together, we simply need to add every entry in one matrix to the entry in the same row and same column in the other matrix. For example:



$$\begin{bmatrix}

a _{11} & a _{12} & a _{13} \\\a _{21} & a _{22} & a _{23} \\\a _{31} & a _{32} & a _{33}

\end{bmatrix} +

\begin{bmatrix}

b _{11} & b _{12} & b _{13} \\\b _{21} & b _{22} & b _{23} \\\b _{31} & b _{32} & b _{33}

\end{bmatrix} =

\begin{bmatrix}

a _{11} + b _{11} & a _{12} + b _{12} & a _{13} + b _{13} \\\a _{21} + b _{21} & a _{22} + b _{22} & a _{23} + b _{23} \\\a _{31} + b _{31} & a _{32} + b _{32} & a _{33} + b _{33}

\end{bmatrix}$$



$$\begin{bmatrix}

1 & 2 & 3 \\\0 & 1 & 2 \\\0 & 0 & 1

\end{bmatrix} +

\begin{bmatrix}

5 & 2 & 1 \\\3 & 3 & 1 \\\4 & 2 & 0

\end{bmatrix} =

\begin{bmatrix}

1 + 5 & 2 + 2 & 3 + 1 \\\0 + 3 & 1 + 3 & 2 + 1 \\\0 + 4 & 0 + 2 & 1 + 0

\end{bmatrix} =

\begin{bmatrix}

6 & 4 & 4 \\\3 & 4 & 3 \\\4 & 2 & 1

\end{bmatrix}$$



\paragraph{Multiplying matrices by scalars}



When multiplying a matrix by a scalar (or number), all we need to do is multiply each entry of the matrix by the scalar. For example:



$$3\begin{bmatrix}

1 & 2 & 3 \\\0 & 1 & 2 \\\0 & 0 & 1

\end{bmatrix} =

\begin{bmatrix}

3 & 6 & 9 \\\0 & 3 & 6 \\\0 & 0 & 3

\end{bmatrix}$$



$$-\frac{1}{2}\begin{bmatrix} 2 & -4 & 8 \\\\ -2 & 1 & -6 \\\\ 0 & 0 & 7 \end{bmatrix} = \begin{bmatrix} -1 & 2 & -4 \\\\ 1 & -\frac{1}{2} & 3 \\\\ 0 & 0 & -\frac{7}{2} \end{bmatrix}$$



\paragraph{Multiplying matrices}



Matrix multiplication is not commutative; that is, for most matrices $A$ and $B$, $AB \neq BA$. Two matrices can only be multiplied together if the number of columns in the first matrix equals the number of rows in the second matrix. For example, we can take the product of a 3-by-2 matrix times a 2-by-5 matrix, but not a 3-by-2 matrix times a 3-by-2 matrix. The product of two matrices has the same number of rows as the first matrix and the same number of columns as the second matrix. For example, a 2-by-3 matrix times a 3-by-2 matrix will result in a 2-by-2 matrix.



The product of two 3-by-3 matrices is

$$\begin{bmatrix}

a _{11} & a _{12} & a _{13} \\\a _{21} & a _{22} & a _{23} \\\a _{31} & a _{32} & a _{33}

\end{bmatrix}

\begin{bmatrix}

b _{11} & b _{12} & b _{13} \\\b _{21} & b _{22} & b _{23} \\\b _{31} & b _{32} & b _{33}

\end{bmatrix} =

\begin{bmatrix}

c _{11} & c _{12} & c _{13} \\\c _{21} & c _{22} & c _{23} \\\c _{31} & c _{32} & c _{33}

\end{bmatrix}$$



where $c _{ij}$ is the dot product of the i-th row of the first matrix and the j-th column of the second matrix; that is, 

$$c _{ij} = a _{i1}b _{1j} + a _{i2}b _{2j} + a _{i3}b _{3j}$$



Here is a 2-by-2 example of matrix multiplication:



$$\begin{bmatrix}

1 & 2 \\\3 & 4

\end{bmatrix}

\begin{bmatrix}

5 & 6 \\\7 & 8

\end{bmatrix} =

\begin{bmatrix}

1(5) + 2(7) & 1(6) + 2(8) \\\3(5) + 4(7) & 3(6) + 4(8)

\end{bmatrix} =

\begin{bmatrix}

5 + 14 & 6 + 16 \\\15 + 28 & 18 + 32

\end{bmatrix}

\begin{bmatrix}

19 & 22 \\\43 & 50

\end{bmatrix}$$



Note that there is no "matrix division". Rather, we multiply one matrix by the inverse of another (for example, $AB^{-1}$, NOT $A/B$) to obtain a "quotient" from two matrices.



\paragraph{Inverse of an n-by-n matrix}



An n-by-n matrix A is the inverse of n-by-n matrix B (and B the inverse of A) if $BA = AB = I$, where I is an identity matrix.



The inverse of an n-by-n matrix can be calculated by creating an n-by-2n matrix which has the original matrix on the left and the identity matrix on the right. Row reduce this matrix and the right half will be the inverse. If the matrix does not row reduce completely (i.e., a row is formed with all zeroes as its entries), it does not have an inverse.



\paragraph{Example 1}



Let $\mathrm{A} = \begin{bmatrix}

  1 \& 4 \& 4 \\\  2 \& 5 \& 8 \\\  3 \& 6 \& 9

\end{bmatrix}$



We begin by expanding and partitioning A to include the identity matrix, and then proceed to row reduce A until we reach the identity matrix on the left-hand side.



$$\begin{bmatrix}

  1 & 4 & 4 & \big| & 1 & 0 & 0 \\\  2 & 5 & 8 & \big| & 0 & 1 & 0 \\\  3 & 6 & 9 & \big| & 0 & 0 & 1

\end{bmatrix}

\rightsquigarrow

\begin{bmatrix}

  1 & 4 & 4 & \big| & 1 & 0 & 0 \\\  0 & -3 & 0 & \big| & -2 & 1 & 0 \\\  0 & -6 & -3 & \big| & -3 & 0 & 1

\end{bmatrix}

\newline

\rightsquigarrow

\begin{bmatrix}

  1 & 4 & 4 & \big| & 1 & 0 & 0 \\\  0 & 1 & 0 & \big| & 2/3 & -1/3 & 0 \\\  0 & -6 & -3 & \big| & -3 & 0 & 1

\end{bmatrix}

\rightsquigarrow$$

$$\begin{bmatrix}

  1 & 0 & 4 & \big| & -5/3 & 4/3 & 0 \\\  0 & 1 & 0 & \big| & 2/3 & -1/3 & 0 \\\  0 & 0 & -3 & \big| & 1 & -2 & 1

\end{bmatrix}

\newline

\rightsquigarrow

\begin{bmatrix}

  1 & 0 & 4 & \big| & -5/3 & 4/3 & 0 \\\  0 & 1 & 0 & \big| & 2/3 & -1/3 & 0 \\\  0 & 0 & 1 & \big| & -1/3 & 2/3 & -1/3

\end{bmatrix}

\rightsquigarrow

\begin{bmatrix}

  1 & 0 & 0 & \big| & -1/3 & -4/3 & 4/3 \\\  0 & 1 & 0 & \big| & 2/3 & -1/3 & 0 \\\  0 & 0 & 1 & \big| & -1/3 & 2/3 & -1/3

\end{bmatrix}$$



The matrix $\mathrm{B} = \begin{bmatrix}

  -1/3 \& -4/3 \& 4/3 \\\  2/3 \& -1/3 \& 0 \\\  -1/3 \& 2/3 \& -1/3

\end{bmatrix}$ is then the inverse of the original matrix A.



\paragraph{Inverse of a Linear Transformation}



We now consider how to represent the inverse of a linear map. We start by recalling some facts about function inverses. Some functions have no inverse, or have an inverse on the left side or right side only.



\paragraph{Definitions}



Where $\pi:\mathbb{R}^3\to \mathbb{R}^2$ is the projection map



$$\begin{pmatrix} x \\\\ y \\\\ z \end{pmatrix}

\mapsto

\begin{pmatrix} x \\\\ y \end{pmatrix}$$



and $\eta:\mathbb{R}^2\to \mathbb{R}^3$ is the embedding



$$\begin{pmatrix} x \\\\ y \end{pmatrix}

\mapsto

\begin{pmatrix} x \\\\ y \\\\ 0 \end{pmatrix}$$



the composition $\pi\circ \eta$ is the identity map on $\mathbb{R}^2$.



$$\begin{pmatrix} x \\\\ y \end{pmatrix}

\stackrel{\eta}{\longmapsto}

\begin{pmatrix} x \\\\ y \\\\ 0 \end{pmatrix}

\stackrel{\pi}{\longmapsto}

\begin{pmatrix} x \\\\ y \end{pmatrix}$$



We say $\pi$ is a \textbf{left inverse map} of $\eta$ or, what is the same thing, that $\eta$ is a \textbf{right inverse map} of $\pi$. However, composition in the other order $\eta\circ \pi$ doesn't give the identity map--- here is a vector that is not sent to itself under $\eta\circ \pi$.



$$\begin{pmatrix} 0 \\\\ 0 \\\\ 1 \end{pmatrix}

\stackrel{\pi}{\longmapsto}

\begin{pmatrix} 0 \\\\ 0 \end{pmatrix}

\stackrel{\eta}{\longmapsto}

\begin{pmatrix} 0 \\\\ 0 \\\\ 0 \end{pmatrix}$$



In fact, the projection $\pi$ has no left inverse at all. For, if $f$ were to be a left inverse of $\pi$ then we would have



$$\begin{pmatrix} x \\\\ y \\\\ z \end{pmatrix}

\stackrel{\pi}{\longmapsto}

\begin{pmatrix} x \\\\ y \end{pmatrix}

\stackrel{f}{\longmapsto}

\begin{pmatrix} x \\\\ y \\\\ z \end{pmatrix}$$



for all of the infinitely many $z$'s. But no function $f$ can send a single argument to more than one value.



(An example of a function with no inverse on either side is the zero transformation on $\mathbb{R}^2$.)



Some functions have a \textbf{two-sided inverse map}, another function that is the inverse of the first, both from the left and from the right. For instance, the map given by $\vec{v}\mapsto 2\cdot \vec{v}$ has the two-sided inverse $\vec{v}\mapsto (1/2)\cdot\vec{v}$. In this subsection we will focus on two-sided inverses. The appendix shows that a function has a two-sided inverse if and only if it is both one-to-one and onto. The appendix also shows that if a function $f$ has a two-sided inverse then it is unique, and so it is called "the" inverse, and is denoted $f^{-1}$. So our purpose in this subsection is, where a linear map $h$ has an inverse, to find the relationship between ${\rm Rep} _{B,D}(h)$ and ${\rm Rep} _{D,B}(h^{-1})$.



A matrix $G$ is a \textbf{left inverse matrix} of the matrix $H$ if $GH$ is the identity matrix. It is a \textbf{right inverse matrix} if $HG$ is the identity. A matrix $H$ with a two-sided inverse is an \textbf{invertible matrix}. That two-sided inverse is called \textbf{the inverse matrix} and is denoted $H^{-1}$.



Because of the correspondence between linear maps and matrices, statements about map inverses translate into statements about matrix inverses.



\paragraph{Lemmas and Theorems}



\begin{itemize}
    \item If a matrix has both a left inverse and a right inverse then the two are equal.
    \item A matrix is invertible if and only if it is nonsingular.
    \item Proof: \textit{(For both results.)} Given a matrix $H$, fix spaces of appropriate dimension for the domain and codomain. Fix bases for these spaces. With respect to these bases, $H$ represents a map $h$. The statements are true about the map and therefore they are true about the matrix.
    \item A product of invertible matrices is invertible--- if $G$ and $H$ are invertible and if $GH$ is defined then $GH$ is invertible and $(GH)^{-1}=H^{-1}G^{-1}$.
    \item Proof: \textit{(This is just like the prior proof except that it requires two maps.)} Fix appropriate spaces and bases and consider the represented maps $h$ and $g$. Note that $h^{-1}g^{-1}$ is a two-sided map inverse of $gh$ since $(h^{-1}g^{-1})(gh) = h^{-1}(\text{id})h = h^{-1}h =\text{id}$ and $(gh)(h^{-1}g^{-1}) = g(\text{id})g^{-1} = gg^{-1} = \text{id}$. This equality is reflected in the matrices representing the maps, as required.
\end{itemize}

Here is the arrow diagram giving the relationship between map inverses and matrix inverses. It is a special case of the diagram for function composition and matrix multiplication.



% [Image: ]



Beyond its place in our general program of seeing how to represent map operations, another reason for our interest in inverses comes from solving linear systems. A linear system is equivalent to a matrix equation, as here.



$$\begin{aligned} \begin{aligned} x _1 & + & x _2 & = & 3 \\\\ 2x _1 & - & x _2 & = & 2 \end{aligned} & \quad \Leftrightarrow \begin{pmatrix} 1 & 1 \\\\ 2 & -1 \end{pmatrix} \begin{pmatrix} x _1 \\\\ x _2 \end{pmatrix} = \begin{pmatrix} 3 \\\\ 2 \end{pmatrix} & \quad\quad\quad \left( \* \right) \end{aligned}$$



By fixing spaces and bases (e.g., $\mathbb{R}^2, \mathbb{R}^2$ and $\mathcal{E} _2,\mathcal{E} _2$), we take the matrix $H$ to represent some map $h$. Then solving the system is the same as asking: what domain vector $\vec{x}$ is mapped by $h$ to the result $\vec{d}\,$? If we could invert $h$ then we could solve the system by multiplying ${\rm Rep} _{D,B}(h^{-1})\cdot{\rm Rep} _{D}(\vec{d})$ to get ${\rm Rep} _{B}(\vec{x})$.



\paragraph{Example 2}



We can find a left inverse for the matrix just given



$$\begin{pmatrix} m  & n  \\\\ p  & q \end{pmatrix} \begin{pmatrix} 1  & 1  \\\\ 2  & -1 \end{pmatrix} = \begin{pmatrix} 1  & 0  \\\\ 0  & 1 \end{pmatrix}$$



by using Gauss' method to solve the resulting linear system.



$$\begin{aligned}

m & + & 2n & & & & & = & 1 \\\m & - & n & & & & & = & 0 \\\ & & & & p & + & 2q & = & 0 \\\ & & & & p & - & q & = & 1

\end{aligned}$$



Answer: $m=1/3$, $n=1/3$, $p=2/3$, and $q=-1/3$. This matrix is actually the two-sided inverse of $H$, as can easily be checked. With it we can solve the system ($*$) above by applying the inverse.



$$\begin{pmatrix} x \\\\ y \end{pmatrix} =\begin{pmatrix} 1/3 & 1/3  \\\\ 2/3 & -1/3 \end{pmatrix} \begin{pmatrix} 3 \\\\ 2 \end{pmatrix} =\begin{pmatrix} 5/3 \\\\ 4/3 \end{pmatrix}$$



\paragraph{Remark}



Why solve systems this way, when Gauss' method takes less arithmetic (this assertion can be made precise by counting the number of arithmetic operations, as computer algorithm designers do)? Beyond its conceptual appeal of fitting into our program of discovering how to represent the various map operations, solving linear systems by using the matrix inverse has at least two advantages.



First, once the work of finding an inverse has been done, solving a system with the same coefficients but different constants is easy and fast: if we change the entries on the right of the system ($*$) then we get a related problem



$$\begin{pmatrix} 1  & 1  \\\\ 2  & -1 \end{pmatrix} \begin{pmatrix} x \\\\ y \end{pmatrix} = \begin{pmatrix} 5 \\\\ 1 \end{pmatrix}$$



with a related solution method.



$$\begin{pmatrix} x \\\\ y \end{pmatrix} = \begin{pmatrix} 1/3  & 1/3  \\\\ 2/3  & -1/3 \end{pmatrix} \begin{pmatrix} 5 \\\\ 1 \end{pmatrix} = \begin{pmatrix} 2 \\\\ 3 \end{pmatrix}$$



In applications, solving many systems having the same matrix of coefficients is common.



Another advantage of inverses is that we can explore a system's sensitivity to changes in the constants. For example, tweaking the $3$ on the right of the system ($*$) to



$$\begin{pmatrix} 1  & 1  \\\\ 2  & -1 \end{pmatrix} \begin{pmatrix} x _1 \\\\ x _2 \end{pmatrix} = \begin{pmatrix} 3.01 \\\\ 2 \end{pmatrix}$$



can be solved with the inverse.



$$\begin{pmatrix} 1/3  & 1/3  \\\\ 2/3  & -1/3 \end{pmatrix} \begin{pmatrix} 3.01 \\\\ 2 \end{pmatrix} = \begin{pmatrix} (1/3)(3.01)+(1/3)(2) \\\\ (2/3)(3.01)-(1/3)(2) \end{pmatrix}$$



to show that $x _1$ changes by $1/3$ of the tweak while $x _2$ moves by $2/3$ of that tweak. This sort of analysis is used, for example, to decide how accurately data must be specified in a linear model to ensure that the solution has a desired accuracy.



We finish by describing the computational procedure usually used to find the inverse matrix.



\begin{itemize}
    \item A matrix is invertible if and only if it can be written as the product of elementary reduction matrices. The inverse can be computed by applying to the identity matrix the same row steps, in the same order, as are used to Gauss-Jordan reduce the invertible matrix.
    \item Proof: A matrix $H$ is invertible if and only if it is nonsingular and thus Gauss-Jordan reduces to the identity. This reduction can be done with elementary matrices
\end{itemize}

$$R _{r}\cdot R _{r-1}\cdot\dots R _{1}\cdot H=I$$.



> This equation gives the two halves of the result.

>

> First, elementary matrices are invertible and their inverses are also elementary. Applying $R _r^{-1}$ to the left of both sides of that equation, then $R _{r-1}^{-1}$, etc., gives $H$ as the product of elementary matrices $H=R _1^{-1}\cdots R _r^{-1}\cdot I$ (the $I$ is here to cover the trivial $r=0$ case).

>

> Second, matrix inverses are unique and so comparison of the above equation with $H^{-1}H=I$ shows that $H^{-1}=R _r\cdot R _{r-1}\dots R _1\cdot I$. Therefore, applying $R _1$ to the identity, followed by $R _2$, etc., yields the inverse of $H$.



\paragraph{Example 3}



To find the inverse of



$$\begin{pmatrix} 1  & 1  \\\\ 2  & -1 \end{pmatrix}$$



we do Gauss-Jordan reduction, meanwhile performing the same operations on the identity. For clerical convenience we write the matrix and the identity side-by-side, and do the reduction steps together.



$$\begin{array}{rcl} \left(\begin{array}{cc|cc} 1  & 1   & 1  & 0  \\\\ 2  & -1  & 0  & 1 \end{array}\right) &\xrightarrow[]{-2\rho _1+\rho _2} &\left(\begin{array}{cc|cc} 1  & 1   & 1  & 0  \\\\ 0  & -3  & -2 & 1 \end{array}\right) \\\\ &\xrightarrow[]{-1/3\rho _2} &\left(\begin{array}{cc|cc} 1  & 1   & 1   & 0     \\\\ 0  & 1   & 2/3 & -1/3 \end{array}\right)  \\\\ &\xrightarrow[]{-\rho _2+\rho _1} &\left(\begin{array}{cc|cc} 1  & 0   & 1/3 & 1/3  \\\\ 0  & 1   & 2/3 & -1/3 \end{array}\right) \end{array}$$



This calculation has found the inverse.



$$\begin{pmatrix} 1  & 1  \\\\ 2  & -1 \end{pmatrix}^{-1} = \begin{pmatrix} 1/3 & 1/3  \\\\ 2/3 & -1/3 \end{pmatrix}$$



\paragraph{Example 4}



This one happens to start with a row swap.



$$\begin{array}{rcl}

\left(\begin{array}{ccc|ccc}

0  & 3  & -1  & 1  & 0  & 0  \\\1  & 0  & 1   & 0  & 1  & 0  \\\1  & -1 & 0   & 0  & 0  & 1

\end{array}\right)

& \xrightarrow[]{\rho _1\leftrightarrow\rho _2}

&   \left(\begin{array}{ccc|ccc}

1  & 0  & 1   & 0  & 1  & 0  \\\0  & 3  & -1  & 1  & 0  & 0  \\\1  & -1 & 0   & 0  & 0  & 1

\end{array}\right)      \\\& \xrightarrow[]{-\rho _1+\rho _3}

&   \left(\begin{array}{ccc|ccc}

1  & 0  & 1   & 0  & 1  & 0  \\\0  & 3  & -1  & 1  & 0  & 0  \\\0  & -1 & -1  & 0  & -1 & 1

\end{array}\right) \\\& \vdots \\\& \xrightarrow[]{}

&   \left(\begin{array}{ccc|ccc}

1  & 0  & 0   & 1/4  & 1/4  & 3/4  \\\0  & 1  & 0   & 1/4  & 1/4  & -1/4 \\\0  & 0  & 1   & -1/4 & 3/4  & -3/4

\end{array}\right)

\end{array}$$



\paragraph{Example 5}



A non-invertible matrix is detected by the fact that the left half won't reduce to the identity.



$$\left(\begin{array}{cc|cc}

1  & 1   & 1  & 0  \\\2  & 2   & 0  & 1

\end{array}\right)

\xrightarrow[]{-2\rho _1+\rho _2}

\left(\begin{array}{cc|cc}

1  & 1   & 1  & 0  \\\0  & 0   & -2 & 1

\end{array}\right)$$



This procedure will find the inverse of a general $n \times n$ matrix. The $2 \times 2$ case is handy.



\paragraph{Resources}



\begin{itemize}
    \item \href{https://www.khanacademy.org/math/algebra-home/alg-matrices/alg-adding-and-subtracting-matrices/v/matrix-addition-and-subtraction-1}{Matrix Addition and Subtraction}, Khan Academy
    \item \href{https://www.khanacademy.org/math/algebra-home/alg-matrices/alg-multiplying-matrices-by-scalars/v/scalar-multiplication}{Multiplying Matrices by Scalars}, Khan Academy
    \item \href{https://www.khanacademy.org/math/algebra-home/alg-matrices/alg-multiplying-matrices-by-matrices/v/matrix-multiplication-intro}{Matrix Multiplication}, Khan Academy
    \item \href{https://openstax.org/books/college-algebra/pages/7-5-matrices-and-matrix-operations}{Matrices and Matrix Operations}, OpenStax
    \item \href{https://www.youtube.com/watch?v=0L90Kkn90J8}{Multiplying a 2-by-3 Matrix and 3-by-2 Matrix}, patrickJMT
\end{itemize}

\paragraph{Licensing}



Content obtained and/or adapted from:



\begin{itemize}
    \item \href{https://en.wikibooks.org/wiki/Linear_Algebra/Inverses}{Inverses, Wikibooks: Linear Algebra} under a CC BY-SA license
    \item \href{https://en.wikibooks.org/wiki/Linear_Algebra/The_Inverse_of_a_Matrix}{The Inverse of a Matrix, Wikibooks: Linear Algebra} under a CC BY-SA license
\end{itemize}

\subsection{Learning Objectives}

\begin{enumerate}
    \item \textbf{Add and Subtract Matrices:}
    \item \textbf{Outcome:} Students will be able to correctly add and subtract matrices of the same dimension.
    \item \textbf{Details:} This includes understanding the requirement for matrices to have identical dimensions and applying the operation element-wise across corresponding entries.
    \item \textbf{Multiply Matrices by Scalars:}
    \item \textbf{Outcome:} Students will be able to multiply a matrix by a scalar value.
    \item \textbf{Details:} This includes multiplying every entry in the matrix by the scalar and simplifying the resulting matrix.
    \item \textbf{Multiply Two Matrices:}
    \item \textbf{Outcome:} Students will be able to multiply two matrices and determine the resulting matrix's dimensions.
    \item \textbf{Details:} This includes understanding the conditions under which two matrices can be multiplied (the number of columns in the first matrix must equal the number of rows in the second matrix) and computing the product by performing dot products for the resulting matrix entries.
\end{enumerate}

\subsection{Assessment Instrument}

\begin{enumerate}
    \item \textbf{Assessment Instrument: Matrix Operations Quiz} \textbf{Format:}
    \item \textbf{Type:} Mixed-format quiz (multiple-choice, short answer, and problem-solving questions)
    \item \textbf{Duration:} 45 minutes
    \item \textbf{Total Marks:} 30 points \textbf{Sections:}
    \item \textbf{Adding and Subtracting Matrices (10 points):}
    \item \textbf{Question 1:} Add the following matrices: $$ \begin{bmatrix} 1 & 2 & 3 \\\\ 4 & 5 & 6 \\\\ 7 & 8 & 9 \end{bmatrix}  +  \begin{bmatrix} 9 & 8 & 7 \\\\ 6 & 5 & 4 \\\\ 3 & 2 & 1 \end{bmatrix} $$
    \item \textbf{Question 2:} Subtract the following matrices: $$ \begin{bmatrix} 4 & 6 & 8 \\\\ 2 & 4 & 6 \\\\ 0 & 2 & 4 \end{bmatrix}  -  \begin{bmatrix} 1 & 3 & 5 \\\\ 0 & 2 & 4 \\\\ 6 & 8 & 10 \end{bmatrix} $$
    \item \textbf{Question 3:} Determine if the following operation is possible, and if so, perform the addition: $$ \begin{bmatrix} 2 & 4 \\\\ 6 & 8  \end{bmatrix} +  \begin{bmatrix} 1 & 3 \\\\ 5 & 7  \end{bmatrix} $$
    \item \textbf{Multiplying Matrices by Scalars (10 points):}
    \item \textbf{Question 4:} Multiply the following matrix by the scalar 3: $$ \begin{bmatrix} 2 & -4 & 1 \\\\ 0 & 5 & -3 \\\\ -2 & 7 & 6 \end{bmatrix} $$
    \item \textbf{Question 5:} Multiply the matrix below by the scalar $-\frac{1}{2}$: $$ \begin{bmatrix} 4 & -8 & 2 \\\\ -6 & 10 & 12 \\\\ 0 & -2 & 3 \end{bmatrix} $$
    \item \textbf{Question 6:} Explain the result when any matrix is multiplied by the scalar 0.
    \item \textbf{Matrix Multiplication (10 points):}
    \item \textbf{Question 7:} Multiply the following matrices: $$ \begin{bmatrix} 1 & 2 \\\\ 3 & 4 \end{bmatrix} \times \begin{bmatrix} 2 & 0 \\\\ 1 & 3 \end{bmatrix} $$
    \item \textbf{Question 8:} Determine if the following multiplication is possible, and if so, compute the product: $$ \begin{bmatrix} 1 & 2 & 3 \\\\ 4 & 5 & 6 \end{bmatrix} \times \begin{bmatrix} 7 & 8 \\\\ 9 & 10 \\\\ 11 & 12 \end{bmatrix} $$
    \item \textbf{Question 9:} Explain why matrix multiplication is not commutative by providing an example with matrices $A$ and $B$ where $AB \neq BA$. \textbf{Scoring Rubric:}
    \item \textbf{Adding and Subtracting Matrices:} 3 points for each correctly performed operation (total 9 points) and 1 point for determining if the operation is possible (Question 3).
    \item \textbf{Multiplying Matrices by Scalars:} 3 points for each correctly multiplied matrix (total 9 points) and 1 point for the explanation in Question 6.
    \item \textbf{Matrix Multiplication:} 3 points for each correctly multiplied matrix (total 9 points) and 1 point for the explanation in Question 9. The quiz assesses students' understanding and application of matrix operations, including addition, subtraction, scalar multiplication, and matrix multiplication. It ensures that they can perform these operations accurately and understand the underlying principles.
\end{enumerate}

%%===============================================================
\section{Solving Systems with Inverses}
\label{sec:solving-systems-with-inverses}
%%===============================================================

\paragraph{Solving a linear system}



There are several algorithms for solving a system of linear equations.



\paragraph{Describing the solution}



When the solution set is finite, it is reduced to a single element. In this case, the unique solution is described by a sequence of equations whose left-hand sides are the names of the unknowns and right-hand sides are the corresponding values, for example $(x=3, \;y=-2,\; z=6)$. When an order on the unknowns has been fixed, for example the alphabetical order the solution may be described as a vector of values, like $(3, \,-2,\, 6)$ for the previous example.



To describe a set with an infinite number of solutions, typically some of the variables are designated as \textbf{free} (or \textbf{independent}, or as \textbf{parameters}), meaning that they are allowed to take any value, while the remaining variables are \textbf{dependent} on the values of the free variables.



For example, consider the following system:



$$ \begin{aligned} x &&\; + \;&& 3y &&\; - \;&& 2z &&\; = \;&& 5 & \\\\ 3x &&\; + \;&& 5y &&\; + \;&& 6z &&\; = \;&& 7 & \end{aligned} $$

The solution set to this system can be described by the following equations:



$$x=-7z-1 \quad \text{ and } \quad y=3z+2\text{.}$$ Here \textit{z} is the free variable, while \textit{x} and \textit{y} are dependent on \textit{z}. Any point in the solution set can be obtained by first choosing a value for \textit{z}, and then computing the corresponding values for \textit{x} and \textit{y}.



Each free variable gives the solution space one degree of freedom, the number of which is equal to the dimension of the solution set. For example, the solution set for the above equation is a line, since a point in the solution set can be chosen by specifying the value of the parameter \textit{z}. An infinite solution of higher order may describe a plane, or higher-dimensional set.



Different choices for the free variables may lead to different descriptions of the same solution set. For example, the solution to the above equations can alternatively be described as follows:



$$y=-\frac{3}{7}x + \frac{11}{7}\;\;\;\;\text{and}\;\;\;\;z=-\frac{1}{7}x-\frac{1}{7}\text{.}$$ Here \textit{x} is the free variable, and \textit{y} and \textit{z} are dependent.



\paragraph{Elimination of variables}



The simplest method for solving a system of linear equations is to repeatedly eliminate variables. This method can be described as follows:



1.  In the first equation, solve for one of the variables in terms of the others.

2.  Substitute this expression into the remaining equations. This yields a system of equations with one fewer equation and one fewer unknown.

3.  Repeat until the system is reduced to a single linear equation.

4.  Solve this equation, and then back-substitute until the entire solution is found.



For example, consider the following system:



$$\begin{aligned}

 x &&\; + \;&& 3y &&\; - \;&& 2z &&\; = \;&& 5 & \\\3x &&\; + \;&& 5y &&\; + \;&& 6z &&\; = \;&& 7 & \\\2x &&\; + \;&& 4y &&\; + \;&& 3z &&\; = \;&& 8 &

\end{aligned}$$

Solving the first equation for \textit{x} gives \textit{x} = 5 + 2\textit{z} - 3\textit{y}, and plugging this into the second and third equation yields



$$ \begin{aligned} -4y & & \; + \; & & 12z & & \; = \; & & -8 & \\\\ -2y & & \; + \; & & 7z & & \; = \; & & -2 & \end{aligned} $$



Solving the first of these equations for \textit{y} yields \textit{y} = 2 + 3\textit{z}, and plugging this into the second equation yields \textit{z} = 2. We now have:



$$

\begin{aligned}

 x &&\; = \;&& 5 &&\; + \;&& 2z &&\; - \;&& 3y & \\\ y &&\; = \;&& 2 &&\; + \;&& 3z && && & \\\ z &&\; = \;&& 2 && && && && &

\end{aligned}

$$

Substituting \textit{z} = 2 into the second equation gives \textit{y} = 8, and substituting \textit{z} = 2 and \textit{y} = 8 into the first equation yields \textit{x} = -15. Therefore, the solution set is the single point $(x, y, z) = (-15, 8, 2)$.



\paragraph{Row reduction}



In \textbf{row reduction} (also known as \textbf{Gaussian elimination}), the linear system is represented as an augmented matrix:



$$\left[\begin{array}{rrr|r}

1 & 3 & -2 & 5 \\\3 & 5 & 6 & 7 \\\2 & 4 & 3 & 8

\end{array}\right]\text{.}$$

This matrix is then modified using elementary row operations until it reaches reduced row echelon form. There are three types of elementary row operations:



> \textbf{Type 1}: Swap the positions of two rows.

> \textbf{Type 2}: Multiply a row by a nonzero scalar.

> \textbf{Type 3}: Add to one row a scalar multiple of another.



Because these operations are reversible, the augmented matrix produced always represents a linear system that is equivalent to the original.



There are several specific algorithms to row-reduce an augmented matrix, the simplest of which are Gaussian elimination and Gauss–Jordan elimination. The following computation shows Gauss–Jordan elimination applied to the matrix above:



$$\begin{aligned}\left[\begin{array}{rrr|r}

1 & 3 & -2 & 5 \\\3 & 5 & 6 & 7 \\\2 & 4 & 3 & 8

\end{array}\right]&\sim

\left[\begin{array}{rrr|r}

1 & 3 & -2 & 5 \\\0 & -4 & 12 & -8 \\\2 & 4 & 3 & 8

\end{array}\right]\sim

\left[\begin{array}{rrr|r}

1 & 3 & -2 & 5 \\\0 & -4 & 12 & -8 \\\0 & -2 & 7 & -2

\end{array}\right]\sim

\left[\begin{array}{rrr|r}

1 & 3 & -2 & 5 \\\0 & 1 & -3 & 2 \\\0 & -2 & 7 & -2

\end{array}\right]

\\\&\sim

\left[\begin{array}{rrr|r}

1 & 3 & -2 & 5 \\\0 & 1 & -3 & 2 \\\0 & 0 & 1 & 2

\end{array}\right]\sim

\left[\begin{array}{rrr|r}

1 & 3 & -2 & 5 \\\0 & 1 & 0 & 8 \\\0 & 0 & 1 & 2

\end{array}\right]\sim

\left[\begin{array}{rrr|r}

1 & 3 & 0 & 9 \\\0 & 1 & 0 & 8 \\\0 & 0 & 1 & 2

\end{array}\right]\sim

\left[\begin{array}{rrr|r}

1 & 0 & 0 & -15 \\\0 & 1 & 0 & 8 \\\0 & 0 & 1 & 2

\end{array}\right].\end{aligned}$$

The last matrix is in reduced row echelon form, and represents the system \textit{x} = -15, \textit{y} = 8, \textit{z} = 2. A comparison with the example in the previous section on the algebraic elimination of variables shows that these two methods are in fact the same; the difference lies in how the computations are written down.



\paragraph{Cramer's rule}



\textbf{Cramer's rule} is an explicit formula for the solution of a system of linear equations, with each variable given by a quotient of two determinants. For example, the solution to the system



$$\begin{aligned}

 x &\; + &\; 3y &\; - &\; 2z &\; = &\; 5 \\\3x &\; + &\; 5y &\; + &\; 6z &\; = &\; 7 \\\2x &\; + &\; 4y &\; + &\; 3z &\; = &\; 8

\end{aligned}$$ is given by



$$ x=\frac {\, \begin{vmatrix}5 & 3 & -2 \\\\ 7 & 5 & 6 \\\\ 8 & 4 & 3\end{vmatrix} \,} {\, \begin{vmatrix}1 & 3 & -2 \\\\ 3 & 5 & 6 \\\\ 2 & 4 & 3\end{vmatrix} \,} ,\;\;\;\; y=\frac {\, \begin{vmatrix}1 & 5 & -2 \\\\ 3 & 7 & 6 \\\\ 2 & 8 & 3\end{vmatrix} \,} {\, \begin{vmatrix}1 & 3 & -2 \\\\ 3 & 5 & 6 \\\\ 2 & 4 & 3\end{vmatrix} \,} ,\;\;\;\; z=\frac {\, \begin{vmatrix}1 & 3 & 5 \\\\ 3 & 5 & 7 \\\\ 2 & 4 & 8\end{vmatrix} \,} {\, \begin{vmatrix}1 & 3 & -2 \\\\ 3 & 5 & 6 \\\\ 2 & 4 & 3\end{vmatrix} \,}. $$ 

For each variable, the denominator is the determinant of the matrix of coefficients, while the numerator is the determinant of a matrix in which one column has been replaced by the vector of constant terms.



Though Cramer's rule is important theoretically, it has little practical value for large matrices, since the computation of large determinants is somewhat cumbersome. (Indeed, large determinants are most easily computed using row reduction.) Further, Cramer's rule has very poor numerical properties, making it unsuitable for solving even small systems reliably, unless the operations are performed in rational arithmetic with unbounded precision.



\paragraph{Matrix solution}



If the equation system is expressed in the matrix form $A\mathbf{x}=\mathbf{b}$, the entire solution set can also be expressed in matrix form. If the matrix \textit{A} is square (has \textit{m} rows and \textit{n}=\textit{m} columns) and has full rank (all \textit{m} rows are independent), then the system has a unique solution given by



$$\mathbf{x}=A^{-1}\mathbf{b}$$



where $A^{-1}$ is the inverse of \textit{A}. More generally, regardless of whether \textit{m}=\textit{n} or not and regardless of the rank of \textit{A}, all solutions (if any exist) are given using the Moore-Penrose pseudoinverse of \textit{A}, denoted $A^+$, as follows:



$$\mathbf{x}=A^+ \mathbf{b} + \left(I - A^+ A\right)\mathbf{w}$$



where $\mathbf{w}$ is a vector of free parameters that ranges over all possible \textit{n}×1 vectors. A necessary and sufficient condition for any solution(s) to exist is that the potential solution obtained using $\mathbf{w}=\mathbf{0}$ satisfy $A\mathbf{x}=\mathbf{b}$ — that is, that $AA^+ \mathbf{b}=\mathbf{b}.$ If this condition does not hold, the equation system is inconsistent and has no solution. If the condition holds, the system is consistent and at least one solution exists. For example, in the above-mentioned case in which \textit{A} is square and of full rank, $A^+$ simply equals $A^{-1}$ and the general solution equation simplifies to



$$\mathbf{x}=A^{-1}\mathbf{b} + \left(I - A^{-1}A\right)\mathbf{w} = A^{-1}\mathbf{b} + \left(I-I\right)\mathbf{w} = A^{-1}\mathbf{b}$$ as previously stated, where $\mathbf{w}$ has completely dropped out of the solution, leaving only a single solution. In other cases, though, $\mathbf{w}$ remains and hence an infinitude of potential values of the free parameter vector $\mathbf{w}$ give an infinitude of solutions of the equation.



\paragraph{Other methods}



While systems of three or four equations can be readily solved by hand (see Cracovian), computers are often used for larger systems. The standard algorithm for solving a system of linear equations is based on Gaussian elimination with some modifications. Firstly, it is essential to avoid division by small numbers, which may lead to inaccurate results. This can be done by reordering the equations if necessary, a process known as \textit{pivoting}. Secondly, the algorithm does not exactly do Gaussian elimination, but it computes the LU decomposition of the matrix \textit{A}. This is mostly an organizational tool, but it is much quicker if one has to solve several systems with the same matrix \textit{A} but different vectors \textbf{b}.



If the matrix \textit{A} has some special structure, this can be exploited to obtain faster or more accurate algorithms. For instance, systems with a symmetric positive definite matrix can be solved twice as fast with the Cholesky decomposition. Levinson recursion is a fast method for Toeplitz matrices. Special methods exist also for matrices with many zero elements (so-called sparse matrices), which appear often in applications.



A completely different approach is often taken for very large systems, which would otherwise take too much time or memory. The idea is to start with an initial approximation to the solution (which does not have to be accurate at all), and to change this approximation in several steps to bring it closer to the true solution. Once the approximation is sufficiently accurate, this is taken to be the solution to the system. This leads to the class of iterative methods. For some sparse matrices, the introduction of randomness improves the speed of the iterative methods.



There is also a quantum algorithm for linear systems of equations.



\paragraph{Resources}



\begin{itemize}
    \item \href{https://mathresearch.utsa.edu/wikiFiles/MAT1053/Solving_Systems_with_Inverses/MAT1053_M10.2Solving_Systems_with_Inverses.pdf}{Solving Systems with Inverses}, Book Chapter
    \item \href{https://mathresearch.utsa.edu/wikiFiles/MAT1053/Solving_Systems_with_Inverses/MAT1053_M10.2Solving_Systems_with_InversesGN.pdf}{Guided Notes}
\end{itemize}

\paragraph{Licensing}



Content obtained and/or adapted from:



\begin{itemize}
    \item \href{https://en.wikipedia.org/wiki/System_of_linear_equations}{System of linear equations, Wikipedia} under a CC BY-SA license
\end{itemize}

\subsection{Learning Objectives}

\begin{enumerate}
    \item Here are three Student Learning Objectives (SLOs) for solving a linear system, following the provided format:
    \item \textbf{Understand Solution Sets of Linear Systems:}
    \item \textbf{Outcome:} Students will be able to identify and classify the solution set of a linear system as either a unique solution, infinitely many solutions, or no solution.
    \item \textbf{Details:} This includes understanding the implications of the rank of the coefficient matrix, the augmented matrix, and the consistency of the system.
    \item \textbf{Apply Gaussian Elimination for Solving Linear Systems:}
    \item \textbf{Outcome:} Students will be able to solve a linear system of equations using Gaussian elimination, including row reduction to reduced row echelon form (RREF).
    \item \textbf{Details:} This includes performing elementary row operations, interpreting the resulting matrix, and back-substituting to find the solution.
    \item \textbf{Use Matrix Methods to Solve Linear Systems:}
    \item \textbf{Outcome:} Students will be able to solve a system of linear equations using matrix operations, including applying the inverse matrix and Cramer’s rule.
    \item \textbf{Details:} This includes understanding when each method is applicable, calculating the determinant, and interpreting the solution in terms of the original system. These SLOs cover the key methods and concepts necessary for understanding and solving linear systems.
\end{enumerate}

\subsection{Assessment Instrument}

\begin{enumerate}
    \item \textbf{Assessment Instrument: Solving Systems with Inverses Quiz} \textbf{Format:}
    \item \textbf{Type:} Mixed-format quiz (multiple-choice, short answer, and problem-solving questions)
    \item \textbf{Duration:} 45 minutes
    \item \textbf{Total Marks:} 30 points \textbf{Sections:}
    \item \textbf{Identifying Solutions of Linear Systems (10 points):}
    \item \textbf{Question 1:} Consider the system of equations: $$\begin{aligned} x + 2y &= 5 \\\\ 3x + 4y &= 11 \end{aligned}$$ Solve the system using the inverse of the coefficient matrix. Show all steps. (5 points)
    \item \textbf{Question 2:} Given the system: $$ \begin{aligned} 2x - y &= 3 \\\\ 4x + z &= 7 \\\\ y + 2z &= 4 \end{aligned} $$ Determine if the system has a unique solution, infinitely many solutions, or no solution using the matrix inverse method. Justify your answer. (5 points)
    \item \textbf{Application of Matrix Inverses (10 points):}
    \item \textbf{Question 3:} If the matrix $ A = \begin{pmatrix} 1 & 2 \\\\ 3 & 5 \end{pmatrix} $ and the inverse $ A^{-1} $ is given by $ A^{-1} = \begin{pmatrix} 5 & -2 \\\\ -3 & 1 \end{pmatrix} $, solve the equation $ A\mathbf{x} = \mathbf{b} $ where $ \mathbf{b} = \begin{pmatrix} 4 \\\\ 10 \end{pmatrix} $. (5 points)
    \item \textbf{Question 4:} Solve the following system of equations using matrix inverses and state the solution as a vector: $$\begin{aligned} 3x + y &= 7 \\\\ 2x + 4y &= 10 \end{aligned}$$ (5 points)
    \item \textbf{Understanding the Inverse Matrix Method (10 points):}
    \item \textbf{Question 5:} Explain why the inverse matrix method cannot be used to solve a system if the determinant of the coefficient matrix is zero. Provide a brief example to illustrate your explanation. (5 points)
    \item \textbf{Question 6:} Consider the system $ A\mathbf{x} = \mathbf{b} $ where $ A $ is a 3x3 matrix. If $ \mathbf{x} = A^{-1}\mathbf{b} $ gives a solution, describe the conditions under which this system has a unique solution. (5 points) \textbf{Scoring Rubric:}
    \item \textbf{Identifying Solutions of Linear Systems:} 5 points for each correctly solved system (total 10 points)
    \item \textbf{Application of Matrix Inverses:} 5 points for each correct application of the matrix inverse method (total 10 points)
    \item \textbf{Understanding the Inverse Matrix Method:} 5 points for each correct explanation and justification (total 10 points) This quiz assesses students' understanding and application of solving linear systems using matrix inverses. It ensures that they can correctly identify solution types, apply matrix operations, and understand the underlying concepts of the inverse matrix method.
\end{enumerate}

%%===============================================================
\section{Sets - Definitions}
\label{sec:sets-definitions}
%%===============================================================

% [Image: A set of polygons in an Euler diagram]



Figure: A set of polygons in an Euler diagram.



In mathematics, a \textbf{set} is a collection of elements. The elements that make up a set can be any kind of mathematical objects: numbers, symbols, points in space, lines, other geometrical shapes, variables, or even other sets. The set with no element is the empty set; a set with a single element is a singleton. A set may have a finite number of elements or be an infinite set. Two sets are equal if they have precisely the same elements.



Sets are ubiquitous in modern mathematics. Indeed, set theory, more specifically Zermelo–Fraenkel set theory, has been the standard way to provide rigorous foundations for all branches of mathematics since the first half of the 20th century.



\paragraph{Origin}



The concept of a set emerged in mathematics at the end of the 19th century. The German word for set, \textit{Menge}, was coined by Bernard Bolzano in his work \textit{Paradoxes of the Infinite}.



% [Image: Passage with a translation of the original set definition of Georg Cantor. The German word Menge for set is translated with aggregate here.]



Figure: Passage with a translation of the original set definition of Georg Cantor. The German word \textit{Menge} for \textit{set} is translated with \textit{aggregate} here.



Georg Cantor, one of the founders of set theory, gave the following definition at the beginning of his \textit{Beiträge zur Begründung der transfiniten Mengenlehre}:



> A set is a gathering together into a whole of definite, distinct objects of our perception or our thought—which are called elements of the set.



Bertrand Russell called a set a \textit{class}:



> When mathematicians deal with what they call a manifold, aggregate, \textit{Menge}, \textit{ensemble}, or some equivalent name, it is common, especially where the number of terms involved is finite, to regard the object in question (which is in fact a class) as defined by the enumeration of its terms, and as consisting possibly of a single term, which in that case \textit{is} the class.



\paragraph{Naïve set theory}



The foremost property of a set is that it can have elements, also called \textit{members}. Two sets are equal when they have the same elements. More precisely, sets \textit{A} and \textit{B} are equal if every element of \textit{A} is an element of \textit{B}, and every element of \textit{B} is an element of \textit{A}; this property is called the \textit{extensionality of sets}.



The simple concept of a set has proved enormously useful in mathematics, but paradoxes arise if no restrictions are placed on how sets can be constructed:



\begin{itemize}
    \item Russell's paradox shows that the "set of all sets that \textit{do not contain themselves}", i.e., $\\{ x | x \text{ is a set and } x \not \in x\\}$, cannot exist.
    \item Cantor's paradox shows that "the set of all sets" cannot exist.
\end{itemize}

Naïve set theory defines a set as any \textit{well-defined} collection of distinct elements, but problems arise from the vagueness of the term \textit{well-defined}.



\paragraph{Axiomatic set theory}



In subsequent efforts to resolve these paradoxes since the time of the original formulation of naïve set theory, the properties of sets have been defined by axioms. Axiomatic set theory takes the concept of a set as a primitive notion. The purpose of the axioms is to provide a basic framework from which to deduce the truth or falsity of particular mathematical propositions (statements) about sets, using first-order logic. According to Gödel's incompleteness theorems however, it is not possible to use first-order logic to prove any such particular axiomatic set theory is free from paradox.



\paragraph{How sets are defined and set notation}



Mathematical texts commonly denote sets by capital letters in italic, such as $A$, $B$, $C$. A set may also be called a \textit{collection} or \textit{family}, especially when its elements are themselves sets.



\paragraph{Roster notation}



\textbf{Roster} or \textbf{enumeration notation} defines a set by listing its elements between curly brackets, separated by commas:



> $A = \\{4, 2, 1, 3\\}$

>

> $B = \\{ \text{blue, white, red} \\}$



In a set, all that matters is whether each element is in it or not, so the ordering of the elements in roster notation is irrelevant (in contrast, in a sequence, a tuple, or a permutation of a set, the ordering of the terms matters). For example, $\\{2, 4, 6\\}$ and $\\{4, 6, 4, 2\\}$ represent the same set.



For sets with many elements, especially those following an implicit pattern, the list of members can be abbreviated using an ellipsis '...'. For instance, the set of the first thousand positive integers may be specified in roster notation as



> $\\{1, 2, 3, ..., 1000\\}$



\paragraph{Infinite sets in roster notation}



An infinite set is a set with an endless list of elements. To describe an infinite set in roster notation, an ellipsis is placed at the end of the list, or at both ends, to indicate that the list continues forever. For example, the set of nonnegative integers is



> $\\{0, 1, 2, 3, 4, ...\\}$



and the set of all integers is



> $\\{..., -3, -2, -1, 0, 1, 2, 3, ...\\}$



\paragraph{Semantic definition}



Another way to define a set is to use a rule to determine what the elements are:



> Let $A$ be the set whose members are the first four positive integers.

> Let $B$ be the set of colors of the French flag.



Such a definition is called a \textit{semantic description}.



\paragraph{Set-builder notation}



Set-builder notation specifies a set as a selection from a larger set, determined by a condition on the elements. For example, a set $F$ can be defined as follows:



> $ F = \\{n \mid n \text{ is an integer, and } 0 \leq n \leq 19\\}$.



In this notation, the vertical bar "|" means "such that", and the description can be interpreted as "$F$ is the set of all numbers $n$ such that $n$ is an integer in the range from 0 to 19 inclusive". Some authors use a colon ":" instead of the vertical bar.



\paragraph{Classifying methods of definition}



Philosophy uses specific terms to classify types of definitions:



\begin{itemize}
    \item An \textit{intensional definition} uses a \textit{rule} to determine membership. Semantic definitions and definitions using set-builder notation are examples.
    \item An \textit{extensional definition} describes a set by \textit{listing all its elements}. Such definitions are also called \textit{enumerative}.
    \item An \textit{ostensive definition} is one that describes a set by giving \textit{examples} of elements; a roster involving an ellipsis would be an example.
\end{itemize}

\paragraph{Membership}



If $B$ is a set and $x$ is an element of $B$, this is written in shorthand as $x \in B$, which can also be read as "\textit{x} belongs to \textit{B}", or "\textit{x} is in \textit{B}". The statement "\textit{y} is not an element of \textit{B}" is written as $y \not \in B$, which can also be read as or "\textit{y} is not in \textit{B}".



For example, with respect to the sets $A = \\{1, 2, 3, 4\\}$, $B = \\{\text{blue, white, red}\\}$, and  $\newline$ $F = \\{ n | n \text{ is an integer, and } 0 \leq n \leq 19 \\}$,



> $4 \in A$ and $12 \in F$; and

>

> $20 \not \in F$ and green $\not \in B$



\paragraph{The empty set}



The \textit{empty set} (or \textit{null set}) is the unique set that has no members. It is denoted $\varnothing$ or $\emptyset$ or $\\{\hspace{0.5em}\\}$ or $\phi$.



\paragraph{Singleton sets}



A \textit{singleton set} is a set with exactly one element; such a set may also be called a \textit{unit set}. Any such set can be written as {\textit{x}}, where \textit{x} is the element. The set {\textit{x}} and the element \textit{x} mean different things; Halmos draws the analogy that a box containing a hat is not the same as the hat.



\paragraph{Subsets}



If every element of set \textit{A} is also in \textit{B}, then \textit{A} is described as being a \textit{subset of B}, or \textit{contained in B}, written \textit{A} $\subseteq$ \textit{B}, or \textit{B} $\supseteq$ \textit{A}. The latter notation may be read \textit{B contains A}, \textit{B includes A}, or \textit{B is a superset of A}. The relationship between sets established by $\subseteq$ is called \textit{inclusion} or \textit{containment}. Two sets are equal if they contain each other: \textit{A} $\subseteq$ \textit{B} and \textit{B} $\subseteq$ \textit{A} is equivalent to \textit{A} = \textit{B}.



If \textit{A} is a subset of \textit{B}, but \textit{A} is not equal to \textit{B}, then \textit{A} is called a \textit{proper subset} of \textit{B}. This can be written \textit{A} ? \textit{B}. Likewise, \textit{B} ? \textit{A} means \textit{B is a proper superset of A}, i.e. \textit{B} contains \textit{A}, and is not equal to \textit{A}.



A third pair of operators $\subset$ and $\supset$ are used differently by different authors: some authors use \textit{A} $\subset$ \textit{B} and \textit{B} $\supset$ \textit{A} to mean \textit{A} is any subset of \textit{B} (and not necessarily a proper subset), while others reserve \textit{A} $\subset$ \textit{B} and \textit{B} $\supset$ \textit{A} for cases where \textit{A} is a proper subset of \textit{B}.



Examples:



\begin{itemize}
    \item The set of all humans is a proper subset of the set of all mammals.
    \item {1, 3} $\subset$ {1, 2, 3, 4}.
    \item {1, 2, 3, 4} $\subseteq$ {1, 2, 3, 4}.
\end{itemize}

The empty set is a subset of every set, and every set is a subset of itself:



\begin{itemize}
    \item $\varnothing$ $\subseteq$ \textit{A}.
    \item \textit{A} $\subseteq$ \textit{A}.
\end{itemize}

\paragraph{Euler and Venn diagrams}



% [Image: A is a subset of B. B is a superset of A.]



Figure: \textit{A} is a subset of \textit{B}. \textit{B} is a superset of \textit{A}.



An Euler diagram is a graphical representation of a collection of sets; each set is depicted as a planar region enclosed by a loop, with its elements inside. If $A$ is a subset of $B$, then the region representing $A$ is completely inside the region representing $B$. If two sets have no elements in common, the regions do not overlap.



A Venn diagram, in contrast, is a graphical representation of $n$ sets in which the $n$ loops divide the plane into $2^n$ zones such that for each way of selecting some of the $n$ sets (possibly all or none), there is a zone for the elements that belong to all the selected sets and none of the others. For example, if the sets are $A$, $B$, and $C$, there should be a zone for the elements that are inside $A$ and $C$ and outside $B$ (even if such elements do not exist).



\paragraph{Special sets of numbers in mathematics}



% [Image: The natural numbers N are contained in the integers Z, which are contained in the rational numbers Q, which are contained in the real numbers R, which are contained in the complex numbers C]



Figure: $\mathbb{N} \subset \mathbb{Z} \subset \mathbb{Q} \subset \mathbb{R} \subset \mathbb{C}$. The natural numbers $\mathbb{N}$ are contained in the integers $\mathbb{Z}$, which are contained in the rational numbers $\mathbb{Q}$, which are contained in the real numbers $\mathbb{R}$, which are contained in the complex numbers $\mathbb{C}$



There are sets of such mathematical importance, to which mathematicians refer so frequently, that they have acquired special names and notational conventions to identify them.



Many of these important sets are represented in mathematical texts using bold (e.g. $\mathbb{Z}$) or blackboard bold (e.g. $\mathbb Z$) typeface. These include



\begin{itemize}
    \item $\mathbf{N}$ or $\mathbb {N}$, the set of all natural numbers: $\mathbb{N}=\{0,1,2,3,...\}$ (often, authors exclude $0$;
    \item $\mathbf{Z}$ or $\mathbb {Z}$, the set of all integers (whether positive, negative or zero): $\mathbb{Z}=\\{...,-2,-1,0,1,2,3,...\\}$;
    \item $\mathbf{Q}$ or $\mathbb {Q}$, the set of all rational numbers (that is, the set of all proper and improper fractions): $\newline$ $\mathbb{Q}=\left\\{\frac {a}{b}\mid a,b\in\mathbb{Z},b\ne0\right\\}$. For example, $-\tfrac{7}{4} \in \mathbb{Q}$ and $5 = \tfrac{5}{1} \in \mathbb{Q}$;
    \item $\mathbf{R}$ or $\mathbb {R}$, the set of all real numbers, including all rational numbers and all irrational numbers (which include algebraic numbers such as $\sqrt2$ that cannot be rewritten as fractions, as well as transcendental numbers such as $\pi$ and $e$;
    \item $\mathbf{C}$ or $\mathbb{C}$, the set of all complex numbers: $\mathbb{C} := \\{ a + bi \mid a,b \in \mathbb{R}\\}$, for example, $1 + 2i \in \mathbb{C}$.
\end{itemize}

Each of the above sets of numbers has an infinite number of elements. Each is a subset of the sets listed below it.



Sets of positive or negative numbers are sometimes denoted by superscript plus and minus signs, respectively. For example, $\mathbf{Q}^+$ represents the set of positive rational numbers.



\paragraph{Functions}



A \textit{function} (or \textit{mapping}) from a set $A$ to a set $B$ is a rule that assigns to each "input" element of $A$ an "output" that is an element of $B$; more formally, a function is a special kind of relation, one that relates each element of $A$ to \textit{exactly one} element of $B$. A function is called



\begin{itemize}
    \item injective (or one-to-one) if it maps any two different elements of $A$ to \textit{different} elements of $B$,
    \item surjective (or onto) if for every element of $B$, there is at least one element of $A$ that maps to it, and
    \item bijective (or a one-to-one correspondence) if the function is both injective and surjective — in this case, each element of $A$ is paired with a unique element of $B$, and each element of $B$ is paired with a unique element of $A$, so that there are no unpaired elements.
\end{itemize}

An injective function is called an \textit{injection}, a surjective function is called a \textit{surjection}, and a bijective function is called a \textit{bijection} or \textit{one-to-one correspondence}.



\paragraph{Cardinality}



The cardinality of a set $S$, denoted $\|S\|$, is the number of members of $S$. For example, if \textit{B} = {blue, white, red}, then $\|B\| = 3$. Repeated members in roster notation are not counted, so



\|{blue, white, red, blue, white}\| = 3, too.



More formally, two sets share the same cardinality if there exists a one-to-one correspondence between them.



The cardinality of the empty set is zero.



\paragraph{Infinite sets and infinite cardinality}



The list of elements of some sets is endless, or \textit{infinite}. For example, the set $\N$ of natural numbers is infinite. In fact, all the special sets of numbers mentioned in the section above, are infinite. Infinite sets have \textit{infinite cardinality}.



Some infinite cardinalities are greater than others. Arguably one of the most significant results from set theory is that the set of real numbers has greater cardinality than the set of natural numbers. Sets with cardinality less than or equal to that of $\N$ are called \textit{countable sets}; these are either finite sets or \textit{countably infinite sets} (sets of the same cardinality as $\N$); some authors use "countable" to mean "countably infinite". Sets with cardinality strictly greater than that of $\N$ are called \textit{uncountable sets}.



However, it can be shown that the cardinality of a straight line (i.e., the number of points on a line) is the same as the cardinality of any segment of that line, of the entire plane, and indeed of any finite-dimensional Euclidean space.



\paragraph{The Continuum Hypothesis}



The Continuum Hypothesis, formulated by Georg Cantor in 1878, is the statement that there is no set with cardinality strictly between the cardinality of the natural numbers and the cardinality of a straight line. In 1963, Paul Cohen proved that the Continuum Hypothesis is independent of the axiom system ZFC consisting of Zermelo–Fraenkel set theory with the axiom of choice. (ZFC is the most widely-studied version of axiomatic set theory.)



\paragraph{Power sets}



The power set of a set $S$ is the set of all subsets of $S$. The empty set and $S$ itself are elements of the power set of $S$, because these are both subsets of $S$. For example, the power set of $\\{1, 2, 3\\}$ is



$$\\{\varnothing, \\{1\\}, \\{2\\}, \\{3\\}, \\{1, 2\\}, \\{1, 3\\}, \\{2, 3\\}, \\{1, 2, 3\\}\\}$$



The power set of a set $S$ is commonly written as $P(S)$ or $2^S$.



If $S$ has $n$ elements, then $P(S)$ has $2^n$ elements. For example, $\\{1, 2, 3\\}$ has three elements, and its power set has $2^3 = 8$ elements, as shown above.



If $S$ is infinite (whether countable or uncountable), then $P(S)$ is uncountable. Moreover, the power set is always strictly "bigger" than the original set, in the sense that any attempt to pair up the elements of $S$ with the elements of $P(S)$ will leave some elements of $P(S)$ unpaired. (There is never a bijection from $S$ onto $P(S)$.)



\paragraph{Partitions}



A partition of a set \textit{S} is a set of nonempty subsets of \textit{S}, such that every element \textit{x} in \textit{S} is in exactly one of these subsets. That is, the subsets are pairwise disjoint (meaning any two sets of the partition contain no element in common), and the union of all the subsets of the partition is \textit{S}.



\paragraph{Basic operations}



There are several fundamental operations for constructing new sets from given sets.



\paragraph{Unions}



% [Image: ]



Figure: The \textbf{union} of $A$ and $B$, denoted $A \cup B$



Two sets can be joined: the \textit{union} of $A$ and $B$, denoted by $A \cup B$, is the set of all things that are members of $A$ or of $B$ or of both.



\textbf{Examples:}



\begin{itemize}
    \item $ \\{1, 2\\} \cup \\{1, 2\\} = \\{1, 2\\}.$
\end{itemize}

\begin{itemize}
    \item $ \\{1, 2\\} \cup \\{2, 3\\} = \\{1, 2, 3\\}.$
\end{itemize}

\begin{itemize}
    \item $ \\{1, 2, 3\\} \cup \\{3, 4, 5\\} = \\{1, 2, 3, 4, 5\\}.$
\end{itemize}

\textbf{Some basic properties of unions:}



\begin{itemize}
    \item $ A \cup B = B \cup A.$
\end{itemize}

\begin{itemize}
    \item $ A \cup (B \cup C) = (A \cup B) \cup C.$
\end{itemize}

\begin{itemize}
    \item $ A \subseteq (A \cup B).$
\end{itemize}

\begin{itemize}
    \item $ A \cup A = A.$
\end{itemize}

\begin{itemize}
    \item $ A \cup \varnothing = A.$
\end{itemize}

\begin{itemize}
    \item $ A \subseteq B$ if and only if $A \cup B = B$.
\end{itemize}

\paragraph{Intersections}



A new set can also be constructed by determining which members two sets have "in common". The \textit{intersection} of \textit{A} and \textit{B}, denoted by $A \cap B$, is the set of all things that are members of both \textit{A} and \textit{B}. If $A \cap B$ = $\varnothing$, then \textit{A} and \textit{B} are said to be \textit{disjoint}.



% [Image: The intersection of A and B]



Figure: The \textbf{intersection} of $A$ and $B$, denoted $A \cap B$.



\textbf{Examples:}



\begin{itemize}
    \item $\\{1, 2\\} \cap \\{1, 2\\} = \\{1, 2\\}.$
    \item $\\{1, 2\\} \cap \\{2, 3\\} = \\{2\\}.$
    \item $\\{1, 2\\} \cap \\{3, 4\\} = \varnothing.$
\end{itemize}

\textbf{Some basic properties of intersections:}



\begin{itemize}
    \item $A \cap B = B \cap A.$
    \item $A \cap (B \cap C) = (A \cap B) \cap C.$
    \item $A \cap B \subseteq A.$
    \item $A \cap A = A.$
    \item $A \cap \varnothing = \varnothing.$
    \item $A \subseteq B$ if and only if $A \cap B = A$.
\end{itemize}

\paragraph{Complements}



% [Image: The relative complement of B in A.]



Figure: The \textbf{relative complement} of \textit{B} in \textit{A}.



% [Image: The complement of A in U]



Figure: The \textbf{complement} of \textit{A} in \textit{U}



% [Image: The symmetric difference of A and B]



Figure: The \textbf{symmetric difference} of \textit{A} and \textit{B}



Two sets can also be "subtracted". The \textit{relative complement} of \textit{B} in \textit{A} (also called the \textit{set-theoretic difference} of \textit{A} and \textit{B}), denoted by $A \setminus B$ (or $A - B$), is the set of all elements that are members of \textit{A}, but not members of \textit{B}. It is valid to "subtract" members of a set that are not in the set, such as removing the element \textit{green} from the set {1, 2, 3}; doing so will not affect the elements in the set.



In certain settings, all sets under discussion are considered to be subsets of a given universal set \textit{U}. In such cases, $U \setminus A$ is called the \textit{absolute complement} or simply \textit{complement} of \textit{A}, and is denoted by $A'$ or $A^C$.



\begin{itemize}
    \item $A' = U \setminus A$
\end{itemize}

\textbf{Examples:}



\begin{itemize}
    \item $\\{1, 2\\} \setminus \\{1, 2\\} = \varnothing.$
    \item $\\{1, 2, 3, 4\\} \setminus \\{1, 3\\} = \\{2, 4\\}.$
    \item If $U$ is the set of integers, $E$ is the set of even integers, and $O$ is the set of odd integers, then $U \setminus E = E' = O.$
\end{itemize}

\textbf{Some basic properties of complements include the following:}



\begin{itemize}
    \item $A \setminus B \neq B \setminus A$ for $A \neq B.$
    \item $A \cup A' = U.$
    \item $A \cap A' = \varnothing.$
    \item $(A')' = A.$
    \item $\varnothing \setminus A = \varnothing.$
    \item $A \setminus \varnothing = A.$
    \item $A \setminus A = \varnothing.$
    \item $A \setminus U = \varnothing.$
    \item $A \setminus A' = A$ and $A' \setminus A = A'.$
    \item $U' = \varnothing$ and $\varnothing' = U.$
    \item $A \setminus B = (A \cap B)'.$
    \item If $A \subseteq B$ then $A \setminus B = \varnothing.$
\end{itemize}

An extension of the complement is the symmetric difference, defined for sets \textit{A}, \textit{B} as



$$A\,\Delta\,B = (A \setminus B) \cup (B \setminus A).$$ For example, the symmetric difference of {7, 8, 9, 10} and {9, 10, 11, 12} is the set {7, 8, 11, 12}. The power set of any set becomes a Boolean ring with symmetric difference as the addition of the ring (with the empty set as neutral element) and intersection as the multiplication of the ring.



\paragraph{Cartesian product}



A new set can be constructed by associating every element of one set with every element of another set. The \textit{Cartesian product} of two sets \textit{A} and \textit{B}, denoted by \textit{A} × \textit{B,} is the set of all ordered pairs (\textit{a}, \textit{b}) such that \textit{a} is a member of \textit{A} and \textit{b} is a member of \textit{B}.



\textbf{Examples:}



\begin{itemize}
    \item $\\{1, 2\\} \times \\{\text{red}, \text{white}, \text{green}\\} = \\{(1, \text{red}), (1, \text{white}), (1, \text{green}), (2, \text{red}), (2, \text{white}), (2, \text{green})\\}.$
    \item $\\{1, 2\\} \times \\{1, 2\\} = \\{(1, 1), (1, 2), (2, 1), (2, 2)\\}.$
    \item $\\{a, b, c\\} \times \\{d, e, f\\} = \\{(a, d), (a, e), (a, f), (b, d), (b, e), (b, f), (c, d), (c, e), (c, f)\\}.$
\end{itemize}

\textbf{Some basic properties of Cartesian products:}



\begin{itemize}
    \item $A \times \varnothing = \varnothing.$
    \item $A \times (B \cup C) = (A \times B) \cup (A \times C).$
    \item $(A \cup B) \times C = (A \times C) \cup (B \times C).$
\end{itemize}

Let \textit{A} and \textit{B} be finite sets; then the cardinality of the Cartesian product is the product of the cardinalities:



\begin{itemize}
    \item \|\textit{A} × \textit{B}\| = \|\textit{B} × \textit{A}\| = \|\textit{A}\| × \|\textit{B}\|.
\end{itemize}

\paragraph{Applications}



Sets are ubiquitous in modern mathematics. For example, structures in abstract algebra, such as groups, fields and rings, are sets closed under one or more operations.



One of the main applications of naive set theory is in the construction of relations. A relation from a domain $A$ to a codomain $B$ is a subset of the Cartesian product $A \times B$. For example, considering the set $S = \\{\text{rock}, \text{paper}, \text{scissors}\\}$ of shapes in the game of the same name, the relation "beats" from $S$ to $S$ is the set $\newline$ $B = \\{(\text{scissors}, \text{paper}), (\text{paper}, \text{rock}), (\text{rock}, \text{scissors})\\}$; thus $x$ beats $y$ in the game if the pair $(x, y)$ is a member of $B$.



Another example is the set $F$ of all pairs $(x, x^2)$, where $x$ is real. This relation is a subset of $\mathbb{R} \times \mathbb{R}$, because the set of all squares is a subset of the set of all real numbers. Since for every $x$ in $\mathbb{R}$, one and only one pair $(x, x^2)$ is found in $F$, it is called a function. In functional notation, this relation can be written as $F(x) = x^2$.



\paragraph{Principle of inclusion and exclusion}



% [Image: The inclusion-exclusion principle is used to calculate the size of the union of sets: the size of the union is the size of the two sets, minus the size of their intersection.]



Figure: The inclusion-exclusion principle is used to calculate the size of the union of sets: the size of the union is the size of the two sets, minus the size of their intersection.



The inclusion–exclusion principle is a counting technique that can be used to count the number of elements in a union of two sets—if the size of each set and the size of their intersection are known. It can be expressed symbolically as



$$|A \cup B| = |A| + |B| - |A \cap B|.$$



A more general form of the principle can be used to find the cardinality of any finite union of sets:



$$ \begin{aligned} \left|A _{1}\cup A _{2}\cup A _{3}\cup\ldots\cup A _{n}\right| =& \left(\left|A _{1}\right|+\left|A _{2}\right|+\left|A _{3}\right|+\ldots\left|A _{n}\right|\right) \\\\ &{} - \left(\left|A _{1}\cap A _{2}\right|+\left|A _{1}\cap A _{3}\right|+\ldots\left|A _{n-1}\cap A _{n}\right|\right) \\\\ &{} + \ldots \\\\ &{} + \left(-1\right)^{n-1}\left(\left|A _{1}\cap A _{2}\cap A _{3}\cap\ldots\cap A _{n}\right|\right) \end{aligned} $$



\paragraph{De Morgan's laws}



Augustus De Morgan stated two laws about sets.



If $A$ and $B$ are any two sets then,





\begin{itemize}
    \item $ (A \cup B)' = A' \cap B'$
\end{itemize}

The complement of $A$ union $B$ equals the complement of $A$ intersected with the complement of $B$.



\begin{itemize}
    \item $ (A \cap B)' = A' \cup B'$
\end{itemize}

The complement of $A$ intersected with $B$ is equal to the complement of $A$ union to the complement of $B$.



\paragraph{Resources}



\begin{itemize}
    \item \href{https://link-springer-com.libweb.lib.utsa.edu/content/pdf/10.1007\%2F978-1-4419-7127-2.pdf}{Course Textbook}, pages 93-100
\end{itemize}

\paragraph{Licensing}



Content obtained and/or adapted from:



\begin{itemize}
    \item \href{https://en.wikipedia.org/wiki/Set_(mathematics}{Set (mathematics), Wikipedia}) under a CC BY-SA license
\end{itemize}

\subsection{Learning Objectives}

\begin{enumerate}
    \item Here are three Student Learning Objectives (SLOs) based on the concepts of set theory and Euler diagrams:
    \item \textbf{Identify and Define Sets:}
    \item \textbf{Outcome:} Students will be able to identify and define a set by listing its elements using roster notation or by describing it using set-builder notation.
    \item \textbf{Details:} This includes distinguishing between finite and infinite sets, as well as understanding and applying the concept of well-defined collections of elements in various contexts.
    \item \textbf{Analyze Subset Relationships:}
    \item \textbf{Outcome:} Students will be able to analyze and determine subset relationships within sets, including proper subsets and the empty set, using Euler diagrams as a visual aid.
    \item \textbf{Details:} This includes understanding and using the subset notation (e.g., $A \subseteq B$) and recognizing when one set is a subset or a proper subset of another, as well as identifying when sets are equal.
    \item \textbf{Apply Set Operations:}
    \item \textbf{Outcome:} Students will be able to apply basic set operations, such as union, intersection, and complement, to solve problems involving sets and visualize these operations using Euler diagrams.
    \item \textbf{Details:} This includes the ability to perform operations like $A \cup B$, $A \cap B$, and $A \setminus B$, understanding the properties of these operations, and using diagrams to represent the results of these operations clearly.
\end{enumerate}

\subsection{Assessment Instrument}

\begin{enumerate}
    \item \textbf{Assessment Instrument: Sets - Definitions Quiz} \textbf{Format:}
    \item \textbf{Type:} Mixed-format quiz (multiple-choice, short answer, and problem-solving questions)
    \item \textbf{Duration:} 45 minutes
    \item \textbf{Total Marks:} 30 points \textbf{Sections:}
    \item \textbf{Identifying and Defining Sets (10 points):}
    \item \textbf{Question 1:} Define the following sets using roster notation or set-builder notation:
    \item a. The set of all even numbers less than 20.
    \item b. The set of natural numbers between 1 and 10, inclusive.
    \item \textbf{Question 2:} Identify whether the following are well-defined sets. Provide a brief explanation:
    \item a. The set of all tall students in the class.
    \item b. The set of all vowels in the English alphabet.
    \item \textbf{Analyzing Subset Relationships (10 points):}
    \item \textbf{Question 3:} Consider the following sets:
    \item $A = \\{2, 4, 6, 8, 10\\}$
    \item $B = \\{4, 8, 12\\}$
    \item \textbf{Question 4:} Determine whether $A \subseteq B$, $B \subseteq A$, or neither. Justify your answer.
    \item \textbf{Question 5:} Draw an Euler diagram to represent the relationship between the following sets:
    \item $C = \\{a, b, c\\}$
    \item $D = \\{a, b, c, d\\}$
    \item $E = \\{d, e\\}$
    \item \textbf{Applying Set Operations (10 points):}
    \item \textbf{Question 6:} Given the sets $X = \\{1, 2, 3, 4\\}$ and $Y = \\{3, 4, 5, 6\\}$, find:
    \item a. $X \cup Y$
    \item b. $X \cap Y$
    \item c. $X \setminus Y$
    \item \textbf{Question 7:} If $Z = \\{3, 5\\}$, is $Z$ a subset of $Y$? Explain your reasoning. \textbf{Scoring Rubric:}
    \item \textbf{Identifying and Defining Sets:} 2.5 points for each correct definition and explanation (total 10 points).
    \item \textbf{Analyzing Subset Relationships:} 2.5 points for each correct identification and justification (total 10 points).
    \item \textbf{Applying Set Operations:} 2.5 points for each correctly solved operation (total 10 points). The quiz will help assess the students' understanding of the fundamental concepts of sets, including their ability to define, identify, and manipulate sets using operations and subset relationships.
\end{enumerate}

%%===============================================================
\section{Sets - Operations}
\label{sec:sets-operations}
%%===============================================================

In mathematics, \textbf{the algebra of sets}, not to be confused with the mathematical structure of \textit{an} algebra of sets, defines the properties and laws of sets, the set-theoretic operations of union, intersection, and complementation and the relations of set equality and set inclusion. It also provides systematic procedures for evaluating expressions, and performing calculations, involving these operations and relations.



Any set of sets closed under the set-theoretic operations forms a Boolean algebra with the join operator being \textit{union}, the meet operator being \textit{intersection}, the complement operator being \textit{set complement}, the bottom being $\varnothing$ and the top being the universe set under consideration.



\paragraph{Fundamentals}



The algebra of sets is the set-theoretic analogue of the algebra of numbers. Just as arithmetic addition and multiplication are associative and commutative, so are set union and intersection; just as the arithmetic relation "less than or equal" is reflexive, antisymmetric and transitive, so is the set relation of "subset".



It is the algebra of the set-theoretic operations of union, intersection and complementation, and the relations of equality and inclusion. For a basic introduction to sets see the article on sets, for a fuller account see naive set theory, and for a full rigorous axiomatic treatment see axiomatic set theory.



\paragraph{The fundamental properties of set algebra}



The binary operations of set union ($\cup$) and intersection ($\cap$) satisfy many identities. Several of these identities or "laws" have well established names.



>    Commutative property:

>    - $A \cup B = B \cup A$

>    - $A \cap B = B \cap A$



> Associative property:

>    - $(A \cup B) \cup C = A \cup (B \cup C)$

>    - $(A \cap B) \cap C = A \cap (B \cap C)$



> Distributive property:

>    - $A \cup (B \cap C) = (A \cup B) \cap (A \cup C)$

>    - $A \cap (B \cup C) = (A \cap B) \cup (A \cap C)$



The union and intersection of sets may be seen as analogous to the addition and multiplication of numbers. Like addition and multiplication, the operations of union and intersection are commutative and associative, and intersection \textit{distributes} over union. However, unlike addition and multiplication, union also distributes over intersection.



Two additional pairs of properties involve the special sets called the empty set Ø and the universe set $U$; together with the complement operator ($A^C$ denotes the complement of $A$. This can also be written as $A'$, read as A prime). The empty set has no members, and the universe set has all possible members (in a particular context).



> Identity :

>    - $A \cup \varnothing = A$

>    - $A \cap U = A$



> Complement :

>    - $A \cup A^C = U$

>    - $A \cap A^C = \varnothing$



The identity expressions (together with the commutative expressions) say that, just like 0 and 1 for addition and multiplication, $\varnothing$ and \textbf{U} are the identity elements for union and intersection, respectively.



Unlike addition and multiplication, union and intersection do not have inverse elements. However the complement laws give the fundamental properties of the somewhat inverse-like unary operation of set complementation.



The preceding five pairs of formulae—the commutative, associative, distributive, identity and complement formulae—encompass all of set algebra, in the sense that every valid proposition in the algebra of sets can be derived from them.



Note that if the complement formulae are weakened to the rule $(A^C)^C = A$, then this is exactly the algebra of propositional linear logic.



\paragraph{The principle of duality}



Each of the identities stated above is one of a pair of identities such that each can be transformed into the other by interchanging $\cup$ and $\cap$, and also $\varnothing$ and \textbf{U}.



These are examples of an extremely important and powerful property of set algebra, namely, the \textbf{principle of duality} for sets, which asserts that for any true statement about sets, the \textbf{dual} statement obtained by interchanging unions and intersections, interchanging \textbf{U} and Ø and reversing inclusions is also true. A statement is said to be \textbf{self-dual} if it is equal to its own dual.



\paragraph{Some additional laws for unions and intersections}



The following proposition states six more important laws of set algebra, involving unions and intersections.



\textbf{PROPOSITION 3}: For any subsets \textit{A} and \textit{B} of a universe set \textbf{U}, the following identities hold:



> idempotent laws:

>    - $A \cup A = A$

>    - $A \cap A = A$



> domination laws:

>    - $A \cup U = U$

>    - $A \cap \varnothing = \varnothing$



> absorption laws:

>    - $A \cup (A \cap B) = A$

>    - $A \cap (A \cup B) = A$



As noted above, each of the laws stated in proposition 3 can be derived from the five fundamental pairs of laws stated above . As an illustration, a proof is given below for the idempotent law for union.



\textit{Proof:}



$$ \begin{aligned} A \cup A & = (A \cup A) \cap U && \text{by the identity law of intersection} \\\\ & = (A \cup A) \cap (A \cup A^C) && \text{by the complement law for union} \\\\ & = A \cup (A \cap A^C) && \text{by the distributive law of union over intersection} \\\\ & = A \cup \varnothing && \text{by the complement law for intersection} \\\\ & = A && \text{by the identity law for union} \end{aligned} $$



The following proof illustrates that the dual of the above proof is the proof of the dual of the idempotent law for union, namely the idempotent law for intersection.



\textit{Proof:}



$$ \begin{aligned} A \cap A &= (A \cap A) \cup \varnothing && \text{by the identity law for union} \\\\ &= (A \cap A) \cup (A \cap A^C) && \text{by the complement law for intersection} \\\\ &= A \cap (A \cup A^C) && \text{by the distributive law of intersection over union} \\\\ &= A \cap U && \text{by the complement law for union} \\\\ &= A && \text{by the identity law for intersection} \end{aligned} $$



Intersection can be expressed in terms of set difference :



$A \cap B = A \setminus (A \setminus B)$



\paragraph{Some additional laws for complements}



The following proposition states five more important laws of set algebra, involving complements.



\textbf{PROPOSITION 4}: Let \textit{A} and \textit{B} be subsets of a universe \textbf{U}, then:



> De Morgan's laws:

>    - $(A \cup B)^C = A^C \cap B^C$

>    - $(A \cap B)^C = A^C \cup B^C$



> double complement or involution law:

>    - ${(A^{C})}^{C} = A$



> complement laws for the universe set and the empty set:

>    - $\varnothing^C = U$

>    - $U^C = \varnothing$



Notice that the double complement law is self-dual.



The next proposition, which is also self-dual, says that the complement of a set is the only set that satisfies the complement laws. In other words, complementation is characterized by the complement laws.



\textbf{PROPOSITION 5}: Let \textit{A} and \textit{B} be subsets of a universe \textbf{U}, then:



> uniqueness of complements:

>    - If $A \cup B = U$, and $A \cap B = \varnothing$, then $B = A^C$



\paragraph{The algebra of inclusion}



The following proposition says that inclusion, that is the binary relation of one set being a subset of another, is a partial order.



\textbf{PROPOSITION 6}: If \textit{A}, \textit{B} and \textit{C} are sets then the following hold:



> reflexivity:

>    - $A \subseteq A$



> antisymmetry:

>    - $A \subseteq B$ and $B \subseteq A$ if and only if $A = B$



> transitivity:

>    - If $A \subseteq B$ and $B \subseteq C$, then $A \subseteq C$



The following proposition says that for any set \textit{S}, the power set of \textit{S}, ordered by inclusion, is a bounded lattice, and hence together with the distributive and complement laws above, show that it is a Boolean algebra.



\textbf{PROPOSITION 7}: If \textit{A}, \textit{B} and \textit{C} are subsets of a set \textit{S} then the following hold:



> existence of a least element and a greatest element:

>    - $\varnothing \subseteq A \subseteq S$



> existence of joins:

>    - $A \subseteq A \cup B$

>    - If $A \subseteq C$ and $B \subseteq C$, then $A \cup B \subseteq C$



> existence of meets:

>    - $A \cap B \subseteq A$

>    - If $C \subseteq A$ and $C \subseteq B$, then $C \subseteq A \cap B$



The following proposition says that the statement $A \subseteq B$ is equivalent to various other statements involving unions, intersections and complements.



\textbf{PROPOSITION 8}: For any two sets \textit{A} and \textit{B}, the following are equivalent:



\begin{itemize}
    \item $A \subseteq B$
\end{itemize}

\begin{itemize}
    \item $A \cap B = A$
\end{itemize}

\begin{itemize}
    \item $A \cup B = B$
\end{itemize}

\begin{itemize}
    \item $A \setminus B = \varnothing$
\end{itemize}

\begin{itemize}
    \item $B^C \subseteq A^C$
\end{itemize}

The above proposition shows that the relation of set inclusion can be characterized by either of the operations of set union or set intersection, which means that the notion of set inclusion is axiomatically superfluous.



\paragraph{The algebra of relative complements}



The following proposition lists several identities concerning relative complements and set-theoretic differences.



\textbf{PROPOSITION 9}: For any universe \textbf{U} and subsets \textit{A}, \textit{B}, and \textit{C} of \textbf{U}, the following identities hold:



\begin{itemize}
    \item $C \setminus (A \cap B) = (C \setminus A) \cup (C \setminus B)$
\end{itemize}

\begin{itemize}
    \item $C \setminus (A \cup B) = (C \setminus A) \cap (C \setminus B)$
\end{itemize}

\begin{itemize}
    \item $C \setminus (B \setminus A) = (A \cap C)\cup(C \setminus B)$
\end{itemize}

\begin{itemize}
    \item $(B \setminus A) \cap C = (B \cap C) \setminus A = B \cap (C \setminus A)$
\end{itemize}

\begin{itemize}
    \item $(B \setminus A) \cup C = (B \cup C) \setminus (A \setminus C)$
\end{itemize}

\begin{itemize}
    \item $(B \setminus A) \setminus C = B \setminus (A \cup C)$
\end{itemize}

\begin{itemize}
    \item $A \setminus A = \varnothing$
\end{itemize}

\begin{itemize}
    \item $\varnothing \setminus A = \varnothing$
\end{itemize}

\begin{itemize}
    \item $A \setminus \varnothing = A$
\end{itemize}

\begin{itemize}
    \item $B \setminus A = A^C \cap B$
\end{itemize}

\begin{itemize}
    \item $(B \setminus A)^C = A \cup B^C$
\end{itemize}

\begin{itemize}
    \item $U \setminus A = A^C$
\end{itemize}

\begin{itemize}
    \item $A \setminus U = \varnothing$
\end{itemize}

\paragraph{Resources}



\begin{itemize}
    \item \href{https://link-springer-com.libweb.lib.utsa.edu/content/pdf/10.1007\%2F978-1-4419-7127-2.pdf}{Course Textbook}, pages 101-115
\end{itemize}

\paragraph{Licensing}



Content obtained and/or adapted from:



\begin{itemize}
    \item \href{https://en.wikipedia.org/wiki/Algebra_of_sets}{Algebra of sets, Wikipedia} under a CC BY-SA license
\end{itemize}

\subsection{Learning Objectives}

\begin{enumerate}
    \item \textbf{Understanding Set Operations:}
    \item \textbf{Outcome:} Students will be able to perform basic set operations including union, intersection, and complement.
    \item \textbf{Details:} This involves recognizing and applying the commutative, associative, and distributive properties in the context of set operations. Students will practice evaluating expressions using these operations, ensuring a solid understanding of how different sets combine and relate to each other.
    \item \textbf{Applying Set Identities:}
    \item \textbf{Outcome:} Students will be able to apply set identities, such as De Morgan's laws and the idempotent laws, to simplify complex set expressions.
    \item \textbf{Details:} Students will learn to use specific set identities to simplify and evaluate expressions. This includes recognizing the duality principle and how it can be used to transform set expressions while maintaining their validity.
    \item \textbf{Exploring Set Inclusion and Order:}
    \item \textbf{Outcome:} Students will be able to analyze and understand the relationship between sets using the concepts of set inclusion and partial order.
    \item \textbf{Details:} Students will explore the properties of reflexivity, antisymmetry, and transitivity in set inclusion. They will learn to determine whether one set is a subset of another and understand how this relationship forms a partial order among sets, linking it to the structure of a Boolean algebra.
\end{enumerate}

\subsection{Assessment Instrument}

\begin{enumerate}
    \item \textbf{Assessment Instrument: Sets - Operations Quiz} \textbf{Format:}
    \item \textbf{Type:} Mixed-format quiz (multiple-choice, short answer, and problem-solving questions)
    \item \textbf{Duration:} 45 minutes
    \item \textbf{Total Marks:} 30 points \textbf{Sections:}
    \item \textbf{Understanding Set Operations (10 points):}
    \item \textbf{Question 1:} Given the sets $ A = \\{1, 2, 3, 4\\} $, $ B = \\{3, 4, 5, 6\\} $, and $ C = \\{5, 6, 7, 8\\} $:
    \item a. Find $ A \cup B $ (2 points)
    \item b. Find $ A \cap C $ (2 points)
    \item c. Find $ B \cup C $ (2 points)
    \item d. Find $ A \cap (B \cup C) $ (4 points)
    \item \textbf{Applying Set Identities (10 points):}
    \item \textbf{Question 2:} Simplify the following expressions using set identities:
    \item a. $ (A \cup B) \cap A^C $ where $ A = \\{x \mid x \text{ is even}\\} $ and $ B = \\{x \mid x > 5\\} $ (4 points)
    \item b. $ A \cup (A \cap B) $ where $ A = \\{1, 2, 3\\} $ and $ B = \\{3, 4, 5\\} $ (3 points)
    \item c. $ (A \cap B)^C $ where $ A = \\{x \mid x \text{ is prime}\\} $ and $ B = \\{x \mid x \text{ is odd}\\} $ (3 points)
    \item \textbf{Exploring Set Inclusion and Order (10 points):}
    \item \textbf{Question 3:} Answer the following questions related to set inclusion and order:
    \item a. Determine whether $ A = \\{2, 4, 6\\} $ is a subset of $ B = \\{2, 4, 6, 8\\} $. Explain your reasoning. (3 points)
    \item b. Prove that $ A \subseteq B $ and $ B \subseteq C $ imply $ A \subseteq C $ using an example where $ A = \\{1, 2\\} $, $ B = \\{1, 2, 3\\} $, and $ C = \\{1, 2, 3, 4\\} $. (4 points)
    \item c. If $ A \subseteq B $, what can you infer about $ A \cup B $ and $ A \cap B $? Provide a brief explanation. (3 points) \textbf{Scoring Rubric:}
    \item \textbf{Understanding Set Operations:} 2 points for each correct answer in sub-questions a, b, c; 4 points for sub-question d (total 10 points).
    \item \textbf{Applying Set Identities:} 4 points for correctly simplifying the first expression, 3 points each for the next two expressions (total 10 points).
    \item \textbf{Exploring Set Inclusion and Order:} 3 points for a correct explanation in sub-question a, 4 points for a correct proof in sub-question b, and 3 points for correct inference in sub-question c (total 10 points). The quiz will help assess the students' understanding of set operations, application of set identities, and their ability to analyze set inclusion and order. It ensures that they can perform operations accurately and reason through relationships between sets.
\end{enumerate}

%%===============================================================
\section{Venn Diagrams}
\label{sec:venn-diagrams}
%%===============================================================

% [Image: Venn diagram showing the uppercase glyphs shared by the Greek, Latin, and Russian alphabets]



Figure: Venn diagram showing the uppercase glyphs shared by the Greek, Latin, and Russian alphabets.



A \textbf{Venn diagram} is a widely used diagram style that shows the logical relation between sets, popularized by John Venn in the 1880s. The diagrams are used to teach elementary set theory, and to illustrate simple set relationships in probability, logic, statistics, linguistics and computer science. A Venn diagram uses simple closed curves drawn on a plane to represent sets. Very often, these curves are circles or ellipses.



Similar ideas had been proposed before Venn. Christian Weise in 1712 (\textit{Nucleus Logicoe Wiesianoe}) and Leonhard Euler (\textit{Letters to a German Princess}) in 1768, for instance, came up with similar ideas. The idea was popularised by Venn in \textit{Symbolic Logic}, Chapter V "Diagrammatic Representation", 1881.



\paragraph{Details}



A Venn diagram may also be called a \textbf{primary diagram}, \textbf{set diagram} or \textbf{logic diagram}. It is a diagram that shows \textit{all} possible logical relations between a finite collection of different sets. These diagrams depict elements as points in the plane, and sets as regions inside closed curves. A Venn diagram consists of multiple overlapping closed curves, usually circles, each representing a set. The points inside a curve labelled \textit{S} represent elements of the set \textit{S}, while points outside the boundary represent elements not in the set \textit{S}. This lends itself to intuitive visualizations; for example, the set of all elements that are members of both sets \textit{S} and \textit{T}, denoted $S \cap T$ and read "the intersection of \textit{S} and \textit{T}", is represented visually by the area of overlap of the regions \textit{S} and \textit{T}.



In Venn diagrams, the curves are overlapped in every possible way, showing all possible relations between the sets. They are thus a special case of Euler diagrams, which do not necessarily show all relations. Venn diagrams were conceived around 1880 by John Venn. They are used to teach elementary set theory, as well as illustrate simple set relationships in probability, logic, statistics, linguistics, and computer science.



A Venn diagram in which the area of each shape is proportional to the number of elements it contains is called an \textbf{area-proportional} (or \textbf{scaled}) \textbf{Venn diagram}.



Venn diagrams are used heavily in the logic of class branch of reasoning.



\paragraph{Example}



% [Image: Sets A (creatures with two legs) and B (creatures that fly)]



Figure: Sets A (creatures with two legs) and B (creatures that fly)



This example involves two sets, A and B, represented here as coloured circles. The orange circle, set A, represents all types of living creatures that are two-legged. The blue circle, set B, represents the living creatures that can fly. Each separate type of creature can be imagined as a point somewhere in the diagram. Living creatures that can fly and have two legs—for example, parrots—are then in both sets, so they correspond to points in the region where the blue and orange circles overlap. This overlapping region would only contain those elements (in this example, creatures) that are members of both set A (two-legged creatures) and set B (flying creatures).



Humans and penguins are bipedal, and so are in the orange circle, but since they cannot fly, they appear in the left part of the orange circle, where it does not overlap with the blue circle. Mosquitoes can fly, but have six, not two, legs, so the point for mosquitoes is in the part of the blue circle that does not overlap with the orange one. Creatures that are not two-legged and cannot fly (for example, whales and spiders) would all be represented by points outside both circles.



The combined region of sets A and B is called the \textit{union} of A and B, denoted by $A \cup B$. The union in this case contains all living creatures that either are two-legged or can fly (or both).



The region included in both A and B, where the two sets overlap, is called the \textit{intersection} of A and B, denoted by $A \cap B$. In this example, the intersection of the two sets is not empty, because there \textit{are} points that represent creatures that are in \textit{both} the orange and blue circles.



\paragraph{Overview}



% [Image: ]



Figure: Intersection of two sets $A \cap B$



% [Image: ]



Figure: Union of two sets $A \cup B$



% [Image: ]



Figure: Symmetric difference of two sets $A \triangle B$



% [Image: ]



Figure: Relative complement of $A$ (left) in $B$ (right), meaning $A^c \cap B = B \setminus A$



% [Image: ]



Figure: Absolute complement of $A$ in $U$, meaning $A^c = U \setminus A$



A Venn diagram is constructed with a collection of simple closed curves drawn in a plane. According to Lewis, the "principle of these diagrams is that classes \[or \textit{sets}\] be represented by regions in such relation to one another that all the possible logical relations of these classes can be indicated in the same diagram. That is, the diagram initially leaves room for any possible relation of the classes, and the actual or given relation, can then be specified by indicating that some particular region is null or is not-null".



Venn diagrams normally comprise overlapping circles. The interior of the circle symbolically represents the elements of the set, while the exterior represents elements that are not members of the set. For instance, in a two-set Venn diagram, one circle may represent the group of all wooden objects, while the other circle may represent the set of all tables. The overlapping region, or \textit{intersection}, would then represent the set of all wooden tables. Shapes other than circles can be employed as shown below by Venn's own higher set diagrams. Venn diagrams do not generally contain information on the relative or absolute sizes (cardinality) of sets. That is, they are schematic diagrams generally not drawn to scale.



Venn diagrams are similar to Euler diagrams. However, a Venn diagram for \textit{n} component sets must contain all $2^n$ hypothetically possible zones, that correspond to some combination of inclusion or exclusion in each of the component sets. Euler diagrams contain only the actually possible zones in a given context. In Venn diagrams, a shaded zone may represent an empty zone, whereas in an Euler diagram, the corresponding zone is missing from the diagram. For example, if one set represents \textit{dairy products} and another \textit{cheeses}, the Venn diagram contains a zone for cheeses that are not dairy products. Assuming that in the context \textit{cheese} means some type of dairy product, the Euler diagram has the cheese zone entirely contained within the dairy-product zone—there is no zone for (non-existent) non-dairy cheese. This means that as the number of contours increases, Euler diagrams are typically less visually complex than the equivalent Venn diagram, particularly if the number of non-empty intersections is small.



The difference between Euler and Venn diagrams can be seen in the following example. Take the three sets:



\begin{itemize}
    \item $A = \\{1,\, 2,\, 5\\}$
    \item $B = \\{1,\, 6\\}$
    \item $C = \\{4,\, 7\\}$
\end{itemize}

The Euler and the Venn diagram of those sets are:



% [Image: Euler diagram]



Figure: Euler diagram



% [Image: Venn diagram]



Figure: Venn diagram



\paragraph{Extensions to higher numbers of sets}



Venn diagrams typically represent two or three sets, but there are forms that allow for higher numbers. Shown below, four intersecting spheres form the highest order Venn diagram that has the symmetry of a simplex and can be visually represented. The 16 intersections correspond to the vertices of a tesseract (or the cells of a 16-cell, respectively).



% [Image: ]



% [Image: ]



% [Image: ]

% [Image: ]

% [Image: ]



For higher numbers of sets, some loss of symmetry in the diagrams is unavoidable. Venn was keen to find "symmetrical figures\...elegant in themselves," that represented higher numbers of sets, and he devised an \textit{elegant} four-set diagram using ellipses (see below). He also gave a construction for Venn diagrams for \textit{any} number of sets, where each successive curve that delimits a set interleaves with previous curves, starting with the three-circle diagram.



% [Image: Venn's construction for four sets]



Figure: Venn's construction for four sets



% [Image: ]



Figure: Venn's construction for five sets



% [Image: ]



Figure: Venn's construction for six sets



% [Image: ]



Figure: Venn's four-set diagram using ellipses



% [Image: ]



Figure: \textbf{Non-example:} This Euler diagram is not a Venn diagram for four sets as it has only 13 regions (excluding the outside); there is no region where only the yellow and blue, or only the red and green circles meet.



% [Image: ]



Figure: Five-set Venn diagram using congruent ellipses in a five-fold rotationally symmetrical arrangement devised by Branko Grünbaum. Labels have been simplified for greater readability; for example, A denotes A n Bc n Cc n Dc n Ec, while BCE denotes Ac n B n C n Dc n E.



% [Image: ]



Figure: Six-set Venn diagram made of only triangles (interactive version)



\paragraph{Edwards–Venn diagrams}



% [Image: ]



Figure: Three sets



% [Image: ]



Figure: Four sets



% [Image: ]



Figure: Five sets



% [Image: ]



Figure: Six sets



Anthony William Fairbank Edwards constructed a series of Venn diagrams for higher numbers of sets by segmenting the surface of a sphere, which became known as Edwards–Venn diagrams. For example, three sets can be easily represented by taking three hemispheres of the sphere at right angles (\textit{x} = 0, \textit{y} = 0 and \textit{z} = 0). A fourth set can be added to the representation, by taking a curve similar to the seam on a tennis ball, which winds up and down around the equator, and so on. The resulting sets can then be projected back to a plane, to give \textit{cogwheel} diagrams with increasing numbers of teeth—as shown here. These diagrams were devised while designing a stained-glass window in memory of Venn.



\paragraph{Other diagrams}



Edwards–Venn diagrams are topologically equivalent to diagrams devised by Branko Grünbaum, which were based around intersecting polygons with increasing numbers of sides. They are also two-dimensional representations of hypercubes.



Henry John Stephen Smith devised similar n-set diagrams using sine curves\[20\] with the series of equations



$$y_i = \frac{\sin\left(2^i x\right)}{2^i} \text{ where } 0 \leq i \leq n-2 \text{ and } i \in \mathbb{N}$$



Charles Lutwidge Dodgson (also known as Lewis Carroll) devised a five-set diagram known as Carroll's square. Joaquin and Boyles, on the other hand, proposed supplemental rules for the standard Venn diagram, in order to account for certain problem cases. For instance, regarding the issue of representing singular statements, they suggest to consider the Venn diagram circle as a representation of a set of things, and use first-order logic and set theory to treat categorical statements as statements about sets. Additionally, they propose to treat singular statements as statements about set membership. So, for example, to represent the statement "a is F" in this retooled Venn diagram, a small letter "a" may be placed inside the circle that represents the set F.



\paragraph{Related concepts}



% [Image: Venn diagram as a truth table]



Figure: Venn diagram as a truth table.



Venn diagrams correspond to truth tables for the propositions $x\in A$, $x\in B$, etc., in the sense that each region of Venn diagram corresponds to one row of the truth table. This type is also known as Johnston diagram. Another way of representing sets is with John F. Randolph's R-diagrams.



\paragraph{Resources}



\begin{itemize}
    \item \href{https://mathresearch.utsa.edu/wikiFiles/MAT1053/Venn_Diagrams_and_Cardinality/MAT1053_M11.2Venn_Diagrams_and_Cardinality.pdf}{Venn Diagrams and Cardinality}, Book Chapter
    \item \href{https://mathresearch.utsa.edu/wikiFiles/MAT1053/Venn_Diagrams_and_Cardinality/MAT1053_M11.2Venn_Diagrams_and_CardinalityGN.pdf}{Guided Notes}
\end{itemize}

\paragraph{Licensing}



Content obtained and/or adapted from:



\begin{itemize}
    \item \href{https://en.wikipedia.org/wiki/Venn_diagram}{Venn diagram, Wikipedia} under a CC BY-SA license
\end{itemize}

\subsection{Learning Objectives}

\begin{enumerate}
    \item \textbf{Interpret Shared Glyphs in Different Alphabets:}
    \item \textbf{Outcome:} Students will be able to identify and classify the uppercase glyphs shared among the Greek, Latin, and Russian alphabets using a Venn diagram.
    \item \textbf{Details:} This includes recognizing the common glyphs between the three alphabets and understanding the visual representation of their relationships through the overlapping regions in a Venn diagram.
    \item \textbf{Analyze Set Relationships Using Venn Diagrams:}
    \item \textbf{Outcome:} Students will be able to describe the logical relationships between the Greek, Latin, and Russian alphabets based on their shared and unique glyphs.
    \item \textbf{Details:} This involves explaining the concepts of intersection and union of sets as depicted in the Venn diagram, including identifying the glyphs that are common to all three alphabets and those that are unique to each alphabet.
    \item \textbf{Apply Set Theory to Linguistics:}
    \item \textbf{Outcome:} Students will be able to apply the principles of set theory to analyze the structure of alphabets and understand the similarities and differences between the Greek, Latin, and Russian alphabets.
    \item \textbf{Details:} This involves using Venn diagrams to model the relationships between different alphabets and to explore how set theory can be used to categorize and analyze linguistic symbols.
\end{enumerate}

\subsection{Assessment Instrument}

\begin{enumerate}
    \item \textbf{Assessment Instrument: Venn Diagrams Quiz} \textbf{Format:}
    \item \textbf{Type:} Mixed-format quiz (multiple-choice, short answer, and problem-solving questions)
    \item \textbf{Duration:} 45 minutes
    \item \textbf{Total Marks:} 30 points \textbf{Sections:}
    \item \textbf{Understanding Shared Glyphs (10 points):}
    \item \textbf{Question 1:} Identify the common glyphs shared between the Greek, Latin, and Russian alphabets as shown in the Venn diagram:
    \item a. Which glyphs are shared between all three alphabets?
    \item b. Which glyphs are shared only between Greek and Latin alphabets?
    \item c. Which glyphs are unique to the Russian alphabet?
    \item \textbf{Question 2:} (Multiple-choice) Which of the following glyphs does \textbf{not} belong to the intersection of all three alphabets?
    \item a. $A$
    \item b. $B$
    \item c. $C$
    \item d. $\Delta$
    \item \textbf{Analyzing Set Relationships (10 points):}
    \item \textbf{Question 3:} (Short answer) Explain the concept of the intersection of sets using the Venn diagram of the Greek, Latin, and Russian alphabets. Provide an example based on the diagram.
    \item \textbf{Question 4:} (Problem-solving) If another alphabet were introduced and overlapped with the existing Venn diagram, how would the regions of intersection change? Describe the new intersections using set notation.
    \item \textbf{Applying Set Theory to Linguistics (10 points):}
    \item \textbf{Question 5:} (Short answer) Using the Venn diagram, describe the union of the Greek and Latin alphabets. How many glyphs are included in this union?
    \item \textbf{Question 6:} (Problem-solving) Create a Venn diagram with three sets representing Greek, Latin, and Russian alphabets. Label each region according to the glyphs found in those regions and calculate the number of glyphs in the symmetric difference between the Latin and Russian alphabets. \textbf{Scoring Rubric:}
    \item \textbf{Understanding Shared Glyphs:} 5 points for correct identification of shared and unique glyphs (total 10 points).
    \item \textbf{Analyzing Set Relationships:} 5 points for accurate explanation and correct problem-solving regarding intersections (total 10 points).
    \item \textbf{Applying Set Theory to Linguistics:} 5 points for correct interpretation of unions and accurate creation and analysis of the Venn diagram (total 10 points). The quiz is designed to assess students' understanding of Venn diagrams in the context of linguistic symbols, their ability to analyze set relationships, and apply set theory principles in real-world scenarios.
\end{enumerate}

%%===============================================================
\section{Cardinality}
\label{sec:cardinality}
%%===============================================================

% [Image: The set S of all Platonic solids has 5 elements. Thus |S|=5.]



Figure: The set $S$ of all Platonic solids has 5 elements. Thus $|S|=5$.



In mathematics, the \textbf{cardinality} of a set is a measure of the "number of elements" of the set. For example, the set $A = \\{2, 4, 6\\}$ contains 3 elements, and therefore $A$ has a cardinality of 3. Beginning in the late 19th century, this concept was generalized to infinite sets, which allows one to distinguish between the different types of infinity, and to perform arithmetic on them. There are two approaches to cardinality: one which compares sets directly using bijections and injections, and another which uses cardinal numbers. The cardinality of a set is also called its \textbf{size}, when no confusion with other notions of size is possible.



The cardinality of a set $A$ is usually denoted $|A|$, with a vertical bar on each side; this is the same notation as absolute value, and the meaning depends on context. The cardinality of a set $A$ may alternatively be denoted by $n(A)$, $\overline{\overline{A}}$, $\operatorname{card}(A)$, or $\\#A$.



\paragraph{History}



Beginning in the late 19th century, this concept was generalized to infinite sets, which allows one to distinguish between the different types of infinity, and to perform arithmetic on them.



\paragraph{Comparing sets}



% [Image: Bijective function from N to the set E of even numbers. Although E is a proper subset of N, both sets have the same cardinality.]



Figure: Bijective function from \textbf{N} to the set \textit{E} of even numbers. Although \textit{E} is a proper subset of \textbf{N}, both sets have the same cardinality.



% [Image: N does not have the same cardinality as its power set P(N)]



Figure: $\mathbf{N}$ does not have the same cardinality as its power set $\mathcal{P}(\mathbf{N})$: For every function $f$ from $\mathbf{N}$ to $\mathcal{P}(\mathbf{N})$, the set $T = \\{n \in \mathbf{N} : n \notin f(n)\\}$ disagrees with every set in the range of $f$, hence $f$ cannot be surjective. The picture shows an example $f$ and the corresponding $T$; \textbf{red}: $n \in f(n) \setminus T$, \textbf{blue}: $n \in T \setminus f(n)$.



While the cardinality of a finite set is just the number of its elements, extending the notion to infinite sets usually starts with defining the notion of comparison of arbitrary sets (some of which are possibly infinite).



\textbf{Definition 1: $\|A\| = \|B\|$}



Two sets $A$ and $B$ have the same cardinality if there exists a bijection (a.k.a., one-to-one correspondence) from $A$ to $B$, that is, a function from $A$ to $B$ that is both injective and surjective. Such sets are said to be \textit{equipotent}, \textit{equipollent}, or \textit{equinumerous}. This relationship can also be denoted $A \approx B$ or $A \sim B$.



For example, the set $E = \\{0, 2, 4, 6, \ldots\\}$ of non-negative even numbers has the same cardinality as the set $\mathbf{N} = \\{0, 1, 2, 3, \ldots\\}$ of natural numbers, since the function $f(n) = 2n$ is a bijection from $\mathbf{N}$ to $E$ (see picture).



\textbf{Definition 2: $\|A\| \leq \|B\|$}



$A$ has cardinality less than or equal to the cardinality of $B$, if there exists an injective function from $A$ into $B$.



\textbf{Definition 3: $\|A\| < \|B\|$}



$A$ has cardinality strictly less than the cardinality of $B$, if there is an injective function, but no bijective function, from $A$ to $B$.



For example, the set $\mathbf{N}$ of all natural numbers has cardinality strictly less than its power set $\mathcal{P}(\mathbf{N})$, because $g(n) = \\{n\\}$ is an injective function from $\mathbf{N}$ to $\mathcal{P}(\mathbf{N})$, and it can be shown that no function from $\mathbf{N}$ to $\mathcal{P}(\mathbf{N})$ can be bijective (see picture). By a similar argument, $\mathbf{N}$ has cardinality strictly less than the cardinality of the set $\mathbf{R}$ of all real numbers. For proofs, see Cantor's diagonal argument or Cantor's first uncountability proof.



If $\|A\| \leq \|B\|$ and $\|B\| \leq \|A\|$, then $\|A\| = \|B\|$ (a fact known as the Schröder–Bernstein theorem). The axiom of choice is equivalent to the statement that $\|A\| \leq \|B\|$ or $\|B\| \leq \|A\|$ for every $A$, $B$.



\paragraph{Cardinal numbers}



In the above section, "cardinality" of a set was defined functionally. In other words, it was not defined as a specific object itself. However, such an object can be defined as follows.



The relation of having the same cardinality is called equinumerosity, and this is an equivalence relation on the class of all sets. The equivalence class of a set \textit{A} under this relation, then, consists of all those sets which have the same cardinality as \textit{A}. There are two ways to define the "cardinality of a set":



1.  The cardinality of a set \textit{A} is defined as its equivalence class under equinumerosity.

2.  A representative set is designated for each equivalence class. The most common choice is the initial ordinal in that class. This is usually taken as the definition of cardinal number in axiomatic set theory.



Assuming the axiom of choice, the cardinalities of the infinite sets are denoted



$$\aleph _0 < \aleph _1 < \aleph _2 < \ldots$$



For each ordinal $\alpha$, $\aleph _{\alpha + 1}$ is the least cardinal number greater than $\aleph _\alpha$.



The cardinality of the natural numbers is denoted aleph-null ($\aleph _0$), while the cardinality of the real numbers is denoted by "$\mathfrak c$" (a lowercase fraktur script "c"), and is also referred to as the cardinality of the continuum. Cantor showed, using the diagonal argument, that ${\mathfrak c} >\aleph _0$. We can show that $\mathfrak c = 2^{\aleph _0}$, this also being the cardinality of the set of all subsets of the natural numbers.



The continuum hypothesis says that $\aleph _1 = 2^{\aleph _0}$, i.e. $2^{\aleph _0}$ is the smallest cardinal number bigger than $\aleph _0$, i.e. there is no set whose cardinality is strictly between that of the integers and that of the real numbers. The continuum hypothesis is independent of ZFC, a standard axiomatization of set theory; that is, it is impossible to prove the continuum hypothesis or its negation from ZFC—provided that ZFC is consistent.



\paragraph{Finite, countable and uncountable sets}



If the axiom of choice holds, the law of trichotomy holds for cardinality. Thus we can make the following definitions:



\begin{itemize}
    \item Any set $X$ with cardinality less than that of the natural numbers, or $\|X\| < \|\mathbf{N}\|$, is said to be a finite set.
    \item Any set $X$ that has the same cardinality as the set of natural numbers, or $\|X\| = \|\mathbf{N}\| = \aleph _0$, is said to be a countably infinite set.
    \item Any set $X$ with cardinality greater than that of the natural numbers, or $\|X\| > \|\mathbf{N}\|$, for example $\|\mathbf{R}\| = \mathfrak{c} > \|\mathbf{N}\|$, is said to be uncountable.
\end{itemize}

\paragraph{Infinite sets}



Our intuition gained from finite sets breaks down when dealing with infinite sets. In the late nineteenth century Georg Cantor, Gottlob Frege, Richard Dedekind and others rejected the view that the whole cannot be the same size as the part. One example of this is Hilbert's paradox of the Grand Hotel. Indeed, Dedekind defined an infinite set as one that can be placed into a one-to-one correspondence with a strict subset (that is, having the same size in Cantor's sense); this notion of infinity is called Dedekind infinite. Cantor introduced the cardinal numbers, and showed—according to his bijection-based definition of size—that some infinite sets are greater than others. The smallest infinite cardinality is that of the natural numbers ($\aleph _0$).



\paragraph{Cardinality of the continuum}



One of Cantor's most important results was that the cardinality of the continuum ($\mathfrak{c}$) is greater than that of the natural numbers ($\aleph _0$); that is, there are more real numbers $\mathbf{R}$ than natural numbers \textbf{N}. Namely, Cantor showed that

$$\mathfrak{c} = 2^{\aleph _0} = \beth _1$$

> (see Beth one) satisfies:



$$2^{\aleph _0} > \aleph _0$$



> (see Cantor's diagonal argument or Cantor's first uncountability proof).



The continuum hypothesis states that there is no cardinal number between the cardinality of the reals and the cardinality of the natural numbers, that is,



$$2^{\aleph _0} = \aleph _1$$



However, this hypothesis can neither be proved nor disproved within the widely accepted ZFC axiomatic set theory, if ZFC is consistent.



Cardinal arithmetic can be used to show not only that the number of points in a real number line is equal to the number of points in any segment of that line, but that this is equal to the number of points on a plane and, indeed, in any finite-dimensional space. These results are highly counterintuitive, because they imply that there exist proper subsets and proper supersets of an infinite set \textit{S} that have the same size as \textit{S}, although \textit{S} contains elements that do not belong to its subsets, and the supersets of \textit{S} contain elements that are not included in it.



The first of these results is apparent by considering, for instance, the tangent function, which provides a one-to-one correspondence between the interval (-½p, ½p) and $\mathbf{R}$ (see also Hilbert's paradox of the Grand Hotel).



The second result was first demonstrated by Cantor in 1878, but it became more apparent in 1890, when Giuseppe Peano introduced the space-filling curves, curved lines that twist and turn enough to fill the whole of any square, or cube, or hypercube, or finite-dimensional space. These curves are not a direct proof that a line has the same number of points as a finite-dimensional space, but they can be used to obtain such a proof.



Cantor also showed that sets with cardinality strictly greater than $\mathfrak c$ exist (see his generalized diagonal argument and theorem). They include, for instance:



\begin{itemize}
    \item the set of all subsets of $\mathbf{R}$, i.e., the power set of $\mathbf{R}$, written $\mathcal{P}(\mathbf{R})$ or $2^{\mathbf{R}}$
    \item the set $\mathbf{R}^{\mathbf{R}}$ of all functions from $\mathbf{R}$ to $\mathbf{R}$
\end{itemize}

Both have cardinality



$$2^\mathfrak {c} = \beth _2 > \mathfrak c$$



> (see Beth two).



The cardinal equalities $\mathfrak{c}^2 = \mathfrak{c},$ $\mathfrak c^{\aleph _0} = \mathfrak c,$ and $\mathfrak c ^{\mathfrak c} = 2^{\mathfrak c}$ can be demonstrated using cardinal arithmetic:



$$\mathfrak{c}^2 = \left(2^{\aleph _0}\right)^2 = 2^{2\times{\aleph _0}} = 2^{\aleph _0} = \mathfrak{c},$$



$$\mathfrak c^{\aleph _0} = \left(2^{\aleph _0}\right)^{\aleph _0} = 2^{{\aleph _0}\times{\aleph _0}} = 2^{\aleph _0} = \mathfrak{c},$$



$$\mathfrak c ^{\mathfrak c} = \left(2^{\aleph _0}\right)^{\mathfrak c} = 2^{\mathfrak c\times\aleph _0} = 2^{\mathfrak c}.$$



\paragraph{Examples and properties}



\begin{itemize}
    \item If $X = \\{a, b, c\\}$ and $Y = \\{\text{apples}, \text{oranges}, \text{peaches}\\}$, then $\|X\| = \|Y\|$ because $\\{(a, \text{apples}), (b, \text{oranges}), (c, \text{peaches})\\}$ is a bijection between the sets $X$ and $Y$. The cardinality of each of $X$ and $Y$ is 3.
    \item If $\|X\| \leq \|Y\|$, then there exists $Z$ such that $\|X\| = \|Z\|$ and $Z \subseteq Y$.
    \item If $\|X\| \leq \|Y\|$ and $\|Y\| \leq \|X\|$, then $\|X\| = \|Y\|$. This holds even for infinite cardinals, and is known as the Cantor–Bernstein–Schroeder theorem.
    \item Sets with cardinality of the continuum include the set of all real numbers, the set of all irrational numbers, and the interval $[0, 1]$.
\end{itemize}

\paragraph{Union and intersection}



If \textit{A} and \textit{B} are disjoint sets, then



$$\left\vert A \cup B \right\vert = \left\vert A \right\vert + \left\vert B \right\vert.$$



From this, one can show that in general, the cardinalities of unions and intersections are related by the following equation:



$$\left\vert C \cup D \right\vert + \left\vert C \cap D \right\vert = \left\vert C \right\vert + \left\vert D \right\vert.$$



\paragraph{Resources}



\begin{itemize}
    \item \href{https://mathresearch.utsa.edu/wikiFiles/MAT1053/Venn_Diagrams_and_Cardinality/MAT1053_M11.2Venn_Diagrams_and_Cardinality.pdf}{Venn Diagrams and Cardinality}, Book Chapter
    \item \href{https://mathresearch.utsa.edu/wikiFiles/MAT1053/Venn_Diagrams_and_Cardinality/MAT1053_M11.2Venn_Diagrams_and_CardinalityGN.pdf}{Guided Notes}
\end{itemize}

\paragraph{Licensing}



Content obtained and/or adapted from:



\begin{itemize}
    \item \href{https://en.wikipedia.org/wiki/Cardinality}{Cardinality, Wikipedia} under a CC BY-SA license
\end{itemize}

\subsection{Learning Objectives}

\begin{enumerate}
    \item Here are three Student Learning Objectives (SLOs) based on the provided context:
    \item \textbf{Understand the Concept of Cardinality:}
    \item \textbf{Outcome:} Students will be able to define the concept of cardinality and distinguish between finite and infinite cardinalities.
    \item \textbf{Details:} This includes understanding how cardinality measures the size of a set, how it applies to both finite and infinite sets, and the notation used to represent cardinality (e.g., $|A|$, $n(A)$). Students will explore examples of both finite sets (such as the set of Platonic solids) and infinite sets (such as the set of natural numbers).
    \item \textbf{Apply Bijection to Compare Sets:}
    \item \textbf{Outcome:} Students will be able to determine whether two sets have the same cardinality by constructing a bijection between them.
    \item \textbf{Details:} This involves identifying when two sets are equipotent by constructing a one-to-one correspondence between the elements of the sets. Students will practice creating bijective functions between finite sets and between infinite sets, such as between the set of natural numbers and the set of even numbers.
    \item \textbf{Analyze Infinite Cardinalities:}
    \item \textbf{Outcome:} Students will be able to explain the difference between countably infinite and uncountably infinite sets and demonstrate the use of Cantor's diagonal argument.
    \item \textbf{Details:} This includes distinguishing between different sizes of infinity, such as $\aleph _0$ for the cardinality of the set of natural numbers and $\mathfrak{c}$ for the cardinality of the real numbers. Students will explore the implications of Cantor's diagonal argument and understand why $\mathfrak{c}$ is greater than $\aleph _0$.
\end{enumerate}

\subsection{Assessment Instrument}

\begin{enumerate}
    \item \textbf{Assessment Instrument: Cardinality Quiz} \textbf{Format:}
    \item \textbf{Type:} Mixed-format quiz (multiple-choice, short answer, and problem-solving questions)
    \item \textbf{Duration:} 45 minutes
    \item \textbf{Total Marks:} 30 points \textbf{Sections:}
    \item \textbf{Understanding Cardinality (10 points):}
    \item \textbf{Question 1:} Define the concept of cardinality and explain its significance in set theory. Provide an example of a set with finite cardinality and another with infinite cardinality. (Short Answer)
    \item \textbf{Question 2:} Multiple Choice: Which of the following statements is true regarding cardinality?
    \item a. The cardinality of the set of natural numbers is larger than the cardinality of the set of even numbers.
    \item b. Two sets have the same cardinality if there exists a bijective function between them.
    \item c. The cardinality of any finite set is undefined.
    \item d. The set of real numbers has a smaller cardinality than the set of natural numbers.
    \item \textbf{Comparing Sets Using Bijection (10 points):}
    \item \textbf{Question 3:} Provide a bijective function between the set of natural numbers $\mathbf{N}$ and the set of even numbers $E = \\{0, 2, 4, 6, \ldots\\}$. Show that this function is both injective and surjective. (Problem-solving)
    \item \textbf{Question 4:} Multiple Choice: If $\|A\| \leq \|B\|$, what does it imply?
    \item a. There is a bijection from $A$ to $B$.
    \item b. There is an injective function from $A$ to $B$.
    \item c. The cardinality of $A$ is greater than that of $B$.
    \item d. $A$ and $B$ must be finite sets.
    \item \textbf{Analyzing Infinite Cardinalities (10 points):}
    \item \textbf{Question 5:} Explain Cantor's diagonal argument and how it demonstrates that the set of real numbers has a greater cardinality than the set of natural numbers. (Short Answer)
    \item \textbf{Question 6:} Multiple Choice: Which of the following best describes the cardinality of the set of real numbers?
    \item a. Countably infinite
    \item b. Finite
    \item c. Uncountably infinite
    \item d. None of the above \textbf{Scoring Rubric:}
    \item \textbf{Understanding Cardinality:} 5 points for Question 1, 5 points for Question 2 (total 10 points)
    \item \textbf{Comparing Sets Using Bijection:} 5 points for Question 3, 5 points for Question 4 (total 10 points)
    \item \textbf{Analyzing Infinite Cardinalities:} 5 points for Question 5, 5 points for Question 6 (total 10 points) The quiz assesses students' understanding of cardinality, their ability to compare sets using bijections, and their comprehension of infinite cardinalities, particularly through Cantor's diagonal argument.
\end{enumerate}

%%===============================================================
\section{Introduction to Probability}
\label{sec:introduction-to-probability}
%%===============================================================

% [Image: The probabilities of rolling several numbers using two dice.]



Figure: The probabilities of rolling several numbers using two dice.



\textbf{Probability} is the branch of mathematics concerning numerical descriptions of how likely an event is to occur, or how likely it is that a proposition is true. The probability of an event is a number between 0 and 1, where, roughly speaking, 0 indicates impossibility of the event and 1 indicates certainty. The higher the probability of an event, the more likely it is that the event will occur. A simple example is the tossing of a fair (unbiased) coin. Since the coin is fair, the two outcomes ("heads" and "tails") are both equally probable; the probability of "heads" equals the probability of "tails"; and since no other outcomes are possible, the probability of either "heads" or "tails" is 1/2 (which could also be written as 0.5 or 50\%).



These concepts have been given an axiomatic mathematical formalization in probability theory, which is used widely in areas of study such as statistics, mathematics, science, finance, gambling, artificial intelligence, machine learning, computer science, game theory, and philosophy to, for example, draw inferences about the expected frequency of events. Probability theory is also used to describe the underlying mechanics and regularities of complex systems.



\paragraph{Terminology of the probability theory}



\textbf{Experiment:} An operation which can produce some well-defined outcomes, is called an Experiment.



> \textbf{Example:} When we toss a coin, we know that either head or tail shows up. So, the operation of tossing a coin may be said to have two well-defined outcomes, namely, (a) heads showing up; and (b) tails showing up.



\textbf{Random Experiment:} When we roll a die we are well aware of the fact that any of the numerals 1,2,3,4,5, or 6 may appear on the upper face but we cannot say that which exact number will show up.



> \textit{Such an experiment in which all possible outcomes are known and the exact outcome cannot be predicted in advance, is called a Random Experiment.}



\textbf{Sample Space:} All the possible outcomes of an experiment as an whole, form the Sample Space.



> \textbf{Example:} When we roll a die we can get any outcome from 1 to 6. All the possible numbers which can appear on the upper face form the Sample Space(denoted by S). Hence, the Sample Space of a dice roll is \textbf{S={1,2,3,4,5,6}}



\textbf{Outcome:} Any possible result out of the Sample Space \textbf{S} for a Random Experiment is called an Outcome.



> \textbf{Example:} When we roll a die, we might obtain 3 or when we toss a coin, we might obtain heads.



\textbf{Event:} Any subset of the Sample Space \textbf{S} is called an Event (denoted by \textbf{E}). When an outcome which belongs to the subset \textbf{E} takes place, it is said that an Event has occurred. Whereas, when an outcome which does not belong to subset \textbf{E} takes place, the Event has not occurred.



> \textbf{Example:} Consider the experiment of throwing a die. Over here the Sample Space \textbf{S={1,2,3,4,5,6}.} Let \textbf{E} denote the event of 'a number appearing less than 4.' Thus the Event \textbf{E={1,2,3}.} If the number 1 appears, we say that Event \textbf{E} has occurred. Similarly, if the outcomes are 2 or 3, we can say Event \textbf{E} has occurred since these outcomes belong to subset \textbf{E.}



\textbf{Trial:} By a trial, we mean performing a random experiment.



> \textbf{Example:} (i) Tossing a \textit{fair} coin, (ii) rolling an unbiased die.



\paragraph{Theory}



Like other theories, the theory of probability is a representation of its concepts in formal terms---that is, in terms that can be considered separately from their meaning. These formal terms are manipulated by the rules of mathematics and logic, and any results are interpreted or translated back into the problem domain.



There have been at least two successful attempts to formalize probability, namely the Kolmogorov formulation and the Cox formulation. In Kolmogorov's formulation (see also probability space), sets are interpreted as events and probability as a measure on a class of sets. In Cox's theorem, probability is taken as a primitive (i.e., not further analyzed), and the emphasis is on constructing a consistent assignment of probability values to propositions. In both cases, the laws of probability are the same, except for technical details.



There are other methods for quantifying uncertainty, such as the Dempster--Shafer theory or possibility theory, but those are essentially different and not compatible with the usually-understood laws of probability.



\paragraph{Applications}



Probability theory is applied in everyday life in risk assessment and modeling. The insurance industry and markets use actuarial science to determine pricing and make trading decisions. Governments apply probabilistic methods in environmental regulation, entitlement analysis, and financial regulation.



An example of the use of probability theory in equity trading is the effect of the perceived probability of any widespread Middle East conflict on oil prices, which have ripple effects in the economy as a whole. An assessment by a commodity trader that a war is more likely can send that commodity's prices up or down, and signals other traders of that opinion. Accordingly, the probabilities are neither assessed independently nor necessarily rationally. The theory of behavioral finance emerged to describe the effect of such groupthink on pricing, on policy, and on peace and conflict.



In addition to financial assessment, probability can be used to analyze trends in biology (e.g., disease spread) as well as ecology (e.g., biological Punnett squares). As with finance, risk assessment can be used as a statistical tool to calculate the likelihood of undesirable events occurring, and can assist with implementing protocols to avoid encountering such circumstances. Probability is used to design games of chance so that casinos can make a guaranteed profit, yet provide payouts to players that are frequent enough to encourage continued play.



Another significant application of probability theory in everyday life is reliability. Many consumer products, such as automobiles and consumer electronics, use reliability theory in product design to reduce the probability of failure. Failure probability may influence a manufacturer's decisions on a product's warranty.



The cache language model and other statistical language models that are used in natural language processing are also examples of applications of probability theory.



\paragraph{Mathematical treatment}



Consider an experiment that can produce a number of results. The collection of all possible results is called the sample space of the experiment, sometimes denoted as $\Omega$. The power set of the sample space is formed by considering all different collections of possible results. For example, rolling a die can produce six possible results. One collection of possible results gives an odd number on the die. Thus, the subset {1,3,5} is an element of the power set of the sample space of dice rolls. These collections are called "events". In this case, {1,3,5} is the event that the die falls on some odd number. If the results that actually occur fall in a given event, the event is said to have occurred.



A probability is a way of assigning every event a value between zero and one, with the requirement that the event made up of all possible results (in our example, the event {1,2,3,4,5,6}) is assigned a value of one. To qualify as a probability, the assignment of values must satisfy the requirement that for any collection of mutually exclusive events (events with no common results, such as the events {1,6}, {3}, and {2,4}), the probability that at least one of the events will occur is given by the sum of the probabilities of all the individual events.



The probability of an event \textit{A} is written as $P(A)$, $p(A)$, or $\text{Pr}(A)$. This mathematical definition of probability can extend to infinite sample spaces, and even uncountable sample spaces, using the concept of a measure.



The \textit{opposite} or \textit{complement} of an event \textit{A} is the event \[not \textit{A}\] (that is, the event of \textit{A} not occurring), often denoted as $A', A^c$, $\overline{A}, A^\complement, \neg A$, or ${\sim}A$; its probability is given by \textit{P}(not \textit{A}) = 1 - \textit{P}(\textit{A}). As an example, the chance of not rolling a six on a six-sided die is 1 -- (chance of rolling a six) = 1 - $\tfrac{1}{6} = \tfrac{5}{6}$.



If two events \textit{A} and \textit{B} occur on a single performance of an experiment, this is called the intersection or joint probability of \textit{A} and \textit{B}, denoted as $P(A \cap B)$.



\paragraph{Independent events}



If two events, \textit{A} and \textit{B} are independent then the joint probability is



$$P(A \text{ and }B) =  P(A \cap B) = P(A) P(B).$$ For example, if two coins are flipped, then the chance of both being heads is $\tfrac{1}{2}\times\tfrac{1}{2} = \tfrac{1}{4}$.



\paragraph{Mutually exclusive events}



If either event \textit{A} or event \textit{B} can occur but never both simultaneously, then they are called mutually exclusive events.



If two events are mutually exclusive, then the probability of \textit{both} occurring is denoted as $P(A \cap B)$ and



$$P(A \text{ and }B) =  P(A \cap B) = 0$$ If two events are mutually exclusive, then the probability of \textit{either} occurring is denoted as $P(A \cup B)$ and



$$ \begin{aligned} P(A\text{ or }B) & = P(A \cup B) \\\\ & = P(A) + P(B) - P(A \cap B) \\\\ & = P(A) + P(B) - 0 \\\\ & = P(A) + P(B) \end{aligned} $$



For example, the chance of rolling a 1 or 2 on a six-sided dice is $P(1\text{ or }2) = P(1) + P(2) = \tfrac{1}{6} + \tfrac{1}{6} = \tfrac{1}{3}.$



\paragraph{Not mutually exclusive events}



If the events are not mutually exclusive then



$$ \begin{aligned} P\left(A \text{ or } B\right) &= P(A \cup B) \\\\ &= P\left(A\right) + P\left(B\right) - P\left(A \text{ and } B\right). \end{aligned} $$



For example, when drawing a single card at random from a regular deck of cards, the chance of getting a heart or a face card $(J,Q,K)$ (or one that is both) is $\tfrac{13}{52} + \tfrac{12}{52} - \tfrac{3}{52} = \tfrac{11}{26}$, since among the 52 cards of a deck, 13 are hearts, 12 are face cards, and 3 are both: here the possibilities included in the "3 that are both" are included in each of the "13 hearts" and the "12 face cards", but should only be counted once.



\paragraph{Conditional probability}



\textit{Conditional probability} is the probability of some event \textit{A}, given the occurrence of some other event \textit{B}. Conditional probability is written $P(A \mid B)$, and is read "the probability of \textit{A}, given \textit{B}". It is defined by



$$P(A \mid B) = \frac{P(A \cap B)}{P(B)}.\,$$



If $P(B)=0$ then $P(A \mid B)$ is formally undefined by this expression. In this case $A$ and $B$ are independent, since $P(A \cap B) = P(A)P(B) = 0$. However, it is possible to define a conditional probability for some zero-probability events using a s-algebra of such events (such as those arising from a continuous random variable).



For example, in a bag of 2 red balls and 2 blue balls (4 balls in total), the probability of taking a red ball is $1/2$; however, when taking a second ball, the probability of it being either a red ball or a blue ball depends on the ball previously taken. For example, if a red ball was taken, then the probability of picking a red ball again would be $1/3$, since only 1 red and 2 blue balls would have been remaining. And if a blue ball was taken previously, the probability of taking a red ball will be $2/3$.



\paragraph{Inverse probability}



In probability theory and applications, \textit{Bayes' rule} relates the odds of event $A _1$ to event $A _2$, before (prior to) and after (posterior to) conditioning on another event $B$. The odds on $A _1$ to event $A _2$ is simply the ratio of the probabilities of the two events. When arbitrarily many events $A$ are of interest, not just two, the rule can be rephrased as \textit{posterior is proportional to prior times likelihood}, $P(A|B)\propto P(A) P(B|A)$ where the proportionality symbol means that the left hand side is proportional to (i.e., equals a constant times) the right hand side as $A$ varies, for fixed or given $B$ (Lee, 2012; Bertsch McGrayne, 2012). In this form it goes back to Laplace (1774) and to Cournot (1843); see Fienberg (2005).



\paragraph{Summary of probabilities}



$$ \begin{array}{|c|l|} \hline \textbf{Event} & \textbf{Probability} \\\\ \hline A & P(A) \in [0, 1] \\\\ \hline \text{not A} & P(A^\complement) = 1 - P(A) \\\\ \hline \text{A or B} & \begin{aligned} P(A \cup B) & = P(A) + P(B) - P(A \cap B) \\\\ P(A \cup B) & = P(A) + P(B) \quad \text{if A and B are mutually exclusive} \end{aligned} \\\\ \hline \text{A and B} & \begin{aligned} P(A \cap B) & = P(A \mid B)P(B) = P(B \mid A)P(A) \\\\ P(A \cap B) & = P(A)P(B) \quad \text{if A and B are independent} \end{aligned} \\\\ \hline \text{A given B} & P(A \mid B) = \frac{P(A \cap B)}{P(B)} = \frac{P(B \mid A)P(A)}{P(B)} \\\\ \hline \end{array} $$ 



\paragraph{Resources}



\begin{itemize}
    \item \href{https://www.khanacademy.org/math/cc-seventh-grade-math/cc-7th-probability-statistics/cc-7th-basic-prob/v/basic-probability}{Introduction to Theoretical Probability}, Khan Academy
\end{itemize}

\paragraph{Licensing}



Content obtained and/or adapted from:



\begin{itemize}
    \item \href{https://en.wikipedia.org/wiki/Probability}{Probability, Wikipedia} under a CC BY-SA license
\end{itemize}

\subsection{Learning Objectives}

\begin{enumerate}
    \item Here are three Student Learning Objectives (SLOs) based on the provided image and the explanation of probability concepts:
    \item \textbf{Calculate Probability for Specific Events:}
    \item \textbf{Outcome:} Students will be able to calculate the probability of rolling a specific sum with two dice.
    \item \textbf{Details:} This includes understanding the total possible outcomes when rolling two dice (36 outcomes) and calculating the probability of rolling specific sums such as 7, 11, or 2 by counting the favorable outcomes for each event.
    \item \textbf{Analyze Probability Distributions:}
    \item \textbf{Outcome:} Students will be able to analyze and interpret the probability distribution of sums obtained by rolling two dice.
    \item \textbf{Details:} This involves identifying the most likely outcomes (such as a sum of 7) and the least likely outcomes (such as a sum of 2 or 12), and explaining why certain sums have higher or lower probabilities based on the number of combinations that produce those sums.
    \item \textbf{Apply Conditional Probability:}
    \item \textbf{Outcome:} Students will be able to apply conditional probability to scenarios involving rolling two dice.
    \item \textbf{Details:} This includes determining the probability of rolling a specific number given that the first die shows a particular value. For example, students might calculate the probability of rolling a sum of 8 given that the first die shows a 3. These objectives are designed to build a foundational understanding of probability as it applies to the specific case of rolling two dice and can be expanded to more complex scenarios as students progress.
\end{enumerate}

\subsection{Assessment Instrument}

\begin{enumerate}
    \item \textbf{Assessment Instrument: Introduction to Probability Quiz} \textbf{Format:}
    \item \textbf{Type:} Mixed-format quiz (multiple-choice, short answer, and problem-solving questions)
    \item \textbf{Duration:} 45 minutes
    \item \textbf{Total Marks:} 30 points \textbf{Sections:}
    \item \textbf{Calculating Probability for Specific Events (10 points):}
    \item \textbf{Question 1:} (Multiple Choice) What is the probability of rolling a sum of 7 with two dice?
    \item a. $\frac{5}{36}$
    \item b. $\frac{7}{36}$
    \item c. $\frac{1}{6}$
    \item d. $\frac{1}{4}$
    \item \textbf{Question 2:} (Short Answer) If you roll two dice, how many possible outcomes will result in a sum of 11? What is the probability of this event?
    \item \textbf{Analyzing Probability Distributions (10 points):}
    \item \textbf{Question 3:} (Problem-Solving) Using the probability distribution for the sums of two dice, explain why a sum of 7 is more likely than a sum of 2.
    \item \textbf{Question 4:} (Short Answer) Determine the probability of rolling a sum greater than or equal to 10 with two dice.
    \item \textbf{Applying Conditional Probability (10 points):}
    \item \textbf{Question 5:} (Problem-Solving) Given that the first die shows a 3, what is the probability that the sum of two dice will be 8?
    \item \textbf{Question 6:} (Multiple Choice) If the first die shows a 6, what is the probability that the total sum will be 9?
    \item a. $\frac{1}{6}$
    \item b. $\frac{1}{12}$
    \item c. $\frac{1}{3}$
    \item d. $\frac{1}{18}$ \textbf{Scoring Rubric:}
    \item \textbf{Calculating Probability for Specific Events:} 5 points for each correct answer (total 10 points)
    \item \textbf{Analyzing Probability Distributions:} 5 points for each correct analysis or calculation (total 10 points)
    \item \textbf{Applying Conditional Probability:} 5 points for each correct application of conditional probability (total 10 points) This quiz assesses students' understanding of basic probability concepts, including calculating probabilities for specific events, analyzing probability distributions, and applying conditional probability. It ensures that students can reason about likelihoods in the context of rolling two dice and understand the underlying principles of probability theory.
\end{enumerate}

%%===============================================================
\section{Expected Value}
\label{sec:expected-value}
%%===============================================================

Let $X$ be a (discrete) random variable with a finite number of finite outcomes $x _1, x _2, \ldots, x _k$ occurring with probabilities $p _1, p _2, \ldots, p _k,$ respectively. The \textbf{expectation} of $X$ is defined as



$$\operatorname{E}[X] =\sum _{i=1}^k x _i\,p _i=x _1p _1 + x _2p _2 + \cdots + x _kp _k.$$



Since $p _1 + p _2 + \cdots + p _k = 1,$ the expected value is the weighted sum of the $x _i$ values, with the probabilities $p _i$ as the weights.



If all outcomes $x _i$ are equiprobable (that is, $p _1 = p _2 = \cdots = p _k$), then the weighted average turns into the simple average. On the other hand, if the outcomes $x _i$ are not equiprobable, then the simple average must be replaced with the weighted average, which takes into account the fact that some outcomes are more likely than others.



% [Image: An illustration of the convergence of sequence averages of rolls of a dice to the expected value of 3.5 as the number of rolls (trials) grows.]



Figure: An illustration of the convergence of sequence averages of rolls of a dice to the expected value of 3.5 as the number of rolls (trials) grows.



\paragraph{Examples}



\begin{itemize}
    \item Let $X$ represent the outcome of a roll of a fair six-sided dice. More specifically, $X$ will be the number of pips showing on the top face of the dice after the toss. The possible values for $X$ are 1, 2, 3, 4, 5, and 6, all of which are equally likely with a probability of 1/6. The expectation of $X$ is
\end{itemize}

$$\operatorname{E}[X] = 1\cdot\frac16 + 2\cdot\frac16 + 3\cdot\frac16 + 4\cdot\frac16 + 5\cdot\frac16 + 6\cdot\frac16 = 3.5.$$



> If one rolls the dice $n$ times and computes the average (arithmetic mean) of the results, then as $n$ grows, the average will almost surely converge to the expected value, a fact known as the strong law of large numbers.



\begin{itemize}
    \item The roulette game consists of a small ball and a wheel with 38 numbered pockets around the edge. As the wheel is spun, the ball bounces around randomly until it settles down in one of the pockets. Suppose random variable $X$ represents the (monetary) outcome of a $ \$1 $ bet on a single number ("straight up" bet). If the bet wins (which happens with probability 1/38 in American roulette), the payoff is $ \$ 35 $; otherwise the player loses the bet. The expected profit from such a bet will be
\end{itemize}

$$\operatorname{E}[\,\text{gain from }\$1\text{ bet}\,] = -\$1 \cdot \frac{37}{38} + \$35 \cdot \frac{1}{38} = -\$\frac{1}{19}.$$



> That is, the bet of $\text{$\\$$}$1 stands to lose $\text{$\\$$}$ $\frac{1}{19}$, so its expected value is –$\text{$\\$$}$ $\frac{1}{19} \approx \$ 0.05.$



\paragraph{Resources}



\begin{itemize}
    \item \href{https://mathresearch.utsa.edu/wikiFiles/MAT1053/Expected_Value/MAT1053_M12.2Expected_Value.pdf}{Expected Value}, Book Chapter
    \item \href{https://mathresearch.utsa.edu/wikiFiles/MAT1053/Expected_Value/MAT1053_M12.2Expected_Value.pdf}{Guided Notes}
\end{itemize}

\paragraph{Licensing}



Content obtained and/or adapted from:



\begin{itemize}
    \item \href{https://en.wikipedia.org/wiki/Expected_value}{Expected Value, Wikipedia} under a CC BY-SA license
\end{itemize}

\subsection{Learning Objectives}

\begin{enumerate}
    \item Here are three Student Learning Objectives (SLOs) based on the provided content:
    \item \textbf{Calculate Expected Value of Discrete Random Variables:}
    \item \textbf{Outcome:} Students will be able to calculate the expected value of a discrete random variable with a finite number of outcomes.
    \item \textbf{Details:} This includes understanding the concept of expected value as a weighted average, applying the formula for expectation, and solving problems where probabilities and outcomes are given. Examples include calculating the expected value for scenarios such as rolling a fair dice or a roulette game.
    \item \textbf{Interpret the Law of Large Numbers:}
    \item \textbf{Outcome:} Students will be able to explain and apply the Law of Large Numbers in the context of discrete random variables.
    \item \textbf{Details:} This includes recognizing the significance of the expected value as the long-run average outcome, understanding how repeated trials of a random experiment lead to convergence towards the expected value, and illustrating this concept with examples such as dice rolls and roulette bets.
    \item \textbf{Differentiate Between Simple and Weighted Averages:}
    \item \textbf{Outcome:} Students will be able to differentiate between simple and weighted averages in the context of probability distributions.
    \item \textbf{Details:} This includes identifying situations where outcomes are equiprobable versus when they are not, and applying the correct type of average (simple or weighted) to calculate the expected value. Students will practice calculating both types of averages and understand the conditions under which each is used.
\end{enumerate}

\subsection{Assessment Instrument}

\begin{enumerate}
    \item \textbf{Assessment Instrument: Expected Value Quiz} \textbf{Format:}
    \item \textbf{Type:} Mixed-format quiz (multiple-choice, short answer, and problem-solving questions)
    \item \textbf{Duration:} 45 minutes
    \item \textbf{Total Marks:} 30 points \textbf{Sections:}
    \item \textbf{Calculating Expected Value (10 points):}
    \item \textbf{Question 1:} (Multiple Choice) A fair six-sided dice is rolled. What is the expected value of the outcome?
    \item a. 2.5
    \item b. 3.5
    \item c. 4.5
    \item d. 5.5
    \item \textbf{Question 2:} (Short Answer) A random variable $X$ has possible outcomes 1, 2, and 3 with probabilities 0.2, 0.5, and 0.3, respectively. Calculate the expected value of $X$.
    \item \textbf{Question 3:} (Problem Solving) Suppose you are playing a game where you flip a fair coin. If it lands heads, you win `$`5; if it lands tails, you lose `$`2. What is the expected value of your winnings?
    \item \textbf{Applying the Law of Large Numbers (10 points):}
    \item \textbf{Question 4:} (Multiple Choice) As the number of trials increases, the average result of rolling a fair six-sided dice should converge to:
    \item a. 1.0
    \item b. 3.5
    \item c. 6.0
    \item d. 7.0
    \item \textbf{Question 5:} (Short Answer) Explain how the Law of Large Numbers applies to the expected value of rolling a six-sided dice multiple times.
    \item \textbf{Question 6:} (Problem Solving) A roulette wheel has 38 pockets, 18 of which are red, 18 black, and 2 green. If you bet `$`1 on red, you win `$`1 if the ball lands on red and lose `$`1 if it lands on any other color. Calculate the expected value of your bet.
    \item \textbf{Differentiating Between Simple and Weighted Averages (10 points):}
    \item \textbf{Question 7:} (Multiple Choice) If all outcomes of a random variable are equally likely, the expected value is equivalent to:
    \item a. The simple average
    \item b. The median
    \item c. The mode
    \item d. The weighted average
    \item \textbf{Question 8:} (Short Answer) Consider a discrete random variable with outcomes 3, 6, and 9, and probabilities 0.1, 0.3, and 0.6, respectively. Calculate the weighted average (expected value) of the outcomes.
    \item \textbf{Question 9:} (Problem Solving) A game involves drawing a card from a deck of 10 cards numbered 1 through 10. The payout is equal to the number on the card drawn, multiplied by 10 cents. What is the expected payout from drawing a card? \textbf{Scoring Rubric:}
    \item \textbf{Calculating Expected Value:} 2 points for each correct answer (total 10 points)
    \item \textbf{Applying the Law of Large Numbers:} 2 points for each correct answer (total 10 points)
    \item \textbf{Differentiating Between Simple and Weighted Averages:} 2 points for each correct answer (total 10 points) This quiz will assess the students' understanding and ability to calculate expected values, apply the Law of Large Numbers, and differentiate between simple and weighted averages. It ensures they can solve problems and interpret the concepts accurately.
\end{enumerate}

%%===============================================================
\section{Conditional Probability}
\label{sec:conditional-probability}
%%===============================================================

In probability theory, \textbf{conditional probability} is a measure of the probability of an event occurring, given that another event (by assumption, presumption, assertion, or evidence) has already occurred. If the event of interest is $A$ and the event $B$ is known or assumed to have occurred, "the conditional probability of $A$ given $B$," or "the probability of $A$ under the condition $B$," is usually written as $P(A \mid B)$ or occasionally $P _{B}(A)$. This can also be understood as the fraction of probability $B$ that intersects with $A$:

$$P(A \mid B) = \frac{P(A \cap B)}{P(B)}.$$



For example, the probability that any given person has a cough on any given day may be only 5\%. But if we know or assume that the person is sick, then they are much more likely to be coughing. For example, the conditional probability that someone unwell is coughing might be 75\%, in which case we would have that $P(\text{Cough}) = 5\%$ and $P(\text{Cough} \mid \text{Sick}) = 75\%$. However, there doesn't have to be a relationship or dependence between $A$ and $B$, and they don't have to occur simultaneously.



$P(A \mid B)$ may or may not be equal to $P(A)$ (the unconditional probability of $A$). If $P(A \mid B) = P(A)$, then events $A$ and $B$ are said to be \textit{independent}: in such a case, knowledge about either event does not alter the likelihood of the other. $P(A \mid B)$ (the conditional probability of $A$ given $B$) typically differs from $P(B \mid A)$.



For example, if a person has dengue fever, they might have a 90\% chance of testing positive for the disease. In this case, what is being measured is that if event $B$ (having dengue) has occurred, the probability of $A$ (testing positive) given that $B$ occurred is 90\%: $P(A \mid B) = 90\%$. Alternatively, if a person tests positive for dengue fever, they may have only a 15\% chance of actually having this rare disease due to high false positive rates. In this case, the probability of the event $B$ (having dengue) given that the event $A$ (testing positive) has occurred is 15\%: $P(B \mid A) = 15\%$. It should be apparent now that falsely equating the two probabilities can lead to various errors of reasoning, which is commonly seen through base rate fallacies.



While conditional probabilities can provide extremely useful information, limited information is often supplied or at hand. Therefore, it can be useful to reverse or convert a conditional probability using Bayes' theorem:

$$P(A \mid B) = \frac{P(B \mid A) \cdot P(A)}{P(B)}.$$

Another option is to display conditional probabilities in a conditional probability table to illuminate the relationship between events.



\paragraph{Definition}



% [Image: Illustration of conditional probabilities with an Euler diagram.]



Figure: Illustration of conditional probabilities with an Euler diagram. The unconditional probability P(\textit{A}) = 0.30 + 0.10 + 0.12 = 0.52. However, the conditional probability $P(A \| B _{1}) = 1$, $P(A \| B _{2}) = 0.12 \div (0.12 + 0.04) = 0.75$, and $P(A\|B _{3}) = 0$.



% [Image: On a tree diagram, branch probabilities are conditional on the event associated with the parent node. (Here, the overbars indicate that the event does not occur.)]



Figure: On a tree diagram, branch probabilities are conditional on the event associated with the parent node. (Here, the overbars indicate that the event does not occur.)



% [Image: Venn Pie Chart describing conditional probabilities]



Figure: Venn Pie Chart describing conditional probabilities.



\paragraph{Conditioning on an event}



\paragraph{Kolmogorov definition}



Given two events $A$ and $B$ from the sigma-field of a probability space, with the unconditional probability of $B$ being greater than zero (i.e., $P(B) > 0$), the conditional probability of $A$ given $B$ is defined to be the quotient of the probability of the joint events $A$ and $B$, and the probability of $B$:



$$P(A \mid B) = \frac{P(A \cap B)}{P(B)},$$



where $P(A \cap B)$ is the probability that both events $A$ and $B$ occur. This may be visualized as restricting the sample space to situations in which $B$ occurs. The logic behind this equation is that if the possible outcomes for $A$ and $B$ are restricted to those in which $B$ occurs, this set serves as the new sample space.



Note that the above equation is a definition—not a theoretical result. We just denote the quantity $\frac{P(A \cap B)}{P(B)}$ as $P(A \mid B)$, and call it the conditional probability of $A$ given $B$.



\paragraph{As an axiom of probability}



Some authors, such as de Finetti, prefer to introduce conditional probability as an axiom of probability:



$$P(A \cap B) = P(A \mid B)P(B)$$



Although mathematically equivalent, this may be preferred philosophically; under major probability interpretations, such as the subjective theory, conditional probability is considered a primitive entity. Moreover, this "multiplication rule" can be practically useful in computing the probability of $A \cap B$ and introduces a symmetry with the summation axiom for mutually exclusive events



$$P(A \cup B) = P(A) + P(B) - P(A \cap B)$$



\paragraph{As the probability of a conditional event}



Conditional probability can be defined as the probability of a conditional event $A_B$. The Goodman–Nguyen–Van Fraassen conditional event can be defined as



$$A _B = \bigcup _{i \ge 1} \left( \bigcap _{j < i} \overline{B} _j, A _i B _i \right)$$



It can be shown that



$$P(A _B)= \frac{P(A \cap B)}{P(B)}$$



which meets the Kolmogorov definition of conditional probability.



\paragraph{Conditioning on an event of probability zero}



If $P(B)=0$, then according to the definition, $P(A|B)$ is undefined.



The case of greatest interest is that of a random variable $Y$, conditioned on a continuous random variable $X$ resulting in a particular outcome $x$. The event $B = \\{ X = x \\}$ has probability zero and, as such, cannot be conditioned on.



Instead of conditioning on $X$ being \textit{exactly} $x$, we could condition on it being closer than distance $\epsilon$ away from $x$. The event $B = \\{ x-\epsilon < X < x+\epsilon \\}$ will generally have nonzero probability and hence, can be conditioned on. We can then take the limit



$$\lim _{\epsilon \to 0} P(A \mid x-\epsilon < X < x+\epsilon).$$



For example, if two continuous random variables $X$ and $Y$ have a joint density $f _{X,Y}(x,y)$, then by L'Hôpital's rule:



$$ \begin{aligned} \lim _{\epsilon \to 0} P(Y \in U \mid x _0-\epsilon < X < x _0+\epsilon) & = \lim _{\epsilon \to 0} \frac{\int _{x _0-\epsilon}^{x _0+\epsilon} \int _U f _{X, Y}(x, y) \,\mathrm{d}y \,\mathrm{d}x}{\int _{x _0-\epsilon}^{x _0+\epsilon} \int _\mathbb{R} f _{X, Y}(x, y) \,\mathrm{d}y \,\mathrm{d}x} \hspace{30em} \\\\ & = \frac{\int _U f _{X, Y}(x _0, y) \,\mathrm{d}y}{\int _\mathbb{R} f _{X, Y}(x _0, y) \,\mathrm{d}y} \end{aligned} $$



The resulting limit is the conditional probability distribution of $Y$ given $X$ and exists when the denominator, the probability density $f _X(x _0)$, is strictly positive.



It is tempting to \textit{define} the undefined probability $P(A \mid X = x)$ using this limit, but this cannot be done in a consistent manner. In particular, it is possible to find random variables $X$ and $W$ and values $x$, $w$ such that the events $\\{X = x\\}$ and $\\{W = w\\}$ are identical but the resulting limits are not.



$$\lim _{\epsilon \to 0} P(A \mid x-\epsilon \le X \le x+\epsilon) \neq \lim _{\epsilon \to 0} P(A \mid w-\epsilon \le W \le w+\epsilon).$$ The Borel–Kolmogorov paradox demonstrates this with a geometrical argument.



\paragraph{Conditioning on a discrete random variable}



Let $X$ be a discrete random variable and its possible outcomes denoted $V$. For example, if $X$ represents the value of a rolled die, then $V$ is the set $\\{ 1, 2, 3, 4, 5, 6 \\}$. Let us assume for the sake of presentation that $X$ is a discrete random variable, so that each value in $V$ has a nonzero probability.



For a value $x$ in $V$ and an event $A$, the conditional probability is given by $P(A \mid X=x)$. Writing



$$c(x, A) = P(A \mid X=x)$$



for short, we see that it is a function of two variables, $x$ and $A$.



For a fixed $A$, we can form the random variable $Y = c(X, A)$. It represents an outcome of $P(A \mid X=x)$ whenever a value $x$ of $X$ is observed.



The conditional probability of $A$ given $X$ can thus be treated as a random variable $Y$ with outcomes in the interval $[0,1]$. From the law of total probability, its expected value is equal to the unconditional probability of $A$.



\paragraph{Partial conditional probability}



The partial conditional probability $P(A\mid B _1 \equiv b _1, \ldots, B _m \equiv b _m)$ is about the probability of event $A$ given that each of the condition events $B _i$ has occurred to a degree $b _i$ (degree of belief, degree of experience) that might be different from 100\%. Frequentistically, partial conditional probability makes sense, if the conditions are tested in experiment repetitions of appropriate length $n$. Such $n$-bounded partial conditional probability can be defined as the conditionally expected average occurrence of event $A$ in testbeds of length $n$ that adhere to all of the probability specifications $B _i \equiv b _i$, i.e.:



$$P^n(A\mid B _1 \equiv b _1, \ldots, B _m \equiv b _m)=

\operatorname E(\overline{A}^n\mid\overline{B}^n _1=b _1, \ldots, \overline{B}^n _m=b _m)$$ Based on that, partial conditional probability can be defined as



$$P(A\mid B _1 \equiv b _1, \ldots, B _m \equiv b _m)

= \lim _{n\to\infty}

P^n(A\mid B _1 \equiv b _1, \ldots, B _m \equiv b _m),$$ where $b _i n \in \mathbb{N}$



Jeffrey conditionalization is a special case of partial conditional probability, in which the condition events must form a partition:



$$P(A\mid B _1 \equiv b _1, \ldots, B _m \equiv b _m)

= \sum^m _{i=1} b _i P(A\mid B _i)$$



\paragraph{Example}



Suppose that somebody secretly rolls two fair six-sided dice, and we wish to compute the probability that the face-up value of the first one is 2, given the information that their sum is no greater than 5.



\begin{itemize}
    \item Let $D _1$ be the value rolled on dice 1.
    \item Let $D _2$ be the value rolled on dice 2.
\end{itemize}

$\textbf{Probability that}$ $D _1 = 2$



Table 1 shows the sample space of 36 combinations of rolled values of the two dice, each of which occurs with probability 1/36, with the numbers displayed in the red and dark gray cells being $D _1 + D _2$.



$D _1 = 2$ in exactly 6 of the 36 outcomes; thus $P(D _1 = 2) = \frac{6}{36} = \frac{1}{6}$:



$$ \begin{array}{|cc|} \hline \hspace{-1em} \begin{array}{c} \begin{array}{c} \\\\ \textbf{+} \\\\ \hline \\\\ \\\\ D _1 \\\\ \\\\ \\\\ \\\\ \end{array} & \hspace{-1em} \begin{array}{|c} \\\\ \\\\ \hline 1 \\\\ \hline 2 \\\\ \hline 3 \\\\ \hline 4 \\\\ \hline 5 \\\\ \hline 6 \\\\ \end{array} \end{array} & \hspace{-2em} \begin{array}{c} D _2 \\\\ \begin{array}{|c|c|c|c|c|c|} \hline 1 & 2 & 3 & 4 & 5 & 6 \\\\ \hline 2 & 3 & 4 & 5 & 6 & 7 \\\\ \hline \textcolor{red}{3} & \textcolor{red}{4} & \textcolor{red}{5} & \textcolor{red}{6} & \textcolor{red}{7} & \textcolor{red}{8} \\\\ \hline 4 & 5 & 6 & 7 & 8 & 9 \\\\ \hline 5 & 6 & 7 & 8 & 9 & 10 \\\\ \hline 6 & 7 & 8 & 9 & 10 & 11 \\\\ \hline 7 & 8 & 9 & 10 & 11 & 12 \\\\ \hline \end{array} \end{array} \\\\ \hline \end{array} $$



$\textbf{Probability that}$ $D _1 + D _2 \leq 5$



Table 2 shows that $D _1 + D _2 \leq 5$ for exactly 10 of the 36 outcomes, thus $P(D _1 + D _2 \leq 5) = \frac{10}{36}$:



$$ \begin{array}{|cc|} \hline \hspace{-1em} \begin{array}{c} \begin{array}{c} \\\\ \textbf{+} \\\\ \hline \\\\ \\\\ D _1 \\\\ \\\\ \\\\ \\\\ \end{array} & \hspace{-1em} \begin{array}{|c} \\\\ \\\\ \hline 1 \\\\ \hline 2 \\\\ \hline 3 \\\\ \hline 4 \\\\ \hline 5 \\\\ \hline 6 \\\\ \end{array} \end{array} & \hspace{-2em} \begin{array}{c} D _2 \\\\ \begin{array}{|c|c|c|c|c|c|} \hline 1 & 2 & 3 & 4 & 5 & 6 \\\\ \hline \textcolor{red}{2} & \textcolor{red}{3} & \textcolor{red}{4} & \textcolor{red}{5} & 6 & 7 \\\\ \hline \textcolor{red}{3} & \textcolor{red}{4} & \textcolor{red}{5} & 6 & 7 & 8 \\\\ \hline \textcolor{red}{4} & \textcolor{red}{5} & 6 & 7 & 8 & 9 \\\\ \hline \textcolor{red}{5} & 6 & 7 & 8 & 9 & 10 \\\\ \hline 6 & 7 & 8 & 9 & 10 & 11 \\\\ \hline 7 & 8 & 9 & 10 & 11 & 12 \\\\ \hline \end{array} \end{array} \\\\ \hline \end{array} $$



$\textbf{Probability that}$ $D _1 = 2$ $\textbf{\textit{given that}}$ $D _1 + D _2 \leq 5$



Table 3 shows that for 3 of these 10 outcomes, $D _1 = 2$.



Thus, the conditional probability $P(D _1 = 2 \| D _1 + D _2 \leq 5) = \frac{3}{10} = 0.3$:



$$ \begin{array}{|cc|} \hline \hspace{-1em} \begin{array}{c} \begin{array}{c} \\\\ \textbf{+} \\\\ \hline \\\\ \\\\ D _1 \\\\ \\\\ \\\\ \\\\ \end{array} & \hspace{-1em} \begin{array}{|c} \\\\ \\\\ \hline 1 \\\\ \hline 2 \\\\ \hline 3 \\\\ \hline 4 \\\\ \hline 5 \\\\ \hline 6 \\\\ \end{array} \end{array} & \hspace{-2em} \begin{array}{c} D _2 \\\\ \begin{array}{|c|c|c|c|c|c|} \hline 1 & 2 & 3 & 4 & 5 & 6 \\\\ \hline \textcolor{red}{2} & \textcolor{red}{3} & \textcolor{red}{4} & \textcolor{red}{5} & 6 & 7 \\\\ \hline \fcolorbox{red}{white}{\textbf{\textcolor{red}{3}}} & \fcolorbox{red}{white}{\textbf{\textcolor{red}{4}}} & \fcolorbox{red}{white}{\textbf{\textcolor{red}{5}}} & 6 & 7 & 8 \\\\ \hline \textcolor{red}{4} & \textcolor{red}{5} & 6 & 7 & 8 & 9 \\\\ \hline \textcolor{red}{5} & 6 & 7 & 8 & 9 & 10 \\\\ \hline 6 & 7 & 8 & 9 & 10 & 11 \\\\ \hline 7 & 8 & 9 & 10 & 11 & 12 \\\\ \hline \end{array} \end{array} \\\\ \hline \end{array} $$



Here, in the earlier notation for the definition of conditional probability, the conditioning event \textit{B} is that $D _1 + D _2 \leq 5$, and the event \textit{A} is $D _1 = 2$. We have

$$P(A\mid B)=\tfrac{P(A \cap B)}{P(B)} = \tfrac{3/36}{10/36}=\tfrac{3}{10},$$ as seen in the table.



\paragraph{Use in inference}



In statistical inference, the conditional probability is an update of the probability of an event based on new information. The new information can be incorporated as follows:



\begin{itemize}
    \item Let \textit{A}, the event of interest, be in the sample space, say (\textit{X},\textit{P}).
    \item The occurrence of the event \textit{A} knowing that event \textit{B} has or will have occurred, means the occurrence of \textit{A} as it is restricted to \textit{B}, i.e. $A \cap B$.
    \item Without the knowledge of the occurrence of \textit{B}, the information about the occurrence of \textit{A} would simply be \textit{P}(\textit{A})
    \item The probability of \textit{A} knowing that event \textit{B} has or will have occurred, will be the probability of $A \cap B$ relative to \textit{P}(\textit{B}), the probability that \textit{B} has occurred.
    \item This results in $P(A|B) = P(A \cap B)/P(B)$ whenever \textit{P}(\textit{B}) \> 0 and 0 otherwise.
\end{itemize}

This approach results in a probability measure that is consistent with the original probability measure and satisfies all the Kolmogorov axioms. This conditional probability measure also could have resulted by assuming that the relative magnitude of the probability of \textit{A} with respect to \textit{X} will be preserved with respect to \textit{B} (cf. a Formal Derivation below).



The wording "evidence" or "information" is generally used in the Bayesian interpretation of probability. The conditioning event is interpreted as evidence for the conditioned event. That is, \textit{P}(\textit{A}) is the probability of \textit{A} before accounting for evidence \textit{E}, and \textit{P}(\textit{A}\|\textit{E}) is the probability of \textit{A} after having accounted for evidence \textit{E} or after having updated \textit{P}(\textit{A}). This is consistent with the frequentist interpretation, which is the first definition given above.



\paragraph{Statistical independence}



Events \textit{A} and \textit{B} are defined to be statistically independent if



$$P(A \cap B) = P(A) P(B).$$



If \textit{P}(\textit{B}) is not zero, then this is equivalent to the statement that



$$P(A\mid B) = P(A).$$



Similarly, if \textit{P}(\textit{A}) is not zero, then



$$P(B\mid A) = P(B)$$



is also equivalent. Although the derived forms may seem more intuitive, they are not the preferred definition as the conditional probabilities may be undefined, and the preferred definition is symmetrical in \textit{A} and \textit{B}.



\textbf{Independent events vs. mutually exclusive events}



The concepts of mutually independent events and mutually exclusive events are separate and distinct. The following table contrasts results for the two cases (provided that the probability of the conditioning event is not zero).



$$ \begin{array}{|c|c|c|} \hline  & \textbf{If statistically independent} & \textbf{If mutually exclusive} \\\\ \hline P(A \mid B) & P(A) & 0 \\\\ \hline P(B \mid A) & P(B) & 0 \\\\ \hline P(A \cap B) & P(A)P(B) & 0 \\\\ \hline \end{array} $$



In fact, mutually exclusive events cannot be statistically independent (unless both of them are impossible), since knowing that one occurs gives information about the other (in particular, that the latter will certainly not occur).



\paragraph{Common fallacies}



> \textit{These fallacies should not be confused with Robert K. Shope's 1978 "conditional fallacy", which deals with counterfactual examples that beg the question}.



\paragraph{Assuming conditional probability is of similar size to its inverse}



% [Image: A geometric visualisation of Bayes' theorem.]



Figure: A geometric visualization of Bayes' theorem. In the table, the values 2, 3, 6, and 9 give the relative weights of each corresponding condition and case. The figures denote the cells of the table involved in each metric, with the probability being the fraction of each figure that is shaded. This shows that

        $\newline$ $P(A \mid B) P(B) = P(B \mid A) P(A),$

        $\newline$ i.e., $P(A \mid B) = \frac{P(B \mid A) P(A)}{P(B)}.$

        $\newline$ Similar reasoning can be used to show that

        $\newline$ $P(\overline{A} \mid B) = \frac{P(B \mid \overline{A}) P(\overline{A})}{P(B)},$

        $\newline$ etc.



In general, it cannot be assumed that $P(A \mid B) \approx P(B \mid A)$. This can be an insidious error, even for those who are highly conversant with statistics. The relationship between $P(A \mid B)$ and $P(B \mid A)$ is given by Bayes' theorem:



$$ \begin{aligned} P(B\mid A) & = \frac{P(A\mid B) P(B)}{P(A)} \\\\ \Leftrightarrow \frac{P(B\mid A)}{P(A\mid B)} & = \frac{P(B)}{P(A)} \end{aligned} $$



That is, $P(A \mid B) \approx P(B \mid A)$ only if $\frac{P(B)}{P(A)} \approx 1$, or equivalently, $P(A) \approx P(B)$.



\paragraph{Assuming marginal and conditional probabilities are of similar size}



In general, it cannot be assumed that $P(A) \approx P(A \mid B)$. These probabilities are linked through the law of total probability:



$$P(A) = \sum _n P(A \cap B _n) = \sum _n P(A\mid B _n)P(B _n).$$



where the events $(B _n)$ form a countable partition of $\Omega$.



This fallacy may arise through selection bias. For example, in the context of a medical claim, let $S _{C}$ be the event that a sequela (chronic disease) $S$ occurs as a consequence of circumstance (acute condition) $C$. Let $H$ be the event that an individual seeks medical help. Suppose that in most cases, $C$ does not cause $S$ (so that $P(S _{C})$ is low). Suppose also that medical attention is only sought if $S$ has occurred due to $C$. From the experience of patients, a doctor may therefore erroneously conclude that $P(S _{C})$ is high. The actual probability observed by the doctor is $P(S _{C} \mid H)$.



\paragraph{Over- or under-weighting priors}



Not taking prior probability into account partially or completely is called \textit{base rate neglect}. The reverse, insufficient adjustment from the prior probability is \textit{conservatism}.



\paragraph{Formal derivation}



Formally, $P(A \mid B)$ is defined as the probability of $A$ according to a new probability function on the sample space, such that outcomes not in $B$ have probability 0 and that it is consistent with all original probability measures.



Let $\Omega$ be a sample space with elementary events $\\{\omega\\}$, and let $P$ be the probability measure with respect to the $\sigma$-algebra of $\Omega$. Suppose we are told that the event $B \subseteq \Omega$ has occurred. A new probability distribution (denoted by the conditional notation) is to be assigned on $\\{\omega\\}$ to reflect this. All events that are not in $B$ will have null probability in the new distribution. For events in $B$, two conditions must be met: the probability of $B$ is one, and the relative magnitudes of the probabilities must be preserved. The former is required by the axioms of probability, and the latter stems from the fact that the new probability measure has to be the analog of $P$ in which the probability of $B$ is one—meaning that every event that is not in $B$ therefore has a null probability. Hence, for some scale factor $\alpha$, the new distribution must satisfy:



1.  $\omega \in B : P(\omega\mid B) = \alpha P(\omega)$

2.  $\omega \not \in B : P(\omega\mid B) = 0$

3.  $\sum _{\omega \in \Omega} {P(\omega\mid B)} = 1.$



Substituting 1 and 2 into 3 to select \textit{a}:



$$\begin{aligned} 1 & = \sum _{\omega \in \Omega} {P(\omega \mid B)} && \\\\ & = \sum _{\omega \in B} {P(\omega \mid B)} + \cancel{\sum _{\omega \not \in B} P(\omega \mid B)} && \quad \text{ using } \sum _{\omega \not \in B} P(\omega \mid B) = 0 \hspace{20em} \\\\ & = \alpha \sum _{\omega \in B} {P(\omega)} && \\\\ & = \alpha \cdot P(B) && \\\\ \Rightarrow \alpha & = \frac{1}{P(B)} \end{aligned} $$



So the new probability distribution is



1.  $\omega \in B: P(\omega\mid B) = \frac{P(\omega)}{P(B)}$

2.  $\omega \not \in B: P(\omega\mid B) = 0$



Now for a general event \textit{A},



$$\begin{aligned} P(A\mid B) & = \sum _{\omega \in A \cap B} {P(\omega \mid B)} + \cancel{\sum _{\omega \in A \cap B^c} P(\omega\mid B)} && \quad \text{ using } \sum _{\omega \in A \cap B^c} P(\omega\mid B) = 0 \hspace{20em} \\\\ & = \sum _{\omega \in A \cap B} {\frac{P(\omega)}{P(B)}} && \\\\ & = \frac{P(A \cap B)}{P(B)} \end{aligned} $$ 



\paragraph{Licensing}



Content obtained and/or adapted from:



\begin{itemize}
    \item \href{https://en.wikipedia.org/wiki/Conditional_probability}{Conditional probability, Wikipedia} under a CC BY-SA license
\end{itemize}

\subsection{Learning Objectives}

\begin{enumerate}
    \item Here are three Student Learning Objectives (SLOs) for the topic of \textbf{conditional probability} in probability theory:
    \item \textbf{Apply Conditional Probability:}
    \item \textbf{Outcome:} Students will be able to calculate the conditional probability of an event given the occurrence of another event.
    \item \textbf{Details:} This includes understanding and applying the formula $P(A \mid B) = \frac{P(A \cap B)}{P(B)}$ to solve problems where students must determine the probability of event $A$ occurring under the condition that event $B$ has already occurred.
    \item \textbf{Interpret Independence of Events:}
    \item \textbf{Outcome:} Students will be able to identify and explain when two events are independent using conditional probability.
    \item \textbf{Details:} This includes recognizing that if $P(A \mid B) = P(A),$ then events $A$ and $B$ are independent, and understanding the implications of independence in the context of real-world scenarios.
    \item \textbf{Apply Bayes' Theorem:}
    \item \textbf{Outcome:} Students will be able to use Bayes' Theorem to update the probability of an event based on new information.
    \item \textbf{Details:} This includes using the formula $P(A \mid B) = \frac{P(B \mid A) \cdot P(A)}{P(B)}$ to compute the revised probability of an event $A$ given evidence $B$, and understanding the applications of this theorem in fields such as medicine, finance, and decision-making.
\end{enumerate}

\subsection{Assessment Instrument}

\begin{enumerate}
    \item \textbf{Assessment Instrument: Conditional Probability Quiz} \textbf{Format:}
    \item \textbf{Type:} Mixed-format quiz (multiple-choice, short answer, and problem-solving questions)
    \item \textbf{Duration:} 45 minutes
    \item \textbf{Total Marks:} 30 points \textbf{Sections:}
    \item \textbf{Understanding Conditional Probability (10 points):}
    \item \textbf{Question 1:} Define conditional probability and provide the formula used to calculate it. (Short Answer, 4 points)
    \item \textbf{Question 2:} Given the following probabilities: $P(A) = 0.4$, $P(B) = 0.5$, and $P(A \cap B) = 0.2$, calculate $P(A \mid B)$. (Problem-Solving, 6 points)
    \item \textbf{Independence of Events (10 points):}
    \item \textbf{Question 3:} Consider two events, $C$ and $D$, where $P(C) = 0.6$ and $P(D) = 0.3$. If $P(C \cap D) = 0.18$, are the events $C$ and $D$ independent? Justify your answer using the appropriate condition for independence. (Problem-Solving, 6 points)
    \item \textbf{Question 4:} If $P(A \mid B) = P(A)$, what does this imply about the relationship between events $A$ and $B$? (Short Answer, 4 points)
    \item \textbf{Application of Bayes' Theorem (10 points):}
    \item \textbf{Question 5:} A medical test for a disease has a 95\% accuracy rate, meaning the probability of testing positive if you have the disease is 0.95. The probability of having the disease is 0.01, and the probability of testing positive overall is 0.05. Using Bayes' Theorem, calculate the probability that a person actually has the disease given that they tested positive. (Problem-Solving, 10 points) \textbf{Scoring Rubric:}
    \item \textbf{Understanding Conditional Probability:}
    \item \textbf{Question 1:} 2 points for the correct definition, 2 points for the correct formula (total 4 points)
    \item \textbf{Question 2:} 2 points for setting up the problem correctly, 4 points for the correct answer (total 6 points)
    \item \textbf{Independence of Events:}
    \item \textbf{Question 3:} 2 points for identifying the correct formula for independence, 4 points for correctly applying it and interpreting the result (total 6 points)
    \item \textbf{Question 4:} 4 points for the correct explanation (total 4 points)
    \item \textbf{Application of Bayes' Theorem:}
    \item \textbf{Question 5:} 4 points for setting up Bayes' Theorem correctly, 6 points for the correct calculation and interpretation (total 10 points) The quiz will assess the students' understanding and application of conditional probability, independence of events, and Bayes' Theorem. It ensures that they can correctly interpret and calculate probabilities in various scenarios, fostering a deeper comprehension of probability theory.
\end{enumerate}

%%===============================================================
\section{Counting Rules}
\label{sec:counting-rules}
%%===============================================================

\paragraph{The Counting Principle}



Before we can delve into the properties of probability and odds, we need to understand the \textbf{Counting Principle}. We use the Counting Principle to determine how many different ways one can choose/do certain events. It is easier to define the Principle through examples:



\paragraph{Independent Events}



Let's say that John is at a deli sub restaurant. There are 3 different breads, 3 different cheeses, 3 different condiments, and 3 different vegetables he can place on his sandwich, assuming that he can only place one from each category on to his sandwich. How many different ways can he set up his sandwich?



Since choosing a cheese doesn't affect the number of choices of vegetables, condiments, or bread, these events are called \textbf{independent events}. For this problem, we will multiply $3\times3\times3\times3$ , so $3^4$ , which is 81. So there are 81 different possible combinations to form this sandwich.



\paragraph{Practice Problems}



1\) Thomas goes to a McDonalds' restaurant and chooses to create his own burger. He has 2 different kinds of cheese, 3 different breads, and 3 different sauces he can choose from, but he can only choose one of each category. How many different ways can he create this burger?



2\) Diane is ordering pizza for her family. There are 4 different possible sizes of the pizza. Also, she has to choose one of 5 toppings to place on the pizza and one of 3 different types of cheese for the pizza. In addition, she must choose one of 3 different kinds of crust. How many different ways can she have her pizza?



3\)



> a) How many 3-digit numbers can be formed from the digits 2, 3, 4, 5, 7 and 9?



> b) How many of these numbers are less than 400?



\paragraph{Answers}



1\) $(2)(3)(3)=18$



2\) $(4)(5)(3)(3)=180$



3\)



> a\) Since there are six available digits, the answer is $(6)(6)(6)=216$



> b\) For the value of the 3-digit number to be less than 400, we only have two choices for the first digit, namely 2 or 3. After that we can choose the other two digits freely. The answer is thus $(2)(6)(6)=72$.



\paragraph{Dependent Events}



Assume that John is now working at a library. He must put 5 books on a shelf in any order. How many different ways can he order the books on the shelf? Unlike the independent events, when John puts a book on the shelf, that eliminates one book from the remaining choices of books to put next on the shelf; thus these are referred to as \textbf{dependent events}. At first, he has 5 different choices, so our first number in the multiplication problem will be 5. Now that one is missing, the number is reduced to 4. Next, it's down to 3, and so on. So, the problem will be



$$(5)(4)(3)(2)(1)$$ However, there is a symbol for this very idea. A $!$ represents the term \textit{factorial}. So, for example, $3!=(3)(2)(1)$ . Factorials are very useful in statistics and probability.



Therefore, the problem can be rewritten as $5!$ , which ends up being equal to 120.



However, not all dependent event problems are that simple. Let's say that there are 10 dogs at a dog competition. How many different ways can one select the Champion AND the Runner-Up? This problem could actually be considered simpler than the last one, but it doesn't involve a factorial. So, how many different ways can the judge determine the Champion? Of course, there are 10 different dogs, so 10 different ways. Next, how many dogs are left to select the Runner-Up from? Well, you removed one, so it's down to 9. Instead of placing a factorial, though, you will only multiply 10 and 9, resulting in 90.



\paragraph{Independent Or Dependent?}



To help you differentiate between the two, we will do a few more examples, but we will have to decide if it is dependent or independent before solving the problem.



Let's say that you are creating a 5-digit garage door opener code (the digits would include 0-9). If there were absolutely no restrictions, would the events be independent of each other or dependent on each other? Of course, there are no restrictions, since you could have five 4's for the code, so to solve the problem, you multiply 10 by itself 5 times, resulting in 100000.



Alternatively, suppose that the first number of the code cannot be 0, and that there can be no repeated numbers whatsoever. Obviously these events are dependent on each other, because there cannot be any repetitions. Let's look at this problem one number at a time.



The first number can be all the numbers except 0, reducing the possible numbers to 9. The second number \textbf{can} be 0 this time, so the possible numbers returns to 10. However, it cannot be a repeat of the previous number, so there are 9 possible choices again. After that, the numbers will reduce by one each time, due to the fact that there cannot be repeats, so the problem would look like this



$$(9)(9)(8)(7)(6)=\frac{(9)(9!)}{(5!)}=27216$$ Now, just one more example. Let's say that you were choosing your schedule for school. There are 8 periods each day, and there are 7 classes to choose from. Nonetheless, you must have a lunch period during the 4th period. We can think of 4th period as non existent because it is in a constant position and therefore does not affect the possible choices. With 7 slots and 7 options, the answer is simply $7!$ .



$$(7!)=5040$$



\paragraph{Review Of The Counting Principle}



So, we use the Counting Principle to determine the different unique ways we can do something, such as a sandwich or a selection of classes. Sometimes, these events will affect each other, such as when you can't choose the same number twice for your garage door code, so they are dependent events. However, other times, one event has no effect on the next event, such as when you have different cheeses and breads to choose for your sandwich, so they are independent events. The Counting Principle is a fundamental mathematical idea and an essential part of probability.



\paragraph{Counting Rules}



\textbf{Rule 1:} If any one of $k$ mutually exclusive and exhaustive events can occur on each of $n$ trials, there are $k^n$ different sequences that may result from a set of such trials.



\begin{itemize}
    \item Example: Flip a coin three times, finding the number of possible sequences. $n=3\ ,\ k=2$ , therefore, $\newline$ $k^n=2^3=8$
\end{itemize}

\textbf{Rule 2:} If $k _1,\dots,k _n$ are the numbers of distinct events that can occur on trials $1,\dots,n$ in a series, the number of different sequences of $n$ events that can occur is $k _1\times\cdots\times k _n$ .



\begin{itemize}
    \item Example: Flip a coin and roll a die, finding the number of possible sequences. Therefore, $(k _1)(k _2)=(2)(6)=12$
\end{itemize}

\textbf{Rule 3:} The number of different ways that $n$ distinct things may be arranged in order is

$$n!=(1)(2)(3)\times\cdots\times(n-1)(n) \text{, where } 0!=1$$

An arrangement in order is called a permutation, so that the total number of permutations of $n$ objects is $n!$ (the symbol $n!$ Is called n-factorial).



\begin{itemize}
    \item Example: Arrange 10 items in order, finding the number of possible ways. Therefore,
\end{itemize}
$$10!=10\cdot9\cdot8\cdot7\cdot6\cdot5\cdot4\cdot3\cdot2\cdot1=3,628,800$$



\textbf{Rule 4:} The number of ways of selecting and arranging $k$ objects from among $n$ distinct objects is: $\frac{n!}{(n-k)!}$ , or as seen on calculators, \[nPr\].



\begin{itemize}
    \item Example: pick 3 things from 10 items, and arrange them in order. Therefore $n=10\ ,\ k=3$ , therefore $\frac{10!}{(10-3)!}=\frac{10!}{7!}=720$
\end{itemize}

\textbf{Rule 5:} The total number of ways of selecting $k$ distinct combinations of $n$ objects, irrespective of order (ie order NOT important), is: $\binom{n}{k}=\frac{n!}{k!(n-k)!}$ or as seen on calculators, \[nCr\].



\begin{itemize}
    \item Example: Pick 3 items from 10 in any order, where $n=10\ ,\ k=3$ , Therefore, $$\binom{10}{3}=\frac{10!}{3!(7!)}=\frac{720}{6}=120$$
\end{itemize}

\paragraph{Resources}



\begin{itemize}
    \item \href{https://www.hamilton.ie/ollie/EE304/Counting.pdf}{Some Simple Counting Rules}, Hamilton Institute
    \item \href{https://www.stat.purdue.edu/\textasciitilde{}huang251/slides4.pdf}{Basic Counting Rules, Permutations, Combinations}, Purdue University
\end{itemize}

\paragraph{Licensing}



Content obtained and/or adapted from:



\begin{itemize}
    \item \href{https://en.wikibooks.org/wiki/Probability/The_Counting_Principle}{The Counting Principle, Wikibooks} under a CC BY-SA license
\end{itemize}

\subsection{Learning Objectives}

\begin{enumerate}
    \item Student Learning Objectives (SLOs) for The Counting Principle
    \item \textbf{Apply the Counting Principle to Determine the Number of Outcomes:}
    \item \textbf{Outcome:} Students will be able to apply the Counting Principle to calculate the total number of possible outcomes for a series of independent events.
    \item \textbf{Details:} This includes recognizing independent events and using multiplication to determine the total number of combinations, as illustrated in scenarios like creating sandwiches or choosing outfits.
    \item \textbf{Differentiate Between Independent and Dependent Events:}
    \item \textbf{Outcome:} Students will be able to distinguish between independent and dependent events in various scenarios and correctly apply the appropriate counting method.
    \item \textbf{Details:} This involves understanding how the selection of one event affects the other, using examples such as arranging books on a shelf or selecting digits for a garage door code.
    \item \textbf{Use Factorials in Counting Dependent Events:}
    \item \textbf{Outcome:} Students will be able to use factorial notation to calculate the number of possible arrangements for dependent events.
    \item \textbf{Details:} This includes solving problems that involve arranging distinct objects in a specific order, recognizing when to apply factorial calculations, and understanding the concept of permutations as related to dependent events.
\end{enumerate}

\subsection{Assessment Instrument}

\begin{enumerate}
    \item \textbf{Assessment Instrument: Counting Rules Quiz} \textbf{Format:}
    \item \textbf{Type:} Mixed-format quiz (multiple-choice, short answer, and problem-solving questions)
    \item \textbf{Duration:} 45 minutes
    \item \textbf{Total Marks:} 30 points \textbf{Sections:}
    \item \textbf{Applying the Counting Principle (10 points):}
    \item \textbf{Question 1:} Determine the number of possible outcomes in the following scenarios:
    \item a. A pizza shop offers 3 types of crusts, 4 types of sauces, and 5 toppings. How many different pizzas can be made if you choose one of each?
    \item b. A password consists of 4 letters followed by 2 digits. If the letters and digits can be repeated, how many different passwords can be created?
    \item c. You roll a 6-sided die twice. How many different outcomes are possible?
    \item d. A student must choose 1 elective from 5 options and 1 sport from 3 options. How many different combinations can the student choose?
    \item \textbf{Independent and Dependent Events (10 points):}
    \item \textbf{Question 2:} Identify whether the events described are independent or dependent, and calculate the number of possible outcomes:
    \item a. A drawer contains 5 red socks and 3 blue socks. Two socks are drawn without replacement. How many different pairs of socks can be drawn?
    \item b. A lottery ticket requires selecting 6 numbers out of 49 without replacement. How many different combinations can be chosen?
    \item c. A card is drawn from a deck, and then a second card is drawn without replacing the first. How many different outcomes are possible?
    \item d. A student chooses a shirt and then a pair of pants from a selection of 4 shirts and 3 pants. How many different outfits can be created?
    \item \textbf{Using Factorials and Permutations (10 points):}
    \item \textbf{Question 3:} Solve the following problems using factorials or permutations:
    \item a. How many ways can 8 books be arranged on a shelf?
    \item b. A teacher has 6 different prizes and wants to give 3 of them to students as 1st, 2nd, and 3rd place awards. How many different ways can the prizes be awarded?
    \item c. How many ways can 4 runners finish a race if all runners finish at different times?
    \item d. A combination lock requires a 3-digit code. How many possible codes can be created if digits can be repeated? \textbf{Scoring Rubric:}
    \item \textbf{Applying the Counting Principle:} 2.5 points for each correct answer (total 10 points)
    \item \textbf{Independent and Dependent Events:} 2.5 points for each correct identification and calculation (total 10 points)
    \item \textbf{Using Factorials and Permutations:} 2.5 points for each correct solution (total 10 points) The quiz assesses students' understanding and application of the Counting Principle, their ability to distinguish between independent and dependent events, and their proficiency in using factorials and permutations to solve counting problems.
\end{enumerate}
